% Framework v3.027: exact 600-cell/Hurwitz return channels, Cartan rank, native clutch, Clifford constraints and conditional physical gates.
% Framework research-mathematical source integration v3.020; E3/E4 pending; no physical occurrence promoted.
% Current integrated source: Paper4B-QuantumResponse-v1_51.tex (PapersMinimal-v3_019; 16 November 2026).
% Current integrated paper source: Paper4B-QuantumResponse-v1_50.tex (v3_017).
% Prior source: Paper4B-QuantumResponse-v1_49.tex. New cross-DAG mathematics is scoped, not physical-occurrence certified.
% Paper-series v2.381: complete algebraic exceptional torsion atlas and all-cocycle twist-stability; original claims retained.
% Paper-series v2.380: unified bibliography, exact exceptional double/length/gauge mathematics in I-B; previous theorem labels retained.
% Paper-series v2.379: full equivariant Dirac/Hodge/eta/torsion atlas in Paper I-B; paper-root references synchronized.
% Paper-series v2.378: 2026 exceptional literature consolidation and reference synchronization. The mathematical additions are located in Paper I-B.
% Paper-series v2.377: reference-only rebind to current Paper I-B v1.1; original mathematics retained except new Paper I-B comparison chapter.
% Framework v2.370: derived exceptional-root/Hurwitz comparison in I-A/I-C; other papers reference-only synchronized.
% Framework v2.368: preserve prior II-B convention/unital fixes and verified II-D corrections; current references synchronized.
% Framework v2.358: current 30-root source reference synchronization; selectively updated III-E, III-F, VI-B, VI-C.
% Framework v2.356: new source-owned conditional results and source-registry synchronization.
% Framework v2.347: exact operational test-resolution tiers, one-body response iff criterion, terminating Majorana hierarchy, odd-transfer witness, cutoff cumulants, basis-free regular-CCR covariance. Prior domains retained.
% Framework v2.320: signed Petersen quantum sectors, mirror-transfer CP no-go, conditional interference and source-visible response; exact finite certificates. Preserves inherited original-source theorem bodies.
% Framework v2.275: exact positive dense-twin word cubature, finite repeatability and eighth-moment classification, nested sampling frames, joint action and POVM-jet exports; preserves v2.274.
% Framework v2.274: two-scalar positive convolution certificates, first-leaked exceptional coset detection, branch-symmetric twin mixing, and finite endpoint moments; all v2.273 branches preserved.
% Framework v2.273: restored missing Problem 3 theorems from v2.268; exact exceptional frames, twin geometry, convolution repeatability and projector/occurrence firewalls; merged on v2.272 without overwriting independent branches.
% Framework v2.272: merged v2.270 covariant native-readout mathematics with v2.271 Problem 4 finite-action Hopf hypercharge/holonomy corrections; no universal bulk-gauge uniqueness claim.
% Independent branches: v2.270 readout/guard exports and v2.270 K16 action-consumed charge closure; all prior scoped assumptions retained.
% Framework v2.267: exact Paper II-A converse/observable certificates, III-F action tests, V-B right orbit and V-C native gauge scope; prior v2.266 preserved.
% Framework v2.266: exact Paper II-C outcome-count/readout converse results and scoped downstream imports; 30-root references synchronized.
% Framework v2.265: II-C coordinate-stage constitution and exact first-symbol host inheritance; prior mathematical domains retained.
% Framework v2.264: II-B reverse and action-bridge integration; preserves prior mathematical domains; active 30-root references synchronized.
% Framework v2.263: Paper I-D adds bridge-controlled root-carry parity and source-relative old-cap cancellation; other root mathematical bodies retained; active 30-root references synchronized.
% Framework v2.262: merges Paper II-F three-gap closure from v2.260 into v2.261, preserving newer II-B/II-D repairs; 30-root references synchronized.
% Framework v2.261: merged Paper II-B GP2-level/geometric repairs; preserves v2.260 II-D repairs; active references synchronized.
% Framework v2.260: Paper II-D signed RS, finite-symbol typing, and conditional curvature/Jordan repairs; references synchronized.
% Framework v2.259: conditional native SU(2) atlas and explicit later discharge; v2.258 mathematical bodies retained elsewhere.
% Framework v2.258: preserves parallel II-A/dynamics update and adds complete readout/frequency closures.
% Framework v2.257: exact minimal complete position readout and arbitrary-preparation frequency law.
% Framework v2.256: exact native inertia and sharp bounded ordered-source closure; preserves concurrent v2.255 work.
% Framework v2.255: preserves parallel detector integration and adds exact finite readout closures.
% Framework v2.254: shared bibliography/reference sync; substantive Pasterski detector-Ward additions only IV-B/IV-C/VII-C.
% Framework v2.253: typed raw coordinate defect in I-A and synchronized VI-A import; active references synchronized.
% Framework v2.252: optional shell-frame source-return calculation in III-G; active references synchronized.
% Framework v2.251: forward common-action execution, joint source jets and scoped coupled control.
% Framework v2.251: exact full-heat variations, fermion source jets and explicit physical completion frontier.
% Framework v2.251: certified interacting source jets, pole stability and scoped completion proofs.
% Framework v2.251: interacting finite-source evaluation and coupled local vacuum classification.
% Framework v2.251: exact finite source execution and native causal common-parent bridge.
% Framework v2.251: constructive parent and scoped source-incidence audit; see release report.
% Framework v2.251: common-action two-jet, corrected source/response closure and synchronized references.
% Framework v2.244: scoped rooted source-order Schur diagnostics and synchronized paper references.
% Framework v2.242: signed rooted phase lift, complete transverse return and exact family cubic.
% Framework v2.241: proved rooted conormal geometry, conditional cap realization and scoped action tests.
% Framework v2.240: observable-minimal action and response certificates integrated; inherited v2.239 geometry retained.
% Framework v2.238: Paper IV restructured; inherited mathematics retained.
\documentclass[12pt,a4paper]{article}
\usepackage{mathtools}
\usepackage{float}
\usepackage{caption}
\usepackage{tikz}
\usepackage{tikz-cd}
\usetikzlibrary{arrows,positioning,calc,decorations.pathmorphing}
\usepackage{microtype}
\usepackage{xltabular}
% Paper-series v2.381: new I-B torsion/modular results, bibliography and current source roots.
% Paper-series v2.380: new exceptional bibliography and synchronized paper references.
% Paper-series v2.379 current 31-paper macro/aux registry.
% Paper-series v2.378: current 31-paper AUX/bibliography registry updated for Paper I-B v1.2 literature additions.
% Framework v2.371: current-source reference registry refreshed; old theorem labels retained.
% Framework v2.370: current 30-paper source/AUX root rebinding; native mathematical macros unchanged.
% Framework v2.368: synchronized merged 30-root registry; mathematical notation unchanged.
% Framework v2.358: active 30-paper source registry synchronized; existing mathematical macros unchanged.
% Framework v2.357: selective new source-owned certificates; active reference roots updated for five revised papers.
% Framework v2.356: active 30-root registry synchronized to the 30 papers in this release.
% Framework v2.275: exact positive dense-twin word cubature, finite repeatability and eighth-moment classification, nested sampling frames, joint action and POVM-jet exports; preserves v2.274.
% Framework v2.274: filename and root-source synchronization.
% Framework v2.273: restored missing Problem 3 theorems from v2.268; exact exceptional frames, twin geometry, convolution repeatability and projector/occurrence firewalls; merged on v2.272 without overwriting independent branches.
% Framework v2.272: merged v2.270 covariant native-readout mathematics with v2.271 Problem 4 finite-action Hopf hypercharge/holonomy corrections; no universal bulk-gauge uniqueness claim.
% Independent branches: v2.270 readout/guard exports and v2.270 K16 action-consumed charge closure; all prior scoped assumptions retained.
% Framework v2.269: synchronized 30 active AUX roots; mathematics and bibliography preserved from Macros-v2_161.
% Framework v2.266: Paper II-C converse-derived outcome resources; II-E, III-C, IV-A/B, V-A and VI-B downstream mathematical closures; 30-root registry synchronized.
% Framework v2.265: II-C coordinate-stage constitution and exact first-symbol host inheritance; prior mathematical domains retained.
% Framework v2.264: paper II-B reverse/audit and IV-A observable-minimal exchange; synchronized external references to 30 active roots.
% Framework v2.263: Paper I-C adds bridge-controlled root-carry parity and source-relative old-cap cancellation; other root mathematical bodies retained; active 30-root references synchronized.
% Framework v2.262: integrates Paper II-F radius, Casimir and readout results; active 30-root references synchronized.
% Framework v2.261: merged Paper II-B GP2-level/geometric repairs; preserves v2.260 II-D repairs; active references synchronized.
% Framework v2.260: Paper II-D signed RS, finite-symbol typing, and conditional curvature/Jordan repairs; references synchronized.
% Framework v2.259: conditional native SU(2) atlas and explicit later discharge; v2.258 mathematical bodies retained elsewhere.
% Framework v2.258: preserves parallel II-A/dynamics update and adds complete readout/frequency closures.
% Framework v2.257: exact minimal complete position readout and arbitrary-preparation frequency law.
% Framework v2.256: exact native inertia and sharp bounded ordered-source closure; preserves concurrent v2.255 work.
% Framework v2.255: preserves parallel detector integration and adds exact finite readout closures.
% Framework v2.254: Pasterski bibliography and synchronized 30-paper external-reference registry.
% Framework v2.253: typed raw coordinate defect in I-A and synchronized VI-A import; active references synchronized.
% Framework v2.252: synchronized source registry; optional III-G shell-frame calculation.
% Framework v2.251: forward common-action execution, joint source jets and scoped coupled control.
% Framework v2.251: exact full-heat variations, fermion source jets and explicit physical completion frontier.
% Framework v2.251: certified interacting source jets, pole stability and scoped completion proofs.
% Framework v2.251: synchronized active source registry; mathematical notation unchanged.
% Framework v2.251: synchronized active source registry; mathematical notation unchanged.
% Framework v2.251: synchronized active source registry; mathematical notation unchanged.
% Framework v2.251: synchronized source registry; no mathematical notation changes.
% Framework v2.244: synchronized 30-paper AUX registry and rooted phase-source action-order tests.
% Framework v2.242: signed rooted phase lift, complete transverse return and exact family cubic.
% Framework v2.241: proved rooted conormal geometry, conditional cap realization and scoped action tests.
% Framework v2.239: all 30 active source roots synchronized; no mathematical macro changes.
% Framework v2.238: Paper IV A/B/C split and 30-root reference registry.
% Framework v2.237: Paper VII-D technical split and 28-root reference registry.
% Framework v2.236: Paper VI-C technical-volume reference added; all current root bindings advanced.
% Framework v2.235: 26-root AUX and bibliography bindings synchronized; mathematical macros unchanged.
% BEGIN INLINE SOURCE: Paper00-TeXBuildCompatibility-v1_0.tex
% Rendering-only repairs for inherited long appendices and bitmap text fonts.
% Integer counters through 26 retain the original labels; later values use AA,
% AB, ... . This file changes no action coefficient, state or theorem premise.
\RequirePackage{alphalph}
\makeatletter
\let\BHPrimitiveAlph\@Alph
\newalphalph{\BHExtendedAlph}{\BHPrimitiveAlph}{26}
\let\@Alph\BHExtendedAlph
\AtBeginDocument{%
  \@ifpackageloaded{microtype}{\microtypesetup{expansion=false}}{}%
}
\makeatother

% END INLINE SOURCE: Paper00-TeXBuildCompatibility-v1_0.tex

% Framework v2.164: current reference/bibliography profile; mathematical definitions retained.
% Framework v2.162: finite metric, functional Hilbert and scoped Ward integration.
% Framework v2.161: inherits Macros-v2_57.tex exactly except the current external-reference binding.
% Macros v2.57 / framework v2.156: current reference registry only; mathematical definitions unchanged.
% Framework v2.143 / Macros v2.55: current external-reference bindings; mathematical notation unchanged.
% Framework v2.110 / Macros v2.54: current external-reference registry and bibliography; mathematical notation inherited.
% Framework v2.67 / Macros-v2_53: adds Born-action-cell and joint-EH readout symbols; absolute unit realization is metrological; MATCH_{J/scal}^{EH} remains the sole G24 normalization bridge.
% Framework v2.35 / Macros-v2_52: grouped paper numbering, active AUX registry, and Born-v6_20 bibliography binding synchronized; notation unchanged.
% Framework v2.26: active-source AUX and bibliography bindings synchronized; notation unchanged.
% Framework v2.09: reviewed early-reciprocity integration; scoped physical imports.
% Framework v2.08: Born--Euler mathematics and conditional complete-action potential criterion.
% Framework v2.07: active AUX registry includes Paper I-B v1.2; shared bibliography Born-v6_18.
% Macros-v2_46 / framework v2.04: mathematical atlas and downstream realization correction.
% Macros-v2_46 / framework v2.03: scoped structure integration and current dependency bindings.
% Macros-v2_46 / framework v2.01: shared source bindings; prior revision comments are historical.
% v2.42 / framework v1.99: Synchronized RG/AS owner and bibliography bindings.
% v2.41 / framework v1.99: active manuscript AUX bindings and geometric-closure bibliography synchronized.
% Release-v1.80 centralized source synchronization for scalar-circle/comparison-band integration.
% Release v1.79: Hurwitz hierarchy and kinetic-normalized action response integration; physical mass matching remains conditional.
% Release v1.78: approved M8 closure-batch integration; physical and analytic domains retained.
% Current v1.78 registry: all active TeX roots are synchronized below.
% Earlier release comments are historical scope records, not current status.
% Release-v1.77 synchronization: native same-owner Noether closure, generator--flow/history typing, analytic-certification firewalls, scalaron/K16, LF/QCD, quantum-Ward, and RAR boundary-owner transport are registered; shared mathematical notation is unchanged.
% Release-v1.76 synchronization: registry points to the v1.76 active paper roots.
% Release-v1.76 synchronization: registers the local A+Phi+g Noether/constraint closure and dynamics-frontier propagation; shared mathematical notation is unchanged.
% Release-v1.75 synchronization: registry points to the v1.75 active paper roots.
% Release-v1.75 synchronization: integrates the joint-Dirac Einstein bulk real form and the local odd-scalar heat-coefficient action on explicit analytic first-jet/interior domains; v1.74 Yang--Mills closure is retained. Fermionic reality, pure-metric curvature-squared/exact-nonlocal terms, noncanonical boundaries, physical vacuum selection, and actual solution-history tokens remain open.
% Release-v1.74 synchronization: registry points to the v1.74 active paper roots and retains the v1.73 theorem labels.
% Release-v1.74 synchronization: canonical Macros-v2_34.tex; generated local nonlinear Yang--Mills action, constrained connection-field phase space, source-free Gauss momentum map, and Legendre closure are integrated on their proved domains. The remaining non-gauge ACT_cl sectors and actual solution-history selection remain open.
% Release-v1.71 synchronization: closes the former JOINT_x/JOINT_5
% connection-occurrence gap by canonical structure-group extension of the
% nonflat Born connection, faithful associated representation, and the local
% cotangent--Clifford operator constructor.  Actual fluctuation members,
% actions, vacua, solutions, calibrations, and analytic QFT domains remain
% separately typed.  The active-AUX registry is advanced once.
% Release-v1.71 canonical macro and active-AUX registry.
% ======================================================================
% Macros-v2_39.tex -- canonical publication macro file
% Release v1.71 retains the v1.70 null-ray/QCD boundaries and additionally
% removes the duplicate structural connection actuality premise.  It does not
% claim global bundle topology, select a gauge/Higgs member, or promote the
% finite construction to an interacting continuum theory.
% Release v1.68 constructs the finite light-front total-fraction operators and
% parton number/momentum operator measures, proves their simplex and Hausdorff-limit
% theorems, and separates positive cut measures from scheme-renormalized
% light-ray distributions.  It also separates finite real Schur poles, stable
% norm-resolvent/Riesz pole-residue limits, and continued resonance
% pole/residue limits.  The Paper-VII-C scattering package is claim-indexed:
% local QFT and spectral claims import no S-matrix or infrared premise, while
% asymptotic amplitudes, LSZ, S-matrix poles, completeness, and cross sections
% import only their named massive or massless/long-range clauses.
% Release v1.67 factorizes BND48_H into a same-action Hessian/Schur
% descendant and still-open typed boundary-trace, critical-lift, and
% refinement-limit conditions.  It proves the finite-collar Dirichlet-to-
% Neumann law sqrt(y)tanh(L sqrt(y)), its nonuniform critical limit, and the
% finite-return ultra-infrared turnover.  The Paper-I-A Hardy constant mode is
% an exact geometric near-hit, not yet the physical metric boundary mode.
% The current local f(R) Hessian remains massive and therefore still refutes
% a RAR prediction on the minimal action branch.
% Release v1.66 derives the transverse depth-six normalizer bridge from the
% already generated C6 Reynolds image and Paper-I-A binary-dihedral normalizer.
% The bridge is the unique least 2D3-stable scalar-containing realization of
% the faithful quotient 2D3/C6=C2, has trace defect six and hence K=48, but is
% transverse to the nested active carrier.  Exact K=44 and K=52 contraction,
% semigroup, and energy-midpoint identities are registered as structural
% readouts rather than projector events.  The active E48 carrier gap remains
% unchanged.  A same-action boundary-symbol/Calderon intertwiner is isolated
% as the necessary and still unproved owner bridge to the static DtN sort;
% representation equivalence alone neither derives RAR nor creates a K=48
% action junction, pole, occupation, or cosmological transition.
% Release v1.65 derives the native E_48-centered carrier mass gap from the
% GP1--GP3-generated active spectrum, distinguishes its carrier-free interval
% from a particle located at E_48, and withdraws the unsupported promotion to
% a positive continuum Yang--Mills mass gap.  The latter remains a valid
% conditional transfer theorem on an enlarged continuum branch, but it is not
% an active framework prediction or an unfinished native-gap proof.
% Release v1.64 synchronizes Derivation Batches 14--22.  It registers the
% GP1--GP3 heat-profile orbit, same-Hessian electroweak response and pole
% ownership, K16/K48 action-response corrections, perturbation state and
% non-Gaussianity typing, finite-vacuum/gap/QCD limit repairs, factorized
% position-area ownership, the actual static metric response, and the ordered
% Wilson/string-tension limits.  Strict continuum convergence, positive
% physical gap, surviving higher cumulants, and positive tension remain open.
% Release v1.61 closes TIMEWIDTH16 as an independent direct-production datum.
% JUNC_{2|7} continues to own the finite intrinsic microscopic profile/width,
% while the actual same-action background/production history determines its
% physical-time pullback, support, and duration.  The current Paper-I-A--VII-C AUX
% registry is advanced once; shared mathematical notation is unchanged.
% Release v1.60 derives the final-octet primitive-remnant depletion, canonical
% Hopf-current representative, exact Hopf descent to dOmega/(4pi), and
% closed-FLRW a^{-3} neutral-current shape.  It retypes K=48 as the onset
% of the final one-unit geometric transfer and records cDM/RAR as distinct
% physical projections of that common neutral parent while preserving
% occurrence, Hilbert-owner, stress-matching, and static-kernel premises.
% Release v1.59 transports the v1.58 exceptional-shell/K16 closure through
% Papers III-F and VI-B--VII-C.  It replaces Paper VI-B's obsolete mixed-Hessian
% output typing by the third-variation pair vertex and Bogoliubov occupation
% support, derives the post-seed Higgs-pair feedback vertex and its exact
% alpha--delta kernel factors, propagates the unique K16/K48 action-junction
% split, identifies the QCD delta count with the upstream 2O shell atom only
% as a common-scalar numerical compatibility, and narrows THERM16 to dynamical
% approach/beta selection because the finite Gibbs/KMS target is already
% unique on a supplied simple observable algebra once its Hamiltonian and
% inverse temperature are supplied; direct-sum central weights remain typed.
% Release v1.58 registers the exceptional nearest-shell closure, the derived
% transverse K=16 action event, the unique interior action junction, the
% type-correct ordered pair vertex, the thin-layer squeeze law, and the
% radiation-like/matter-like kinematic split.  REG_0, transfer magnitude,
% physical layer width, and thermalization remain explicitly conditional.
% Release v1.57 registers the active-depth-six obstruction, the native
% K=48 carrier desert, the typed 7|9 owner reduction, and the dynamical
% reheating handoff across Papers III-B--VI-B.
% Release v1.56 registers the least-support n_2 inflation dressing of the
% persistent n_7 Laurent response, the analytic global-convexity and positive
% endpoint-measure proofs, the derived inflation core and K=16 coefficient
% freeze, and the no-particle dark-sector split into static RAR plus the
% conditional homogeneous effective-curvature response at reflected K=48.
% There is no n_6 action, scalaron-DM mass, or second Einstein--Hilbert term.
% Release v1.55 registers the protected-octet resolved/residual reflection
% q_sp+Z_W=8, its j<->8-j and K<->64-K scale covariance, and the derived
% rank-one-exit quadratic grade C_1^2/U=C_2 in Papers III-C, III-E, and VIII.
% The reflection is a typed trace/scale theorem, not a time reversal or a
% physical event-occurrence theorem; reheating transfer remains conditional.
% Release v1.54 registers the exact positive-part Laurent indices, the
% active-domain/terminal-endpoint action split, the no-second-EH firewall, and
% the q_DM=6/K_DM=48 scalaron readout typing across Papers III-D/E/F, VI,
% VIII, and IX.
% Release v1.53 registers the common nonnegative Laurent obstacle projection
% for n_1^+ and n_7^+, discharges the three former exit/freeze assumptions,
% and advances the Paper-III-D/E/F, VIII, and IX consumers.  Shared notation
% and bibliography ownership are unchanged.
% Release v1.51 registers the PQTI derivation batch and its updated
% Paper-III-A/C/D/E/F, VIII, and IX consumers.  Shared notation and
% bibliography ownership are unchanged.
% Release v1.50 advances the Paper-III-B AUX target to v1.40 after the
% rigorous audit repair; shared notation and bibliography ownership are unchanged.
% Release v1.49 advances the Paper-II-E/F AUX target to v3.14 after the
% minimal audit repair; shared notation and bibliography ownership are unchanged.
% Release v1.46 synchronizes Paper II-A's minimal frozen constructor-origin
% theorem and its five dependent AUX targets.  Shared notation is unchanged.
% Release v1.45 updates the Paper-II-B AUX target to v2.23 after the
% GP3-rigidity and appendix-placement repair; shared notation is unchanged.
% Release v1.44 advances the centralized Paper-I-A--VII-C AUX registry after the
% Paper-I-A v5.16 protected-completion scope repair.  Every current
% paper imports this one file; no paper-level opt-in or \externaldocument list
% is permitted.  The bibliography remains centralized here through Born-v6_13.
% Shared bibliography, theorem infrastructure, colors, notation alignment,
% and current Paper-I-A--VII-C external-reference registration.
% This file centralizes common imports, bibliography sorting, colors, theorem
% infrastructure, and the notation shared by the complete series.
% v2.9 (release v1.43) made the Paper-I-A--VII-C AUX registry automatic and
% centralized the bibliography and all external-reference ownership.
% v2.8 release-registry refresh (v1.33) registers the rooted generated-parent
% cascade and the POL/REG/free-law/11+4/Born-group closure revisions.  No
% shared mathematical notation or Paper-II-A fundamental principle is changed.
% v2.8 release-registry refresh (v1.32) registers the universal generated-
% continuation revisions in Papers II-A, III-A, III-C, IV, and VIII.  No
% shared mathematical notation or Paper-II-A fundamental principle is changed.
% v2.8 release-registry refresh (v1.30) registers the CRP4/projective-state
% source revisions and their direct registry-only consumers.  No shared
% mathematical notation or Paper-II-A axiom is changed.
% v2.7 release-registry refresh (v1.27) updates only the current external
% Paper-I-A--VII-C basenames after the placement/assumption synchronization; no
% shared mathematical notation or Paper-II-A axiom is changed.
% v2.7 synchronizes the geometric defect-repair revision: Paper-II-A v3.42,
% Paper-III-C v1.26, Paper-III-D v1.27, Paper-III-E v1.32, and Paper-III-F
% v1.34.  It registers the common trace-defect/spectral-filtration theorem,
% its owner-typed response pullback, and the reduced GRED and TERM interfaces.
% No shared mathematical notation is changed.
% v2.6 synchronizes Paper-III-E v1.31's cone-type repair, minimal-response
% classification, and explicit classical action-role boundary, together with
% Paper-III-F v1.32's downstream response-atlas synchronization.  No shared
% mathematical notation is changed.
% v2.5 synchronizes the Paper-III-D v1.26 faithful-trace-duality proof,
% exact ordered-owner joint locus, and typed alternative-presentation
% bookkeeping repair, together with its Paper-III-E and closure-frontier
% references.  No shared mathematical notation is changed.
% v2.4 synchronizes the Paper-III-C v1.25 orbit/process/owner repair, its
% downstream Papers III-B--F, V, VI, VIII, and IX revisions, and Born-v6_13.
% v2.3 synchronizes the Paper-II-A v3.41 witness/support/principal inheritance
% layer, the Paper-III-A v1.24 joint-realizability/readout firewall, all dependent
% host-rate revisions, and the current Paper-I-A--VII-B source registry.
% v2.2 synchronizes the generated active-realization, determinant-depth,
% dimension-weighted hierarchy, and typed cosmological-event interfaces.
% v2.1 synchronizes the exceptional-support interface revisions: Paper I-A
% v5.9, Paper II-A v3.30, Paper II-E/F v3.7, and Paper IV v4.16.
% v2.0 introduces the stage-safe G0 occurrence notation, synchronizes the
% active Paper-I-A--IV source registry (current rigorous Paper-III closure:
% III-A v1.17, III-B v1.22, III-C--E v1.17, III-F v1.14, IV v4.14), and makes
% external-reference self-skip depend on logical paper role rather than a
% versioned job name.
% ======================================================================

% ======================================================================
% macros.tex — Shared macros for the Born Geometry series
% ======================================================================

% ======================================================================
% 1. PACKAGES, BIBLIOGRAPHY, COLORS, AND UNIT DEFINITIONS
% ======================================================================
% Core packages shared across the complete Paper-I-A--VII-C series.  Loading them here makes
% the macro file a genuine series preamble while remaining compatible with existing
% paper preambles that load the same packages before \input.
\usepackage{amsmath}
\usepackage{amssymb}
\usepackage{amsthm}
\usepackage{mathrsfs}
\usepackage{array}
\usepackage{booktabs}
\usepackage{xcolor}
\makeatletter
\@ifpackageloaded{geometry}{}{\usepackage{geometry}\geometry{margin=1in}}
\makeatother
\usepackage{mdframed}

% Shared support packages required by the macros below and by most papers.
\usepackage{ragged2e}
\usepackage{longtable}
\usepackage{tabularx}
\usepackage{enumitem}
\usepackage[most]{tcolorbox}

% Series-wide numeric bibliography convention using the portable BibTeX path.
% All Papers I-A--VII-C use the ordinary \cite interface.
\usepackage[numbers,sort&compress]{natbib}
\providecommand{\printbibliography}{%
  \bibliographystyle{plainnat}%
  \bibliography{Born-v6_59}%
}

% Shared colors.  \providecolor avoids clobbering a paper-specific override.
\providecolor{postulatecolor}{RGB}{230,242,255}
\providecolor{compare}{RGB}{230,245,230}
\providecolor{definitioncolor}{RGB}{230,242,255}

\providecolor{openriskcolor}{RGB}{255,235,235}
\providecolor{bridgecolor}{RGB}{240,240,255}
\providecolor{gold}{RGB}{255,215,0}

% ======================================================================
% 2. CUSTOM ENVIRONMENTS
% ======================================================================
% Common table column types used in all papers.  Guarded so that a paper may
% define its own versions before loading this file.
\makeatletter
\@ifundefined{NC@find@L}{%
  \newcolumntype{L}[1]{>{\RaggedRight\arraybackslash}p{#1}}%
}{}
\@ifundefined{NC@find@C}{%
  \newcolumntype{C}[1]{>{\Centering\arraybackslash}p{#1}}%
}{}
\@ifundefined{NC@find@Y}{%
  \newcolumntype{Y}{>{\RaggedRight\arraybackslash}X}%
}{}
\makeatother

% ======================================================================
% 3. THEOREM STYLES AND ENVIRONMENTS
% ======================================================================
% Custom underlined theorem style (useful for remarks or notes)
\newtheoremstyle{underlined}
  {3pt}   % Space above
  {3pt}   % Space below
  {}      % Body font (empty means normal/roman text, typical for remarks)
  {}      % Indent amount
  {}      % Theorem head font (can be \itshape, \bfseries, etc.)
  {.}     % Punctuation after theorem head
  { }     % Space after theorem head
  {\underline{\thmname{#1}\thmnumber{ #2}}\thmnote{ (#3)}} % Theorem head formatting

% Standard plain theorem style environments.  Each environment is guarded so
% the shared macro file can be loaded after a paper-specific theorem setup.
\theoremstyle{plain}
\makeatletter
\@ifundefined{theorem}{\newtheorem{theorem}{Theorem}[section]}{}
\@ifundefined{lemma}{\newtheorem{lemma}{Lemma}[section]}{}
\@ifundefined{proposition}{\newtheorem{proposition}{Proposition}[section]}{}
\@ifundefined{corollary}{\newtheorem{corollary}{Corollary}[section]}{}
\makeatother

% Definition style environments (typically non-italicized body)
\theoremstyle{definition}
\makeatletter
\@ifundefined{definition}{\newtheorem{definition}{Definition}[section]}{}
\@ifundefined{postulate}{\newtheorem{postulate}{Postulate}[section]}{}
\@ifundefined{closure}{\newtheorem{closure}{Closure}[section]}{}
\makeatother

% Applying the custom underlined style
\theoremstyle{underlined}
\makeatletter
\@ifundefined{remark}{\newtheorem{remark}{Remark}[section]}{}
\@ifundefined{convention}{\newtheorem{convention}{Convention}[section]}{}
\@ifundefined{principle}{\newtheorem{principle}{Principle}[section]}{}
\makeatother

% Framed proof environment.
% Every \begin{proof} ... \end{proof} is automatically framed/shaded.
% This avoids mdframed/trivlist nesting, which can drop the top border after
% theorem-like environments.
\makeatletter
\renewenvironment{proof}[1][\proofname]{%
  \par\addvspace{6pt}%
  \noindent
  \begin{tcolorbox}[%
    enhanced,
    breakable,
    colback=black!4,
    colframe=black!35,
    boxrule=0.35pt,
    arc=1pt,
    left=6pt,
    right=6pt,
    top=4pt,
    bottom=4pt,
    boxsep=0pt,
    before skip=0pt,
    after skip=6pt,
    parbox=false
  ]%
  \pushQED{\qed}%
  \normalfont\itshape #1\@addpunct{.}\enspace
  \normalfont\ignorespaces
}{%
  \unskip\nobreak\hfill\popQED
  \end{tcolorbox}%
  \par\addvspace{2pt}%
}
\makeatother

\mdfdefinestyle{postulatebox}{
  backgroundcolor=yellow!10,
  linecolor=black,
  linewidth=0.4pt,
  skipabove=\baselineskip,
  skipbelow=\baselineskip,
  nobreak=true
}

% ======================================================================
% PROOF-READABILITY HELPERS
% ======================================================================

\newcommand{\ProofPlan}[1]{%
  \par\smallskip
  \noindent\textbf{Proof plan. }#1%
  \par\smallskip
}

\newlist{proofsteps}{enumerate}{2}
\setlist[proofsteps,1]{
  label=\textbf{(\alph*)},
  ref=\alph*,
  leftmargin=2.2em,
  itemsep=0.45\baselineskip,
  topsep=0.35\baselineskip
}
\setlist[proofsteps,2]{
  label=\textbf{(\roman*)},
  ref=\roman*,
  leftmargin=2.0em,
  itemsep=0.25\baselineskip,
  topsep=0.2\baselineskip
}


% ======================================================================
% 4. NUMBER SYSTEMS AND ALPHABETS
% ======================================================================
\providecommand{\RR}{\mathbb{R}} % Real numbers
\providecommand{\PP}{\mathbb{P}} % Projective space / Probability
\providecommand{\CC}{\mathbb{C}} % Complex numbers
\providecommand{\ZZ}{\mathbb{Z}} % Integers
\providecommand{\NN}{\mathbb{N}} % Natural numbers
\providecommand{\HH}{\mathbb{H}} % Quaternions / Hyperbolic space
\providecommand{\LFRayGeom}{\mathsf{LF\!\mbox{-}RAY\!\mbox{-}GEOM}}
\providecommand{\LFRep}{\mathsf{LF\!\mbox{-}REP}}
\providecommand{\LFState}{\mathsf{LF\!\mbox{-}STATE}}
\providecommand{\LFLim}{\mathsf{LF\!\mbox{-}LIM}}
\providecommand{\LFScat}{\mathsf{LF\!\mbox{-}SCAT}}
% Number systems / projective spaces
\providecommand{\CP}{\mathbb{CP}}
\providecommand{\RP}{\mathbb{RP}}

% ======================================================================
% 5. GROUPS AND LIE THEORY
% ======================================================================
\providecommand{\SO}{\ensuremath{\mathrm{SO}}}     % Special Orthogonal group
\providecommand{\SU}{\ensuremath{\mathrm{SU}}}     % Special Unitary group
\providecommand{\Spin}{\ensuremath{\mathrm{Spin}}} % Spin group
\providecommand{\CS}{\ensuremath{\mathrm{CS}}}     % Chern-Simons (capitalized)
\providecommand{\cs}{\ensuremath{\mathrm{cs}}}     % Chern-Simons (lowercase)

% ======================================================================
% 6. COMMON MANIFOLDS AND TOPOLOGY
% ======================================================================
\providecommand{\Sone}{S^1}                 % Standard circle

% ======================================================================
% 7. MATH OPERATORS (Declared via \DeclareMathOperator)
% ======================================================================
\DeclareMathOperator{\End}{End}   % Endomorphisms
\DeclareMathOperator{\Aut}{Aut}   % Automorphisms
\DeclareMathOperator{\spec}{spec} % Lowercase spectrum
\DeclareMathOperator{\tr}{tr}     % Lowercase trace (if you already have \Tr)
\DeclareMathOperator{\supp}{supp} % Support
\DeclareMathOperator{\TrA}{Tr_{A}} % Partial trace over subsystem A
\DeclareMathOperator{\TrB}{Tr_{B}} % Partial trace over subsystem B
% Lie/algebra operators
\DeclareMathOperator{\Ad}{Ad}
\DeclareMathOperator{\Sym}{Sym}
\DeclareMathOperator{\su}{su}
% ======================================================================
% 8. STANDARD OPERATOR-LIKE SYMBOLS (Safe definitions via \providecommand)
% ======================================================================
\providecommand{\vol}{\operatorname{vol}}   % Volume form
\providecommand{\Vol}{\operatorname{Vol}}   % Total volume
\providecommand{\Ric}{\operatorname{Ric}}   % Ricci tensor
\providecommand{\Rm}{\operatorname{Rm}}     % Riemann tensor
\providecommand{\Tr}{\operatorname{Tr}}     % Uppercase trace
\providecommand{\Spec}{\operatorname{Spec}} % Uppercase spectrum
\providecommand{\diag}{\operatorname{diag}} % Diagonal matrix
\providecommand{\rank}{\operatorname{rank}} % Rank
\providecommand{\scal}{\mathcal R} % Scalar curvature
\providecommand{\Cl}{\operatorname{Cl}}     % Clifford algebra
\providecommand{\diam}{\operatorname{diam}} % Diameter

% ======================================================================
% 9. FUNCTIONS, TRIGONOMETRY, AND CALCULUS
% ======================================================================

% Codifferential naming conventions
% (Avoid \ddag: commonly defined by physics.sty and symbol packages)
\providecommand{\codiff}{\mathrm{d}^{\dagger}} % Codifferential operator

% Quantum mechanics bra-ket notation (needed for basis tables)
\providecommand{\ket}[1]{\lvert #1 \rangle}
\providecommand{\bra}[1]{\langle #1 \rvert}
\providecommand{\braket}[2]{\langle #1 \mid #2 \rangle}

% ======================================================================
% 10. PHYSICAL SCALES AND CONSTANTS (Paper 1c / Born Geometry)
% ======================================================================
\providecommand{\Rp}{\ensuremath{R_{\mathrm P}}}     % matched common EH curvature scale (not defined from observed G)
\providecommand{\hBorn}{\ensuremath{h_{\mathrm B}}} % physical representative 2\pi\hbar of one primitive Born scalar action cell
\providecommand{\UEHjoint}{\ensuremath{U_{J}}} % joint-Dirac EH curvature-coordinate scale 8 kappa_J
\providecommand{\Gread}{\ensuremath{G_{\rm scal}}} % scalar-action coordinate c_op^3/(h_B U); physical matching explicit
\providecommand{\GreadJoint}{\ensuremath{G_{J,\mathrm{raw}}}} % raw ordinary-trace coordinate = G_scal/96
\providecommand{\Rds}{R_{\mathrm{dS}}}    % de Sitter Curvature
\providecommand{\Mpl}{M_{\mathrm{Pl}}}    % Planck mass
\providecommand{\lp}{\ell_{\mathrm{Pl}}}  % Planck length
\providecommand{\Pmax}{P_{\mathrm{max}}}  % Maximum momentum
\providecommand{\Mplbar}{\bar{M}_{\mathrm{Pl}}} % Reduced Planck mass
\providecommand{\eff}{\mathrm{eff}}       % Subscript for effective quantities
\providecommand{\Cosmo}{\ensuremath{\Lambda}}                % Cosmological constant (use \Cosmo in source)
\providecommand{\LCDM}{\ensuremath{\Cosmo\mathrm{CDM}}}  % \LambdaCDM (typesets with \Cosmo)
\providecommand{\Lcut}{\ensuremath{\Lambda_{\mathrm{cut}}}}  % Cutoff scale (brace when used as script)
\providecommand{\Rcs}{\ensuremath{R_{\mathrm{cs}}}}      % Configuration-space radius
\providecommand{\CosmoEff}{\ensuremath{\Cosmo_{\eff}}}   % Effective cosmological constant

% Disambiguation: \Rmom is the momentum-space (S^3) radius. 
% Reserve plain R for the 4D curvature scalar in f(R).
\providecommand{\Rmom}{\ensuremath{R_{\mathrm{mom}}}} 
\providecommand{\Punit}{P_{0}} % Momentum unit set by \Rmom (= \hbar/\Rmom)
\providecommand{\Rvec}{\ensuremath{R_{\mathrm{vec}}}} % invariant vector-space radius in the reduced round sector
% Specific numeric values and expressions
\providecommand{\Vlogcap}{L_{\mathrm{Hopf}}^{(1)}}
\providecommand{\Hopf}{\operatorname{Hopf}}
% Scalar phase (S^1) sector + universal cover notation (Paper I-A)
\providecommand{\tper}{T_\star}      % Time period/scale
\providecommand{\Eper}{E_\star}      % Energy period/scale
\providecommand{\ttilde}{\tilde t}   % Universal cover time
\providecommand{\Etilde}{\tilde E}   % Universal cover energy

% ======================================================================
% 10A. SERIES-WIDE LOGICAL / NOTATIONAL DISAMBIGUATIONS
% ======================================================================
% Use these symbols to avoid collisions among: (i) the Paper--III
% no-insertion/geometric-closure class, (ii) Paper--I-A trimer diagonals,
% (iii) Paper--V-A/V-C family-space diagonals, (iv) scalar-index running,
% and (v) protected shell-clock labels.
\providecommand{\Rprim}{\ensuremath{\mathsf R_{\mathrm{prim}}^{\Sigma}}} % primitive finite structural readouts
\providecommand{\Rloop}{\ensuremath{\mathsf R_{\circ}^{\Sigma}}} % closed-loop readout role
\providecommand{\Rbdry}{\ensuremath{\mathsf R_{\partial}^{\Sigma}}} % boundary-conditioned readout role
\providecommand{\Rgen}{\ensuremath{\mathsf R_{\mathrm{fin}}^{\Sigma}}} % finite generated closure of primitive roles
\providecommand{\PhysOcc}{\operatorname{PhysOcc}}
\providecommand{\Actual}{\operatorname{Actual}}
\providecommand{\GeomLic}{\operatorname{GeomLic}}
\providecommand{\InvDist}{\operatorname{InvDist}}
\providecommand{\Dcl}{\ensuremath{\mathsf{D}_1}}      % unmarked positive Dirac-even closed spectral-readout form class
\providecommand{\Dheat}{\ensuremath{\mathsf D_{\mathrm{heat}}}} % one-parameter heat-semigroup subclass
\providecommand{\DtriI}{\ensuremath{D^{\mathrm{tri}}_1}} % Paper--I-A cyclic-trimer diagonal
\providecommand{\DtriII}{\ensuremath{D^{\mathrm{tri}}_2}}
\providecommand{\DfamI}{\ensuremath{D^{\mathrm{fam}}_1}} % Family-space diagonal selector
\providecommand{\DfamII}{\ensuremath{D^{\mathrm{fam}}_2}}
\providecommand{\alpharun}{\ensuremath{\alpha_{\mathrm{run}}}} % scalar spectral-index running
\providecommand{\Ksh}{\ensuremath{K_{\mathrm{sh}}}} % protected shell-clock coordinate
\providecommand{\ddef}{\ensuremath{\delta_{\mathrm{def}}}} % universal defect scalar after Paper--III realization

% ======================================================================
% 10B. SERIES-WIDE LOGICAL STATUS TAGS
% ======================================================================
% Upright small-caps wrapper for audit/status labels.  It prevents LaTeX from
% trying to use unavailable small-caps italic font shapes inside theorem bodies
% such as TU/lmr/m/scit.
\providecommand{\AuditLabel}[1]{}
\renewcommand{\AuditLabel}[1]{{\normalfont\scshape #1}}

% These tags are bookkeeping labels for consistency audits. They do not
% create new mathematical hypotheses; they record how a datum is being used
% in the proof chain.
\providecommand{\StatusPrimitive}{}
\renewcommand{\StatusPrimitive}{\AuditLabel{Primitive}}
\providecommand{\StatusDerived}{}
\renewcommand{\StatusDerived}{\AuditLabel{Derived}}
\providecommand{\StatusImported}{}
\renewcommand{\StatusImported}{\AuditLabel{Imported}}
\providecommand{\StatusConditional}{}
\renewcommand{\StatusConditional}{\AuditLabel{Conditional}}
\providecommand{\StatusBridge}{}
\renewcommand{\StatusBridge}{\AuditLabel{Bridge-derived}}
\providecommand{\StatusComparison}{}
\renewcommand{\StatusComparison}{\AuditLabel{Comparison-only}}

% ======================================================================
% 10C. SERIES-WIDE CLAIM-WORDING AUDIT LABELS
% ======================================================================
% These are textual audit labels used to keep modal language consistent
% across the series.  They introduce no new mathematical assumptions.
\providecommand{\ClaimForced}{}
\renewcommand{\ClaimForced}{\AuditLabel{forced}}
\providecommand{\ClaimUnique}{}
\renewcommand{\ClaimUnique}{\AuditLabel{unique}}
\providecommand{\ClaimCanonical}{}
\renewcommand{\ClaimCanonical}{\AuditLabel{canonical}}
\providecommand{\ClaimMinimal}{}
\renewcommand{\ClaimMinimal}{\AuditLabel{minimal}}
\providecommand{\ClaimClosed}{}
\renewcommand{\ClaimClosed}{\AuditLabel{closed}}
\providecommand{\ClaimPrediction}{}
\renewcommand{\ClaimPrediction}{\AuditLabel{prediction}}

% ======================================================================
% 10D. SERIES-WIDE THEOREM-CHAIN AUDIT LABELS
% ======================================================================
% These labels mark the direction of load-bearing theorem chains. They are
% audit annotations only; they introduce no new mathematical assumptions.
\providecommand{\ChainInput}{}
\renewcommand{\ChainInput}{\AuditLabel{Chain input}}
\providecommand{\ChainOutput}{}
\renewcommand{\ChainOutput}{\AuditLabel{Chain output}}
\providecommand{\ChainGuard}{}
\renewcommand{\ChainGuard}{\AuditLabel{Chain guard}}
\providecommand{\ChainNoBackflow}{}
\renewcommand{\ChainNoBackflow}{\AuditLabel{No backflow}}

% ======================================================================
% 10E. SERIES-WIDE LOCAL THEOREM-AUDIT LABELS
% ======================================================================
% These labels mark a statement-level dependency audit around individual
% theorem/definition nodes. They are bookkeeping labels only.
\providecommand{\LocalAudit}{}
\renewcommand{\LocalAudit}{\AuditLabel{Local audit}}
\providecommand{\UsesOnly}{}
\renewcommand{\UsesOnly}{\AuditLabel{Uses only}}
\providecommand{\Exports}{}
\renewcommand{\Exports}{\AuditLabel{Exports}}
\providecommand{\DoesNotUse}{}
\renewcommand{\DoesNotUse}{\AuditLabel{Does not use}}

% ======================================================================
% 10F. SERIES-WIDE NUMERICAL-ROW AUDIT LABELS
% ======================================================================
% These labels classify individual rows in numerical, benchmark, status,
% and comparison tables.  They are audit annotations only and do not change
% any theorem, value, formula, or comparison number.
\providecommand{\RowAudit}{}
\renewcommand{\RowAudit}{\AuditLabel{Numerical-row audit}}
\providecommand{\RowTheorem}{}
\renewcommand{\RowTheorem}{\AuditLabel{Theorem-fixed row}}
\providecommand{\RowBridge}{}
\renewcommand{\RowBridge}{\AuditLabel{Bridge-fixed row}}
\providecommand{\RowRepresentative}{}
\renewcommand{\RowRepresentative}{\AuditLabel{Representative row}}
\providecommand{\RowComparison}{}
\renewcommand{\RowComparison}{\AuditLabel{Comparison row}}
\providecommand{\RowCalibrated}{}
\renewcommand{\RowCalibrated}{\AuditLabel{Calibrated row}}
\providecommand{\RowEstimate}{}
\renewcommand{\RowEstimate}{\AuditLabel{Estimate row}}
\providecommand{\RowOpen}{}
\renewcommand{\RowOpen}{\AuditLabel{Open row}}
\providecommand{\RowNoSelector}{}
\renewcommand{\RowNoSelector}{\AuditLabel{No selector}}

% ======================================================================
% 10G. TABLE HELPERS
% ======================================================================
% Table helpers used by Paper V-C.  The unused readout-label preset has been
% removed; only active table helpers remain.
\providecommand{\gpsection}[1]{%
  \addlinespace[0.55em]
  \multicolumn{3}{@{}l}{\textsc{#1}}\\[-0.25em]
  \midrule
}
\providecommand{\mathcell}[1]{%
  \(\begin{gathered}#1\end{gathered}\)%
}

\newcommand{\DerBy}[1]{%
  \begin{tabular}[t]{@{}l@{}}#1\end{tabular}%
}

\newcommand{\TableNotes}[1]{%
  \par\medskip
  \begin{minipage}{0.94\linewidth}
  \footnotesize\emph{Notes.} #1
  \end{minipage}\par
}
% ======================================================================
% 10H. PAPER VI-A / QCD NOTATION
% ======================================================================
% Shared symbols used by Paper VI-A and downstream QCD/Yang--Mills handoffs.
\providecommand{\Pqcd}{\mathfrak P_{\rm VIII}^{\rm QCD}}
\providecommand{\Iqcd}{\mathfrak I_{\rm QCD}^{(7)}}
\providecommand{\muQCD}{\mu_{\rm QCD}^{\rm int}}
\providecommand{\LambdaQCD}{\Lambda_{\rm QCD}}
\providecommand{\LambdaBH}{\Lambda_{\rm BH}^{\rm QCD}}
\providecommand{\HQCD}{\mathcal H_{\rm QCD}^{\rm BH}}
\providecommand{\HQCDfin}{H_{\rm QCD}^{\rm fin}}
\providecommand{\Hraw}{\mathcal H_{\rm raw}}
\providecommand{\HphysQCD}{\mathcal H_{\rm phys}^{\rm QCD}}
\providecommand{\Vcol}{V_C}
\providecommand{\Ps}{P_{\rm sing}}
\providecommand{\PiOct}{\Pi_8}
\providecommand{\MSbar}{\overline{\rm MS}}
\providecommand{\YM}{\mathrm{YM}}
\providecommand{\Ktri}{K_{\triangle}}
\providecommand{\KYM}{K_{\YM}^{\rm BH}}
\providecommand{\aBH}{a_{\rm BH}}
\providecommand{\epsBH}{\varepsilon_{\rm BH}}
\newcommand{\RowConstructive}{\textnormal{\textsc{Constructive}}}
\newcommand{\RowObligation}{\textnormal{\textsc{Bridge obligation}}}
\newcommand{\RowTarget}{\textnormal{\textsc{Target row}}}
\newcommand{\RowValidation}{\textnormal{\textsc{Validation row}}}
\providecommand{\slashed}[1]{\not{#1}}

% ======================================================================
% 10I. PAPER VI-B / CONTINUOUS-QFT NOTATION
% ======================================================================
% Shared symbols used by Paper VI-B and any downstream continuous-QFT or
% Yang--Mills handoff paper.  All definitions are guarded so individual
% papers may override them locally before loading this file if needed.
\providecommand{\A}{\mathfrak A}
\providecommand{\Hphys}{\mathcal H_{\rm phys}}
\providecommand{\QCD}{\mathrm{QCD}}
\providecommand{\BH}{\mathrm{BH}}
\providecommand{\loc}{\mathrm{loc}}

\providecommand{\RBH}{\mathcal R_{\rm BH}}
\providecommand{\DG}{\mathfrak D_G}
\providecommand{\PW}{\mathcal{P\!W}} % Peter--Weyl spectral-symbol subspace
\providecommand{\Toe}{\operatorname{Toe}} % Toeplitz compression/representative
\providecommand{\Hol}{\operatorname{Hol}} % Holonomy

% ======================================================================
% 11. LATE-LOADED PACKAGES (Must be at the end of the preamble)
% ======================================================================
% Load late packages only if a paper has not already loaded them, avoiding
% option clashes in paper-specific preambles.
\makeatletter
\@ifpackageloaded{hyperref}{}{\usepackage{hyperref}}
\makeatother
% --- Series-wide disambiguations (Pass 3) ---
\providecommand{\diracdelta}{\delta_{\rm D}} % Dirac delta distribution (delta reserved for hierarchy)
\providecommand{\Hg}{H_{\mathrm g}} % Hodge harmonic projector (keep H for Hubble)
\providecommand{\Gr}{G_{\mathrm r}} % Green operator (keep G for Newton)
\providecommand{\FD}{\mathrm{D}} % functional/variational derivative symbol (delta reserved)
\providecommand{\Id}{\mathbf{1}} % identity/Kronecker tensor (avoid \delta_{ij})


% BEGIN INLINE SOURCE: Paper00-ActiveExternalReferences-v2_130.tex
% Framework v2.191: current reference binding; existing definitions retained.
% Framework v2.190: current22-root registry with native clock-evolution proofs.
% Framework v2.189: exact current22-root external registry, including retained v2.188 hardening owners.
% Framework v2.187: current registry synchronization only; source release v2.186 preserved.
% Framework v2.186: authoritative current 22-paper registry.
% Framework v2.182: all current 22 logical paper AUX bindings.
% Framework v2.170: current I-B v1.43 external AUX binding.
% Framework v2.167: exact current 22-root registry after CP--BL2 integration.
% Framework v2.166: exact current 22-root compact-release AUX registry.
% Framework v2.165: authoritative current root bindings.
% Framework v2.164: exact current 22 manuscript AUX bindings.
% Framework v2.162: finite metric, functional Hilbert and scoped Ward integration.
% Framework v2.161: current 22 manuscript AUX bindings; source only.
% Framework v2.156 / registry v2.113: synchronized current 22 manuscript roots; source-only registry.
% Framework v2.143 / registry v2.111: active manuscript AUX basenames; no generated AUX or PDF included.
% ======================================================================
% Central external-reference imports, numbers only, no PDF jumping
% ======================================================================

% xr-hyper is loaded centrally so that papers may retain their existing
% \usepackage{xr-hyper} lines without redefining \externaldocument after this
% registry is installed.  If plain xr was loaded earlier, retain it.
\makeatletter
\@ifpackageloaded{xr-hyper}{}{%
  \@ifpackageloaded{xr}{}{\usepackage{xr-hyper}}%
}
\makeatother

% Make external-document registration idempotent.  The registry identifies a
% logical paper by its Paper<number><optional letter> prefix.  Prefix-qualified
% imports remain independent.  Current paper roots contain no local
% \externaldocument calls; this guard also makes accidental duplicates harmless.
\ExplSyntaxOn
\prop_new:N \g_bh_external_refs_prop
\cs_new_eq:NN \BHOriginalExternalDocument \externaldocument

\cs_new_protected:Npn \bh_external_document:nnnn #1#2#3#4
 {
  \tl_set:Nn \l_tmpa_tl {#3}
  \regex_replace_once:nnN
   { \A (?: xref/ )? (Paper[0-9]+[A-F]?) .* \Z }
   { \1 }
   \l_tmpa_tl
  \tl_set:Nx \l_tmpb_tl { #1 | \tl_use:N \l_tmpa_tl }
  \prop_if_in:NVTF \g_bh_external_refs_prop \l_tmpb_tl
   {
    \PackageInfo{BornHopf}
     {Skipping~duplicate~external~refs~for~#3}
   }
   {
    \prop_gput:NVn \g_bh_external_refs_prop \l_tmpb_tl {#3}
    \IfNoValueTF{#4}
     { \BHOriginalExternalDocument[#1][#2]{#3} }
     { \BHOriginalExternalDocument[#1][#2]{#3}[#4] }
   }
 }

\RenewDocumentCommand{\externaldocument}{O{} O{} m o}
 { \bh_external_document:nnnn {#1}{#2}{#3}{#4} }

\NewDocumentCommand{\BHExternalRefs}{m}
 {
  % Do not import the aux file of the logical paper currently being compiled,
  % even when the job and registry use different version suffixes.
  \tl_set:Nx \l_tmpa_tl { \c_sys_jobname_str }
  \regex_replace_once:nnN
   { \A (Paper[0-9]+[A-F]?) .* \Z }
   { \1 }
   \l_tmpa_tl
  \tl_set:Nn \l_tmpb_tl { #1 }
  \regex_replace_once:nnN
   { \A (Paper[0-9]+[A-F]?) .* \Z }
   { \1 }
   \l_tmpb_tl
  \str_if_eq:eeTF { \tl_use:N \l_tmpa_tl } { \tl_use:N \l_tmpb_tl }
   {
    \PackageInfo{BornHopf}
     {Skipping~external~refs~for~#1~because~this~is~the~current~document}
   }
   {
    \IfFileExists{xref/#1.aux}
     { \externaldocument[][nocite]{xref/#1}[#1.pdf] }
     {
      \IfFileExists{#1.aux}
       { \externaldocument[][nocite]{#1} }
       {
        \PackageInfo{BornHopf}
         {External~aux~file~#1.aux~not~found;~no~labels~imported}
       }
     }
   }
 }
\ExplSyntaxOff

% External references are automatic.  This file is the sole owner of both the
% bibliography and the current Paper-I-A--VII-D AUX registry (30 roots including I-C).  Missing auxiliaries
% are informational on the first pass and are populated by the source-only
% two-pass series build.  Basenames deliberately omit temporary download/copy
% suffixes such as ``(1)'' or ``-1''.
  % Paper I-A: intrinsic S3 and exceptional group introduction.
  \BHExternalRefs{Paper1A-Geometry-v5_175}

  % Paper I-B: complete exceptional geometry and finite spectral structures.
  \BHExternalRefs{Paper1B-ExceptionalGeometry-v1_5}
  % Paper I-C: reciprocal product M8 geometry.
  \BHExternalRefs{Paper1D-TechnicalMaterial-v1_65}
  % Paper I-D: terminal technical proofs and calculations.
  \BHExternalRefs{Paper1C-M8Geometry-v1_137}

\ifdefined\BHMathOnly
  % Papers I-A and I-B import only the mathematical atlas labels.
\else
  % Paper II: five physical-mathematical main papers plus one technical proof volume.
  \BHExternalRefs{Paper2A-GeometricConstitution-v3_221}
  \BHExternalRefs{Paper2B-MomentumCarrier-v2_144}
  \BHExternalRefs{Paper2C-ReciprocalLocalization-v2_169}
  \BHExternalRefs{Paper2D-SymmetrySpectra-v2_141}
  % The former II-F export labels remain in II-E; current II-F is the technical volume.
  \BHExternalRefs{Paper2E-KinematicExports-v3_138}
  % Paper II-F: complete technical proof and calculation owner for II-A--II-E.
  \BHExternalRefs{Paper2F-TechnicalMaterial-v1_67}

  % Paper III: Dynamics and Intrinsic Geometric Constitution (six components).
  \BHExternalRefs{Paper3A-GeneratedDynamics-v1_176}
  \BHExternalRefs{Paper3B-SpectralAction-v1_187}
  \BHExternalRefs{Paper3C-ProtectedHierarchy-v1_187}
  \BHExternalRefs{Paper3D-RunningResponse-v1_165}
  \BHExternalRefs{Paper3E-GravityAction-v1_197}
  \BHExternalRefs{Paper3F-CompleteDynamics-v1_263}
  % Paper III-G: complete preserved technical proof owner.
  \BHExternalRefs{Paper3G-TechnicalMaterial-v1_53}

  % Paper IV-A: generated finite quantum mechanics.
  \BHExternalRefs{Paper4A-FiniteQuantumMechanics-v4_214}
  % Paper IV-B: complete finite quantum processes and physical response.
  \BHExternalRefs{Paper4B-QuantumResponse-v1_49}
  % Paper IV-C: extended proofs and finite technical calculations.
  \BHExternalRefs{Paper4C-TechnicalMaterial-v1_49}

  % Paper V-A: The Standard Model.
  \BHExternalRefs{Paper5A-FiniteMatterGeometry-v3_209}

  % Paper V-B: Finite Spectrum and Flavor Closure.
  \BHExternalRefs{Paper5B-TerminalSpectrum-v5_199}

  % Paper V-C: Finite Interaction Boundaries of the Standard Model.
  \BHExternalRefs{Paper5C-InteractionResponse-v2_177}
  % Paper V-D: preserved technical proof and calculation owner.
  \BHExternalRefs{Paper5D-TechnicalMaterial-v1_53}

  % Paper VI-A: Gravity and Closed Cosmology.
  \BHExternalRefs{Paper6A-GravityCosmology-v2_260}

  % Paper VI-B: Geometric Dark Sector and Cosmological Readouts.
  \BHExternalRefs{Paper6B-CosmologyReadouts-v1_206}

  % Paper VI-C: complete technical proof and comparison owner.
  \BHExternalRefs{Paper6C-TechnicalMaterial-v1_50}

  % Paper VII-A: Finite QCD Kernel and Nonperturbative Closure.
\BHExternalRefs{Paper7A-QCDKernel-v1_172}

  % Paper VII-B: Hadronic, Thermal, and Nuclear QCD Readouts.
  \BHExternalRefs{Paper7B-QCDReadouts-v1_157}

  % Paper VII-C: Continuous Quantum Field Theory.
  \BHExternalRefs{Paper7C-ContinuousQFT-v2_202}

  % Paper VII-D: complete technical material for VII-A through VII-C.
  \BHExternalRefs{Paper7D-TechnicalMaterial-v1_53}

\fi

% END INLINE SOURCE: Paper00-ActiveExternalReferences-v2_130.tex


\hypersetup{hidelinks}
\setlength{\emergencystretch}{3em}
\title{\textbf{Born--Hopf Geometry in a Closed Universe IV-B: Quantum Processes and Observable Response} \\[0.5em]\large Complete Finite Action, Source-Cyclic Poles, Charged Transport, and Physical Readout}
\author{Peter Kugelmann\thanks{\texttt{pkugelmann@proton.me}}}
\date{Version 1.55 --- 20 November 2026 \\ Framework v3.066}
% Cross-DAG current-series import for the new source-owned mathematical bridge.
\BHExternalRefs{Paper1B-ExceptionalGeometry-v1_12}
\BHExternalRefs{Paper2C-ReciprocalLocalization-v2_170}

\begin{document}
\maketitle
\begin{abstract}
Paper IV-B develops the finite action-generated quantum process and its physical readout from the existing Born--Hopf finite quantum algebra and dynamical sources. Exact clock evolution, positive-frequency sectors, finite field/Fock branches, charged source resolvents and source-cyclic pole certificates are retained with their state, action, preparation, detector and limiting hypotheses. A finite determinant root, algebraic channel or supplied Gaussian comparison is not automatically a measured particle, cosmological occupation or continuum scattering limit. Paper IV-C carries the complete auxiliary calculations and negative comparison tests. The v2.254 extension gives an exact finite detector-charge covariance identity with a conditional same-action Ward-residual bound and a metric-generated collar flux test; no physical soft-charge or asymptotic memory claim is inferred.
\end{abstract}
\setcounter{tocdepth}{2}
\tableofcontents
\clearpage
\noindent\textit{Original-source reference convention.} In unchanged inherited passages, Part~I refers to the finite quantum-core material now owned by Paper~IV-A. Former appendix references may resolve to the technical companion Paper~IV-C; their theorem labels and hypotheses remain unchanged.

\section{One-particle modes, finite Fock and charged field sectors}
\label{p4b:sec:v238-b1}
\noindent\textit{Source and claim boundary.} Finite mode assignments, statistics, fermionic CAR and bounded-boson calculations are conditional on the inherited one-particle, exchange and occupation data. Finite field presentations do not imply a continuum QFT limit.

\paragraph{Finite-QFT Bridge and Field-Theoretic Presentations}
\label{p4:part:finite-qft-bridge}

\paragraph{Orientation.}
This part packages finite mode quantization and finite-field presentation theorems as bridge results over the closed quantum core.  The statements in this part are exact finite-sector or stated finite-window/interface results; continuum-looking language enters only through the declared presentation or compatible-refinement domains.



\subsection{Finite mode quantization and the QFT bridge}
\label{p4:sec:part2-qft-bridge}

\paragraph{Status of the bridge.}
The material in this section licenses finite QFT language on stated finite-window or compatible-refinement theorem domains.  Continuum, scattering, thermal, and field-theoretic statements below are finite-interface statements unless a compatible large-window or Born--Hopf refinement hypothesis is explicitly stated.




\paragraph*{Finite-QFT bridge dictionary and closure map}
\label{p4:sec:qft-bridge-dictionary}

\begingroup
\small
\setlength{\LTleft}{0pt}
\setlength{\LTright}{0pt}
\setlength{\tabcolsep}{5pt}
\renewcommand{\arraystretch}{1.16}
\emergencystretch=2em
\hbadness=10000
\captionsetup{hypcap=false}

\captionof{table}[QFT bridge dictionary]{QFT bridge dictionary.}
\label{p4:tab:qft-bridge-dictionary}

\begin{longtable}{@{}
  >{\raggedright\arraybackslash}p{0.27\linewidth}
  >{\raggedright\arraybackslash}p{\dimexpr0.73\linewidth-2\tabcolsep\relax}
@{}}

\toprule
\textbf{Finite Born/QM structure}
& \textbf{Field representative, closure, and boundary} \\
\midrule
\endfirsthead

\toprule
\textbf{Finite Born/QM structure}
& \textbf{Field representative, closure, and boundary} \\
\midrule
\endhead

\midrule
\multicolumn{2}{r}{\emph{Continued on next page}}\\
\bottomrule
\endfoot

\bottomrule
\endlastfoot

Hopf \(U(1)\) fiber and projective classes from Paper~I-A
&
\textit{Representative:} comparison between the oriented geometric Hopf
line and central quantum ray phase.
\textit{Finite/interface theorem:}
Theorem~\ref{p4:thm:hopf-ray-phase} proves ordinary projective phase
redundancy and supplies the degree-one associated character
\(\chi_1:U(1)_{\rm Hopf}\to U(1)_{\rm ray}\).
\textit{Boundary:} this compares associated lines; it does not make the
generally noncentral Hopf action on the full carrier equal to scalar
multiplication. \\

Reeb and Hopf charge generator \(Q_\xi\)
&
\textit{Representative:} charge-sector decomposition and superselection.
\textit{Finite/interface theorem:} Corollary~\ref{p4:cor:hopf-reeb-superselection}, using
Theorem~\ref{p4:thm:charge-superselection}, shows that Reeb weights become
exact or effective finite superselection blocks when accessible observables
commute with \(Q_\xi\) and channels are \(U_\xi\)-covariant.  Channel
covariance implies that the charge-dephasing conditional expectation
intertwines the channel; the converse is false in general.  Intertwining
alone suffices for preservation of the dephased subspace.
\textit{Boundary:} the Reeb-weight ladder realizing a later physical sector is
left to the sector papers. \\

Distinguished finite one-particle block \(\mathcal K\)
&
\textit{Representative:} one-particle mode carrier.
\textit{Finite/interface theorem:} the finite Fock construction and
Theorem~\ref{p4:thm:finite-second-quantization} show that the field carrier is a
repackaging of finite tensor and sector structure, not a new quantization axiom.
\textit{Boundary:} a one-particle assignment is required; identifying a
geometric block with a physical particle sector belongs to later sector
papers. \\

Statistics-selected tensor powers of \(\mathcal K\)
&
\textit{Representative:} finite Bose and Fermi Fock sectors.
\textit{Finite/interface theorem:} Theorem~\ref{p4:thm:finite-second-quantization} gives
symmetric and antisymmetric occupation spaces together with finite creation and
annihilation maps inside finite Hilbert space.
\textit{Boundary:} the physical statistics selector is supplied or derived on
a later spin--statistics domain, and \(N_{\max}\) is separate bosonic
theorem-domain data.  Exact infinite bosonic CCR representations are continuum
limits, not native finite objects. \\

One-body operators \(A\in\End(\mathcal K)\) and mode polynomials
&
\textit{Representative:} second-quantized observables and finite field
Hamiltonians.
\textit{Finite/interface theorem:} Theorem~\ref{p4:thm:finite-second-quantization} gives
\(d\Gamma(A)\) and candidate self-adjoint quadratic and quartic generators.
\textit{Boundary:} actual Hamiltonian meaning requires the displayed
interaction/\(\mathsf{DYN}\) domain; fully renormalized interacting
continuum QFT is not claimed here. \\

Positive spatial operator \(A\ge0\) on the Paper~III compact slice, scalar
cover \(\ttilde\), and stable isolated block
&
\textit{Representative:} positive-frequency realization of the upstream
local free dispersion.
\textit{Finite/interface theorem:} Theorem~\ref{p4:thm:finite-one-particle-dispersion} and
Remark~\ref{p4:rem:lorentzian-noncircular} select the positive finite
Hamiltonian on a stable intrinsic block.
\textit{Boundary:} Paper~III already owns the GP1 cone, Lorentzian sign, and
local second-order mode law; Paper~IV does not feed an input back into that
derivation. \\

Lowest spatial eigenvalue \(\lambda_0\) on a stable block
&
\textit{Representative:} physical threshold and rest-energy interpretation of
spectral gaps.
\textit{Finite/interface theorem:} Corollary~\ref{p4:cor:spectral-gap-threshold} shows that
geometric spectral gaps become finite one-particle energy thresholds once the
dispersion theorem applies.
\textit{Boundary:} the geometric origin and particle identification of each gap
belong to Papers~I-A--III and later sector papers. \\

Finite positive Euclidean spectral kernels
&
\textit{Representative:} cyclic real-time reconstruction and finite
reflection positivity.
\textit{Finite/interface theorem:} Theorems~\ref{p4:thm:finite-euclidean-hamiltonian-reconstruction}
and~\ref{p4:thm:finite-os-positivity} construct a cyclic reconstruction
space, prove bosonic reflection positivity on the stated Wick domain, and
classify the finite graded-fermionic branch by positivity of its complete
reflection matrix.
\textit{Boundary:} identifying the cyclic reconstruction with a pre-existing
carrier requires the separate same-carrier semigroup factorization; a
positive degree-one fermionic corner is not sufficient, and full interacting
continuum Wightman reconstruction remains outside Paper~IV. \\

Nested finite mode spaces with compatible embeddings
&
\textit{Representative:} continuum and local field-algebra bridge.
\textit{Finite/interface theorem:} Theorem~\ref{p4:thm:inductive-qft-bridge} gives a
well-defined inductive field-algebra limit for the directed finite algebras.
\textit{Boundary:} the algebra embeddings, compatible states, localizers,
and dynamics must be supplied.  Paper~II's nested symbol spaces do not by
themselves construct those maps or an interacting relativistic continuum
theory. \\

Finite local supports and compatible finite generators
&
\textit{Representative:} algebraic local-net structure and cutoff
microcausality.
\textit{Finite/interface theorem:} Theorems~\ref{p4:thm:compatible-local-net}
and~\ref{p4:thm:approximate-microcausality} distinguish exact-support
local/graded-local pre-nets from bandlimited localizers and bound cutoff
commutators for bounded CAR/truncated observables or Weyl words by the
supplied spectral-leakage and dynamical-locality estimates.  The even CAR
subnet is ordinarily local.
\textit{Boundary:} spectral projection does not preserve exact support, and
full Haag--Kastler classification and interacting continuum microcausality
remain outside scope. \\

Finite Gaussian and quasi-free data with smeared test functions
&
\textit{Representative:} quasi-free field state and weak operator-valued
distribution language.
\textit{Finite/interface theorem:} Theorem~\ref{p4:thm:quasifree-reconstruction} and
Proposition~\ref{p4:prop:weak-ovd-smeared-fields} recover ordinary
field-correlation language on exact CAR or regular Weyl-CCR inductive
systems and their algebraic GNS cores.
\textit{Boundary:} occupation-truncated bosonic matrix algebras provide only
degree-controlled approximate moments under an independent
\(N_{\max}\)-tail limit; non-quasi-free constructive QFT remains beyond this
paper. \\

Compatible finite coarse-graining maps
&
\textit{Representative:} finite coarse-graining bookkeeping, threshold
matching, and running representatives.
\textit{Finite/interface theorem:}
Theorem~\ref{p4:thm:renormalization-semigroup} separates the exact
projective family of CPTP state maps from typed coefficient-transfer and
threshold maps.
\textit{Boundary:} coefficient cocycles require a functorial closed-basis
matching prescription; model-specific beta functions and precision fits are
supplied only in sector papers. \\

Time-dependent finite mode spaces \(\mathcal K_t\)
&
\textit{Representative:} curved-background and Bogoliubov particle-production
interface.
\textit{Finite/interface theorem:} Theorem~\ref{p4:thm:curved-background-finite-bridge} shows
that a supplied Hermitian bundle, connection, quadratic dynamics, and
positive-frequency splitting support Bogoliubov transport.
\textit{Boundary:} Paper~II does not supply the required many-body,
local-observable, state, and propagator ladder; full cosmological likelihood
propagation is not performed here. \\

Quadratic scalar and tensor perturbation blocks on a Paper~VI-A background
&
\textit{Representative:} inflationary perturbation and scalaron one-particle
mode language.
\textit{Finite/interface theorem:} Theorem~\ref{p4:thm:finite-cosmological-perturbation-bridge}
gives stable oscillator and Bogoliubov evolution after positivity of the
full physical quadratic form, or a separately supplied compatible
positive-energy complex structure and stable positive-frequency splitting,
has been supplied.
\textit{Boundary:} positive kinetic energy alone is insufficient, and the
truncated-boson branch needs separate occupation-tail control.  Continuum
correlators additionally require compatible embeddings, states, dynamics,
and a stated convergence topology; numerical CMB and large-scale-structure
propagation remains outside Paper~IV. \\

Time-dependent scalaron or reheating background coupled to finite matter blocks
&
\textit{Representative:} reheating and particle-production rate interface.
\textit{Finite/interface theorem:} Theorem~\ref{p4:thm:finite-reheating-interface} shows that
finite Dyson and Bogoliubov evolution yields transition probabilities, while
a fixed product environment assignment yields CPTP reduction.
\textit{Boundary:} moving carriers require a connection or coherent
trivialization, correlated initial states require a CP assignment or
restricted domain; weak-coupling rates need the displayed reservoir limit
and an invariant occupation algebra, while exact branching ratios require
sector couplings supplied outside Paper~IV. \\

Finite transfer maps from initial perturbation data to late observables
&
\textit{Representative:} Paper~VI-B transfer and readout objects:
\(\mathcal T_8(a)\), \(\mathcal T_g\), \(\kappa_g\), \(\Xi_a\), and
\(\Xi_d,\Xi_*,\Xi_g,\Xi_{\rm BBN}\).
\textit{Finite/interface theorem:} Theorem~\ref{p4:thm:finite-transfer-kernel-status} shows that
these objects are finite transfer-kernel or finite-determinant outputs, not new
constants of the geometry.
\textit{Boundary:} CPTP state transfer, CP trace-nonincreasing selected
branches, positive effects, generic response functionals, and determinant
kernels are distinct typed objects; their numerical evaluation requires the
cosmological propagation pipeline. \\

Finite weak-coupling thermal blocks and heavy neutral modes
&
\textit{Representative:} kinetic finite-temperature QFT interface for
baryogenesis.
\textit{Finite/interface theorem:} Theorem~\ref{p4:thm:minimal-finite-temperature-seesaw-interface}
gives a finite thermal-QFT support theorem for the later reduced baryogenesis
calculation.
\textit{Boundary:} full density-matrix thermal leptogenesis and all plasma
corrections remain downstream. \\

\end{longtable}
\TableNotes{The entries are read in the same spirit as the constitutive-mapping
table of Paper~II, but with finite-QFT bridge status.  Paper~II records how
intrinsic kinematic structures represent physical notions before dynamics;
this table records which one-particle, statistics, field, locality, state,
Hamiltonian, or limiting data must first be supplied and which finite theorem
then applies.  No constitutive datum is hidden by calling the bridge a
postulate-free consequence: each entry is a finite representation theorem, a
role-conditioned theorem, or an explicitly delimited interface/open-residue
statement.  The standard continuum terminology is used in the customary
Osterwalder--Schrader, Haag--Kastler/LSZ, Wick/quasi-free, and KMS senses.
Unless stated otherwise, units are chosen with \(\hbar=k_B=1\), with physical
units restored by replacing \(e^{-itH}\) with \(e^{-itH/\hbar}\) and reading
\(\beta=(k_BT)^{-1}\) in thermal factors.}
\endgroup

Consequently, every QFT notion used later in the series is assigned one of
three statuses here: supplied constitutive/interface data, an exact finite
representation theorem on those data, or a derived open residue whose
continuum, curved-background, scattering, reservoir, or thermodynamic limit
still requires the displayed hypotheses.  Finite dimensionality closes the
second class; it does not create the first or discharge the third.

\begin{proposition}[Finite-QFT dependency ledger]
\label{p4:prop:native-qft-status}
Within the finite Born quantum framework developed in this paper, the following objects are not additional quantum-kinematic postulates once the corresponding finite sector data are supplied:
\begin{enumerate}
\item finite candidate mode carriers, including Paper-I-A Hurwitz-link internal
supports only on an actualized or conditionally generated
\(\mathsf{Supp}_{2T}\) domain, and---after a
one-particle assignment, Bose/Fermi selector and, for bosons, a generated occupation boundary have been
supplied---finite Fock
sectors, creation and annihilation maps, and finite second quantization;
\item smeared fields and quasi-free correlators on exact CAR or regular
Weyl-CCR systems, degree-controlled truncated-boson moments, weak
field-distribution language on an inductive GNS core, and bosonic or
explicitly gauge-invariant graded finite Euclidean reflection positivity;
\item finite Dyson expansions, finite channel scattering amplitudes, finite
Gibbs/KMS identities, finite Bogoliubov transport, exact state
coarse-graining maps, and---after a functorial closed-basis prescription is
supplied---typed coefficient-transfer and threshold maps.
\end{enumerate}
Their finite calculus is a consequence of structures already established in Part~I and in
the bridge theorems of this section: Hilbert-space, spectral, tensor,
positivity, and completely positive map structures.  The complementary continuum layer consists of the infinite-volume interacting theory, exact point-field microcausality before a compatible limit, exact global bosonic CCR as a limiting idealization, and model-specific numerical running or transport coefficients.
\end{proposition}


\begin{proof}
Item~1 is Theorem~\ref{p4:thm:finite-second-quantization}, together with the
\(\mathsf{Supp}_{2T}\)-relative Hurwitz-link support import
Theorem~\ref{p4:thm:hurwitz-link-finite-support-import} when that core role is
supplied, and Proposition~\ref{p4:prop:no-exact-ccr}, which
identifies the finite bosonic boundary defect and prevents the inconsistent
assumption of exact finite-dimensional CCR.  Item~2 is supplied by
Theorems~\ref{p4:thm:finite-euclidean-hamiltonian-reconstruction},
\ref{p4:thm:finite-os-positivity}, \ref{p4:thm:inductive-qft-bridge},
\ref{p4:thm:compatible-local-net}, \ref{p4:thm:approximate-microcausality},
and \ref{p4:thm:quasifree-reconstruction}, together with
Proposition~\ref{p4:prop:weak-ovd-smeared-fields}.  Item~3 follows from
finite-dimensional self-adjoint evolution, Theorem~\ref{p4:thm:dyson-finite},
Proposition~\ref{p4:prop:framework-scattering-window},
Theorems~\ref{p4:thm:moller-bridge},
\ref{p4:thm:quasilocal-thermodynamic-bridge},
\ref{p4:thm:curved-background-finite-bridge},
\ref{p4:thm:renormalization-semigroup}, and the finite open-system closure of
Theorems~\ref{p4:thm:microscopic-bath-gkls}
and~\ref{p4:thm:minimal-finite-temperature-seesaw-interface}.  Each cited
result is either a finite algebraic theorem or an explicitly labelled
compatibility/limit theorem.  The excluded objects all require an additional
finite-support, infinite-volume, continuum, or model-specific construction
not supplied by the finite data alone.
\end{proof}


\begin{remark}[Finite-QFT theorem domain]
\label{p4:rem:qft-postulate-boundary}
The GP1--GP3-native mathematical data used in Paper~IV are the imported
finite carriers and their generated finite test algebras.  At \(m_p=1\) the
latter is the exact six-block algebra of
Theorem~\ref{p4:thm:full-finite-observable-algebra}; the full carrier
presentation is physically used as a graded field algebra only on
\(\mathsf{FENV}\), and as a full event algebra when carrier resolution is
generated or on an independent \(\mathsf{ENV}\) root.  The state--effect
evaluation already has its generated event/probability-role meaning;
\(\mathsf{PROB}\) selects only actual preparation/readout tokens.  Tensor form follows once \(\mathsf{SUB}\)
certifies two occurrences with no cross-relation; CP/Kraus form follows once
\(\mathsf{PROC}\) supplies an actual stable map; and inner/GKLS form applies
to the corresponding supplied \(\mathsf{DYN}\) domains.
Generic statistics and one-particle assignment remain the explicit
\(\mathsf{STAT}\) and \(\mathsf{1P}\) selector families; actual Hamiltonians
remain under \(\mathsf{DYN}\), and nongenerated proper accessibility
restrictions under \(\mathsf{ACC}\).  Fields and localizers remain supplied
sector/interface data.  Point fields, continuum path
integrals, global CCR, infinite-volume vacua, and continuum scattering
matrices enter only after the corresponding compatibility and limit
hypotheses have been stated.
\end{remark}




\paragraph{Support supplied to the rest of the series}
\label{p4:sec:qft-support-map}

The same classification also identifies exactly what Paper~IV contributes downstream.  It identifies the finite theorem-domain conditions under which later papers may use field-theoretic language without adding continuum QFT as an extra foundation.




\subsection{One-particle mode spaces and finite Fock sectors}
\label{p4:sec:part2-fock}

In the finite Born framework, first retain the candidate one-particle atlas
\[
\begin{aligned}
 \mathfrak A_{\rm 1P}:=\{\operatorname{Ran}P:\;&P\text{ is a generated
 spectral/charge/chiral/grading/symmetry projector}\\
 &\text{satisfying the stated stability requirements}\}.
\end{aligned}
\]
Theorem~\ref{p4:thm:stationary-positive-frequency-mode-space} later derives
a canonical positive-frequency mode Hilbert structure on stationary stable
linear-field branches, but that mathematical mode structure does not by itself
supply particle semantics.  In the generic finite atlas, therefore supply a
\emph{one-particle assignment} selecting one intrinsic finite block
\[
  \mathcal K=\operatorname{Ran}P
\]
as the \(n=1\) mode carrier.  The projector \(P\) may be generated by the
constituted finite spectral, charge, chirality, grading, or symmetry
projectors already present in the framework, but those data alone do not
give the block particle meaning.  Typical candidates are joint eigenspaces
or finite spectral windows of constituted commuting data, possibly with an internal factorization
\[
\mathcal K\cong \mathcal K_{\mathrm{ext}}\otimes \mathcal K_{\mathrm{int}}.
\]
Write \(d=\dim \mathcal K\).  Thus the word ``one-particle'' refers to the
\(n=1\) occupation sector assigned to an intrinsic finite mode block, not to
an assumed classical point particle.  Permutation covariance on tensor powers
of \(\mathcal K\) gives the complete Young-sector decomposition of
Theorem~\ref{p4:thm:permutation-sector-decomposition}.  The Bose and Fermi
constructions below begin only after the trivial- or sign-sector selector,
respectively, has been supplied.  The finite second-quantized bridge packages
the selected sectors as mode-occupation spaces.

\begin{proposition}[Bare-algebra naturality does not select a one-particle block]
\label{p4:prop:bare-algebra-one-particle-no-selector}
If \(P\in M_d(\CC)\) is invariant under every inner automorphism, then
\(P=0\) or \(P=I\).  For a finite direct-sum algebra, an
automorphism-natural projection is therefore at most a central sum of blocks
unless additional generated geometry distinguishes a subspace.  In
particular, the bare finite algebra does not select a nontrivial generic
\(\mathsf{1P}\) carrier.
\end{proposition}

\begin{proof}
The equality \(uPu^*=P\) for every unitary \(u\) places \(P\) in the center
of \(M_d(\CC)\).  A scalar projection is \(0\) or \(I\).  Applying this
blockwise proves the direct-sum statement.  A generated isolated spectral
projection can evade the obstruction because the generating operator and
the intrinsic spectral characterization are additional constituted data;
particle semantics still requires the corresponding excitation-role theorem.
\end{proof}

\begin{theorem}[Stationary positive-frequency mode structure and zero-mode split]
\label{p4:thm:stationary-positive-frequency-mode-space}
Let a constituted stationary stable real linear scalar-field role have a finite
mode block \(\mathcal K\) with positive self-adjoint frequency operator
\(H\ge0\).  Set
\[
  P_{\rm osc}:=\mathbf1_{(0,\infty)}(H),
  \qquad
  \mathcal K_{\rm osc}:=P_{\rm osc}\mathcal K,
  \qquad
  \mathcal K=\ker H\oplus\mathcal K_{\rm osc}.
\]
On the real Cauchy-data space
\(\Gamma_{\rm osc}=\mathcal K_{\rm osc,\mathbb R}\oplus
\mathcal K_{\rm osc,\mathbb R}\), write \(\Omega=H/\hbar\) and
\[
  \mathcal G(q,p)=(p,-\Omega^2q),
  \qquad
  J:=-\mathcal G(-\mathcal G^2)^{-1/2}.
\]
Then
\[
  J(q,p)=(-\Omega^{-1}p,\Omega q),
  \qquad J^2=-I,
  \qquad [J,\mathcal G]=0.
\]
For the canonical Cauchy symplectic form
\[
 \sigma((q,p),(q',p'))=\langle q,p'\rangle-\langle p,q'\rangle,
\]
one has \(\sigma(v,Jv)>0\) for \(v\ne0\).  Hence the oscillatory
positive-frequency mode Hilbert structure is canonically generated from the
stationary quadratic field law.  A zero-frequency kernel, when present, is a
separate nonoscillatory role.  This theorem constructs a mode Hilbert space; it
does not by itself assign one-particle semantics in the sense of
\(\mathsf{1P}\).
\end{theorem}

\begin{proof}
On \(\mathcal K_{\rm osc}\), \(\Omega>0\), so the displayed functional
calculus is defined.  Direct multiplication gives \(J^2=-I\), and because
\(J\) is a function of \(\mathcal G\), it commutes with the stationary
evolution.  Finally
\[
 \sigma((q,p),J(q,p))
 =\langle q,\Omega q\rangle+\langle p,\Omega^{-1}p\rangle>0.
\]
The spectral projector \(P_{\rm osc}\) and the kernel split are intrinsic
functional-calculus outputs.
\end{proof}


\begin{theorem}[Hessian-generated Kähler and projective mode geometry]
\label{p4:thm:v245-hessian-kahler-projective-geometry}
On the oscillatory domain of
Theorem~\ref{p4:thm:stationary-positive-frequency-mode-space}, define
\[
 g(v,w):=\sigma(v,Jw).
\]
For \(v=(q,p)\) and \(w=(q',p')\),
\[
 \boxed{
 g(v,w)=\langle q,\Omega q'\rangle
       +\langle p,\Omega^{-1}p'\rangle.}
\]
Then \(g\) is symmetric and positive definite, and
\[
 g(Jv,Jw)=g(v,w),\qquad \sigma(Jv,Jw)=\sigma(v,w).
\]
Hence \((\Gamma_{\rm osc},\sigma,J,g)\) is a compatible linear Kähler
space.  With the positive-frequency coordinate already generated in the
stationary quadratic branch,
\[
 z(v)=\frac1{\sqrt{2\hbar}}
 \left(\Omega^{1/2}q+i\Omega^{-1/2}p\right),
\]
one has
\[
 \boxed{
 \langle z(v),z(w)\rangle
 =\frac1{2\hbar}\bigl(g(v,w)+i\sigma(v,w)\bigr).}
\]
Thus the mode-space Hermitian metric is generated by the same physical
Hessian and canonical symplectic form that generate propagation.

On the normalized projective quotient, use the induced convention
\[
 \omega_{\rm FS}(\delta z,\delta w)
 =2\hbar\,\Im\langle\delta z,\delta w\rangle.
\]
For self-adjoint \(A\) define \(f_A([z])=\langle z,Az\rangle\).  Its
Hamiltonian vector field is
\[
 X_A=-\frac{i}{\hbar}(A-f_AI)z,
\]
and therefore
\[
 \boxed{\{f_A,f_B\}_{\rm FS}
 =\frac1{i\hbar}\langle z,[A,B]z\rangle.}
\]
The Schrödinger kinetic one-form on normalized vectors is
\[
 \boxed{\Theta_Q=i\hbar\langle z,dz\rangle,}
\]
namely \(\hbar\) times the Berry connection in this sign convention.
This theorem concerns the generated positive-frequency field-mode geometry;
it does not replace the independently constituted algebraic state space or
assign one-particle semantics by itself.
\end{theorem}

\begin{proof}
Substituting \(J(q',p')=(-\Omega^{-1}p',\Omega q')\) into the canonical
symplectic form gives the displayed expression for \(g\), and positivity
follows from
\(g(v,v)=\|\Omega^{1/2}q\|^2+\|\Omega^{-1/2}p\|^2\).
The identities with \(J\) follow by direct substitution.  Expanding the
Hermitian product of the two \(z\)-coordinates gives the real part \(g/(2\hbar)\)
and imaginary part \(\sigma/(2\hbar)\).  On the normalized horizontal tangent
space, differentiating \(f_A\) and solving
\(df_A(\cdot)=\omega_{\rm FS}(X_A,\cdot)\) gives the displayed \(X_A\);
substitution of \(X_A,X_B\) yields the commutator bracket.  Finally
\(\langle z,dz\rangle\) is purely imaginary on \(\langle z,z\rangle=1\), so
the Schrödinger kinetic term is the stated projective connection one-form.
\end{proof}
% BEGIN C01 V2174 MODE MOMENT

% BEGIN INLINE SOURCE: Paper4-OccupationMoment-Main-v1_0.tex
% v2.174: a consequence of the existing stationary positive-Hessian domain.


\subsection{Creation and annihilation maps; exact CAR and truncated CCR}
\label{p4:sec:part2-car-ccr}

On the fermionic Fock space define the creation and annihilation maps by
\begin{equation}
\label{p4:eq:fermion-ladder}
c_i^\dagger\ket{\dots,n_i,\dots}
=
(-1)^{\sum_{k<i} n_k}(1-n_i)\ket{\dots,n_i+1,\dots},
\qquad
c_i=(c_i^\dagger)^*.
\end{equation}
They satisfy the exact canonical anticommutation relations
\begin{equation}
\label{p4:eq:car}
\{c_i,c_j\}=0,
\qquad
\{c_i,c_j^\dagger\}=\delta_{ij}\,\mathbf 1.
\end{equation}
So finite fermionic mode quantization is exact inside the supplied
one-particle, fermionic-sector carrier.

For bosonic sectors with occupation cutoff $N_{\max}$, define creation and annihilation maps by the usual occupation-number formulas whenever the resulting state stays inside the cutoff space, and set them to zero on the top-occupation boundary.
Away from that boundary they obey the standard bosonic ladder rules, while at the boundary their commutator acquires a finite defect supported on the highest-occupation subspace.
This is not an accident but a structural necessity.


\begin{proposition}[Finite carriers have no exact CCR]
\label{p4:prop:no-exact-ccr}
\leavevmode\par\noindent
No finite-dimensional Hilbert space admits operators \(a,a^\dagger\) satisfying
\begin{equation}
\label{p4:eq:no-finite-ccr}
[a,a^\dagger]=\mathbf 1
\end{equation}
exactly.  Consequently, bosonic mode operators in the finite Born framework
must be realized with an occupation cutoff or as an effective local limit.
\end{proposition}


\begin{proof}
If $[a,a^\dagger]=\mathbf 1$ on a finite-dimensional space, taking traces gives
\[
\Tr([a,a^\dagger])=0=\Tr(\mathbf 1)=\dim \mathcal H,
\]
a contradiction.
Hence exact bosonic CCR require either an infinite-dimensional carrier or an effective limit in which the finite boundary defect becomes negligible.
\end{proof}


\begin{theorem}[Exact free-vacuum covariance on a finite occupation carrier]
\label{p4:thm:v245-finite-free-vacuum-covariance}
Let \(\mathcal K_{\rm osc}\) be a finite positive-frequency mode block with
\(\Omega>0\), and let a finite symmetric occupation carrier contain the
vacuum and the complete one-particle sector.  Define truncated bosonic ladder
maps by the usual formulas below the top occupation boundary, with
\(a(f)|0\rangle=0\).  Then the vacuum-to-one-particle relation is exact,
\[
 a(f)a^\dagger(g)|0\rangle=\langle f,g\rangle|0\rangle,
\]
independently of the top-boundary CCR defect.  For real mode profiles define
\[
 Q(f)=\sqrt{\frac{\hbar}{2}}
 \left[a(\Omega^{-1/2}f)+a^\dagger(\Omega^{-1/2}f)\right],
\]
\[
 P(f)=-i\sqrt{\frac{\hbar}{2}}
 \left[a(\Omega^{1/2}f)-a^\dagger(\Omega^{1/2}f)\right].
\]
Then
\[
 \boxed{\frac12\langle0|\{Q(f),Q(g)\}|0\rangle
 =\frac{\hbar}{2}\langle f,\Omega^{-1}g\rangle,}
\]
\[
 \boxed{\frac12\langle0|\{P(f),P(g)\}|0\rangle
 =\frac{\hbar}{2}\langle f,\Omega g\rangle,}
\qquad
 \boxed{\frac12\langle0|\{Q(f),P(g)\}|0\rangle=0.}
\]
Equivalently,
\[
 \boxed{C_{QQ}=\frac{\hbar}{2\Omega}
 =\frac{\hbar^2}{2H}},\qquad H=\hbar\Omega.
\]
For the corresponding free Euclidean evolution,
\[
 \boxed{C_E(\tau)=\frac{\hbar}{2\Omega}e^{-|\tau|\Omega}
 =\hbar G_E(\tau).}
\]
Thus exact global CCR are unnecessary for the free-vacuum two-point
normalization on a finite occupation carrier.  The theorem computes the
vacuum covariance if that free ground state is the state under discussion; it
does not actualize that state physically.
\end{theorem}

\begin{proof}
Only vacuum and one-particle matrix elements enter the two-point functions, so
the top-occupation boundary is never reached.  Expanding the anticommutators
and using \(a(f)|0\rangle=0\) and
\(a(f)a^\dagger(g)|0\rangle=\langle f,g\rangle|0\rangle\) gives the three
displayed covariances.  Free Euclidean evolution multiplies a one-particle
mode by \(e^{-|\tau|\Omega}\), giving the final kernel and the equality with
the previously derived Euclidean Green operator.
\end{proof}
% BEGIN C01 V2174 GRADED BOUNDARY AND ORDERING

% BEGIN INLINE SOURCE: Paper4-GradedOccupationBoundary-Main-v1_0.tex
% v2.174: recover a boundary from the already supplied graded field port.


\subsection{Finite paired fermions and their clock action}

\begin{proposition}[Paired Berezin determinant line]
\label{p4:prop:paired-berezin-determinant-line}
Let \(V\) be a generated finite Hermitian space with the licensed
field and independent dual Grassmann roles.  The Riesz map
\(R_h:V\to V^*\), \(R_h(v)(w)=h(v,w)\), is antilinear, and
dual evaluation is canonical.  The top exterior pairing
\[
 \det(V^*)\otimes\det(V)\longrightarrow\CC
\]
is basis-independent.  Fix the paired Berezin ordering so that
the Gaussian with \(D=I\) has integral one.  For
\(D\in\operatorname{End}V\),
\[
 \int[d\bar\psi\,d\psi]_{\mathrm{paired}}\,
       e^{-\bar\psi D\psi}=\det D.
\]
\end{proposition}
\begin{proof}
Riesz duality follows from positivity and nondegeneracy of the
finite Hermitian form.  Under a change of basis, the determinants
of the dual and vector transformations cancel, proving naturality
of the top exterior pairing.  In the Gaussian expansion only the
degree-\(n\) term in \(\bar\psi D\psi\), \(n=\dim V\), saturates
all \(2n\) generators.  Anticommutation gives the signed
permutation sum for \(\det D\).  The stated ordering fixes its
constant sign at \(D=I\).  No finite spectral real structure is
needed for this paired construction.
\end{proof}


\begin{corollary}[Fixed-carrier gauge Jacobian cancellation]
\label{p4:cor:v489-fixed-carrier-gauge-jacobian-cancellation}
On the fixed finite carrier of Proposition~\ref{p4:prop:paired-berezin-determinant-line},
let a unitary gauge/basis transformation act by
\[
 \psi\mapsto U\psi,
 \qquad
 \bar\psi\mapsto \bar\psi U^{-1}.
\]
The vector and independent-dual Grassmann Jacobians are inverse.  Hence
\[
 \boxed{[d\bar\psi\,d\psi]_{\rm paired}
 \ \text{is a gauge/basis scalar on the fixed finite carrier}.}
\]
No additional determinant-line phase selector is required to define this
paired vectorlike finite measure.
\end{corollary}
\begin{proof}
For Grassmann variables the change \(\psi'=U\psi\) contributes the inverse
Jacobian \(\det U^{-1}\), while the independent dual change
\(\bar\psi'=\bar\psi U^{-1}\) contributes \(\det U\).  Their product is one.
Equivalently, this is the basis cancellation already built into the canonical
pairing \(\det(V^*)\otimes\det(V)\to\CC\).
\end{proof}

\begin{remark}[Fixed-carrier measure versus chiral reconstruction]
\label{p4:rem:v489-fixed-versus-chiral-line}
A nontrivial determinant-line bundle is a different problem: it appears when a
later continuum/chiral presentation replaces the fixed carrier by a
configuration-dependent subspace such as \(A\mapsto P_-(A)V\).  Curvature,
holonomy and global integrability of that varying line remain legitimate
chiral-regulator/continuum certificates, but they are not prerequisites for
the finite paired vectorlike measure above.
\end{remark}

\begin{theorem}[Global parallel-line and measure-ray trivialization criterion]
\label{p4:thm:parallel-line-trivialization}
Let $L\to\mathcal A$ be a Hermitian complex line bundle over a path-connected
admitted configuration component, equipped with a unitary connection $\nabla$.
Assume
\[
 F_\nabla=0
\]
and that the holonomy of $\nabla$ is trivial around every piecewise smooth
closed loop in $\mathcal A$.  Then $L$ admits a nowhere-zero global parallel
section.  Any two unit-norm global parallel sections differ by one constant
phase.  Consequently the associated positive measure ray, when this line is
the complete generated measure line, has no further local phase/trivialization
selector.

Conversely, a nowhere-zero global parallel section forces trivial holonomy.
Thus nontrivial holonomy is an obstruction to global parallel
trivialization rather than a missing selector.
\end{theorem}

\begin{proof}
Choose a base point $a_0$ and a nonzero vector $v_0\in L_{a_0}$.  Define
$s(a)$ by parallel transport of $v_0$ along any path from $a_0$ to $a$.
Flatness makes transport homotopy invariant locally, while trivial holonomy
makes the value independent of the chosen path globally.  Hence $s$ is a
nowhere-zero parallel section.  The ratio of two parallel unit sections is a
parallel $U(1)$-valued function and is therefore constant on the connected
base.  Conversely, parallel transport of a global parallel section around a
loop returns the same vector, so the holonomy is trivial.
\end{proof}

\begin{remark}[Chiral/global-measure firewall]
\label{p4:rem:parallel-line-chiral-firewall}
Theorem~\ref{p4:thm:parallel-line-trivialization} is a criterion, not a proof
that the physical chiral determinant line of a later paper has a flat
connection or trivial holonomy on its complete admissible gauge-field family.
Those curvature, locality, gauge-compatibility and global-integrability
hypotheses must be established on that actual owner.
\end{remark}

\begin{corollary}[Doubled Pfaffian presentation]
\label{p4:cor:doubled-pfaffian-presentation}
\leavevmode\par\noindent
In the block ordering \((e_1,\ldots,e_n,f_1,\ldots,f_n)\),
\[
 A_D=\begin{pmatrix}0&D\\-D^T&0\end{pmatrix},
 \qquad
 \operatorname{Pf}A_D=(-1)^{n(n-1)/2}\det D.
\]
\end{corollary}
\begin{proof}
At \(D=I\), interleaving each \(e_j,f_j\) pair uses
\(n(n-1)/2\) transpositions.  For general \(D\), the congruence
by \(\operatorname{diag}(D,I)\) multiplies the Pfaffian by
\(\det D\).  The polynomial identity extends through singular
\(D\).  This is a doubled determinant presentation; it does not
select the phase of a chiral or unpaired determinant.
\end{proof}

\begin{theorem}[Finite fermion clock real form and charge law]
\label{p4:thm:finite-fermion-clock-real-form}
Let the same finite analytic field/dual role carry a Hermitian
kinetic operator \(K\), with interior support or compatible
transported endpoint data.  Then
\[
 I_{E,F}=\int d\tau\,\bar\psi(\partial_\tau+K)\psi
 \ \longmapsto\
 -iS_{L,F},\qquad
 S_{L,F}=\int dt\,\bar\psi(i\partial_t-K)\psi
\]
under \(\tau=it\).  On its finite Fermi Fock space,
\[
 [d\Gamma(K),d\Gamma(Q)]=d\Gamma([K,Q]).
\]
A time-dependent one-particle charge has Heisenberg derivative
\(d\Gamma(\partial_tQ+i[K,Q])\).
\end{theorem}
\begin{proof}
The derivative transforms as \(\partial_\tau=-i\partial_t\)
and \(d\tau=i\,dt\); multiplication gives the displayed pair.
The endpoint condition permits the same integrations by parts
on both slices.  For the number-preserving bilinears
\(d\Gamma(K)=\sum a_i^\dagger K_{ij}a_j\), the canonical
anticommutation relations cancel the quartic terms in the
commutator, leaving \(\sum a_i^\dagger[K,Q]_{ij}a_j\).
The Heisenberg equation supplies the last formula.  A conserved
Noether charge additionally uses the symmetry and endpoint
conditions of the actual claim-complete action owner.
\end{proof}

These results use the finite one-particle, field/dual and
Fermi-statistics branch already constituted here.  They do not
construct a Majorana mass, global determinant phase, chiral
regulator, interacting measure or full quantum PCT completion.
% END M8 v1.78 FER01


% BEGIN M8 v1.78 FER02
% Results: B09-LINE-REALITY, B09-FAITHFUL-OPERATOR, B09-REAL-GRADING-IDENTITIES, B09-UNITARY-NATURALITY, B09-FAMILY-TYPING


\section{Clock polarity, positive-frequency and physical transfer}
\label{p4b:sec:v238-b2}
\noindent\textit{Source and claim boundary.} Finite energy orientation, Lorentzian free-mode dispersion and stationary source response are imported only on their proved upstream action/clock domains. In particular a formal clock solution does not select an actual history.

\subsection{Stationary energy and complete scalar-source response}
\label{p4:sec:stationary-source-energy}
\begin{theorem}[Finite energy/source variational certificate]
\label{p4:thm:finite-energy-source-certificate}
On an already licensed finite physical Hilbert space, use a fixed canonical
kinetic frame and the complete stationary Hermitian family
$H(h)=H-hO$. Let $N>0$, $E_0(h)=\lambda_{\min}(H(h))$, and
$p(h)=[E_0(0)-E_0(h)]/N$. For a physical density matrix $\rho$, put
\[
 q=\operatorname{Tr}(\rho O)/N,\qquad
 \delta=\operatorname{Tr}\{\rho[H-E_0(0)]\}/N\ge0.
\]
The pressure is convex, $p(0)=0$, and for every $h>0$,
\[
 \boxed{p(h)\ge hq-\delta,\qquad
 p'_+(h)\ge q-\delta/h.}
\]
At a differentiable sourced ground branch,
$p'(h)=\langle O\rangle_h/N$. At crossings use the convex one-sided
derivative; no unique state is selected by this notation.
\end{theorem}
\begin{proof}
The variational principle gives
$E_0(h)\le\operatorname{Tr}\rho(H-hO)$.
The ground energy is an infimum of affine functions, hence concave.
Rearrangement gives the first bound; convexity gives
$p'_+(h)\ge[p(h)-p(0)]/h$, proving the second. Feynman--Hellmann
gives the differentiable-branch formula.
\end{proof}

\begin{corollary}[Exact finite constrained-energy duality]
\label{p4:cor:finite-constrained-energy-duality}
Minimize over all density matrices of the supplied physical carrier:
\[
 I(q)=\min_{\substack{\rho\succeq0,\ \operatorname{Tr}\rho=1\\
                 \operatorname{Tr}(\rho O)=Nq}}
       \operatorname{Tr}\{\rho[H-E_0(0)]\}/N,
\]
with $I=+\infty$ outside the feasible interval. Then $I$ is convex and
\[
 \boxed{p(h)=\sup_q\{hq-I(q)\}.}
\]
A restricted trial class gives a variational bound, not this exact equality.
\end{corollary}
\begin{proof}
Each feasible slice of the finite density-matrix set is compact, so the
minimum is attained. Mixing minimizers proves convexity. Taking the
infimum of $\operatorname{Tr}\rho(H-hO)$ over the whole physical
density-matrix set and grouping by $q$ gives the equality.
\end{proof}

This is a stationary specialization of
Theorem~\ref{p4:thm:v4104-finite-interacting-source-response}, after
the complete kinetic and interaction transport of
Theorem~\ref{p3f:thm:v1148-finite-metric-constructor}.
It supplies no initial state, physical occurrence or limiting sequence.
Keep the source-dependent vacuum-energy factor if a transfer representative
subtracts $E_0(h)$ separately at each source. Discarding that factor changes
the physical derivative. A source-independent reference subtraction is
harmless. Nonlinear couplings, moving frames and source-dependent reductions
consume all existing contacts and jets from the complete source theorem.
Paper~VII-A applies this finite certificate to its owned scalar dictionary
and states the additional ordered-limit domain explicitly.

% END INLINE SOURCE: Paper4-StationarySourceEnergy-Main-v1_0.tex



% BEGIN INLINE SOURCE: Paper4-FunctionalMetricWard-Main-v1_0.tex
% Framework v2.162: interacting functional Hilbert and canonical Ward transfer.
\paragraph{Functional Hilbert response with metric derivatives}
\label{p4:sec:v4106-functional-hilbert-response}
Use the energy-complete licensed finite source of
Theorem~\ref{p4:thm:v4104-finite-interacting-source-response}, with all
consumed constraints, projectors, ordering and central energies retained.
The source histories and complete finite generator have the differentiability
needed below on the compact interval $I$. This extends the existing
retarded inverse; it introduces no new finite response owner.

\begin{theorem}[Exact functional Hilbert insertion and interacting response]
\label{p4:thm:v4106-functional-hilbert-response}
Let $H=H(t,x,x^{(1)},\ldots,x^{(s)})$ be the full self-adjoint
many-body finite generator and $J_{ar}=\partial H/\partial x^{a(r)}$.
Along the supplied history define
\[
 \mathfrak d_tA=\dot A+\frac{i}{\hbar}[H,A],\qquad
 \mathscr LX=\dot X+\frac{i}{\hbar}[H,X].
\]
The on-equation action/Hilbert insertion and its boundary operator are
\[
 \mathcal O_a=\sum_{r=0}^s(-\mathfrak d_t)^rJ_{ar},\qquad
 K[v]=\sum_a\sum_{r=1}^s\sum_{\ell=0}^{r-1}
            (-\mathfrak d_t)^\ell J_{ar}v^{a(r-1-\ell)}.
\]
They obey $\delta H[v]=\mathcal O_av^a+\mathfrak d_tK[v]$.
Thus the source part of the state-action variation is
$-\int\operatorname{Tr}(\rho\mathcal O_a)v^a\,dt
 -[\operatorname{Tr}(\rho K[v])]_{t_i}^{t_f}$.
For $\sigma=\delta\rho$, put
\[
 \tau=\sigma+\frac{i}{\hbar}[K[v],\bar\rho],\qquad
 \mathcal K_a[v]=\delta\mathcal O_a[v]
                       +\frac{i}{\hbar}[K[v],\mathcal O_a].
\]
The exact full density tangent and force response are
\[
 \mathscr L\tau=-\frac{i}{\hbar}[\mathcal O_b,\bar\rho]v^b+q,
 \qquad
 \tau_i=\sigma_i+\frac{i}{\hbar}[K[v](t_i),\rho_i],
\]
\begin{align*}
 \delta F_a(t)&=\operatorname{Tr}(\bar\rho\mathcal K_a[v](t))
       +\int_{t_i}^t\Pi^{R,\mathcal O}_{ab}(t,u)v^b(u)\,du
       +\operatorname{Tr}(\mathcal O_a(t)\tau_{\rm ind}(t)),\\
 \Pi^{R,\mathcal O}_{ab}(t,u)&=-\frac{i}{\hbar}\Theta(t-u)
  \operatorname{Tr}\!\left(\bar\rho(u)
  [U(t,u)^\dagger\mathcal O_a(t)U(t,u),\mathcal O_b(u)]\right),
\end{align*}
where $F_a=\operatorname{Tr}(\rho\mathcal O_a)$, $U$ is the full
many-body propagator, $q$ is an independently owned tangent, and
\[
 \tau_{\rm ind}=U(t,t_i)\tau_iU(t,t_i)^\dagger+\mathscr L_R^{-1}q,\qquad
 (\mathscr L_R^{-1}f)(t)=\int_{t_i}^tU(t,u)f(u)U(t,u)^\dagger\,du.
\]
No many-body correlation hierarchy is replaced by one-body closure.
\end{theorem}
\begin{proof}
The operator derivative satisfies the scalar product rule. Repeated
integration by parts gives
$J_{ar}v^{a(r)}=(-\mathfrak d_t)^rJ_{ar}v^a+
\mathfrak d_t\sum_{\ell=0}^{r-1}(-\mathfrak d_t)^\ell
J_{ar}v^{a(r-1-\ell)}$. Sum over $a,r$ and use
$d\operatorname{Tr}(\rho A)/dt=\operatorname{Tr}(\rho\mathfrak d_tA)$.
Differentiate the density equation. The Jacobi identity and
$\mathscr L\bar\rho=0$ imply
$\mathscr L(i[K,\bar\rho]/\hbar)=i[\mathfrak d_tK,\bar\rho]/\hbar$,
which proves the equation for $\tau$.
Solve it with the existing finite retarded inverse. Substitution of
$\sigma=\tau-i[K,\bar\rho]/\hbar$ into $\delta F_a$ and trace cyclicity
give $\mathcal K$ and the displayed kernel, including its sign.
All insertions are Hermitian. Interactions and central terms remain in
$H,U,\mathcal O,\mathcal K$ throughout.
\end{proof}
The inherited finite-time bounds are
\[
 \|\mathscr L_R^{-1}f\|_{C(I;1)}\le |I|\|f\|_{C(I;1)},\qquad
 |\Pi^{R,\mathcal O}_{ab}|\le
               \frac2\hbar\|\mathcal O_a\|\|\mathcal O_b\|.
\]
These bounds consume neither a spectral gap nor a stationary commutator
inverse. They do not establish infinite-time relaxation or uniform
continuum bounds. Interior variations have vanishing boundary jets;
otherwise their actual endpoint term and shifted initial value are retained.
For $s=0$, this theorem is exactly the existing instantaneous-source theorem.

\begin{corollary}[Moving-carrier curvature contact]
\label{p4:cor:v4106-moving-carrier-contact}
For the actual affine source-frame word
$H=\mathcal E-i\hbar\Gamma_0-i\hbar\Gamma_a\dot x^a$,
with $\Gamma_\mu^\dagger=-\Gamma_\mu$, define
\begin{align*}
 \mathcal D_aA&=\partial_aA+[\Gamma_a,A],\\
 \mathcal F_{ab}&=\partial_a\Gamma_b-\partial_b\Gamma_a+
                                      [\Gamma_a,\Gamma_b],\\
 \mathcal F_{a0}&=\partial_a\Gamma_0-\partial_t\Gamma_a+
                                      [\Gamma_a,\Gamma_0].
\end{align*}
Then
\[
 \mathcal O_a=\mathcal D_a\mathcal E-i\hbar\mathcal F_{a0}
                          -i\hbar\mathcal F_{ab}\dot x^b,\quad
 \tau=\sigma+[\Gamma_av^a,\bar\rho],\quad
 \mathcal K_a[v]=(\mathcal D_b\mathcal O_a)v^b
                          -i\hbar\mathcal F_{ab}\dot v^b.
\]
The local velocity contact is therefore $-i\hbar\mathcal F_{ab}$.
For an orthonormal embedding $J$, $P=JJ^\dagger$,
its one-body curvature and CAR lift are
\[
 \gamma_a=J^\dagger J_a,\qquad
 f_{ab}=J^\dagger[P_a,P_b]J,\qquad
 \Gamma_a=d\Gamma(\gamma_a),\quad\mathcal F_{ab}=d\Gamma(f_{ab}).
\]
\end{corollary}
\begin{proof}
Insert $J_{a1}=-i\hbar\Gamma_a$ in the theorem and expand its commutators
and total time derivatives. This gives the first three formulas.
Differentiating $J^\dagger J=I$ and $P=JJ^\dagger$ gives the curvature
formula; the exact CAR Lie homomorphism lifts it.
\end{proof}
A parallel frame along one history can still have nonzero source curvature.
A pure full-unitary change of fixed-rank coordinates is flat, while a moving
subspace can have the displayed curvature. Any projector jets must come
from the existing licensed carrier or isolated-cluster word; these formulas
do not license an arbitrary deformed Dirac-square cutoff or an added
counterdiabatic Hamiltonian. The inherited Riesz/Kato calculus computes
jets of a licensed constant-rank cluster, rather than selecting that cluster.

\paragraph{Canonical quantum transfer of the stress Ward identity}
\label{p4:sec:v4106-canonical-stress-ward}
\begin{theorem}[Finite canonical Ward defect and conditional zero-defect transfer]
\label{p4:thm:v4106-canonical-stress-ward}
Let $R_\xi$ be an interior-supported constituted symmetry tangent of the
complete marked owner. Retain its field, clock, constraint, occupation and
carrier transport. A transformation of an independently fixed physical mark
is outside this symmetry domain unless its variation is included.
For a licensed canonical lift between the finite fibers, with
$V_{\epsilon\xi}=I-i\epsilon G_\xi/\hbar+O(\epsilon^2)$, define
\[
 \mathfrak D_\xi=\delta H[R_\xi]-\dot G_\xi+
                     \frac{i}{\hbar}[G_\xi,H]
               =\delta H[R_\xi]-\mathfrak d_tG_\xi.
\]
If the actual full family satisfies
\[
 H[x^\xi]=V_\xi H[x]V_\xi^\dagger+i\hbar\dot V_\xi V_\xi^\dagger,
\]
then $\mathfrak D_\xi=0$. In general the exact finite identity is
\[
 \int\operatorname{Tr}(\rho\mathcal O_a)R_\xi^a\,dt
 =\int\operatorname{Tr}(\rho\mathfrak D_\xi)\,dt+
 [\operatorname{Tr}(\rho(G_\xi-K[R_\xi]))]_{t_i}^{t_f}.
\]
For the inherited convention
$\delta S=\tfrac12\int d\mu_{\rm act}\,T^{\mu\nu}\delta g_{\mu\nu}+\cdots$,
compact support gives
\[
 \int d\mu_{\rm act}\,\xi^\nu
       (\nabla_\mu\langle T^\mu{}_\nu\rangle-\mathcal F_\nu^{\rm ext})
       =\int\operatorname{Tr}(\rho\mathfrak D_\xi)\,dt.
\]
Here the action measure includes its actual clock-unit conversion and
$\mathcal F^{\rm ext}$ includes every retained external force or mark.
Zero defect yields conserved total stress on the complete coupled equations
when no force is omitted. Linear response additionally consumes its first source variation:
\begin{align*}
 \delta\mathfrak D_\xi[v]={}&\delta^2H[v,R_\xi]
       +\delta H[\delta R_\xi[v]]-\frac d{dt}\delta G_\xi[v]\\
 &+\frac{i}{\hbar}\bigl([\delta G_\xi[v],H]+[G_\xi,\delta H[v]]\bigr)=0.
\end{align*}
\end{theorem}
\begin{proof}
Differentiate canonical covariance to get
$\delta H=\dot G_\xi-i[G_\xi,H]/\hbar$.
Use the functional Hilbert identity and
$\operatorname{Tr}(\rho\mathfrak d_tG)=
d\operatorname{Tr}(\rho G)/dt$ to obtain the endpoint formula.
Varying the metric by its Lie derivative and integrating by parts gives
the stress formula with the retained external force. Differentiation of
the defect gives the final equation, including $\delta R_\xi[v]$.
\end{proof}
The state family is transported as $\rho^\xi=V_\xi\rho V_\xi^\dagger$.
No invariant-state premise is used for this on-equation source-history
identity. Paper~VII-C's fixed-state product Ward theorem retains its own
invariant-state antecedent. A paired complex finite-unitary Berezin
Jacobian cancels on that branch; it does not certify all chiral, Majorana
or continuum anomaly domains.

For $g$ on the CAR one-body space, a genuinely inherited quartic tensor
obeys
$[d\Gamma(g),\widehat V_4(\mathsf V)]
 =\widehat V_4([g^{(2)},\mathsf V])$.
This follows by commuting $d\Gamma(g)$ with all four ladder slots.
It tests the actual transported tensor; it does not define the action
jets from a desired commutator. Normal/time symmetry generators can be
many-body operators. On the FLRW scalar background,
$\xi^0=TQ$, $\xi^i=LD^iQ$ give
$R_\xi=(T'+\mathcal HT,L'-T,-\mathcal HT,L)$.
Its next metric variation includes products such as $Q^2$ and their
generated directions. The full parent jets must therefore precede
restriction to a four-port slice or a gauge leaf.

The theorem supplies a calculable certificate, not automatic covariance of
every round coefficient window. Appendix~\ref{p4:app:v4106-cutoff-ward-defect}
gives the exact leakage defect and an interacting empty-guard example.
Scalar projectability additionally requires the actual full state,
background and response word to carry the appropriate isometry covariance;
isotropy of one stress moment alone is insufficient.

% END INLINE SOURCE: Paper4-FunctionalMetricWard-Main-v1_0.tex


% Framework v2.189: positive complete-source Ward and transport construction.

% BEGIN INLINE SOURCE: Paper4-HomogeneousTimeWard-Main-v1_0.tex
% Framework v2.189 / Paper IV v4.118: complete native passive homogeneous clock Ward.
% Closes complete passive homogeneous time Ward on the licensed finite domain.
% Preserves the actual clock and source domain; no general local Ward promotion.


\subsection{Complete passive-clock pullback and homogeneous finite time Ward}
\label{p4:sec:v4118-passive-time}

Use the complete action/quantization word of
Theorem~\ref{p3f:thm:v1148-finite-metric-constructor}, before fixing lapse
or deleting its source, clock, constraint, ordering, projector, field,
central or boundary data. Its finite full generator is self-adjoint and
has the differentiability required by
Theorem~\ref{p4:thm:v4106-functional-hilbert-response}.
The actual clock-unit conversion belongs to the action measure.

\begin{theorem}[Native passive homogeneous clock pullback]
\label{p4:thm:v4118-passive-clock-pullback}
Let $s=\phi(t)$ be a smooth orientation-preserving coordinate change of
the same interval, with $\lambda=\phi'>0$ and $\phi(t)=t$ near both
endpoints. Pull back the \emph{complete existing history} and all its
retained marks and finite fibres. This is a passive relabelling: the proper
clock form is $N_\phi dt=\phi^*(Nds)$, so
\[
 N_\phi=\lambda N\circ\phi,\qquad
 x_\phi=x\circ\phi
\]
for scalar configurations. Temporal covector/connection coefficients
pull back with $\lambda$; all derivative jets use the full chain rule.
The inherited occupation/constraint/ordering and field maps are the
pullback of the same maps, rather than newly selected members.
Then, in the pulled-back canonical frame, the complete finite generator
obeys
\[
 \boxed{H_\phi(t)=\lambda(t)H(\phi(t)).}
\]
The physical proper-time calibration, occurrence, state, occupation
boundary and action coefficients are unchanged. This does not license
an independently reparametrized calibrated physical history.
\end{theorem}
\begin{proof}
For the actual first-order fermion parent, $G_p$ and $J$ are pulled back
by composition, while its action coefficients $K_p,B_p$ are temporal
one-forms and acquire $\lambda$. The finite constructor therefore gives
\[
 G_\phi=G\circ\phi,\quad S_\phi=S\circ\phi,
 \quad W_\phi=W\circ\phi,\quad Z_\phi=Z\circ\phi,
 \quad K_\phi=\lambda K\circ\phi,
 \quad B_\phi=\lambda B\circ\phi.
\]
Here $\dot J_\phi=\lambda\dot J\circ\phi$ and
$\dot S_\phi=\lambda\dot S\circ\phi$ retain the complete embedding
and positive-root connection terms. Consequently
\[
 e_\phi=\lambda e\circ\phi,\qquad
 \gamma_\phi=\lambda\gamma\circ\phi.
\]
Every interaction slot uses the same pulled-back $Z$; its action density,
including the actual central energy, is a temporal one-form. Fixed actual
operator products, tensor contractions and ordering pull back by
composition. A moving inherited projector $P(t)$ becomes $P(\phi(t))$;
its time jet acquires $\lambda$, so its connection is retained. No
claim that the projector is fixed is used.

For the actual scalar restriction-before-Legendre word write
\[
 L=\tfrac12\dot z^TM\dot z-\dot z^TCz
   +\tfrac12z^TDz-\tfrac12z^TKz-U.
\]
Change variables in the same action. The coefficients satisfy
\[
 M_\phi=(M\circ\phi)/\lambda,\quad C_\phi=C\circ\phi,
 \quad (D_\phi,K_\phi,U_\phi)=
 \lambda(D,K,U)\circ\phi.
\]
Thus $p_\phi=M_\phi\dot z_\phi-C_\phi z_\phi=p\circ\phi$, and
direct substitution in the actual Legendre word proves
$H_{\Phi,\phi}=\lambda H_\Phi\circ\phi$. It retains $C,D$, their
derivative-image leakage, curvature and connection terms. No scalar
dilation, finite exact CCR or alternative occupation projector is invoked.
Boundary functionals pull back by composition; their derivatives are
one-forms, and the intrinsic canonical endpoint is unchanged.

The native Yang--Mills, scalar, fermion and constraint action constituents
are densities of the same complete local covariant action. The same
change-of-variable calculation transports their actually consumed
covariant background/source jets, including curvature and frame jets,
with the full chain rule before the licensed restriction and canonical
construction. This does not establish a new higher-field-derivative
Legendre, auxiliary-variable or quantization extension. It is not legitimate to keep a
coordinate-frequency cutoff, external profile or clock coefficient fixed
while pretending that its physical mark was pulled back. The stated
complete pullback changes only the presentation of the same actual word.
\end{proof}

\begin{theorem}[Constructed finite homogeneous time-Ward lift]
\label{p4:thm:v4118-homogeneous-time-ward}
Let $U(t,s)$ be the unitary full many-body propagator of the existing
finite generator, using its inherited carrier connection. In a common
coherent frame define
\[
 V_\phi(t):=U(\phi(t),t).
\]
Then this is a canonical lift between the corresponding finite fibres,
and the \emph{already independently pulled-back} generator satisfies
\[
 H_\phi=V_\phi HV_\phi^\dagger
              +i\hbar\dot V_\phi V_\phi^\dagger.
\]
For $\phi_\epsilon(t)=t+\epsilon\xi(t)+O(\epsilon^2)$,
\[
 G_\xi=\xi H,\qquad
 \delta_\xi H=\xi\dot H+\dot\xi H,
 \qquad\boxed{\mathfrak D_\xi=0.}
\]
The full first source variation also obeys
$\delta\mathfrak D_\xi[v]=0$, including the actual
$\delta R_\xi[v]$, all functional contacts and all central energies.
\end{theorem}
\begin{proof}
Finite self-adjoint full evolution provides the unique fibre-to-fibre
propagator. Its two endpoint derivatives give
\[
 i\hbar\dot V_\phi=
 \lambda H(\phi(t))V_\phi-V_\phi H(t).
\]
The constructor theorem, rather than this equality, first identified
$H_\phi=\lambda H\circ\phi$. Multiplying by $V_\phi^\dagger$
therefore proves its actual canonical covariance. Infinitesimally
$V_{\phi_\epsilon}=I-i\epsilon\xi H/\hbar+O(\epsilon^2)$.
The Ward defect of
Theorem~\ref{p4:thm:v4106-canonical-stress-ward} is
\[
 \delta_\xi H-\frac{d}{dt}(\xi H)
                +\frac{i}{\hbar}[\xi H,H]=0.
\]
This is an off-shell constructor identity throughout the pulled-back
source family. Differentiation gives
\[
 \delta^2H[v,R_\xi]+\delta H[\delta R_\xi[v]]
 =\xi\frac{d}{dt}\delta H[v]+\dot\xi\delta H[v],
 \qquad\delta G_\xi[v]=\xi\delta H[v].
\]
The two varied commutators
$[\xi\delta H,H]+[\xi H,\delta H]$ cancel exactly, proving the
full first-variation identity. This uses no one-body, Gaussian,
stationary-state or occupation closure. The state family is
$\rho_\phi=V_\phi\rho V_\phi^\dagger=\rho\circ\phi$ on the
same solution. Interior endpoint support removes the actual
$G_\xi-K[R_\xi]$ boundary term; otherwise retain it.
\end{proof}

\begin{corollary}[Complete homogeneous source work]
\label{p4:cor:v4118-homogeneous-source-work}
Put $F_a=\operatorname{Tr}(\rho\mathcal O_a)$ for the complete functional
insertions of Theorem~\ref{p4:thm:v4106-functional-hilbert-response}.
For scalar source configurations $q^A$ and temporal one-form source
coefficients $w^b$, the complete pulled-back time tangent is
\[
 R_\xi N=(\xi N)^{\boldsymbol\cdot},\quad
 R_\xi q^A=\xi\dot q^A,\quad
 R_\xi w^b=(\xi w^b)^{\boldsymbol\cdot}.
\]
Its exact on-state-equation Ward identity is
\[
 \boxed{N\dot F_N=\sum_A\dot q^AF_A-\sum_b w^b\dot F_{w^b}.}
\]
Every actual independent external source is retained in these work terms.
For the homogeneous scalar radius coordinate $q^0=\log a$, with
$F_N=\mathcal V\rho$ and $F_{q^0}=-3N\mathcal V p$, this becomes
\[
 \dot\rho+3\frac{\dot a}{a}(\rho+p)
 =\frac1{N\mathcal V}
   \left(\sum_{A\ne0}\dot q^AF_A
                 -\sum_b w^b\dot F_{w^b}\right).
\]
On the complete coupled equations it is the homogeneous same-action
energy exchange; it inserts no independent kinetic or transfer law.
\end{corollary}
\begin{proof}
The zero-defect functional identity gives
$\int F_aR_\xi^a\,dt=0$ for interior support. Integrate the
$\dot\xi$ terms and use arbitrariness of $\xi$ to obtain
$-N\dot F_N+\sum_A\dot q^AF_A-\sum_b w^b\dot F_{w^b}=0$.
Substitute the inherited lapse/scale source conversion and
$\dot{\mathcal V}=3\mathcal V\dot a/a$.
\end{proof}

\paragraph{Actual closure scope.}
This closes the passive homogeneous time-Ward lift and its complete
finite-source work identity on the current licensed finite action domain.
The physical proper-time clock calibration is preserved; an independent
rate change of its calibrated cover is not licensed by this proof.
The spatially constant tangent $\xi(t)$ does not prove Ward identities for
$\xi(t,y)=T(t)Q_n(y)$, spatial diffeomorphisms, the full $n>2$ scalar
Einstein-response Ward/projectability problem, continuum/chiral anomaly
domains, or an actual coupled history. Differentiating this global time
identity in a spatial harmonic does not recover the missing local
identity. Such application claims keep their actual consumed domains.

% END INLINE SOURCE: Paper4-HomogeneousTimeWard-Main-v1_0.tex



% BEGIN INLINE SOURCE: Paper4-TransportedEmbeddingCovariance-Main-v1_0.tex
% Framework v2.189 / Paper IV v4.118: exact same-word transported-embedding covariance.


\subsection{The action moment of the positive-frequency modes}
\begin{proposition}[Occupation phase moment and charged coordinate bracket]
\label{p4:prop:v2174-occupation-phase-moment}
On the oscillatory classical domain of
Theorem~\ref{p4:thm:v245-hessian-kahler-projective-geometry}, keep
$\sigma=\sum_jdq_j\wedge dp_j$ and $\Omega>0$. The common phase action
$z\mapsto e^{-i\theta}z$ has the action-valued moment
\begin{equation}
 \mu_{\rm occ}=\frac12(q^T\Omega q+p^T\Omega^{-1}p)
       =\hbar\|z\|^2,\qquad
 \xi_{\rm occ}=(\Omega^{-1}p,-\Omega q),\qquad
 \iota_{\xi_{\rm occ}}\sigma=d\mu_{\rm occ}.
 \label{p4:eq:v2174-occupation-moment}
\end{equation}
For the common charged two-real-mode constituent, use the inherited
$z_+,z_-$ and $\chi\ne0$. Then
\begin{align}
 \mu_Q&=\hbar(|z_+|^2-|z_-|^2),&
 X_{\rm cl}&=\chi(\bar z_-+z_+),\nonumber\\
 \{\mu_{\rm occ},X_{\rm cl}\}
     &=-i\chi(\bar z_--z_+),&
 \{\mu_Q,X_{\rm cl}\}&=iX_{\rm cl},\qquad
 \{\mu_{\rm occ},\mu_Q\}=0.
 \label{p4:eq:v2174-charged-phase-brackets}
\end{align}
\end{proposition}
\begin{proof}
The already generated coordinate
$z=(\Omega^{1/2}q+i\Omega^{-1/2}p)/\sqrt{2\hbar}$ gives the first
identity and $\{z_j,\bar z_k\}=-i\delta_{jk}/\hbar$.
Contracting the displayed vector field with $\sigma$ gives $d\mu_{\rm occ}$;
its action on $z$ is $-iz$. The Leibniz rule gives both charged brackets
and the commuting moments. In particular
$\{\mu_{\rm occ},\{\mu_{\rm occ},X_{\rm cl}\}\}=-X_{\rm cl}$.
\end{proof}
This classical stationary domain is not inferred from the quantum Gibbs
ground theorem. Normalizing a projective mode direction removes its
amplitude; it does not determine the full field phase space or its actual
finite quantum lift. That lift retains the inherited ordering, occupation
and vacuum allocation word. Appendix~\ref{p4:app:v2174-occupation-comparisons}
records the neutral-ray comparison and the exact scope of the alternative
capacity routes.

% END INLINE SOURCE: Paper4-OccupationMoment-Main-v1_0.tex

% END C01 V2174 MODE MOMENT



\begin{remark}[Mode structure versus one-particle semantics]
The preceding theorem removes the mathematical positive-frequency splitting on
its stationary stable field domain.  It does not declare that an admissible
field-mode carrier is physically the \(n=1\) particle sector.  The generic
\(\mathsf{1P}\) interface remains unless a field/species theorem supplies
that semantic role, as happens independently on later visible-matter branches.
\end{remark}

For fermionic statistics one has the finite Fock space
\begin{equation}
\label{p4:eq:fermionic-fock}
\mathcal F_-(\mathcal K)
:=
\bigoplus_{n=0}^{d} \Lambda^n \mathcal K,
\end{equation}
which is automatically finite-dimensional.

\begin{theorem}[GP1 bosonic occupation capacity and generated-boundary closure]
\label{p4:thm:gp1-bosonic-occupation-boundary}
Assume a finite nonzero one-particle block \(\mathcal K\) and an actually
realized bosonic occupation role.  Let \(S_{\rm occ}\subseteq\mathbb N_0\)
be the set of occupation degrees retained at a finite realized stage.  GP1
forces
\[
 \boxed{|S_{\rm occ}|<\infty,\qquad
 N_*:=\max S_{\rm occ}<\infty.}
\]
This existential finite-capacity conclusion does not determine either the
numerical value of \(N_*\) or the initial-sublevel form of the retained set.

If, in addition, tensor degree is the constituted occupation grammar
\[
 \ell_{\rm occ}(\operatorname{Sym}^n\mathcal K)=n
\]
and the registered generated-boundary condition
\[
 \mathsf P_{\rm gbd}(\mathrm{occ},N_{\max})
\]
is realized or discharged by a separate theorem, then GP2 exactness/no-omission
inside that generated sublevel gives
\begin{equation}
\label{p4:eq:bosonic-fock-cutoff}
 \boxed{
 \mathcal F_+^{(\le N_{\max})}(\mathcal K)
 =\bigoplus_{n=0}^{N_{\max}}\operatorname{Sym}^n\mathcal K.}
\end{equation}
Its finite dimension is
\[
 \dim\mathcal F_+^{(\le N_{\max})}(\mathcal K)
 =\sum_{n=0}^{N_{\max}}\binom{n+d-1}{n}
 =\binom{N_{\max}+d}{N_{\max}}.
\]
GP1 therefore supplies finiteness of the realized bosonic occupation envelope;
it does not supply the numerical/generated occupation boundary.  The residual
bosonic occupation realization family can therefore be written as
\[
 \mathfrak A_{\rm OCC}
 :=\{S_{\rm occ}\subset\mathbb N_0:0<|S_{\rm occ}|<\infty\},
\]
or, on the standard initial-sublevel branch, by the still-undetermined
\(N_{\max}\).  A generated-boundary theorem or explicit boundary datum is
required to choose/derive that member.
\end{theorem}

\begin{proof}
For every \(n\ge0\) and \(d=\dim\mathcal K>0\),
\[
 \dim\operatorname{Sym}^n\mathcal K
 =\binom{n+d-1}{n}>0.
\]
If \(S_{\rm occ}\) were infinite, the retained occupation/composition envelope
would contain infinitely many nonzero mutually graded sectors and hence be
infinite-dimensional, contrary to GP1's finite total retained
composition/occupation envelope.  Thus \(S_{\rm occ}\) is finite.  The
second statement is exactly Paper~II-A's generated-sublevel mechanism:
\(\mathsf P_{\rm gbd}\) supplies the intrinsic depth while GP2 retains the
complete generated sublevel.  The final dimension is the hockey-stick
binomial identity.
\end{proof}

\begin{corollary}[Occupation cutoff uniqueness under generated saturation]
\label{p4:cor:occupation-cutoff-saturation}
Let \(d=\dim\mathcal K>0\).  Suppose a separate geometric capacity theorem
proves both that the complete bosonic carrier is exactly the full initial
sublevel through its top degree and that its carrier dimension is the
independently generated integer \(D\).  Then its top degree \(N_{\max}\), if
it exists, is the unique nonnegative integer satisfying
\[
 \boxed{D=\binom{N_{\max}+d}{d}.}
\]
Thus the numerical \(\mathsf{OCC}\) member is derived on this saturation
branch.  A dimensional fit or the weaker embedding inequality
\(\binom{N+d}{d}\le D\) does not discharge \(\mathsf{OCC}\).
\end{corollary}

\begin{proof}
The generated-saturation premise and
Theorem~\ref{p4:thm:gp1-bosonic-occupation-boundary} give the displayed
equality.  Moreover,
\[
 \binom{N+1+d}{d}-\binom{N+d}{d}
 =\binom{N+d}{d-1}>0,
\]
so the dimension is strictly increasing in \(N\) and at most one cutoff can
satisfy it.  The conclusion uses the theorem that the carrier is the entire
generated sublevel with no omitted degrees or unused states; dimension
matching alone supplies no such theorem.
\end{proof}

For the standard finite bosonic Fock branch used below, the theorem domain
therefore supplies or derives \(\mathsf P_{\rm gbd}(\mathrm{occ},N_{\max})\),
not the bare fact that some finite capacity exists.  After the statistics
selector and this occupation boundary are fixed, both fermionic and bosonic
constructions stay inside the finite Born framework.  The resulting
many-particle or effective field-theoretic carrier is a selected direct sum of
already admitted tensor powers, rather than a new carrier ontology.

Choosing an orthonormal mode basis $\{e_i\}_{i=1}^d$ gives the usual occupation labels.
For fermions the basis vectors are
\[
\ket{n_1,\dots,n_d},
\qquad
n_i\in\{0,1\},
\]
and for bosons with cutoff $N_{\max}$ they satisfy
\[
\ket{n_1,\dots,n_d},
\qquad
n_i\in\mathbb N,
\quad
\sum_i n_i\le N_{\max}.
\]
This is the finite mode-space form of many-body quantum mechanics and the starting point of the QFT bridge used later.



\subsection{Finite dispersion on the GP1 Lorentzian branch}
\label{p4:sec:finite-one-particle-dispersion}

Paper~III derives the sharp operational cone and its Lorentzian conformal sign
before a field Hamiltonian is chosen.  On every stationary stable sector for
which Paper~III-F already owns the local action, stationary background,
physical reduction, and positive quadratic Hessian,
Theorem~\ref{p3f:thm:v244-stationary-physical-hessian-propagation} now also
generates the local second-order operator itself.  On the primitive scalar
branch this is consistent with
\[
\widetilde E^2=c^2\|p\|_h^2+\mu^2c^4
\]
from Theorem~\ref{p3:thm:p3-p4-lorentzian-chain-closure}.  The present
subsection has the strictly downstream task of taking the positive spectral
square root, constructing the positive-frequency presentation, and
interpreting isolated spectral bottoms as finite energy thresholds.  Sectors
without an owned action/background/Hessian retain the older supplied-free-law
interface.  Nothing below selects or closes the Paper-III Lorentzian sign.

The finite operator algebra supplies intrinsic spectral, charge, chirality,
grading, and symmetry projectors.  These projectors carve the finite carrier
into blocks.  A block becomes a one-particle block precisely when the finite
Born/Fock construction uses it as the one-quantum mode space and the
reconstructed free generator leaves it invariant.  This does not assume a
continuum particle ontology.  The existence of a local second-order field sector is therefore generated on
the action-owned stationary stable branches just named and remains a stated
domain condition elsewhere; finite functional calculus alone cannot select a
relativistic dispersion function.

\begin{definition}[Intrinsic finite spectral block]
\label{p4:def:intrinsic-finite-spectral-block}
Let \(\mathcal A_{\rm fin}\) be the constituted finite operator algebra generated by the Paper~III compact spatial operator, the Paper~II/Paper~III charge and symmetry operators, and the finite Born event calculus of the present paper.  An \emph{intrinsic finite spectral block} is a finite-dimensional subspace
\[
  \mathcal K=\operatorname{Ran}P
\]
where \(P\in\mathcal A_{\rm fin}\) is obtained by finite functional calculus and finite intersections of constituted spectral, charge, chirality, grading, or symmetry projectors.  No external particle label, fitted sector projector, continuum mass shell, or phenomenological threshold is part of the definition.

The block is \emph{one-particle admissible} when the finite Fock construction of Section~\ref{p4:sec:part2-fock} assigns \(\mathcal K\) as its \(n=1\) mode space.  It is \emph{stable} for a one-particle generator \(H_1\) if \([P,H_1]=0\), and \emph{isolated} if the finite spectrum of \(H_1|_{\mathcal K}\) is separated from the complementary finite spectrum by a positive distance.
For a scalar free-law theorem, it is a \emph{common reducing block} only when
\([P,A]=0\) for the normalized spatial operator \(A\); this extra condition is
required before \(A|_{\mathcal K}\) or a common \(A\)--\(H_1\) mode
decomposition is used.
\end{definition}

\begin{remark}[Reeb polarization and positive frequency]
\label{p4:rem:reeb-polarization-not-time-frequency}
\leavevmode\par\noindent
The CR/Szeg\H{o} polarization of Paper~I-A is naturally one-sided with respect to
the Hopf--Reeb boundary geometry, but it is not by itself the positive-frequency
condition used below.  On the primitive local free-law domain,
Theorem~\ref{p4:thm:finite-one-particle-dispersion} instead defines
\(H_{\rm PDE}\) as the unique positive square root of the action-generated
or otherwise validly imported Lorentzian free operator.  An independently supplied Euclidean kernel
represents that same generator only when the whitened same-carrier criterion
of Lemma~\ref{p4:lem:positive-frequency-generator-from-euclidean-block} is
satisfied; on the forward Hamiltonian-generated branch this criterion holds
identically by
Theorem~\ref{p4:thm:hamiltonian-generated-euclidean-branch}.  This keeps
boundary holomorphicity, Hopf charge, positive-frequency selection, and
Euclidean reconstruction logically separate.
\end{remark}

\begin{lemma}[Whitened same-carrier Hamiltonian criterion]
\label{p4:lem:positive-frequency-generator-from-euclidean-block}
Let \(\mathcal K=\operatorname{Ran}P\) be an intrinsic finite spectral block
and let
\(C_{\mathcal K}:[0,\infty)\to\End(\mathcal K)\) be a norm-continuous
Euclidean two-point kernel.  Restrict to the support of
\(C_{\mathcal K}(0)\) and assume that \(C_{\mathcal K}(0)>0\) there.  Define
\[
 S_\tau:=
 C_{\mathcal K}(0)^{-1/2}
 C_{\mathcal K}(\tau)
 C_{\mathcal K}(0)^{-1/2}.
\]
Then the following are equivalent:
\begin{enumerate}[label=(\roman*),leftmargin=2.2em]
\item there is a positive self-adjoint
\(H_{\mathcal K}\in\End(\mathcal K)\) such that
\[
 C_{\mathcal K}(\tau)
 =C_{\mathcal K}(0)^{1/2}
 e^{-\tau H_{\mathcal K}/\hbar}
 C_{\mathcal K}(0)^{1/2};
\]
\item \((S_\tau)_{\tau\ge0}\) is a norm-continuous semigroup of positive
self-adjoint contractions,
\begin{equation}
\label{p4:eq:one-particle-euclidean-spectral-block}
 S_0=\mathbf1,\qquad
 S_{\tau+\sigma}=S_\tau S_\sigma,\qquad
 0\le S_\tau\le\mathbf1.
\end{equation}
\end{enumerate}
When these conditions hold, the generator is unique and
\[
 \boxed{S_\tau=e^{-\tau H_{\mathcal K}/\hbar}.}
\]
Only on this semigroup-factorization domain does the Euclidean kernel
produce a Hamiltonian on the original block.  Reflection positivity without
this condition still gives the cyclic realization of
Theorem~\ref{p4:thm:finite-euclidean-hamiltonian-reconstruction}, but its
Hilbert space may be larger than \(\mathcal K\).
Accordingly, this equivalence is the finite
\((\mathrm{WICK}_{\rm same})\) certificate.

The unwhitened quotient
\[
 T_\tau:=C_{\mathcal K}(\tau)C_{\mathcal K}(0)^{-1}
 =C_{\mathcal K}(0)^{1/2}S_\tau
  C_{\mathcal K}(0)^{-1/2}
\]
is only similar to the contraction semigroup.  Requiring \(T_\tau\) itself
to be positive and self-adjoint in the original inner product is not
necessary.  It is a sufficient special case when
\([H_{\mathcal K},C_{\mathcal K}(0)]=0\), in which case
\(S_\tau=T_\tau\) and the covariance also admits the left factorization
\(C_{\mathcal K}(\tau)=e^{-\tau H_{\mathcal K}/\hbar}
C_{\mathcal K}(0)\).
\end{lemma}

\begin{proof}
If (i) holds, whitening on both sides gives
\(S_\tau=e^{-\tau H_{\mathcal K}/\hbar}\), which proves (ii).
Conversely, a norm-continuous finite-dimensional semigroup has a bounded
generator \(G\) with \(S_\tau=e^{\tau G}\).  Self-adjointness makes \(G\)
self-adjoint, and contractivity gives \(G\le0\).  Set
\(H_{\mathcal K}:=-\hbar G\ge0\) and undo the whitening to obtain (i).
Differentiation at \(\tau=0\) proves uniqueness.  The formula for \(T_\tau\)
is immediate; it is self-adjoint in the original inner product exactly on
the commuting special branch stated above.

The additional condition is necessary.  A scalar output kernel
\[
C(\tau)=e^{-E_1\tau}+e^{-E_2\tau},
\qquad E_1\ne E_2,
\]
can be recorded on a one-dimensional output block, but its minimal cyclic
reconstruction has two independent positive-energy channels.  Positivity
therefore does not identify the reconstruction generator with an operator
on the original one-dimensional block.
\end{proof}

\begin{definition}[Imported calibrated local scalar free-law domain]
\label{p4:def:minimal-uninserted-free-scalar-cover-domain}
The \emph{local scalar free-law domain} on an intrinsic block
\(\mathcal K\) is the Paper-III domain of
Definition~\ref{p3:def:paper-iv-finite-field-input}: it realizes a real,
passive, time-translation invariant, spatially isotropic local finite-order
cone-faithful scalar propagation role whose characteristic cone is contained in
the sharp GP1 cone, together with a nonnegative invariant scalar lower-order
free term.  Paper~III,
Theorem~\ref{p3:thm:minimal-local-scalar-cone-order}, derives second order as
the least nondegenerate scalar principal realization on this domain.  Its
\emph{primitive unimpeded subdomain} is the Paper-III-E branch whose real
scalar support is the generated primitive unlabelled neutral tensor-unit
descendant and whose principal quadratic relation is an undeformed same-owner
descendant of the operational metric relation.  Paper~III-E,
Theorem~\ref{p3e:thm:metric-principal-descendant-saturation}, then derives
that its characteristic cone exhausts the sharp operational boundary.  The spatial
operator is the restriction to a common reducing block of the normalized
nonnegative round scalar Laplacian \(A\) fixed in Paper~III, so
\([P,A]=0\), \(\sigma_2(A)=h^{-1}\), and the smooth parent satisfies
\(A\mathbf1=0\).

For use in Paper~IV, the imported causal datum is explicitly the triple
\((h,c,\mathcal C_c^+)\), where \(h\) is the positive round spatial metric,
\(c>0\), and
\begin{equation}
\label{p4:eq:explicit-imported-gp1-cone}
 \mathcal C_c^+
 :=
 \left\{(\tau,\xi)\in\mathbb R\oplus T_xX:
 \tau\ge0,
 \ h(\xi,\xi)\le c^2\tau^2\right\}.
\end{equation}
Thus ``contained in the sharp GP1 cone'' means containment in
\(\mathcal C_c^+\), and ``exhausts the sharp GP1 boundary'' means equality
of characteristic boundaries with
\(h(\xi,\xi)=c^2\tau^2\), \(\tau\ge0\).  The existence and shape of this
triple are upstream theorem outputs of the clocked-Hopf composite and its
least physical doctrine, followed by closed convexification.  The interior of
the cone is a mathematical causal envelope, not a claim that every interior
history occurs.  Numerical representatives for \(R_x\) and
\((\mathrm{CAL}_t):T_t>0\), and hence \(c=2\pi R_x/T_t\), remain upstream
calibration data; Paper~IV uses
only the displayed cone and does not infer those calibrations from finite
functional calculus.

This remains a theorem-domain condition on the existence/locality of the free
field sector and on its lower-order/rest data.  It is not inferred from finite
functional calculus: on a finite spectrum, \(A^2\) and every function
\(f(A)\) are just as basis-free as \(A\).  Once the local isotropic
cone-faithful propagation role is constituted, Paper~III derives the least
second-order principal structure; on the undeformed same-owner branch the
Paper-II-A/III-E inheritance theorem fixes the primitive coefficient.
\end{definition}

\begin{lemma}[Finite Hamiltonian square on the imported causal local domain]
\label{p4:lem:second-order-free-collapse}
Let \(A\ge0\) be the Paper-III normalized positive spatial operator restricted
to a common reducing intrinsic finite block \(\mathcal K\).  On the imported
causal local second-order free scalar domain, any positive-frequency
one-particle generator that implements this same local free law has the
affine square
\begin{equation}
\label{p4:eq:affine-collapse-second-order-free-sector}
  H_1^2=\hbar^2v^2 A+\beta\mathbf 1_{\mathcal K},
  \qquad
  0<v\le c,
  \quad
  \beta\ge0.
\end{equation}
For the primitive unimpeded undeformed same-owner branch, the inherited
principal relation gives \(v=c\).  Effective media or interacting sectors may have \(v<c\)
and are not identified with that primitive branch.
\end{lemma}

\begin{proof}
Paper~II operator rigidity and spatial isotropy show that an isotropic
passive local second-order scalar law has spatial principal part \(v^2A\).
Containment of its characteristic cone in
\eqref{p4:eq:explicit-imported-gp1-cone} gives \(v\le c\).  On the primitive
unimpeded branch, the same-owner/no-deformation hypotheses give \(v=c\) by
Lemma~\ref{p4:lem:scalar-cover-speed-normalization}.  Time-translation
invariance and the scalar lower-order clause leave only
\(\beta\mathbf1_{\mathcal K}\).  Restoring \(E=\hbar\omega\) gives the
coefficient \(\hbar^2v^2\).  The assumption \(v>0\) is the nondegenerate
propagating branch, and the GP1 cone gives \(v\le c\).  Finally
\(\beta\ge0\) is the stated positive-rest/stability condition; it does not
follow merely from \(H_1^2\ge0\) on a block whose spectrum of \(A\) is bounded
away from zero.
\end{proof}

\begin{lemma}[Principal-owner inheritance fixes the primitive coefficient]
\label{p4:lem:scalar-cover-speed-normalization}
On the primitive unimpeded branch of
Definition~\ref{p4:def:minimal-uninserted-free-scalar-cover-domain}, suppose
the positive-frequency generator satisfies
\[
H_1^2=\hbar^2v^2 A+\beta\mathbf 1_{\mathcal K}.
\]
Then \(v=c\), and hence
\[
H_1^2=\hbar^2c^2A+\beta\mathbf1_{\mathcal K}.
\]
\end{lemma}

\begin{proof}
The principal relation
\(
\omega^2=v^2\lVert k\rVert_h^2
\)
has future characteristic boundary
\[
 \partial\mathcal C_v^+
 =
 \{(\tau,\xi):\tau\ge0,\ h(\xi,\xi)=v^2\tau^2\}.
\]
By the primitive-unimpeded clause, the scalar support is the unlabelled neutral
tensor-unit descendant and the displayed principal relation is an undeformed
descendant of the operational principal owner.  Paper~II-A,
Theorem~\ref{p2a:thm:undeformed-host-principal-inheritance}, therefore makes
its characteristic set equal to the operational boundary in
\eqref{p4:eq:explicit-imported-gp1-cone}.  At \(\tau=1\), the two spatial
sections are the \(h\)-spheres of radii \(v\) and \(c\); equality gives
\(v^2=c^2\).  Positivity gives \(v=c\).
The term \(\beta\mathbf1\) is zeroth order and therefore does not enter the
principal symbol or its characteristic boundary.  No contact/Reeb
normalization is used in this comparison.
\end{proof}


\begin{corollary}[Inherited operational speed on finite free blocks]
\label{p4:cor:finite-inherited-operational-speed}
On the primitive unimpeded branch of
Definition~\ref{p4:def:minimal-uninserted-free-scalar-cover-domain},
Lemma~\ref{p4:lem:scalar-cover-speed-normalization} and Paper~III-E's
metric-principal descendant theorem imply that the speed denoted \(c\) in the
finite scalar cover is the operational speed,
\[
 \boxed{v=c=c_{\rm op}},
\]
and hence
\[
 H_1^2=\hbar^2c_{\rm op}^2A+\beta\mathbf1_{\mathcal K}.
\]
Paper~IV therefore introduces no second finite-quantum propagation-speed
normalization on this branch.  The generic cone-contained effective branch
with \(v<c_{\rm op}\) remains outside this corollary.
\end{corollary}

\begin{theorem}[Positive-frequency realization of the upstream finite dispersion]
\label{p4:thm:finite-one-particle-dispersion}
Let \(S^3_x(R_x)\) denote the smooth round reciprocal-position parent
geometry supplied by Papers~II--III.  The label \(m_x=1\) denotes its
primitive realized reciprocal-position spectral frontier,
\(\mathcal P_x(1)=E_0^x\oplus E_1^x\); it is a finite harmonic-resolution
window on the already smooth reciprocal role, not a position-space cell set,
not an effectively smooth large-\(m_x\) presentation, and not the compact
momentum carrier \(m_p=1\) or the Paper-II finite Toeplitz momentum carrier.
Let \(A\) be the normalized nonnegative scalar Laplacian on this
smooth parent, and let
\(\mathcal K=\operatorname{Ran}P\) be an intrinsic finite spectral block in
the sense of
Definition~\ref{p4:def:intrinsic-finite-spectral-block}.  Suppose that it is
one-particle admissible, is a common reducing block for \(A\), so that the
finite spatial operator is \(A|_{\mathcal K}=PAP|_{\mathcal K}\), and lies in the
primitive unimpeded subdomain of
Definition~\ref{p4:def:minimal-uninserted-free-scalar-cover-domain}.  Define
the positive-frequency realization of the upstream local free law by
\begin{equation}
\label{p4:eq:upstream-positive-frequency-definition}
H_{\rm PDE}
:=
\sqrt{\hbar^2c^2A+\mu^2c^4\mathbf1_{\mathcal K}},
\qquad \mu\ge0.
\end{equation}
Write \(H_1:=H_{\rm PDE}\) on this theorem domain.
Then \(H_{\rm PDE}\) is the unique positive generator whose square implements
the imported law, the block is stable for it, and
\begin{equation}
\label{p4:eq:one-particle-hamiltonian-square}
  H_{\rm PDE}^2
  =
  \hbar^2c^2 A+\mu^2c^4\mathbf 1_{\mathcal K}
\end{equation}
holds.  If \(A\phi_\lambda=\lambda\phi_\lambda\), then
\begin{equation}
\label{p4:eq:finite-one-particle-dispersion}
  E_\lambda^2
  =
  \hbar^2c^2\lambda+\mu^2c^4,
  \qquad
  E_\lambda\ge0.
\end{equation}
Equivalently, in principal-symbol form,
\begin{equation}
\label{p4:eq:finite-one-particle-mass-shell}
  E^2
  =
  c^2\|p\|^2+\mu^2c^4
\end{equation}
inside the scalar cover.  The form and coefficient in this relation are
imported from the Papers-II--III local free-law theorem; Paper~IV selects the
positive square root on the stable finite quantum block.

If the same block is also reconstructed from a Euclidean kernel, that
reconstruction realizes this dispersion only on the explicit compatibility
domain
\begin{equation}
\label{p4:eq:same-free-law-euclidean-compatibility}
C_{\mathcal K}(\tau)
=
C_{\mathcal K}(0)^{1/2}
 e^{-|\tau|H_{\rm PDE}/\hbar}
C_{\mathcal K}(0)^{1/2},
\qquad
C_{\mathcal K}(0)>0.
\end{equation}
On this domain
Lemma~\ref{p4:lem:positive-frequency-generator-from-euclidean-block}
reconstructs \(H_1=H_{\rm PDE}\).  Reflection positivity of an arbitrary
kernel without \eqref{p4:eq:same-free-law-euclidean-compatibility} reconstructs
some positive generator, but does not force the Paper-III dispersion.
\end{theorem}

\begin{proof}
The block \(\mathcal K\) is the range of an intrinsic projector generated by
the constituted finite operator algebra, not an externally fitted particle
label.  Since \([P,A]=0\), the normalized spatial operator restricts to
\(\mathcal K\) and its positive finite functional calculus is well defined.
On the primitive local second-order domain, generated from the physical
action Hessian when the hypotheses of
Theorem~\ref{p3f:thm:v244-stationary-physical-hessian-propagation} hold and
otherwise retained as the explicitly supplied interface, Paper~III,
Theorem~\ref{p3:thm:p3-p4-lorentzian-chain-closure}, gives
\[
\omega_\lambda^2
=c^2\lambda+\frac{\mu^2c^4}{\hbar^2}.
\]
Equivalently, Lemmas~\ref{p4:lem:second-order-free-collapse} and
\ref{p4:lem:scalar-cover-speed-normalization} give the same square with
\(\beta=\mu^2c^4\).  The positive finite spectral theorem selects the unique
positive square root \(H_{\rm PDE}\) in
\eqref{p4:eq:upstream-positive-frequency-definition}; because it is a
function of \(A|_{\mathcal K}\), it commutes with \(P\).
Application to an eigenvector of \(A\) proves
\eqref{p4:eq:finite-one-particle-dispersion}; the principal-symbol expression
is its leading local relation.  No step feeds back into the prior derivation
of the operational cone or Lorentzian sign.

On the Euclidean compatibility domain, whitening
\eqref{p4:eq:same-free-law-euclidean-compatibility} gives exactly
\[
 C_{\mathcal K}(0)^{-1/2}C_{\mathcal K}(\tau)
 C_{\mathcal K}(0)^{-1/2}
 =e^{-|\tau|H_{\rm PDE}/\hbar}.
\]
Hence the same-carrier criterion reconstructs exactly \(H_{\rm PDE}\)
without any commutation premise between \(C_{\mathcal K}(0)\) and
\(H_{\rm PDE}\).  When they do commute, the covariance additionally admits
the simpler left factorization
\(C_{\mathcal K}(\tau)=e^{-|\tau|H_{\rm PDE}/\hbar}C_{\mathcal K}(0)\).
Without the same-law condition, positive Euclidean spectral exponents are
independent data, which proves the final scope statement.
\end{proof}

\begin{theorem}[Pure scalar-factor positive-frequency representation]
\label{p4:thm:pure-scalar-factor-representation}
On the stationary theorem domain of
Theorem~\ref{p4:thm:finite-one-particle-dispersion}, let
\(\tau_s:(\widetilde t,x)\mapsto(\widetilde t+s,x)\) denote the pure
scalar-cover translation.  Its natural action on scalar solutions is
\[
  (T_s\psi)(\widetilde t,x):=\psi(\widetilde t+s,x).
\]
The positive-frequency solution role is invariant under \(T_s\), and on the
finite block the induced representation is
\begin{equation}
\label{p4:eq:pure-scalar-factor-unitary}
  U_E(s)=\exp\left(-\frac{i}{\hbar}sH_{\rm PDE}\right).
\end{equation}
Thus the pure scalar-factor representation, including its orientation and scale,
is theorem-level once the stationary positive-frequency free-law block is
fixed.  This statement does not assert that \(U_E\) is the actual physical
\(\mathsf{DYN}\) member, nor does it actualize a particular state history.
\end{theorem}

\begin{proof}
For an eigenmode \(H_{\rm PDE}\phi_a=E_a\phi_a\), the positive-frequency
solution is \(\psi_a(\widetilde t,x)=e^{-iE_a\widetilde t/\hbar}\phi_a(x)\).
Hence
\[
 T_s\psi_a=e^{-iE_as/\hbar}\psi_a.
\]
Spectral calculus on the finite block therefore gives
\(U_E(s)=\sum_a e^{-iE_as/\hbar}P_a=e^{-isH_{\rm PDE}/\hbar}\).
Because \(\tau_s\) is the pure scalar-factor action, no spatial-isometry
generator is inserted.  The target is the scalar-factor representation role;
Theorem~\ref{p2a:cor:declared-output-actualization-boundary} prevents its
automatic promotion to an actual DYN/history token.
\end{proof}

\begin{corollary}[Stationary free-law kinetic law is generated]
\label{p4:cor:stationary-free-law-dyn-law-closure}
On the stationary stable same-principal-owner domain of
Theorems~\ref{p4:thm:finite-one-particle-dispersion} and
\ref{p4:thm:pure-scalar-factor-representation}, the kinetic generator and
scalar-translation law are generated:
\[
 H_{\rm PDE}
 =\sqrt{\hbar^2c_{\rm op}^2A+\mu^2c_{\rm op}^4I},
 \qquad
 U_E(s)=e^{-isH_{\rm PDE}/\hbar}.
\]
Thus no further Hamiltonian member is selected on this branch.  An actual
\(\mathsf{DYN}\) token, initial state, and trajectory remain separate G0
occurrence data, and the conclusion does not extend to generic interacting or
effective sectors.
\end{corollary}


\begin{theorem}[Stationary quadratic action as a complex-time real-form pair]
\label{p4:thm:stationary-quadratic-real-form-pair}
On the stationary stable finite free-law domain of
Theorem~\ref{p4:thm:finite-one-particle-dispersion}, let
\[
 \Omega:=H_{\rm PDE}/\hbar\ge0
\]
on the real mode space \(\mathcal K_{\mathbb R}\), and extend its real inner
product complex bilinearly to \(\mathcal K_{\mathbb C}\).  For a field germ
\(q\) holomorphic in a sector joining the positive Euclidean and oriented
Lorentzian time axes, define the quadratic densities
\begin{align}
\label{p4:eq:stationary-quadratic-real-form-actions}
 \mathcal L_E^{(2)}(q)
 &:=\frac12\Bigl(
   (\partial_\tau q,\partial_\tau q)
   +(\Omega q,\Omega q)\Bigr),\\
 \mathcal L_L^{(2)}(q)
 &:=\frac12\Bigl(
   (\partial_t q,\partial_t q)
   -(\Omega q,\Omega q)\Bigr).
\end{align}
Under the oriented real-form restriction
\begin{equation}
\label{p4:eq:stationary-quadratic-wick-map}
 \tau=it,
 \qquad
 d\tau=i\,dt,
 \qquad
 \partial_\tau=-i\partial_t,
\end{equation}
one has the local action-form identity
\begin{equation}
\label{p4:eq:stationary-quadratic-action-form-identity}
 \mathcal L_E^{(2)}\,d\tau
 \longmapsto
 -i\,\mathcal L_L^{(2)}\,dt.
\end{equation}
Consequently, on every contour domain whose endpoint and boundary terms are
transported or vanish,
\begin{equation}
\label{p4:eq:stationary-quadratic-weight-identity}
 S_E^{(2)}\longmapsto-iS_L^{(2)},
 \qquad
 e^{-S_E^{(2)}/\hbar}\longmapsto e^{iS_L^{(2)}/\hbar}.
\end{equation}
The Euclidean Euler operator
\begin{equation}
\label{p4:eq:stationary-quadratic-euler-pair}
 K_E=-\partial_\tau^2+\Omega^2
\end{equation}
continues to the Lorentzian free-law operator
\[
 K_L=\partial_t^2+\Omega^2,
\]
and its positive-frequency solutions are generated by
\(U_E(t)=e^{-itH_{\rm PDE}/\hbar}\).  Hence the Euclidean quadratic action,
the Lorentzian quadratic action, their Euler operators, and their
positive-time dynamics are co-realized real-form children of the same
complex-time free-law parent on this theorem domain.

The conclusion is limited to stationary stable quadratic bosonic blocks (and
to gauge physical polarizations after a separately valid gauge reduction).
It does not continue the exact nonlocal spectral trace as an action, fix the
Einstein--Hilbert real-form sign, treat the indefinite conformal mode or odd
fermionic sector, prove interacting infinite-volume reconstruction, or
actualize a physical \(\mathsf{DYN}\) member or history.
\end{theorem}

\begin{proof}
Substitution of \eqref{p4:eq:stationary-quadratic-wick-map} into the
complex-bilinear Euclidean density gives
\[
 \frac12\left(
 (-i)^2(\partial_tq,\partial_tq)+(\Omega q,\Omega q)
 \right)i\,dt
 =-i\,\mathcal L_L^{(2)}\,dt,
\]
which proves \eqref{p4:eq:stationary-quadratic-action-form-identity} and,
under the stated boundary control, the integrated and exponential identities.
The same substitution sends
\(-\partial_\tau^2+\Omega^2\) to
\(\partial_t^2+\Omega^2\).  By
Theorems~\ref{p4:thm:pure-scalar-factor-representation} and
\ref{p4:thm:complex-time-reciprocal-dynamics-square}, the positive-frequency
solution operation is the imaginary-axis child of the same positive generator.
Paper-II-A rooted incidence and evaluation-dual co-realization therefore retain
the complete restricted diagram.  The excluded sectors require additional
real-form, gauge, sign, or continuum data not present in this quadratic
constructor.
\end{proof}

\begin{theorem}[Action-generated positive-frequency coordinate and Schr\"odinger form]
\label{p4:thm:v244-action-generated-schrodinger-form}
On the oscillatory block of
Theorem~\ref{p4:thm:stationary-quadratic-real-form-pair}, let
\(p=\dot q\), \(\Omega>0\), and
\begin{equation}
\label{p4:eq:v244-positive-frequency-coordinate}
 \boxed{
 z:=\frac1{\sqrt{2\hbar}}
 \left(\Omega^{1/2}q+i\Omega^{-1/2}p\right).}
\end{equation}
Then the real quadratic equation
\(\ddot q+\Omega^2q=0\) is equivalent to
\begin{equation}
\label{p4:eq:v244-schrodinger-from-real-action}
 \boxed{i\hbar\dot z=Hz,\qquad H:=\hbar\Omega.}
\end{equation}
Moreover
\[
 \langle z,Hz\rangle
 =\frac12\left(\|p\|^2+\|\Omega q\|^2\right),
\]
and the real canonical action
\[
 S_{\rm can}=\int dt\left[
 \langle p,\dot q\rangle
 -\frac12\left(\|p\|^2+\|\Omega q\|^2\right)\right]
\]
differs only by the endpoint term
\(\frac12[\langle q,p\rangle]_{t_i}^{t_f}\) from
\begin{equation}
\label{p4:eq:v244-schrodinger-action}
 \boxed{
 S_z=\int dt\left[
 \frac{i\hbar}{2}
 \left(\langle z,\dot z\rangle-\langle\dot z,z\rangle\right)
 -\langle z,Hz\rangle\right].}
\end{equation}
Thus the first-order complex Schr\"odinger form is a presentation of the same
stationary real quadratic dynamics; it is not a second independent wave law.
This theorem does not by itself actualize a \(\mathsf{DYN}\) token or assign
one-particle semantics.
\end{theorem}

\begin{proof}
Differentiate \eqref{p4:eq:v244-positive-frequency-coordinate} and use
\(\dot p=-\Omega^2q\):
\[
 \dot z
 =\frac1{\sqrt{2\hbar}}
 \left(\Omega^{1/2}p-i\Omega^{3/2}q\right)
 =-i\Omega z,
\]
which gives \eqref{p4:eq:v244-schrodinger-from-real-action}.  The energy
identity follows by expanding the Hermitian norm; the mixed terms cancel.
Similarly,
\[
 \frac{i\hbar}{2}
 \left(\langle z,\dot z\rangle-\langle\dot z,z\rangle\right)
 =\frac12\left(\langle p,\dot q\rangle-\langle q,\dot p\rangle\right),
\]
and integration of the second term by parts yields the stated endpoint
difference.  Variation of \(S_z\) with fixed endpoints gives the same
first-order equation.
\end{proof}

\begin{corollary}[Generated finite retarded, advanced and Euclidean Green operators]
\label{p4:cor:v244-generated-green-operators}
For \(K_L=\partial_t^2+\Omega^2\), define by finite functional calculus
\[
 s_t(\lambda)=
 \begin{cases}
 \sin(t\lambda)/\lambda,&\lambda\ne0,\\
 t,&\lambda=0.
 \end{cases}
\]
Then the Cauchy solution is
\[
 q(t)=\cos[\Omega(t-t_0)]q_0+s_{t-t_0}(\Omega)p_0,
\]
and
\[
 \boxed{G_R(t)=\Theta(t)s_t(\Omega),\qquad
 G_A(t)=-\Theta(-t)s_t(\Omega)}
\]
satisfy \(K_LG_{R/A}=\delta(t)I\) distributionally.  Their difference is the
causal propagator \(E(t)=s_t(\Omega)\).  On a strictly positive block, the
Euclidean inverse is
\[
 \boxed{G_E(\tau)=\frac1{2\Omega}e^{-|\tau|\Omega}},
 \qquad
 (-\partial_\tau^2+\Omega^2)G_E=\delta(\tau)I.
\]
These propagators are generated from the same finite physical Hessian and do
not constitute an additional transfer-law occurrence.
\end{corollary}

% BEGIN M8 v1.78 FER01
% Results: B03-DUAL, B03-PAIRED-MEASURE, B03-PFAFFIAN, B03-CLOCK, B03-NOETHER


\section{Complete action-generated finite clock evolution}
\label{p4b:sec:v238-b3}
\noindent\textit{Source and claim boundary.} The complete finite action and its already constituted generator control this sector. The mathematically generated evolution law and source variations are kept separate from physical initial-state and readout incidence.

\subsection{Spherical finite-source work, ports and causal response}
\label{p4:sec:spherical-dirac-source}

\begin{proposition}[Exact coordinate-source work]
\label{p4:prop:coordinate-source-work}
Let an already licensed complete finite source have coordinate-only proper
energy \(h(A,B,\phi)\) in an actual flat canonical source frame, with full
state \(\dot\sigma=-i[h,\sigma]/\hbar\). The clock lift is
\(H_t=Nh\); a retained moving frame additionally carries its actual affine
connection term. Put \(V_3=\Omega_z AB^2\), and define the physical ports
\[
 \rho=\frac{\operatorname{Tr}(\sigma h)}{V_3},\quad
 p_z=-\frac A{V_3}\operatorname{Tr}(\sigma h_{,A}),\quad
 p_\perp=-\frac B{2V_3}\operatorname{Tr}(\sigma h_{,B}),\quad
 J_a=\frac{\operatorname{Tr}(\sigma h_{,\phi_a})}{V_3}.
\]
Then, for every full finite state,
\begin{equation}
 \boxed{\dot\rho+u(\rho+p_z)+2v(\rho+p_\perp)
 =J_a\dot\phi_a,\qquad u=\dot A/A,\quad v=\dot B/B.}
 \label{p4:eq:coordinate-source-work}
\end{equation}
\end{proposition}
\begin{proof}
The source ports are the lapse and metric/scalar derivatives of the complete
action energy, with the actual temporal clock lift. The state commutator has
zero expectation against \(h\). Differentiate
\(\rho V_3=\operatorname{Tr}(\sigma h)\), and use
\(\dot V_3/V_3=u+2v\). This proves the identity with all interaction and
central derivatives retained; there is no Wick or one-body closure.
\end{proof}
In native scalar--metric action units these ports are divided by \(\hbar\)
at proper-time \(c=1\), as stated in Paper III-F. Moving frames use their
covariant functional contacts; additional field/open-source work is retained
on both sides. This exact coordinate-source identity is not an assertion of
an anomaly-free general interacting Ward identity.



\paragraph{Canonical product-band generator}
\begin{proposition}[Complete free Dirac source on the spherical product]
\label{p4:prop:spherical-free-dirac-source}
Use an actual admitted complete conjugate-sign free band of the inherited
metric/spin datum
\(A^2dz^2+B^2d\Omega_2^2\), with its owned finite embedding and axial
generator value \(k\). On the unit two-sphere the spin Dirac eigenvalues
are \(\pm\nu\), \(\nu=\ell+1\geq1\), each sign with multiplicity
\(2\nu\) \cite{BaerSphereDiracSpectrum}. No new position carrier, axial
periodicity, band or mass calibration is inserted.
On a fixed inherited invariant scalar mass sector \(\mu_D(\phi)\), the
canonical generator is
\begin{equation}
 \boxed{h_D=P_z\Gamma_z+P_\Omega\Gamma_\Omega+\mu_D\beta,
 \quad P_z=\hbar c_{\rm op}k/A,\quad
 P_\Omega=\hbar c_{\rm op}\nu/B,}
 \label{p4:eq:spherical-free-dirac-generator}
\end{equation}
where the three fixed Hermitian Clifford matrices square to \(I\) and
anticommute pairwise. Consequently
\begin{equation}
 \boxed{E^2=P_z^2+P_\Omega^2+\mu_D^2,\qquad W(h_D)=2E.}
 \label{p4:eq:spherical-free-dirac-width}
\end{equation}
The complete many-body word is \(h=d\Gamma(h_D)+\varepsilon_0I\), with
the actual central allocation retained. For the exact one-body expectation
matrix \(C\) of the full finite state, its physical ports are
\begin{align}
 \rho_D&=\frac{\operatorname{tr}(Ch_D)+\varepsilon_0}{V_3},\nonumber\\
 p_{z,D}&=\frac{P_z\operatorname{tr}(C\Gamma_z)-A\varepsilon_{0,A}}{V_3},
 &p_{\perp,D}&=\frac{P_\Omega\operatorname{tr}(C\Gamma_\Omega)
 -B\varepsilon_{0,B}}{2V_3},\nonumber\\
 J_{D,a}&=\frac{\operatorname{tr}(C\beta)\mu_{D,a}
 +\varepsilon_{0,\phi_a}}{V_3}.
 \label{p4:eq:spherical-free-dirac-ports}
\end{align}
\end{proposition}
\begin{proof}
The raw symmetric kinetic density scales as \(AB^2\). Pull back with the
full embedding \((AB^2)^{-1/2}J_0\), before an instantaneous energy
projection. The canonical kinetic form is \(I\); the symmetric embedding
connection cancels, and the homogeneous coframe has no rotational clock
connection. The product Dirac symbol separates into the axial and sphere
blocks, yielding the generator. The Clifford relations give \(h_D^2=E^2I\).
The physical source ports follow by differentiating the complete many-body
energy and retaining \(\varepsilon_0\). Quadratic expectations close
exactly by the CAR Lie homomorphism even for a non-quasi-free full state.
\end{proof}
When an inherited mass/reduction has additional metric jets, retain them in
the complete functional ports. Mass matrices use the full inherited operator;
the scalar-sector display applies only when that sector is actually invariant.
A symmetry certificate is required for a pointwise homogeneous metric
solution. Without it, these are homogeneous action-average ports.

\begin{corollary}[Positive constituent and anisotropic scaling]
\label{p4:cor:spherical-free-positive-constituent}
For an instantaneous positive-energy occupation of size \(n\), the kinetic
constituent satisfies
\[
 \rho_{\rm kin}=nE/V_3,\qquad
 p_{z,\rm kin}=nP_z^2/(V_3E),\qquad
 p_{\perp,\rm kin}=nP_\Omega^2/(2V_3E).
\]
That occupation condition need not persist when the direction of \(h_D\)
changes. On an actually retained massless \(k=0\) positive member, the
direction is fixed, so the occupation is constant and
\[
 \rho_{\rm kin}\propto(AB^3)^{-1},\qquad p_{z,\rm kin}=0,
 \qquad p_{\perp,\rm kin}=\rho_{\rm kin}/2.
\]
Full stress still includes the central and other consumed sectors.
\end{corollary}
\begin{proof}
Use the Clifford eigenstate expectation \(\langle\Gamma_i\rangle=d_i/E\)
in the preceding ports. For \(\mu_D=k=0\), generators at different times
commute, so the positive occupation is preserved and \(E\propto B^{-1}\).
\end{proof}

\begin{corollary}[Mass work and the complete central trace]
\label{p4:cor:spherical-dirac-trace-source-work}
If the actual mass word is linear homogeneous in the scalar field, then
\begin{equation}
 \boxed{T_D+\operatorname{Tr}_F(\Phi J_D)
 =-\frac{\varepsilon_0+A\varepsilon_{0,A}+B\varepsilon_{0,B}
 -\phi_a\varepsilon_{0,\phi_a}}{V_3},}
 \label{p4:eq:spherical-dirac-trace-source-work}
\end{equation}
where \(T_D=-\rho_D+p_{z,D}+2p_{\perp,D}\).
A conformally homogeneous central energy of joint degree \(-1\) contributes
zero; a constant vacuum density \(\varepsilon_0=V_3\rho_{\rm vac}\)
contributes \(-4\rho_{\rm vac}\). Actual mass offsets retain their residual.
\end{corollary}
\begin{proof}
The full operator identity
\(h_D+A h_{D,A}+B h_{D,B}-\phi_a h_{D,\phi_a}=0\)
cancels both the kinetic and homogeneous mass terms in the trace, for every
full state. Euler homogeneity then evaluates the stated central cases.
\end{proof}
This keeps the vacuum/scalar source which disappears from commutators but
still contributes to the native curvature trace.



\paragraph{Evaluated causal free-source block}
\begin{proposition}[Two-orientation retarded Dirac kernel]
\label{p4:prop:spherical-dirac-retarded-kernel}
On the preceding invariant scalar Clifford sector, put
\(d=(P_z,P_\Omega,\mu_D)\), \(Z=-i\Gamma_z\Gamma_\Omega\beta\), and
\(P_\eta=(I+\eta Z)/2\), \(\eta=\pm1\). For
\(s_\eta=\operatorname{tr}(CP_\eta\Gamma)\), the exact quadratic readout
evolution and its rotation are
\[
 \dot s_\eta=\frac{2\eta}{\hbar}d\times s_\eta,\qquad
 \dot R_\eta(t,s)=\frac{2\eta}{\hbar}[d(t)]_\times R_\eta(t,s),
 \qquad R_\eta(s,s)=I.
\]
For \(d_a=\partial_a d\), the inherited complete retarded kernel is
\begin{equation}
 \boxed{\Pi^R_{ab}(t,s)=\Theta(t-s)\sum_{\eta=\pm1}\frac{2\eta}{\hbar}
 \big[(R_\eta(t,s)^Td_a(t))\times d_b(s)\big]\cdot s_\eta(s).}
 \label{p4:eq:spherical-dirac-retarded-kernel}
\end{equation}
Its contact is \(\operatorname{Tr}(\sigma h_{,ab})\), including all central
and mass derivatives. In particular
\(h_{D,AA}=2P_z\Gamma_z/A^2\) and
\(h_{D,BB}=2P_\Omega\Gamma_\Omega/B^2\) on this member.
\end{proposition}
\begin{proof}
The central involution \(Z\) fixes the two Clifford orientations, with
\([a\cdot\Gamma,b\cdot\Gamma]=2i\eta(a\times b)\cdot\Gamma\)
on its \(\eta\) block. The exact CAR commutator gives the readout evolution.
Heisenberg transport of the later observable replaces \(d_a(t)\) by
\(R_\eta(t,s)^Td_a(t)\). Substitute this into
\(-i\Theta(t-s)\operatorname{Tr}(\sigma(s)[D_a(t,s),D_b(s)])/\hbar\)
from Theorem~\ref{p4:thm:v4104-finite-interacting-source-response}.
\end{proof}
The full state is retained; these port coordinates do not replace it.
Mixing mass sectors use the full unitary formula rather than an unlicensed
invariant projector. This evaluates a supplied finite source block, not the
complete higher-curvature metric response or an infinite-time memory bound.

% END INLINE SOURCE: Paper4-SphericalDiracSource-Main-v1_0.tex



% BEGIN INLINE SOURCE: Paper4-DynamicTestsTensorResponse-Main-v1_0.tex
% Framework v2.163: realizations of existing finite dynamics/response owners.
\paragraph{Owned tests and common finite dynamics}
\label{p4:sec:owned-dynamic-tests}
\begin{corollary}[Common-flow certificate on generated tests]
\label{p4:cor:common-flow-dynamic-tests}
Let $B_*$ be the algebra of
Corollary~\ref{p2a:cor:least-dynamic-test-algebra} in a licensed finite
host, invariant under all Heisenberg generators $i[H(t),\cdot]/\hbar$.
If another licensed Hamiltonian satisfies
$[H'(t)-H(t),B_*]=0$ for all times, their restricted Heisenberg
propagators agree. States agreeing on $B_*$ initially give the same
expectation histories for its tests. The same conclusion holds for a
declared invariant linear test space when only linear expectations are
consumed.
\end{corollary}
\begin{proof}
The restricted generators coincide pointwise. Uniqueness of the finite
linear initial-value problem gives identical propagators on the invariant
space. Pairing these propagated tests with the equal initial restricted
functionals proves the expectation statement.
\end{proof}
For time-dependent generators, the algebra is closed under the complete
family, not a single instantaneous member. This is a certificate for a
common readout; it does not select an unresolved full action. Energy and
stress consumers retain their complete source-jet packet even when two
Hamiltonians differ only by a metric-dependent central energy.

\paragraph{Finite causal TT response of the actual Dirac member}
\label{p4:sec:finite-tt-dirac-response}
\begin{corollary}[TT application of the finite source response]
\label{p4:cor:finite-tt-source-response}
On a licensed free bilinear CAR source, let $C(t)$ be its actual one-body
covariance, $U(t,s)$ its propagator, and $D_A^{(h)}$ the source columns of
Theorem~\ref{p3f:thm:tt-dirac-vertex-contact}. In a fixed canonical frame
with no linear derivative vertex the induced tangent and force are
\begin{align*}
 \dot\sigma&=-\frac{i}{\hbar}[H_D,\sigma]
       -\frac{i}{\hbar}[D_B^{(h)},C]h_B,\\
 \delta F_A^{(h)}&=\operatorname{tr}(CQ_{AB}^{(h)})h_B
                         +\operatorname{tr}(D_A^{(h)}\sigma),\\
 K^R_{AB}(t,s)&=-\frac{i}{\hbar}\Theta(t-s)
  \operatorname{tr}\bigl(C(s)
      [U(t,s)^\dagger D_A^{(h)}(t)U(t,s),D_B^{(h)}(s)]\bigr).
\end{align*}
The retarded tangent has its actual initial value; an independently owned
initial tangent contributes its transported force. All expectation kernels
on the Killing-supported branch vanish, while the complete contact remains.
For an interacting or moving-carrier source, use the full many-body
functional insertion and contact of
Theorem~\ref{p4:thm:v4106-functional-hilbert-response} instead.
\end{corollary}
\begin{proof}
Apply Theorem~\ref{p4:thm:v4104-finite-interacting-source-response} to the
actual TT columns. For bilinear CAR operators,
$[d\Gamma(A),d\Gamma(B)]=d\Gamma([A,B])$ and
$\langle d\Gamma(A)\rangle=\operatorname{tr}(CA)$, so the full response
reduces exactly to these one-body equations without Gaussian factorization.
Variation of the vertex retains its $Q$ contact. The annihilation identities
in III-F prove the Killing-band assertion. The general functional theorem
retains derivative contacts, moving-frame curvature and central jets on
its own domain.
\end{proof}
No new response owner, instantaneous population reset, continuum limit or
universal interacting Ward certificate is introduced. Paper~VI-A couples
this actual source to the same-action gravitational TT block.

% END INLINE SOURCE: Paper4-DynamicTestsTensorResponse-Main-v1_0.tex



% BEGIN INLINE SOURCE: Paper4-FiniteActionEvolution-Main-v1_0.tex
% Framework v2.190: canonical Euler/CAR compatibility and complete propagation.


\subsection{Complete finite action-generated clock evolution}
\label{p4:sec:v2190-finite-action-evolution}

The constructor here consumes the complete normalized action/quantization
word, not only its terminal masses or its quadratic block.  In particular
Theorem~\ref{p3f:thm:v1148-finite-metric-constructor} gives
\[
 \widehat H=d\Gamma(e)-i\hbar d\Gamma(\gamma)
            +\widehat V_{\rm inherited,can}+\varepsilon_0I,
 \quad e=WKW,\quad
 \gamma=WBW+\tfrac12(W\dot S-\dot SW).
\]
Its field and dual-field embeddings, every interaction slot, ordering,
constraint/occupation projector and central term are retained.  A finite
CAR source is finite dimensional.  A joint bosonic word is finite only
when its actual inherited occupation space is finite; mode rank alone
does not meet this domain.

\begin{theorem}[Canonical finite CAR Euler law fixes its full derivation]
\label{p4:thm:v2190-car-euler-derivation}
Let the existing canonical finite CAR generators be \(a_k,a_k^\dagger\),
and let \(Q\) be the retained linear normal-symbol quantization map.  Its
creation slots precede annihilation slots in the word's fixed order.  Let
\(h(\bar c,c,t)\) be the complete even Grassmann normal symbol of the
actual Hermitian \(H(t)=Qh(t)\), including any inherited lower-order and
central terms created by the actual ordering.  For left Grassmann
derivatives,
\[
 \boxed{[a_k,H]=Q(\partial^L_{\bar c_k}h),\qquad
 [a_k^\dagger,H]=Q(\partial^L_{c_k}h).}
\]
Suppose the already licensed same-clock infinitesimal field-reapplication
relation is the image of the normalized canonical Euler word
\[
 i\hbar\dot c_k=\partial^L_{\bar c_k}h,\qquad
 i\hbar\dot{\bar c}_k=\partial^L_{c_k}h.
\]
Then its derivation on the full pure CAR algebra is uniquely
\[
 \boxed{\mathcal D_t(O)=\frac{i}{\hbar}[H(t),O].}
\]
There is no independent Heisenberg-law selector on that same declared
quantized-Euler domain.  The center of a Hamiltonian representation is
invisible to this derivation but remains fixed by the complete action.
\end{theorem}
\begin{proof}
For an even ordered monomial \(M=x_1\cdots x_{2r}\) of odd CAR slots,
moving \(a_k\) through the word gives
\[
 [a_k,M]=\sum_{j=1}^{2r}(-1)^{j-1}
 x_1\cdots\{a_k,x_j\}\cdots x_{2r}.
\]
The only nonzero anticommutators are creation-slot contractions.  This is
exactly the left \(\bar c_k\)-derivative under \(Q\); using
\(a_k^\dagger\) instead contracts the annihilation slots with the same
sign rule.  Linearity proves both identities.  No multiplicativity of
the Grassmann symbol map is assumed.

Varying \(i\hbar\bar c\dot c-h\) with the left-derivative convention
gives the two canonical Euler equations; the symmetrized kinetic form
differs only by its retained endpoint term.  The commutator identities
therefore fix the infinitesimal operation on all canonical generators.
Ordinary Leibniz reapplication fixes it on their entire generated algebra.
The inner derivation displayed has precisely those values, so is unique.
Hermiticity preserves the star operation and all CAR relations.  Two
Hamiltonians giving it differ by the center; the action's inherited
\(\varepsilon_0\) is not thereby discarded.
\end{proof}

For operator coefficients in an additional commuting finite sector, the
two slot identities still hold in the full joint word.  The CAR slots alone
do not determine dynamics in that sector's commutant.  Every other retained
sector and its literal finite commutation/projector boundary terms must
therefore remain in the complete generator and target.  In particular no
exact finite canonical CCR relation is inserted.  The theorem establishes
compatibility of an already licensed Euler/quantization/reapplication
diagram; it does not declare that diagram physical solely because its
normal-symbol formula exists.

\begin{theorem}[Unique propagator of the complete finite action generator]
\label{p4:thm:v2190-complete-finite-propagator}
On the actual finite Hilbert space of the complete licensed word, let
\(H(t)=H(t)^\dagger\) be continuous on a compact clock interval.  When
its physical propagation target is the declared output of that word,
there is a unique unitary propagator
\[
 i\hbar\partial_tU(t,s)=H(t)U(t,s),\quad U(s,s)=I,
 \qquad U(t,s)U(s,r)=U(t,r).
\]
It is a co-realized descendant of the same complete parent; a generated
parent family retains the corresponding compatible propagator family.
For constant \(H\), \(U(t,s)=e^{-i(t-s)H/\hbar}\).  The full state and
pulled-back observable laws are
\[
 \rho(t)=U(t,s)\rho(s)U(t,s)^\dagger,
 \quad\dot\rho=-\frac{i}{\hbar}[H,\rho],
 \qquad O_H(t)=U(t,s)^\dagger O U(t,s).
\]
For a time-dependent generator,
\(\partial_tO_H=(i/\hbar)U^\dagger[H(t),O]U\) for a fixed
Schr\"odinger-picture \(O\).  The time dependence does not in general
define a one-parameter group.
\end{theorem}
\begin{proof}
The Volterra series for the matrix initial-value problem has nth-term
bound \((M|t-s|/\hbar)^n/n!\), where \(M\) bounds \(\|H\|\) on the
compact interval.  Uniform convergence gives a solution.  Iterating the
integral difference equation proves uniqueness.  Differentiating
\(U^\dagger U\) proves unitarity; uniqueness proves composition.  The
state and observable formulas follow by differentiation.  Rooted
co-realization applies after the complete word and its propagation target
are certified, as in Theorem~\ref{p2a:thm:v2190-clock-evolution-closure}.
This is not an extra occurrence of an already generated law.
\end{proof}

\begin{corollary}[Complete frame covariance and exact constraint transport]
\label{p4:cor:v2190-frame-constraint-transport}
For an actual differentiable unitary frame \(R(t)\) of the finite space,
the coefficient change \(\Psi'=R^\dagger\Psi\) gives
\[
 H'=R^\dagger HR-i\hbar R^\dagger\dot R,
 \qquad U'(t,s)=R(t)^\dagger U(t,s)R(s).
\]
For a field frame its actual Fock lift is used, with every interaction slot
transformed; a bilinear-only replacement is insufficient.  The induced
connection is the same kinetic/embedding contribution retained by the
complete finite action constructor.

For a \(C^1\) fixed-rank orthogonal projector \(P(t)\), exact transported
support is equivalent to
\[
 \boxed{P(t)=U(t,s)P(s)U(t,s)^\dagger
 \quad\Longleftrightarrow\quad
 \dot P+\frac{i}{\hbar}[H,P]=0.}
\]
For a constant projector this is \([H,P]=0\).  Compression alone is not
a proof that restricted trajectories equal parent trajectories.
\end{corollary}
\begin{proof}
Differentiate \(\Psi'=R^\dagger\Psi\) and the displayed propagator.
For the projector, differentiation proves necessity.  Conversely
\(\partial_t(U^\dagger P U)=0\) proves sufficiency.  Since \(U\) is
unitary and the rank is fixed, equality of transported ranges is equality
of their orthogonal projectors.  All retained state, source and projector
jets remain part of the same word.
\end{proof}

The joint geometric/backreaction representation is
Theorem~\ref{p3f:thm:v2190-joint-characteristic-flow}.  It requires the
actual source-corrected Legendre matrix, not source normalization alone.
The present finite result neither supplies a physical initial state nor
forces a clock turn to be a period of an interacting Hamiltonian.  It
uses the inherited elapsed clock cover and calibration.  The actual
electron cyclic pole, local Ward transport and consumed continuum limits
remain their existing application evaluations; they are not replaced by
a new universal dynamics principle.  Diagnostics for noncommuting roots,
CAR ordering and moving projectors are in
Appendix~\ref{p4:app:v2190-evolution-diagnostics}.

% END INLINE SOURCE: Paper4-FiniteActionEvolution-Main-v1_0.tex


\paragraph{Module 3: finite thermal/KMS blocks and kinetic limits}
\label{p4:sec:module-finite-thermal-kms}

For a self-adjoint finite generator \(H_N\) supplied on a thermal theorem
domain, the canonical finite Gibbs ensemble is the density matrix
\begin{equation}
\label{p4:eq:finite-thermal-gibbs-native}
\rho_{\beta,N}=Z_{\beta,N}^{-1}e^{-\beta H_N},
\qquad
Z_{\beta,N}=\Tr(e^{-\beta H_N}).
\end{equation}
This is the unique KMS density on a simple full matrix algebra.  On a direct
sum it is one central mixture of the block Gibbs densities; KMS equilibrium
alone leaves their probabilities free, as made explicit in
Proposition~\ref{p4:prop:v236-thermal-atlas}.
Thermal QFT language is licensed by finite trace identities first and by infinite-volume KMS limits only after a compatible limiting dynamics is shown to exist.

\begin{theorem}[Exact finite KMS identity]
\label{p4:thm:finite-thermal-kms-kinetic-interface}
Let \((A_N,H_N)\) be a finite Heisenberg system with Gibbs state \(\rho_{\beta,N}\) and dynamics \(\alpha_z(A)=e^{izH_N}Ae^{-izH_N}\) for real or imaginary time parameter \(z\) on analytic finite matrices.  Then for all \(A,B\in A_N\),
\begin{equation}
\label{p4:eq:finite-kms-native}
\omega_{\beta,N}\!\left(A\alpha_{t+i\beta}(B)\right)
=
\omega_{\beta,N}\!\left(\alpha_t(B)A\right),
\qquad t\in\mathbb R,
\qquad
\omega_{\beta,N}(X)=\Tr(\rho_{\beta,N}X),
\end{equation}
where
\(\alpha_{i\beta}(B)=e^{-\beta H_N}Be^{\beta H_N}\).
At \(t=0\) this reduces to
\(\omega_{\beta,N}(A\alpha_{i\beta}(B))
=\omega_{\beta,N}(BA)\).
\end{theorem}


\begin{proof}
Because \(\alpha_{t+i\beta}=\alpha_t\alpha_{i\beta}\) and
\(e^{-\beta H_N}\) commutes with \(e^{itH_N}\), cyclicity of the trace gives
\[
\begin{aligned}
Z_{\beta,N}\,
\omega_{\beta,N}\!\left(A\alpha_{t+i\beta}(B)\right)
&=
\Tr\!\left(
e^{-\beta H_N}A e^{itH_N}e^{-\beta H_N}
B e^{-itH_N}e^{\beta H_N}
\right)\\
&=
\Tr\!\left(Ae^{itH_N}e^{-\beta H_N}Be^{-itH_N}\right)\\
&=
\Tr\!\left(e^{-\beta H_N}\alpha_t(B)A\right)
=
Z_{\beta,N}\,
\omega_{\beta,N}\!\left(\alpha_t(B)A\right).
\end{aligned}
\]
\end{proof}

\begin{remark}[Boundary between finite KMS data and kinetic rates]
\label{p4:rem:finite-kms-kinetic-boundary}
For a fixed finite bath, inserting its spectral resolution produces a
pure-point thermal spectral \emph{measure}, not an integrable mixing
correlation and not by itself a dissipative transition-rate matrix.  Exact
finite reduction gives the generally recurrent CPTP family of
Theorem~\ref{p4:thm:exact-finite-system-environment-reduction}.
A GKLS kinetic generator is obtained only on the mixing-reservoir,
repeated-probe, or simultaneous large-bath/weak-coupling domains of
Theorem~\ref{p4:thm:microscopic-bath-gkls}, where the reservoir branches use
Fourier transforms of limiting correlations.  Boltzmann or density-matrix
kinetic equations are therefore conditional reduced/interface dynamics.  A
closed scalar population equation further requires invariance of the selected
commutative occupation algebra under the adjoint GKLS generator, as in
Theorem~\ref{p4:thm:finite-reheating-interface}; it is not a consequence of
the finite KMS identity alone and is not a primitive continuum postulate.
\end{remark}

\paragraph{Module 4: finite-window scattering and LSZ boundary}
\label{p4:sec:module-finite-scattering-lsz}

Scattering is finite channel transport before it is an asymptotic continuum \(S\)-matrix.  Choose resolved in/out channel projectors \(P_{\mathrm{in}}\) and \(P_{\mathrm{out}}\) and a finite time window \([-T,T]\).  For normalized channel states \(a\in P_{\mathrm{in}}\mathcal H_N\) and \(b\in P_{\mathrm{out}}\mathcal H_N\), define
\begin{equation}
\label{p4:eq:finite-channel-s-matrix-native}
S^{(N,T)}_{ba}
=
\langle b,P_{\mathrm{out}}U_N(T,-T)P_{\mathrm{in}}a\rangle .
\end{equation}


\begin{theorem}[Finite-window scattering; LSZ as large-window interface]
\label{p4:thm:finite-channel-scattering-lsz-interface}
\leavevmode\par\noindent
For finite projectors \(P_{\mathrm{in}},P_{\mathrm{out}}\) and finite unitary
evolution \(U_N(T,-T)\), the channel matrix
\(S^{(N,T)}=P_{\mathrm{out}}U_N(T,-T)P_{\mathrm{in}}\) is a contraction:
\begin{equation}
\label{p4:eq:finite-scattering-contraction}
(S^{(N,T)})^*S^{(N,T)}\le P_{\mathrm{in}}.
\end{equation}
The missing probability is exactly leakage into unresolved or unobserved
channels,
\begin{equation}
\label{p4:eq:finite-channel-leakage}
\|P_{\mathrm{in}}a\|^2-\|S^{(N,T)}a\|^2
=
\|(\mathbf 1-P_{\mathrm{out}})U_N(T,-T)P_{\mathrm{in}}a\|^2.
\end{equation}
If a compatible large-\(m_x\) reciprocal-position/window refinement exists
in which the channel wave packets separate and the finite amplitudes are
Cauchy, their limits define the corresponding scattering amplitudes.
This statement alone is not LSZ reduction.  LSZ is obtained only on the
further interface domain on which:
\begin{enumerate}[label=(\alph*), leftmargin=2em]
\item the limiting theory is local or graded-local, translation covariant,
      has an invariant vacuum and satisfies the spectrum condition, with
      stable isolated mass shells and Haag--Ruelle/LSZ asymptotic fields for
      the chosen channels;
\item each external channel is represented by an isolated simple
      one-particle pole with nonzero residue;
\item the chosen interpolating field has nonzero overlap with that
      one-particle state;
\item the connected time-ordered correlators and their amputations converge
      in a topology supporting the on-shell limit; and
\item external-leg normalization, boundary terms, and the interchange of
      the large-window and on-shell limits are controlled.
\end{enumerate}
On that domain the standard LSZ theorem identifies the amputated-correlator
limits with the channel-amplitude limits.  LSZ is not a primitive axiom of
Paper~IV and not a flow of \(m_p=1\).
\end{theorem}


\begin{proof}
The contraction inequality follows from \(P_{\mathrm{out}}\le\mathbf 1\) and unitarity:
\[
(S^{(N,T)})^*S^{(N,T)}=P_{\mathrm{in}}U^*P_{\mathrm{out}}UP_{\mathrm{in}}\le P_{\mathrm{in}}U^*UP_{\mathrm{in}}=P_{\mathrm{in}}.
\]
Equation~\eqref{p4:eq:finite-channel-leakage} is Pythagoras for the
orthogonal splitting
\(P_{\mathrm{out}}\oplus(\mathbf 1-P_{\mathrm{out}})\) after unitary
evolution.  A Cauchy family of finite channel amplitudes has the stated
limit by completeness of \(\mathbb C\) (or of the finite channel operator
space).  The additional locality, spectrum, stable-mass-shell,
asymptotic-field, pole, residue, amputation, and limit-interchange
hypotheses display the standard Haag--Ruelle/LSZ theorem domain specialized
to these channels; they do not follow from amplitude convergence or
wave-packet separation.  Applying
LSZ only after those hypotheses are verified proves the final conditional
statement.
\end{proof}

\begin{theorem}[External fermion line and generated closed completion]
\label{p4:thm:external-fermion-line-generated-closure}
Let \(H\) be a finite generated fermion carrier and let \(D:H\to H\) be its
generated kinetic operator.  For an external field
\(\psi\in H\) and dual field \(\bar\psi\in H^*\),
\[
 \boxed{S_F[\bar\psi,\psi;D]=\bar\psi(D\psi)\in\Rbdry.}
\]
For ordinary external vectors this is complex-valued.  In a fermionic
path-integral presentation the fields take values in a generated Grassmann
algebra and the quadratic expression lies in its even part; a numerical
scalar appears only after the licensed Berezin closure.

If the exterior-power constructor is generated and \(n=\dim H\), then
\[
 \bigwedge\nolimits^nD
 =\det(D)\,1_{\wedge^nH},
 \qquad
 \boxed{\det D
 =\operatorname{Tr}_{\wedge^nH}(\wedge^nD)\in\Rloop.}
\]
For \(D>0\),
\[
 \boxed{\log\det D=\operatorname{Tr}\log D.}
\]
For an antisymmetric Majorana form \(A\), a Pfaffian closure additionally
requires generated fermionic order/orientation data and obeys
\(\operatorname{Pf}(A)^2=\det A\).
\end{theorem}

\begin{proof}
The first formula is exactly the evaluation
\[
 H^*\otimes H\xrightarrow{\,1\otimes D\,}H^*\otimes H
 \xrightarrow{\operatorname{ev}}\mathbb C
\]
with the external field and dual-field ports retained, so the Paper-II-A
normal-form theorem places it in \(\Rbdry\).  Choose a basis
\(e_1,\ldots,e_n\).  By the definition of the top exterior power,
\[
 (\wedge^nD)(e_1\wedge\cdots\wedge e_n)
 =De_1\wedge\cdots\wedge De_n
 =\det(D)(e_1\wedge\cdots\wedge e_n).
\]
Because \(\wedge^nH\) is one-dimensional, its ordinary trace is the displayed
determinant.  This is a generated closed completion, not a retyping of the
external line.  If \(D>0\), diagonalization gives
\[
 \log\det D=\log\prod_i\lambda_i
 =\sum_i\log\lambda_i=\operatorname{Tr}\log D.
\]
The Pfaffian identity is the finite antisymmetric exterior closure; its sign
requires the stated order/orientation datum.
\end{proof}

\begin{remark}[Finite and continuum status of fermionic closures]
Determinant and Pfaffian outputs are exact finite closures, but their physical
use remains conditioned by the occurring field/process and Berezin
interfaces.  For nonpositive complex spectra, \(\log D\) needs a branch.  On
an infinite parent, log determinants require the displayed zeta/heat or other
regularization and the Euclidean-to-Hamiltonian reconstruction hypotheses.
Those data are not supplied by the finite closure theorem.
\end{remark}

\paragraph{Module 5: Euclidean kernels, finite Gaussian integrals, and OS status}
\label{p4:sec:module-finite-euclidean-path-integral}

\begin{theorem}[Finite reference-integration ray]
\label{p4:thm:finite-reference-integration-ray}
Let $V$ be an occurring finite-dimensional real bosonic mode space, regarded
with its generated additive-group structure.  Its translation-invariant Haar
measure is unique up to multiplication by a positive constant; hence the
configuration role canonically determines the positive integration ray
\[
 [d\mu_V].
\]
For an occurring compact group configuration factor $G$, normalized Haar
measure is unique.  Relation-free finite products carry the corresponding
product integration ray.  On the paired finite fermion owner, the generated
top-exterior/Berezin functional supplies the algebraic integration line.
Any additional configuration-dependent density, Jacobian, constraint factor,
or phase must be generated by the actual quotient/field-space/fermion owner;
it is not supplied by finite dimensionality alone.
\end{theorem}

\begin{proof}
Existence and uniqueness up to positive scale are the finite-dimensional Haar
theorem for the locally compact additive group $V$.  Compact Haar measure is
uniquely normalized by total mass one.  Product uniqueness follows from the
relation-free product grammar.  The fermionic statement is algebraic and is
proved below by the paired Berezin/top-exterior construction.  None of these
arguments constructs a measure on an infinite or gauge-redundant metric
configuration space.
\end{proof}

\begin{theorem}[Exponential Euclidean-weight rigidity]
\label{p4:thm:euclidean-weight-rigidity}
Let $w:\mathbb R\to\mathbb R_{>0}$ be a continuous scalar weight assigned to
a dimensionless Euclidean action coordinate $\sigma$.  Suppose
\[
 w(0)=1,
 \qquad
 w(\sigma_1+\sigma_2)=w(\sigma_1)w(\sigma_2)
\]
for all $\sigma_1,\sigma_2$ in a \emph{determining composition domain}:
either every real pair is actually composable, or there is a dense additive
subgroup $D\subset\mathbb R$ whose every pair is composable and on which the
printed law holds.  Then there is a unique $c\in\mathbb R$ with
\[
 \boxed{w(\sigma)=e^{c\sigma}}.
\]
On the generated Euclidean damping branch $c<0$.  If the already constituted
action unit is $\sigma=S_E/\hbar$ with unit damping rate, then
\[
 \boxed{w(S_E)=e^{-S_E/\hbar}}.
\]
If that action-rate calibration has not been generated, the theorem stops at
$e^{-\beta S_E}$ with one positive inverse-action scale $\beta$.
When the available relation-free compositions do \emph{not} determine such a
real or dense additive domain, only the character identities on those
composable pairs follow; no global exponential rigidity or independently
selected damping coefficient is inferred.
\end{theorem}

\begin{proof}
Positivity allows $f=\log w$.  On a full-real composition domain,
$f(x+y)=f(x)+f(y)$ for all reals, and continuity gives $f(x)=cx$.
Alternatively, on the stated dense additive subgroup $D$, $f|_D$ is
additive.  For any $x,y\in\mathbb R$, choose $d_n\to x$ and $e_n\to y$
in $D$; continuity passes $f(d_n+e_n)=f(d_n)+f(e_n)$ to $x,y$.
Thus the full-real continuous Cauchy equation holds, giving $f(x)=cx$,
with $c$ uniquely fixed by $w$.  Damping of nonzero positive actions fixes
$c<0$ only on the admitted monotone damping branch.  The action-unit
calibration is a separate input; neither it nor dense actual-action
composability follows from continuity of $w$ alone.
\end{proof}

\begin{definition}[Generated finite Euclidean quantum ray]
\label{p4:def:generated-finite-euclidean-quantum-ray}
Whenever an occurring finite configuration owner supplies the reference
integration ray of Theorem~\ref{p4:thm:finite-reference-integration-ray}, an
integrable claim-complete Euclidean action $S_E$, and either the determining-composition and action-unit branches of
Theorem~\ref{p4:thm:euclidean-weight-rigidity} or an independently
\emph{generated} explicit weight $e^{-S_E/\hbar}$, define
\[
 \boxed{
 [\rho](A)
 :=\left[
   \int A\,e^{-S_E/\hbar}\,d\mu_0
 \right].}
\]
The brackets denote the positive functional ray: multiplying $d\mu_0$ by a
positive constant or shifting $S_E\mapsto S_E+C$ rescales the functional and
does not change normalized nongravitational expectations.  Claims that
consume an absolute vacuum/gravitational action do not quotient by that
constant and retain the corresponding zero-order datum.
\end{definition}

\begin{corollary}[Finite generated quantum-constructor closure]
\label{p4:cor:finite-generated-quantum-constructor}
On the stated finite positive/integrable branch, once the configuration
integration ray, Euclidean action, and action unit are generated, \emph{and}
the determining composition-domain certificate is supplied (or the exact
exponential weight itself is generated), the positive state/action ray of
Definition~\ref{p4:def:generated-finite-euclidean-quantum-ray} is a generated
descendant and requires no independent partition-weight selector.  Restricted
nondense composability by itself does not establish this uniqueness claim.  This statement does not construct a gravitational path measure or
a global chiral determinant phase.
\end{corollary}

\begin{corollary}[Relative Gaussian and paired-Grassmann determinant weights]
\label{p4:cor:v2140-relative-determinant-weights}
Let $H_0,H_1$ be positive definite quadratic Hessians on the same occurring
finite real bosonic mode space $V$, evaluated with one representative of the
integration ray of Theorem~\ref{p4:thm:finite-reference-integration-ray}.
Then
\begin{equation}
\label{p4:eq:v2140-relative-boson-determinant}
 \boxed{
 \frac{\int_V e^{-\frac12\langle x,H_1x\rangle}\,d\mu_V(x)}
      {\int_V e^{-\frac12\langle x,H_0x\rangle}\,d\mu_V(x)}
 =\left(\frac{\det H_1}{\det H_0}\right)^{-1/2}.}
\end{equation}
For invertible $D_0,D_1$ on one paired finite fermion owner,
Proposition~\ref{p4:prop:paired-berezin-determinant-line} gives
\begin{equation}
\label{p4:eq:v2140-relative-grassmann-determinant}
 \boxed{
 \frac{\int[d\bar\psi\,d\psi]_{\rm paired}e^{-\bar\psi D_1\psi}}
      {\int[d\bar\psi\,d\psi]_{\rm paired}e^{-\bar\psi D_0\psi}}
 =\frac{\det D_1}{\det D_0}.}
\end{equation}
Hence a real bosonic quadratic block enters a relative effective action with
weight $+\tfrac12\Tr\log H$, while a paired complex Grassmann block enters
with weight $-\Tr\log D$.  A complex boson is two real bosonic blocks and
therefore has weight $+1$.  These determinant powers are fixed by the finite
integration owners and are not independent loop coefficients.
\end{corollary}

\begin{proof}
Choose an orthonormal basis diagonalizing $H_j$.  Each one-dimensional real
Gaussian contributes the common measure normalization times
$\sqrt{2\pi/\lambda}$, so the common normalization and $(2\pi)^{\dim V/2}$
cancel in the ratio and give
\eqref{p4:eq:v2140-relative-boson-determinant}.  The Grassmann identity is the
ratio of the two exact Berezin determinants from
Proposition~\ref{p4:prop:paired-berezin-determinant-line}.  Taking minus the
logarithm of the corresponding partition-weight ratios gives the stated
effective-action exponents.
\end{proof}

The Euclidean language used in spectral actions and path-integral notation can also be kept finite.  Let \(C_N\ge0\) be a positive covariance on a finite real mode space \(K_N\), and let \(J\in K_N\).  The normalized finite Gaussian generating functional is
\begin{equation}
\label{p4:eq:finite-gaussian-generating-functional}
Z_N[J]
=
\exp\!\left(\frac12\langle J,C_NJ\rangle\right).
\end{equation}
When \(C_N\) has a kernel, this formula is understood as the algebraic Gaussian generating functional on the support of \(C_N\), or equivalently after quotienting null directions. A nondegenerate finite Lebesgue Gaussian integral requires \(C_N>0\) on the integrated mode space.
For fermionic Grassmann modes with antisymmetric quadratic form \(D_N\), the
corresponding finite Berezin integral gives the exact finite
Pfaffian/determinant closures of
Theorem~\ref{p4:thm:external-fermion-line-generated-closure}.  A Pfaffian
sign uses the generated fermionic order/orientation data; continuum log
determinants retain their regularization and reconstruction hypotheses.  Thus
the formal path integral is replaced by finite Gaussian linear algebra plus
finite perturbation theory from Module~2.


\begin{theorem}[Finite Euclidean kernels, perturbative coefficients, and integral boundary]
\label{p4:thm:finite-euclidean-path-integral-status}
Let \(C_N\ge0\) be a finite covariance and let \(Z_N[J]\) be
\eqref{p4:eq:finite-gaussian-generating-functional}.  Then all Euclidean
\(n\)-point functions obtained by differentiating \(Z_N[J]\) are finite and
are determined by \(C_N\).  If the covariance lies in the bosonic/even or
graded-fermionic reflection-positive domain of
Theorem~\ref{p4:thm:finite-os-positivity}, the corresponding Euclidean
correlators have the cyclic finite reconstruction of
Theorem~\ref{p4:thm:finite-euclidean-hamiltonian-reconstruction}.
For every polynomial interaction \(V_N\), expansion in its coupling constants
has finite Gaussian coefficients at each fixed perturbative order.

An exact bosonic Lebesgue integral
\[
Z_{N,V}[J]
:=
\int_{\operatorname{supp}C_N}
\exp\!\left(
-\frac12\langle x,C_N^{-1}x\rangle
-V_N(x)+\langle J,x\rangle
\right)\,dx
\]
is asserted only when its absolute integrability is verified.  A sufficient
condition is the coercive bound
\[
\Re\!\left(\frac12\langle x,C_N^{-1}x\rangle+V_N(x)\right)
\ge a\|x\|^p-b
\]
for some \(a>0\), \(p>1\), and \(b\in\mathbb R\) outside a compact set.
Finiteness of
the mode space alone does not imply this condition.  Fermionic Berezin
integrals remain finite algebraic expressions because Grassmann degree is
finite.
\end{theorem}


\begin{proof}
The Gaussian derivative identities follow from differentiating the finite
exponential \eqref{p4:eq:finite-gaussian-generating-functional}.  OS
positivity is the finite positivity condition used by the reconstruction
theorem below.  At each fixed perturbative order, Wick differentiation of a
polynomial interaction produces finitely many Gaussian moments, hence a
finite coefficient.  For the exact bosonic integral, the displayed coercive
bound dominates the absolute value of the integrand by
\(\exp(b-a\|x\|^p+\|J\|\|x\|)\), which is integrable in finite dimension.
Without an integrability condition the conclusion is false; for example,
\(V_N(x)=-x^4\) makes the real integral diverge.  A Berezin integral instead
terminates algebraically because every Grassmann generator is nilpotent.
\end{proof}

\begin{remark}[Finite interaction versus geometry-selected vertex]
\label{p4:rem:finite-interaction-versus-geometric-vertex}
The theorem above proves finite quantum admissibility: once a finite carrier and a
self-adjoint finite interaction polynomial are specified, perturbation theory is
finite algebra.  A stronger statement is needed to call such a vertex
intrinsic to the constituted Paper-I-A round geometry.  In that case its
tensor coefficient must descend
from the Paper-I-A scalar/vector/spinor triple-overlap recoupling rules, or from a
separately declared exceptional-core algebra such as the
\(\mathsf{Supp}_{2T}\)-relative Hurwitz/trimer modules.
Generic finite coefficients are allowed effective quantum data, but they are not
automatically geometry-selected vertices.
\end{remark}


\paragraph{Hurwitz-link finite internal support import}
\label{p4:sec:hurwitz-link-finite-internal-support}

The upstream modules show that finite mode spaces can support finite fields,
interactions, kernels, and channel amplitudes once the relevant carrier and
the corresponding one-particle, field, interaction, state, and process
theorem-domain data have been supplied.
Papers~V-A--V-B \cite{Kugelmann2026f,Kugelmann2026g} use a more rigid source for
some of those carriers: the Paper-I-A Hurwitz-link module on an actualized or
conditionally generated core-support role \(\mathsf{Supp}_{2T}\).
Paper~IV does not derive that occurrence from the level-one carrier and does
not identify the physical Standard-Model sectors.  Its role is to record
that, once the unmarked core role occurs, Paper~I-A generates its complete
unmarked Hurwitz-link and trimer descendants and the Hilbert/Fock, transport,
and gauge-projection theorems of the constituted quantum roles apply to them.


\begin{theorem}[\(\mathsf{Supp}_{2T}\)-relative Hurwitz supports as candidate finite quantum carriers]
\label{p4:thm:hurwitz-link-finite-support-import}
Assume an occurrence of the unmarked Paper-I-A core-support role
\(\mathsf{Supp}_{2T}\), either directly under G0 or as the generated output
of the Paper-II-A point-support-to-core bridge in
Corollary~\ref{p2:cor:core-support-inheritance-no-duplicate-witness}.
This antecedent is not selected by \(m_p=1\).  For calculation, choose any
\(\sigma=(H,v,\delta)\) in the complete representative orbit, with
\(H\in[2T]\), \(v\in H\), and \(\delta=(u_+,u_-)\) an oriented body
diagonal of the Hurwitz cube link \(N(v)\).  No member of this orbit is
physically selected by the unmarked role.
On this domain, Proposition~\ref{p1:prop:hurwitz-link-cube} and
Theorem~\ref{p1:thm:hurwitz-walsh-clifford} give
\[
N(v)\cong(\ZZ_2)^3,
\qquad
(\RR^{N(v)},\star_{\sigma_-})\cong\Cl(0,3),
\]
so the full and scalar-reduced internal carriers have dimensions
\[
8,
\qquad
7.
\]
The chosen calculation representative and
Theorem~\ref{p1:thm:hurwitz-diagonal-trimer} give the decomposition
\[
N(v)=\{u_+,u_-\}\sqcup T_+\sqcup T_-,
\qquad
|T_+|=|T_-|=3,
\]
and the two distinct generated descendants
\[
  V(C_6^{\rm link})=T_+\sqcup T_-,
  \qquad
  G_\delta:=\pi^{-1}(\operatorname{Stab}^{+}(\delta))
      \cong C_6^{\rm grp}.
\]
Here \(\pi:N_{SU(2)}(H)\cong2O\to S_4\) is the projective quotient,
\(C_6^{\rm link}\) is the induced six-cycle graph, and
\(C_6^{\rm grp}\) is the cyclic binary lift of the oriented diagonal
stabilizer.
Thus \(T_+\sqcup T_-\) is the vertex set of the native axis-removed
staggered antiprism.  The finite complex supports
\[
C_{\rm HLC}:=\CC^{T_+},
\qquad
F_{\rm HLC}:=\CC^{T_-}
\]
are therefore legitimate candidate internal mode carriers for the finite
quantum calculus of Paper~IV.  If a later sector supplies a one-particle
assignment for either support, then
\[
\End(C_{\rm HLC})\cong C_{\rm HLC}\otimes C_{\rm HLC}^*,
\qquad
\dim \End(C_{\rm HLC})=9,
\]
and all projector, statistics-selected Fock, mixing, transport,
block-elimination, and gauge-projection constructions of Paper~IV apply to
that assigned support on their separately stated theorem domains, without a
further quantum-representation axiom.  The one-particle assignment and the
\(\mathsf{Supp}_{2T}\) occurrence antecedent remain explicit; the choice of
\(\sigma\) does not.
\end{theorem}


\begin{proof}
On the declared core-support domain, Paper~I-A,
Proposition~\ref{p1:prop:hurwitz-link-cube} and
Theorems~\ref{p1:thm:hurwitz-walsh-clifford}
and~\ref{p1:thm:hurwitz-diagonal-trimer}, prove the cube-link,
Walsh--Clifford, and antiprism-trimer identities for every representative.
Corollary~\ref{p1:cor:generated-c6-stabilizer-link-bridge} proves that the
complete \(C_6^{\rm grp}\) stabilizer and \(C_6^{\rm link}\) graph families
are equivariant descendants of the unmarked core.  The resulting finite
supports become finite-dimensional Hilbert spaces after complexification and
the supplied counting inner product.  Their generated finite operator
algebras carry the Part-I physical state/effect/event calculus immediately.
Once token-level \(\mathsf{PROB}\), \(\mathsf{SUB}\), and \(\mathsf{PROC}\)
supply the relevant preparations/readouts, occurrences, and
actual maps, the tensor and CP-map representations apply without requiring a
full-envelope premise.
After a one-particle assignment, a statistics selector and, for bosons,
\(N_{\max}\) have been supplied,
Theorem~\ref{p4:thm:finite-second-quantization} gives the corresponding Fock
sectors.  The transported-sector theorems of
Section~\ref{p4:sec:part2-canonical-transported-blocks} apply to a supplied
transported rank-three projector built from \(C_{\rm HLC}\) or
\(F_{\rm HLC}\), and
Theorem~\ref{p4:thm:finite-gauge-constraint-closure} applies to a supplied
compact unitary gauge action.  Thus \(\mathsf{Supp}_{2T}\) supplies the
unmarked finite support and its generated descendants; Paper~IV supplies the
conditional finite quantum representations acting on them.  Equivariance
makes the calculation independent of the chosen \(\sigma\).
\end{proof}


\begin{corollary}[Quantum state, occupation and Gauss transport of the sampled presentation]
\label{p4:cor:v204-quantum-sampled-transport}
Use Paper~II-C's realized $T:W\to V$ on any existing finite quantum carrier.
States and effects transport by $\rho\mapsto T\rho T^*$ and
$E\mapsto TET^*$, preserving $\operatorname{Tr}(\rho E)$ and positivity.
For an already assigned one-particle role and the existing statistics and
occupation envelope of Theorem~\ref{p4:thm:finite-second-quantization},
\[
 \Gamma_{\le N}(T)=\bigoplus_{n=0}^N T^{\otimes n}|_{\mathrm{Sym/Alt}}
\]
is unitary onto the corresponding image sectors and
\[
 \Gamma(T)d\Gamma(A)\Gamma(T)^*=d\Gamma(TAT^*).
\]
If $P_G$ is the existing Haar Gauss projector, the transported gauge action
has projector $\widetilde P_G=\Gamma(T)P_G\Gamma(T)^*$ and the same
physical invariant subspace. Existing species projections and their sum,
constraint commutators, spectra and total-generator identities are preserved.
\end{corollary}
\begin{proof}
Apply Paper~1B, Corollary~\ref{p1b:cor:tensor-transport}
\cite{Kugelmann2026Paper1B}, to Paper~II-C's realized isometry and the
statistics and finite envelope already supplied by
Theorem~\ref{p4:thm:finite-second-quantization}. On every tensor sector,
$d\Gamma(A)$ is the sum of the transported one-body operators, giving
the displayed identity. The same corollary transports the Haar projector
and its invariant subspace. States, effects and all units are understood
on the sampled image, whose ambient identity is $TT^*$ (and its tensor
counterparts). This is the existing quantum construction in its new
presentation.
\end{proof}
The transport does not assign a one-particle interpretation to a carrier
which did not already have one, or alter its statistics or occupation bound.
Once assigned by the existing framework, these are inherited data rather
than new assumptions introduced by sampling.

\begin{remark}[No scalene carrier triangle]
\label{p4:rem:no-scalene-carrier-triangle}
The Paper-I-A/Paper-IV import does not use a Euclidean triangle with side lengths \(7,8,9\).  The numbers \(7,8,9\) enter as finite carrier dimensions: reduced Walsh--Clifford septet, full Walsh--Clifford octet, and rank-three pair carrier \(\dim\End(C_{\rm HLC})\).  The only triangles used here are the Hurwitz-link trimers \(T_+\) and \(T_-\) inside the axis-removed antiprism.
\end{remark}


\begin{corollary}[Imported finite operator algebras on level-one and core supports]
\label{p4:cor:hurwitz-finite-operator-algebras}
On the realized level-one carrier, and on a separately actualized or
conditionally generated \(\mathsf{Supp}_{2T}\) domain where item~(ii) is
invoked, Paper~I-A
supplies the following finite operator algebras as mathematical carriers for
the present calculus:
\begin{enumerate}[label=(\roman*), leftmargin=2em]
\item the protected block \(\mathcal K_{\mathsf B}\), whose full operator algebra has the Clifford--Fierz basis
\[
\End(\mathcal K_{\mathsf B})
=
\langle 1,\Gamma_a,\Gamma_{ab},\Gamma_5\Gamma_a,\Gamma_5\rangle,
\]
as in Proposition~\ref{p1:prop:math-protected-clifford-fierz};
\item on the \(\mathsf{Supp}_{2T}\) domain, each representative oriented
trimer \(T\),
whose rank-three operator algebra is the Weyl algebra
\[
M_3(\CC)=\operatorname{span}_{\CC}\{X^aZ^b:a,b\in\ZZ_3\},
\qquad
ZX=\omega XZ,
\]
as in Proposition~\ref{p1:prop:math-trimer-weyl-a2};
\item the derived/co-realized carrier \(\mathcal H_1\), whose coordinate-coupling graph is the projected graph of Proposition~\ref{p1:prop:math-h1-coordinate-coupling-graph} and Paper~II, Proposition~\ref{p2:prop:h1-coordinate-coupling-graph}.
\end{enumerate}
Consequently these finite algebras may be used as one-particle mode spaces,
accessible subalgebras, transported sectors, or sector projectors after the
displayed support and physical-role selectors have been supplied.
\end{corollary}


\begin{proof}
Each item is a finite-dimensional complex Hilbert carrier or operator algebra
supplied upstream by Paper~I-A or Paper~II; item~(ii) additionally uses the
actualized or conditionally generated \(\mathsf{Supp}_{2T}\) role.
Part~I proves the finite state/effect, tensor, and CP calculus for every
generated finite direct-sum algebra.  It also proves separately that a
carrier-resolving \(\mathsf{ENV}\) extension, when supplied, forces
\(\End(\mathcal H)\).  Sections~\ref{p4:sec:part2-qft-bridge}--
\ref{p4:sec:part2-internal-thresholds} prove the corresponding conditional
Fock, mixing, block-elimination, and sector-threshold constructions.
Thus the present corollary imports only finite mathematical supports and applies the already-proved finite quantum calculus to them.
\end{proof}

\paragraph{Module 6: finite gauge constraints and internal connections}
\label{p4:sec:module-finite-gauge-constraints}

Gauge language enters Paper~IV only as finite sector bookkeeping.  The actual
Standard-Model algebra and representations are selected in Papers~V-A--V-B; on
the \(\mathsf{Supp}_{2T}\) domain, a representative Hurwitz-link trimer \(T_+\)
supplies at most a finite rank-three support on which their later
gauge-closure argument can act.  Nevertheless, once a compact unitary action
has itself been supplied, its finite gauge constraint has theorem-level
status.  Put \(A_N:=\End(\mathcal H_N)\) for the resolved finite observable
algebra on this theorem domain.  Let a compact group \(G\) act through a
finite-dimensional unitary representation \(U(g)\) on \(\mathcal H_N\);
finite groups are included as the special case in which Haar averaging is a
normalized finite sum.  Define the physical projector by Haar averaging
\begin{equation}
\label{p4:eq:finite-gauge-projector}
P_G=\int_G U(g)\,dg,
\end{equation}
with the normalized Haar measure; for finite groups this is the normalized sum over \(G\).  The gauge-invariant observable algebra is
\begin{equation}
\label{p4:eq:finite-gauge-invariant-algebra}
A_N^G=\{A\in A_N: U(g)AU(g)^*=A\text{ for all }g\in G\}.
\end{equation}


\begin{theorem}[Finite gauge-constraint closure]
\label{p4:thm:finite-gauge-constraint-closure}
For any finite-dimensional unitary representation \(U:G\to\mathcal U(\mathcal H_N)\) of a compact group \(G\):
\begin{enumerate}
\item \(P_G\) is an orthogonal projector onto the invariant physical subspace \(\mathcal H_N^G\).
\item \(A_N^G\) is a finite \(*\)-subalgebra, and \(A\in A_N^G\) preserves \(\mathcal H_N^G\).
\item If \(H_N\in A_N^G\), then \(e^{-itH_N}\mathcal H_N^G\subseteq\mathcal H_N^G\) and gauge constraints are stable under finite dynamics.
\item If \(\delta_\theta(A)=i[Q_\theta,A]\) is the infinitesimal action of a charge generator, then invariant states satisfy the finite Ward identity \(\omega(\delta_\theta(A))=0\) whenever \([\rho,Q_\theta]=0\).
\end{enumerate}
Thus gauge projection, Gauss constraints, and Ward identities are finite algebraic consequences of a unitary internal representation, not additional continuum-field postulates.
\end{theorem}


\begin{proof}
The Haar average is self-adjoint and idempotent by invariance of Haar measure, so it is the orthogonal projector onto the invariant vectors.  The commutant of a unitary group action is closed under products and adjoints, hence is a finite \(*\)-subalgebra.  If \(H_N\) commutes with the action, so does its exponential, and the invariant subspace is preserved.  Finally, \(\omega(\delta_\theta(A))=i\Tr(\rho[Q_\theta,A])=i\Tr([\rho,Q_\theta]A)=0\) by cyclicity of the trace.
\end{proof}


\begin{remark}[Finite gauge closure and the local Noether-II packet]
\label{p4:rem:v176-finite-gauge-noether-typing}
The theorem above is the finite representation-theoretic counterpart of the
local gauge constraint packet derived later from the common
\(A+\Phi+g\) action in Paper~III-F.  The two are not identified by notation
alone: a representation/compression or reconstruction interface is required
to transport a field-theoretic moment map into a particular finite carrier.
\end{remark}


\paragraph{Module 7: compactness correction to bosonic CCR}
\label{p4:sec:module-ccr-defect}

The exact finite correction to bosonic CCR is important because it prevents the framework from making an inconsistent finite-dimensional postulate.  For one truncated bosonic oscillator with basis \(\{|0\rangle,\ldots,|n_{\max}\rangle\}\), define
\[
a|n\rangle=\sqrt n\,|n-1\rangle,
\qquad
a^\dagger|n\rangle=\begin{cases}
\sqrt{n+1}\,|n+1\rangle,&n<n_{\max},\\
0,&n=n_{\max}.
\end{cases}
\]


\begin{proposition}[Explicit finite CCR boundary defect]
\label{p4:prop:explicit-finite-ccr-boundary-defect}
On the truncated oscillator above,
\begin{equation}
\label{p4:eq:explicit-finite-ccr-defect}
[a,a^\dagger]
=
\mathbf 1-(n_{\max}+1)|n_{\max}\rangle\langle n_{\max}|.
\end{equation}
For a finite set of modes with occupation cutoffs, the commutator has the form
\begin{equation}
\label{p4:eq:multi-mode-ccr-defect}
[a_i,a_j^\dagger]=\delta_{ij}\mathbf 1-B_{ij,N},
\end{equation}
where \(B_{ij,N}\) is supported on the occupation boundary of the truncated carrier.  Hence exact global bosonic CCR are impossible at finite dimension, but on states \(\rho\) with negligible boundary weight,
\begin{equation}
\label{p4:eq:ccr-defect-expectation-bound}
\bigl|\Tr(\rho([a_i,a_j^\dagger]-\delta_{ij}\mathbf 1))\bigr|
\le \|B_{ij,N}\|\,\Tr(\rho P_{\partial N}).
\end{equation}
\end{proposition}


\begin{proof}
For \(n<n_{\max}\), \(aa^\dagger|n\rangle=(n+1)|n\rangle\) and \(a^\dagger a|n\rangle=n|n\rangle\), so the commutator is \(|n\rangle\).  At the top state, \(aa^\dagger|n_{\max}\rangle=0\) and \(a^\dagger a|n_{\max}\rangle=n_{\max}|n_{\max}\rangle\), so the commutator eigenvalue is \(-n_{\max}\), which is exactly \(1-(n_{\max}+1)\).  The multimode statement is the same calculation on each occupation boundary, and the expectation bound is \(|\Tr(\rho B)|\le\|B\|\Tr(\rho P_{\partial N})\) for boundary-supported \(B\).
\end{proof}

\begin{remark}[Finite-domain meaning of the CCR defect]
\label{p4:rem:ccr-defect-meaning}
The spatial unfolding label \(m_x\) and the bosonic occupation cutoff
\(N_{\max}\) control different structures.  Increasing \(m_x\) enlarges the
reciprocal-position window or the number of resolved spatial modes.  It does
not change the one-mode identity
\[
[a,a^\dagger]
=
\mathbf1-(N_{\max}+1)|N_{\max}\rangle\langle N_{\max}|
\]
and therefore does not by itself remove the occupation-boundary defect.
For a fixed family of observables, the defect becomes negligible only after
a separate estimate such as
\[
\|B_{ij,N_{\max}}\|\,
\Tr(\rho_{N_{\max}}P_{\partial N_{\max}})
\longrightarrow0,
\]
or after a degree-controlled occupation-cutoff limit
\(N_{\max}\to\infty\).  A large-\(m_x\) spatial limit and an
\(N_{\max}\)-limit may be taken together, but neither implies the other.
Later scalaron, perturbation, and plasma oscillator presentations therefore
require both their spatial-window hypothesis and an independent
low-occupation or occupation-tail estimate.
\end{remark}

\paragraph{Module 8: finite Wick and quasi-free calculus}
\label{p4:sec:module-finite-wick-calculus}

The Wick/Feynman calculus also has a native finite form.  For fermionic CAR generators \(c_i,c_i^\dagger\), a gauge-invariant quasi-free state is determined by its two-point covariance
\[
C_{ij}=\omega(c_i^\dagger c_j),
\qquad 0\le C\le\mathbf 1.
\]
For bosonic truncated sectors, Wick factorization is exact before compression in the regular oscillator algebra for regular quasi-free states.  After finite occupation compression it remains valid only under the corresponding compressed-quasi-free / low-boundary-weight hypothesis, up to the boundary defect of Module~7.


\begin{theorem}[Finite Wick/quasi-free calculus]
\label{p4:thm:finite-wick-quasifree-calculus}
On a finite CAR algebra, every gauge-invariant quasi-free state with covariance \(C\) satisfies the exact determinant formula
\begin{equation}
\label{p4:eq:finite-fermion-wick-determinant}
\omega\!\left(c_{i_1}^\dagger\cdots c_{i_k}^\dagger c_{j_k}\cdots c_{j_1}\right)
=
\det\bigl(C_{i_pj_q}\bigr)_{p,q=1}^k,
\end{equation}
and odd correlation functions vanish.

For the bosonic branch, let \(\omega\) be a regular quasi-free state on the
untruncated Fock representation over a finite mode space.  Let \(Q_N\) be the
projection onto total occupation at most \(N\), put
\[
p_N:=\omega(Q_N)>0,
\qquad
\omega_N(X):=p_N^{-1}\omega(Q_NXQ_N),
\]
and let \(a_N^\#=Q_Na^\#Q_N\).  For a word \(W\) of degree \(r\), let \(W_N\)
be the same word in the truncated operators and define
\[
R_{N,r}:=Q_N-Q_{N-r},
\qquad
D_{N,W}:=W_N-Q_NWQ_N,
\]
with \(Q_{N-r}=0\) when \(N<r\).  Then \(D_{N,W}Q_{N-r}=0\), and
\begin{equation}
\label{p4:eq:finite-boson-wick-boundary-error}
\begin{aligned}
|\omega_N(W_N)-\omega_{\mathrm{Wick}}(W)|
\le{}&
\frac{\|D_{N,W}\|}{p_N}\sqrt{\omega(R_{N,r})}\\
&+
\frac{\sqrt{\omega(\mathbf1-Q_N)}}{p_N}
\left(
\sqrt{\omega(W^*W)}
+
\sqrt{\omega(Q_NWW^*Q_N)}
\right)\\
&+
\frac{1-p_N}{p_N}\,|\omega(W)|.
\end{aligned}
\end{equation}
Here \(\omega_{\mathrm{Wick}}(W)=\omega(W)\).  Consequently the bosonic Wick
moments are recovered only when the degree-\(r\) boundary-layer term and the
discarded-tail terms, including their cutoff-dependent norm factors, tend
to zero.  Exact agreement without a limit holds for the corresponding
low-sector matrix elements; equality of compressed-state expectations also
requires absence of discarded state weight.
\end{theorem}


\begin{proof}
The CAR determinant formula is the defining finite quasi-free property,
equivalently obtained by diagonalizing \(C\) and multiplying the one-mode
generating functions.  Odd functions vanish by gauge invariance.

For a bosonic word of degree \(r\), no intermediate creation step can leave
the cutoff when the initial vector lies below occupation \(N-r\).  Hence
\(D_{N,W}Q_{N-r}=0\), and state Cauchy--Schwarz gives
\[
|\omega_N(D_{N,W})|
\le
\frac{\|D_{N,W}\|}{p_N}\sqrt{\omega(R_{N,r})}.
\]
Moreover,
\[
W-Q_NWQ_N=(\mathbf1-Q_N)W+Q_NW(\mathbf1-Q_N).
\]
Applying Cauchy--Schwarz to these two terms and then accounting for the
normalization \(p_N^{-1}\) gives the remaining terms in
\eqref{p4:eq:finite-boson-wick-boundary-error}.

A top-level projector alone cannot control the error.  For example, a
one-mode squeezed vacuum has support only in even occupation sectors.  With
an odd cutoff \(N\), its top-state weight is zero while its discarded even
tail is nonzero.  The former top-boundary-only bound would therefore have
zero right-hand side and a nonzero left-hand side.
\end{proof}

\begin{remark}[Wick calculus support]
\label{p4:rem:finite-wick-downstream-support}
\leavevmode\par\noindent
Module~8 supplies the finite theorem-level status of propagator contractions,
Gaussian integrals, determinant and Pfaffian factors, and leading perturbative
kernels.  Papers~IV--V-C may therefore use Wick-expanded two-point and loop
language as finite-mode bookkeeping, while any claim about an exact interacting
continuum constructive limit remains outside Paper~IV.
\end{remark}



% BEGIN INLINE SOURCE: Paper4-CARSourceVariance-Main-v1_0.tex
% Framework v2.164: exact bilinear variance beside the existing Wick calculus.
\paragraph{Bilinear source variance and its non-Wick return}
\label{p4:sec:car-source-variance}
\begin{theorem}[Exact fixed-source CAR variance and sharp Wick bound]
\label{p4:thm:car-source-variance}
In a fixed canonical finite CAR chart let $O=d\Gamma(S)+cI$,
$S=S^\dagger$ numerical, and $\gamma_{ij}=\omega(c_j^\dagger c_i)$.
Thus $\gamma$ is the transpose of the covariance displayed in
Theorem~\ref{p4:thm:finite-wick-quasifree-calculus}, and $0\le\gamma\le I$.
Set
\[
 \kappa_{ij;kl}=\omega(c_i^\dagger c_jc_k^\dagger c_l)
 -\gamma_{ji}\gamma_{lk}-\gamma_{li}(\delta_{jk}-\gamma_{jk}).
\]
The connected variance satisfies
\[
 \boxed{\operatorname{Var}_\omega(O)=W+\Xi,\quad
 W=\operatorname{Tr}[\gamma S(I-\gamma)S],\quad
 \Xi=\sum_{i,j,k,l}S_{ij}S_{kl}\kappa_{ij;kl},}
\]
and
\[
 \boxed{0\le W\le\tfrac12\operatorname{Tr}S^2
                 \le\tfrac12d_f\|S\|^2.}
\]
On the number-invariant quasi-free branch $\kappa=0$.
For a state with anomalous pairings, include their contractions before
calling the remainder a connected four-point cumulant.
\end{theorem}
\begin{proof}
Expand the variance and subtract the disconnected two-point product.
The stated CAR contraction sums to $W$, and its remainder sums to $\Xi$.
Positivity of $W$ follows from the squared Hilbert--Schmidt norm of
$(I-\gamma)^{1/2}S\gamma^{1/2}$. In a diagonal covariance basis,
Hermiticity symmetrizes the coefficient of $|S_{ij}|^2$ to
$[c_i(1-c_j)+c_j(1-c_i)]/2\le1/2$ for $0\le c_i,c_j\le1$.
This proves the bound, which an off-diagonal unit source on a two-mode
projector covariance saturates. The Wick statement follows from the
already owned finite quasi-free determinant formula.
\end{proof}
The numerical, fixed-frame source is an essential hypothesis.
Noncentral operator coefficients use the complete joint products, and
source-dependent kinetic frames retain their derivatives.
If $O$ commutes with the physical singlet projector, raw CAR products
can evaluate moments on singlet-supported states. Otherwise retain
every intermediate projector. Projection or reduction can create a
non-Wick return in the consumed state; quasi-free input does not guarantee
quasi-free output.

\begin{corollary}[Bounded-source variance route]
\label{p4:cor:bounded-source-connected-return}
If $N_L\to\infty$, $d_{f,L}/N_L$ and $\|S_L\|$ are bounded, then
$W_L=O(N_L)$ and
\[
 \liminf_L\operatorname{Var}(O_L)/N_L^2>0
 \quad\Longleftrightarrow\quad
 \liminf_L\Xi_L/N_L^2>0.
\]
\end{corollary}
\begin{proof}
The bound makes $W_L/N_L^2\to0$; use the exact decomposition.
\end{proof}
Bounded mode density and source norm must be proved on the consumed
refinement/readout branch. Compact momentum alone supplies neither
estimate. This corollary concerns the quadratic-variance certificate,
not every possible symmetry-breaking mechanism.

% END INLINE SOURCE: Paper4-CARSourceVariance-Main-v1_0.tex




\subsection{Release v4.100: exact finite graded effective-average-action response}
\label{p4:sec:v4100-exact-graded-eaa}

\begin{theorem}[Exact finite graded modified-Legendre flow]
\label{p4:thm:v4100-exact-graded-eaa}
Let the occurring finite carrier contain real bosonic coordinates $\phi$ and
paired complex Grassmann coordinates $(\bar\psi,\psi)$ with a convergent
bosonic measure and the finite Berezin measure already constituted in this
paper.  Add regulator terms
\[
 \Delta S_k
 =\frac12\phi^T R_{B,k}\phi+\bar\psi R_{F,k}\psi
\]
and compatible bosonic/Grassmann sources.  On a domain where the modified
graded Legendre map is nonsingular, the exact finite effective-average action
obeys
\begin{equation}
\label{p4:eq:v4100-exact-graded-eaa}
 \boxed{
 \partial_t\Gamma_k
 =\frac12\operatorname{STr}\!\left[
   (\Gamma_k^{(2)}+\mathcal R_k)^{-1}\partial_t\mathcal R_k
 \right],\qquad t=\log k .}
\end{equation}
The supertrace convention is fixed so that one real bosonic block contributes
$+\tfrac12\Tr(G_B\dot R_B)$, one paired complex Grassmann block contributes
$-\Tr(G_F\dot R_F)$, and a paired complex ghost/Jacobian block has the same
negative determinant sign.  Mixed blocks are handled by the inverse graded
Hessian, not by adding independent loop weights.
\end{theorem}

\begin{proof}
Differentiate the finite logarithmic generating integral at fixed mean
superfield.  The source derivatives cancel against those of the graded
Legendre transform.  The remaining bosonic covariance is the inverse
regulated bosonic Hessian, while the paired Berezin derivative gives the
negative fermionic/ghost covariance with the determinant powers already fixed
by the finite Gaussian/Berezin integrals of this paper.  Collecting the blocks
in the doubled superfield convention gives
\eqref{p4:eq:v4100-exact-graded-eaa}.  No loop expansion has been used.
\end{proof}

\begin{remark}[Exact finite identity versus continuum and chiral measure]
\label{p4:rem:v4100-graded-eaa-scope}
Equation~\eqref{p4:eq:v4100-exact-graded-eaa} is exact on the admitted finite
paired branch.  It does not by itself prove an infinite-dimensional
interacting continuum limit, a globally well-defined chiral determinant
phase, BRST/background identities for an arbitrary regulator, or a pole
certificate.  Those remain separately typed analytic/measure obligations.
\end{remark}
% END EXACT-GRADED-EAA v4.100


% BEGIN GENERATED-FINITE-PAIR-CELL-HANDOFF v4.101


\section{Scattering, finite-temperature and kinetic source interfaces}
\label{p4b:sec:v238-b4}
\noindent\textit{Source and claim boundary.} Finite scattering, resonance, damping and thermal-kinetic representations have their distinct channel, scale, state, Markov and limiting hypotheses. The complete action does not automatically yield an asymptotic scattering operator or thermal history.

\subsection{Scattering, resonances, and decay channels}
\label{p4:sec:part2-scattering}

Scattering requires a more careful framework translation than bound-state spectroscopy.
On the microscopic compact momentum side the series-derived compact sector is minimal, $m_p=1$, so there is no literal infinite asymptotic momentum continuum.
The standard scattering picture therefore belongs to the large reciprocal-position-side unfolding, where one can distinguish incoming, interaction, and outgoing channel sectors over long but finite flight times before finite-size recurrences set in.
Accordingly, the framework gives an exact \emph{finite-window scattering theory}; the usual infinite-volume $S$-matrix language is recovered only as an effective idealization of that large-\(m_x\) position-side regime.



\paragraph{Finite-window channel scattering}
\label{p4:sec:part2-finite-window-scattering}

Let $\mathcal K$ be a large reciprocal-position sector and suppose
\[
\mathcal K = \mathcal K_{\mathrm{in}} \oplus \mathcal K_{\mathrm{int}} \oplus \mathcal K_{\mathrm{out}},
\]
where the interaction Hamiltonian is decomposed as
\[
H=H_0+V,
\]
with $H_0$ preserving the incoming and outgoing channel sectors and $V$ supported near the interaction region.
Choose a completely positive trace-preserving preparation map
$J_{\mathrm{in}}$ from incoming channel states into
\(\End(\mathcal K)\), and let $P_{\mathrm{out}}$ project onto the outgoing sector.
For a flight time $\tau$ after the packet has crossed the interaction zone but before revival from the finite boundaries, define the finite-window scattering map
\begin{equation}
\label{p4:eq:finite-window-scattering}
\mathcal S_{\tau}(\rho_{\mathrm{in}})
:=
P_{\mathrm{out}}U(\tau)J_{\mathrm{in}}(\rho_{\mathrm{in}})U(\tau)^*P_{\mathrm{out}},
\qquad
U(\tau)=e^{-i\tau H/\hbar}.
\end{equation}
If $\{Q_\beta\}$ is a channel POVM on $\mathcal K_{\mathrm{out}}$, the detected channel probabilities are
\begin{equation}
\label{p4:eq:scattering-probabilities}
p_\beta(\tau)=\Tr\!\bigl(Q_\beta\,\mathcal S_{\tau}(\rho_{\mathrm{in}})\bigr).
\end{equation}
When these probabilities stabilize on a plateau of times \(\tau_{\mathrm{sc}}\ll \tau \ll \tau_{\mathrm{rev}}\), they define the effective scattering data of the experiment.

\begin{proposition}[Finite-window channel instrument and first-order rate]
\label{p4:prop:framework-scattering-window}
In the finite Born framework, channel scattering is exactly a CP trace-nonincreasing operation of the form \eqref{p4:eq:finite-window-scattering}. Together with the complementary failure/no-exit operation it forms a CPTP instrument. Conditional on successful exit into the outgoing sector, the normalized state update is the usual nonlinear selective update $\rho\mapsto \Phi_\tau(\rho)/\Tr\Phi_\tau(\rho)$, not a linear CPTP channel.
For pure incoming and outgoing channel states $\ket{\alpha}$ and $\ket{\beta}$ of $H_0$, the first-order transition amplitude is
\begin{equation}
\label{p4:eq:first-born-amplitude}
\mathcal A^{(1)}_{\beta\leftarrow\alpha}(\tau)
=
-\frac{i}{\hbar}\int_0^{\tau}
e^{it(E_\beta-E_\alpha)/\hbar}\,
\bra{\beta}V\ket{\alpha}\,dt.
\end{equation}
If \(E_\beta\ne E_\alpha\), this amplitude is bounded and oscillatory:
\[
\left|\mathcal A^{(1)}_{\beta\leftarrow\alpha}(\tau)\right|
\le\frac{2|\langle\beta,V\alpha\rangle|}
{|E_\beta-E_\alpha|}.
\]
It need not converge to zero.  What vanishes is the time-normalized
off-shell amplitude and probability rate,
\(\mathcal A^{(1)}(\tau)/\tau\to0\) and
\(|\mathcal A^{(1)}(\tau)|^2/\tau\to0\).  The familiar on-shell delta
selection is a distributional large-window/continuum statement, not an
ordinary limit of the finite amplitude.
\end{proposition}

\begin{proof}
Complete positivity follows because \(J_{\rm in}\) is CP and
\eqref{p4:eq:finite-window-scattering} is obtained by unitary conjugation
and projection compression, both CP operations.  Trace nonincrease follows
from \(P_{\rm out}\le\mathbf1\) and trace preservation of the preparation and
unitary steps.  Replacing \(P_{\rm out}\) by
\(\mathbf1-P_{\rm out}\) supplies the complementary outcome; the two CP maps
sum to a trace-preserving instrument.
The first-order formula is the Dyson expansion of the interaction-picture propagator from Section~\ref{p4:sec:part2-timedep-perturbation} between free-channel eigenstates.
For \(\Delta E=E_\beta-E_\alpha\ne0\), direct integration gives
\[
\mathcal A^{(1)}_{\beta\leftarrow\alpha}(\tau)
=
\frac{\langle\beta,V\alpha\rangle}{\Delta E}
\bigl(1-e^{i\tau\Delta E/\hbar}\bigr),
\]
which proves the bound and the two normalized limits.
\end{proof}

\begin{remark}[Relation to textbook \texorpdfstring{$S$}{S}-matrix language]
\label{p4:rem:scattering-smatrix}
The usual $S$-matrix is obtained only after an additional asymptotic
construction in which the position-side flight region and boundary-return
time diverge and the free phases, local generators, channel embeddings, and
tails converge in the stated sense.  This is a large-\(m_x\)
reciprocal-position interface limit, not a limit in which the compact
momentum truncation flows away from \(m_p=1\).
The framework does not require that idealization at the fundamental level; all directly measurable scattering statements are already encoded by the finite-window channel probabilities \eqref{p4:eq:scattering-probabilities}.
\end{remark}

For the asymptotic bridge below the large-window parameter is the
reciprocal-position unfolding from Paper~III.  On the nondecreasing
\(\mathsf{SF}_{\Lambda_F}\) branch it obeys
Corollary~\ref{p4:cor:exact-reciprocal-position-resolution}; in particular,
for \(m_x\ge2\),
\[
\ell_F\sqrt{m_x(m_x+2)}
\le R_x
<\ell_F\sqrt{(m_x+1)(m_x+3)}.
\]
Outside that branch, the same scattering construction may use a separately
supplied compatible trajectory \(R_x(m_x)\to\infty\), but it may not infer
one from the \(\Dcl\) action profile alone.  The compact momentum carrier
remains fixed at \(m_p=1\).  Let
\[
H_m=H_{0,m}+V_m
\]
be a compatible directed family of finite channel Hamiltonians.  Assume that
\(V_m\) is uniformly bounded and supported in a fixed interaction region
\(B_{r_0}\), whose physical size does not grow with \(m_x\).  For an
open-channel energy window \(I\) separated from thresholds and from the
bound-sector spectrum, denote by \(\mathcal K^{\mathrm{open}}_{m,I}\) the
corresponding incoming packet core.  Assume separately that this core has a
positive outward group/escape-speed lower bound \(v_*(I)>0\).  Spectral
separation from threshold alone does not exclude a stationary point of the
dispersion.  A large-window scattering
bridge additionally requires the dynamical-locality estimate
\((\mathrm{DLOC}_{\rm pkt})\) stated below; a bound on the dispersion derivative
alone is not an exact finite-propagation theorem.

\begin{lemma}[Two-time return suppression under \((\mathrm{DLOC}_{\rm pkt})\)]
\label{p4:lem:derived-scattering-plateau}
Put \(U_m(t):=e^{-itH_m/\hbar}\) and
\(P_0:=P_{B_{r_0}}\).  Assume on the compatible open packet core the
dynamical-locality bound
\begin{equation}
\label{p4:eq:scattering-quasilocal-propagation}
\left\|
(\Id-P_{B_{r+c|t|+\ell}})U_m(t)P_{B_r}
\right\|
\le \varepsilon_m(\ell,t),
\end{equation}
and define its monotone tail envelope by
\[
\overline\varepsilon_m(\ell,T)
:=
\sup_{\substack{\ell'\ge\ell\\ |t|\le T}}
\varepsilon_m(\ell',t).
\]
Assume this envelope tends to zero in the stated joint
separation/refinement regime.  Let \(Q_{\rm far}^{(m)}\) be a supplied
far-boundary localization projector whose support has one-way separation
\[
d_m
:=
d\!\left(B_{r_0},\operatorname{supp}Q_{\rm far}^{(m)}\right)
=\pi R_x(m_x)-O(r_0)
\]
from the interaction region.  The nominal round-trip cone scale is
\[
\tau_{\rm cone}(m_x):=\frac{2d_m}{c}
=\frac{2\pi R_x(m_x)-O(r_0)}{c}
\longrightarrow\infty.
\]
This is a cone-traversal scale, not an exact first-return time.

For \(0\le s\le\tau\), define the sequential far-detection/return Kraus
operator
\begin{equation}
\label{p4:eq:two-time-return-kraus}
K_{\rm ret}^{(m)}(\tau,s)
:=
P_0U_m(\tau-s)Q_{\rm far}^{(m)}U_m(s)P_0 .
\end{equation}
Then \(\rho\mapsto K_{\rm ret}^{(m)}(\tau,s)\rho
K_{\rm ret}^{(m)}(\tau,s)^*\) is a CP trace-nonincreasing instrument branch.
Whenever \(c\tau/2<d_m\),
\begin{equation}
\label{p4:eq:two-time-return-bound}
\boxed{\;
\sup_{0\le s\le\tau}
\|K_{\rm ret}^{(m)}(\tau,s)\|
\le
\overline\varepsilon_m\!\left(d_m-\frac{c\tau}{2},
\frac{\tau}{2}\right).
\;}
\end{equation}
Consequently, for every normalized input state supported in \(P_0\),
\[
\Tr\!\left(
K_{\rm ret}^{(m)}(\tau,s)\rho
K_{\rm ret}^{(m)}(\tau,s)^*
\right)
\le
\overline\varepsilon_m\!\left(d_m-\frac{c\tau}{2},
\frac{\tau}{2}\right)^2.
\]
If \(\tau_{\rm sc}\ll\tau_m\) and
\(\tau_m/R_x(m_x)\to0\), the right-hand side vanishes uniformly in
\(s\).  Any finite detector-time instrument whose return branches use
\eqref{p4:eq:two-time-return-kraus} at
\(\{s_1,\ldots,s_{J_m}\}\) has total return probability at most
\[
J_m\,
\overline\varepsilon_m\!\left(d_m-\frac{c\tau_m}{2},
\frac{\tau_m}{2}\right)^2,
\]
and therefore also vanishes whenever that product tends to zero.  A
continuous-time detector requires its own detector ancilla or instrument
measure; it is not represented by an unrelated single-time projection.
\end{lemma}

\begin{proof}
The two legs of \eqref{p4:eq:two-time-return-kraus} satisfy
\[
\begin{aligned}
\|Q_{\rm far}^{(m)}U_m(s)P_0\|
&\le
\overline\varepsilon_m(d_m-cs,s),\\
\|P_0U_m(\tau-s)Q_{\rm far}^{(m)}\|
&=
\|Q_{\rm far}^{(m)}U_m(s-\tau)P_0\|\\
&\le
\overline\varepsilon_m(d_m-c(\tau-s),\tau-s),
\end{aligned}
\]
whenever the displayed residual separations are positive.  The equality in
the second line uses unitarity.  Since all other factors are contractions,
\[
\|K_{\rm ret}^{(m)}(\tau,s)\|
\le
\min\!\left\{
\|Q_{\rm far}^{(m)}U_m(s)P_0\|,
\|P_0U_m(\tau-s)Q_{\rm far}^{(m)}\|
\right\}.
\]
One of \(s\) and \(\tau-s\) is at most \(\tau/2\).  Monotonicity of the
envelope therefore gives \eqref{p4:eq:two-time-return-bound}.  A single Kraus
operator defines a CP trace-nonincreasing map, and
\(\Tr(K\rho K^*)\le\|K\|^2\Tr\rho\) proves the probability bound.
For a finite instrument, outcome probabilities add.

If \(\tau_m/R_x(m_x)\to0\), then
\(d_m-c\tau_m/2\to\infty\) on the refinement scale, so the assumed tail
envelope tends to zero.  Because the estimate has tails, it gives
suppression rather than an exact zero-return interval.  Nor is
\(\tau_{\rm cone}\) a lower bound on every spectral recurrence time: a finite
quasiperiodic state can recur without a packet circumnavigating the spatial
boundary.
\end{proof}


\begin{lemma}[Strong-limit criterion for finite-window scattering]
\label{p4:lem:finite-window-scattering-cauchy-criterion}
Let \(S_m(\tau_m):\mathcal D_I\to\mathcal H_{\rm out}\) be uniformly bounded
finite-window channel operators with one common complete target Hilbert
space.  The large-window
operator
\[
S_{\rm as}\psi=\lim_{m_x\to\infty}S_m(\tau_m)\psi
\]
exists strongly on \(\mathcal D_I\) if and only if the vectors
\(\{S_m(\tau_m)\psi\}_m\) are Cauchy for every \(\psi\in\mathcal D_I\).
If \(\mathcal D_I\) is dense in the open-channel energy subspace, the
same uniformly bounded family extends the limit uniquely and strongly to
that subspace.
\end{lemma}

\begin{proof}
For each fixed \(\psi\), the target Hilbert space is complete, so
pointwise convergence is equivalent to the Cauchy property of
\(\{S_m(\tau_m)\psi\}_m\).  Uniform boundedness gives a common norm bound.
If \(\psi_n\in\mathcal D_I\) tends to any vector in the open-channel
subspace, the uniform bound makes \(S_{\rm as}\psi_n\) Cauchy and
independent of the approximating sequence.  This gives the unique strong
extension.
\end{proof}


\begin{theorem}[Conditional large-window M\o ller/cross-section criterion]
\label{p4:thm:moller-bridge}
Let \(m_x\to\infty\) be the reciprocal-position unfolding with
\[
\ell_F\sqrt{m_x(m_x+2)}
\le R_x
<\ell_F\sqrt{(m_x+1)(m_x+3)},
\]
on the homogeneous sharp-frontier branch, or with a separately supplied
compatible \(R_x(m_x)\to\infty\) trajectory on another displayed branch,
while \(m_p=1\) remains fixed.
Let \(H_m=H_{0,m}+V_m\) be a compatible directed family of finite channel Hamiltonians with \(V_m\) localized in a fixed interaction region.
Let \(I\) be a compact open-channel energy window separated from thresholds
and bound-sector spectrum, restrict to the compatible open packet core in
\(I\), and let
\(j_{\rm in}^{(m)}\) be its coherently supplied isometric embedding into the
\(m\)-th finite carrier.  Let
\[
j_{\rm out}^{(m)}:
\mathcal K^{(m)}_{\rm out}\longrightarrow\mathcal H_{\rm out}
\]
be coherent isometric embeddings of all finite outgoing spaces into one
common complete outgoing Hilbert space.  Work on the large-window bridge
domain where the
supplied directed finite generators agree on fixed local regions up to errors
vanishing on the compatible packet core, the channel embeddings are coherent
along the refinement system.
For every plateau choice
\[
\tau_{\mathrm{sc}}\ll \tau_m\ll \tau_{\mathrm{cone}}(m_x)
\]
with \(\tau_m/R_x(m_x)\to0\), the finite-window channel operators
\[
S_m(\tau_m)
:=
j_{\rm out}^{(m)}
P^{(m)}_{\mathrm{out}}
e^{i\tau_mH_{0,m}/\hbar}
e^{-i\tau_mH_m/\hbar}
j^{(m)}_{\mathrm{in}}
\]
obey, for every \(\psi\) in the packet core, the componentwise physical-error
bound
\[
\begin{aligned}
\|S_m(\tau_m)\psi-S_n(\tau_n)\psi\|
\le{}&
\epsilon^{\mathrm{loc}}_{m,n}(\psi)
+\epsilon^{\mathrm{leak}}_m(\psi)
+\epsilon^{\mathrm{leak}}_n(\psi)\\
&+\epsilon^{\mathrm{ret}}_m(\psi)
+\epsilon^{\mathrm{ret}}_n(\psi)
+\epsilon^{\mathrm{tail}}_{m,n}(\psi),
\end{aligned}
\]
where every term tends to zero as \(m,n\to\infty\), uniformly on bounded
subsets of the packet core.  For compactly supported \(V_m\), the tail term
is controlled by the stated local/leakage estimates; for short-range
interactions use instead the pre-return Cook condition displayed in the
proof.  Under these hypotheses the operators
converge strongly on the open-channel energy subspace:
\[
S_{\mathrm{as}}
:=
\mathrm{s}\!-\!\lim_{m_x\to\infty}S_m(\tau_m).
\]
Thus, on the domain with the displayed compatibility and error-control hypotheses,
the asymptotic M\o ller scattering operator is the large-\(m_x\) limit of
finite-window channel transport with the free phase removed.  GP1 supplies
the cone; the stated \((\mathrm{DLOC}_{\rm pkt})\) estimate supplies approximate
return suppression for the two-time instrument branches of
Lemma~\ref{p4:lem:derived-scattering-plateau}.  Neither one by itself proves
the local-compatibility, leakage, Cook-tail, embedding, or detector and flux
hypotheses; once their displayed errors vanish, Cauchy convergence follows.
If, in addition, the outgoing detector POVMs refine to angular cells
\(\Delta\Omega\) and the incident flux admits a stable normalization, the
finite-window cell probabilities define differential cross sections.  Their
large-window and fine-cell limit agrees with the textbook cross section
whenever that final detector and flux limit exists.

The large-window hypotheses are not Paper-IV premises for finite scattering.
They are the conditions handled by the continuous-QFT handoff for admissible
Born--Hopf refinement systems and enter here only as the theorem domain of the
scattering interface.
\end{theorem}


\begin{proof}
\ProofPlan{We prove the large-window scattering bridge by using the plateau input, establishing a strong Cauchy estimate for finite channel operators, controlling compact-support or short-range tails, and passing to detector-cell probabilities and the fine-cell cross-section limit.}
\begin{proofsteps}

\item \textbf{Large-window plateau input.}
Lemma~\ref{p4:lem:derived-scattering-plateau} proves that the nominal
cone-traversal scale diverges and that each stated sequential
far-detection/return branch is suppressed on the open packet core in the
joint limit.  A different detector history must supply its own instrument
estimate.  The remaining substantive interface burden is the displayed
local/leakage/return/tail control; once those bounds vanish, the strong
Cauchy property is a theorem rather than a separate hypothesis.
Take two unfoldings \(m,n\) and a packet-core vector
\(\psi\in\mathcal D_I\).
On pre-traversal windows, the GP1 operational cone together with the
supplied local-generator compatibility and
\((\mathrm{DLOC}_{\rm pkt})\) leakage estimate makes
the packet see the same fixed interaction region, the same compatible local
free dynamics, and the same compact momentum carrier \(m_p=1\); increasing
\(m_x\) only moves the reciprocal-position boundary outward.

\item \textbf{Cauchy estimate for finite channel operators.}
By the theorem-domain local propagation and embedding estimates, the
difference satisfies
\[
\begin{aligned}
\bigl\|S_m(\tau_m)\psi-S_n(\tau_n)\psi\bigr\|
\le{}&
\epsilon^{\mathrm{loc}}_{m,n}
+
\epsilon^{\mathrm{leak}}_m+
\epsilon^{\mathrm{leak}}_n
+
\epsilon^{\mathrm{ret}}_m+
\epsilon^{\mathrm{ret}}_n
+
\epsilon^{\mathrm{tail}}_{m,n} .
\end{aligned}
\]
Here \(\epsilon^{\mathrm{loc}}_{m,n}\to0\) by compatibility of the supplied
directed finite generators on every fixed local region;
\(\epsilon^{\mathrm{leak}}_m\to0\) by the finite spectral-cutoff refinement;
\(\epsilon^{\mathrm{ret}}_m\to0\) by
Lemma~\ref{p4:lem:derived-scattering-plateau} for the stated finite
detector-time instrument (or by the independently supplied return error for
another detector model); and for compactly supported
\(V_m\), \(\epsilon^{\mathrm{tail}}_{m,n}\) is controlled by the same
\((\mathrm{DLOC}_{\rm pkt})\) leakage after the packet clears \(B_{r_0}\).
It tends to zero
but need not vanish identically.

\item \textbf{Short-range/Cook tail control.}
For short-range rather than compactly supported interactions, require the
uniform pre-return Cook-tail condition: for some \(0<\theta<1\),
\[
\lim_{T\to\infty}\;
\limsup_{m\to\infty}
\int_T^{\theta\tau_{\rm cone}(m_x)}
\|V_m e^{-itH_{0,m}/\hbar}j_{\rm in}^{(m)}\psi\|\,dt
=0
\qquad(\psi\in\mathcal D_I).
\]
An individual finite recurrent system is not required to have an integrable
tail out to \(+\infty\); the vanishing is a uniform joint
large-window/pre-return statement.
By the theorem-domain hypothesis all displayed errors tend to zero; hence
\(\{S_m(\tau_m)\psi\}_m\) is Cauchy for every packet-core vector \(\psi\).
Lemma~\ref{p4:lem:finite-window-scattering-cauchy-criterion} then gives the strong limit on the packet core and extends it strongly from the dense packet core to the open-channel energy subspace by uniform boundedness.


\item \textbf{Detector-cell probabilities and fine-cell limit.}
For a fixed detector cell \(Q_{\Delta\Omega}^{(n)}\), boundedness gives
\[
\langle \psi,S_m(\tau_m)^*Q_{\Delta\Omega}^{(n)}S_m(\tau_m)\psi\rangle
\longrightarrow
\langle \psi,S_{\mathrm{as}}^*Q_{\Delta\Omega}^{(n)}S_{\mathrm{as}}\psi\rangle.
\]
If the incident flux normalization also converges, normalized finite-window cell probabilities have a large-window limit.
Passing next to the fine-cell limit \(\Delta\Omega\to0\) yields the differential cross section whenever that detector/flux refinement exists.
\end{proofsteps}

\end{proof}

This theorem is an interface theorem derived on its natural finite-window domain.
The derivation applies to localized or Cook-short-range interactions on open-channel packet states away from thresholds and bound sectors.
Asymptotic completeness, metastable resonance structure, threshold singularities, and fully renormalized continuum QFT belong to the corresponding continuum or sector refinement layers.
The large-window limit is a reciprocal-position/harmonic unfolding limit and is compatible with fixed compact momentum carrier \(m_p=1\); it is not a dynamical momentum-\(m\) limit.



\paragraph{Feshbach reduction and resonances}
\label{p4:sec:part2-feshbach-resonances}

Resonances arise when an internal quasi-bound sector couples weakly to open channels.
Let
\[
\mathcal K = P\mathcal K \oplus Q\mathcal K,
\]
with $P$ projecting onto open channels and $Q$ onto an internal sector.
Whenever $z$ lies outside the spectrum of $PHP$, the exact finite Feshbach operator on the internal sector is
\begin{equation}
\label{p4:eq:feshbach-effective}
H^{\mathrm{eff}}_Q(z)
=
QHQ + QHP\,(z-PHP)^{-1}PHQ.
\end{equation}
Its eigenvalue branches determine the displaced quasi-bound levels seen in the channel amplitudes.

\begin{theorem}[Exact finite Feshbach reduction and conditional resonance width]
\label{p4:thm:finite-resonance-reduction}
Suppose $Q\mathcal K$ contains an isolated level $\ket{r}$ of $QHQ$ near the open-channel energies of $PHP$.
For $z\notin\mathrm{spec}(QHQ)$ such that the Schur complement is invertible, equivalently $z\in\rho(H)\cap\rho(QHQ)$, the reduced channel resolvent has the Schur-complement form
\begin{equation}
\label{p4:eq:resolvent-resonance}
P(z-H)^{-1}P
=
\bigl(z-PHP-PHQ(z-QHQ)^{-1}QHP\bigr)^{-1}.
\end{equation}
Together with the internal-sector effective operator above, this gives a
meromorphic finite energy dependence whose exact poles, for self-adjoint
\(H\), lie on the real spectrum.  Large finite-window response may occur
near an isolated eigenvalue or avoided crossing, equivalently near a
solution of the effective level equation
\begin{equation}
\label{p4:eq:effective-level-equation}
z\in\mathrm{spec}\bigl(H_Q^{\mathrm{eff}}(z)\bigr).
\end{equation}
No nonzero decay width follows in the exact finite Hermitian problem.  If an
additional large-window/continuum limiting-absorption domain supplies a
boundary value of the self-energy
\[
\Sigma(E+i0)=\Delta(E)-\frac{i}{2}\Gamma(E),
\qquad \Gamma(E)\ge0,
\]
that varies slowly near an isolated level, then the analytically continued
channel response has the Breit--Wigner profile to leading order.  This last
statement is conditional on that continuum or reservoir interface.
\end{theorem}

\begin{proof}
Equation \eqref{p4:eq:resolvent-resonance} is the Schur-complement identity
for the block decomposition defined by \(P\) and \(Q\).  The finite spectral
theorem gives a rational resolvent with real poles at the eigenvalues of the
self-adjoint matrix \(H\).  Near an isolated eigenvalue its rank-one
contribution can dominate finite-window response, but it supplies no
imaginary part and hence no irreversible width.  On the additional
limiting-absorption domain, substituting the displayed slowly varying
boundary value into the scalar isolated-level Schur complement produces a
denominator
\[
E-E_r-\Delta(E_r)+i\Gamma(E_r)/2,
\]
whose squared modulus is the leading Breit--Wigner profile.
\end{proof}

\begin{remark}[Status of finite resonance reduction]
\label{p4:rem:finite-resonance-status}
The exact content of
Theorem~\ref{p4:thm:finite-resonance-reduction} is the finite
Schur-complement/Feshbach identity specialized to the present channel
decomposition \cite{Feshbach1958,Feshbach1962}.  A complex pole, exponential
decay law, or nonzero width belongs to the explicitly stated continuum,
reservoir, or coarse-grained interface.
\end{remark}

\begin{remark}[Finite and continuum resonance domains]
\label{p4:rem:resonance-meaning}
In the exact finite theory a ``resonant'' sector means an internal level
causing strong but recurrent finite-window response.  Only after a
large-\(m_x\) limiting-absorption, reservoir, or irreversible coarse-graining
limit supplies a continuous spectral boundary value may one use complex-pole
and width language in the usual sense.
\end{remark}



\paragraph{Decay channels and open-system reduction}
\label{p4:sec:part2-decay-channels}

Once unresolved channel sectors are traced out, one obtains exact CPTP
reduced maps of the type proved in Part~I.  They become a time-homogeneous
GKLS process only on an additional Markov-semigroup or reservoir domain.
Let $Q$ project onto a metastable internal sector and let $P=\mathbf 1-Q$ collect the unresolved escape channels.
If a supplied reduced evolution including decay products is a
norm-continuous CPTP semigroup, its generator is GKLS.  If one retains only
the surviving metastable \(Q\)-block, the corresponding Markovian restriction
is generally trace-decreasing:
\begin{equation}
\label{p4:eq:resonance-gkls}
\dot\rho_Q
=
-i[H_{\mathrm{eff}},\rho_Q]
-
\frac12\sum_k\{L_k^*L_k,\rho_Q\}
+
\text{possible recycling terms internal to }Q.
\end{equation}
The operators \(L_k\) encode the available escape channels.  For a single
metastable mode \(\ket{r}\) with no recycling back into \(Q\), this Markovian
model gives the population law
\begin{equation}
\label{p4:eq:exponential-decay}
\Tr(\rho_Q(t)\ket{r}\bra{r})
=e^{-\Gamma t}\Tr(\rho_Q(0)\ket{r}\bra{r}),
\qquad
\Gamma=\sum_k \bra{r}L_k^*L_k\ket{r}.
\end{equation}
The exact finite closed model remains recurrent.  Exponential decay is the
output of the stated irreversible reduced-semigroup interface, not a
consequence of finite Feshbach reduction alone.

\begin{remark}[Bridge between scattering and spectroscopy]
\label{p4:rem:scattering-spectroscopy-bridge}
The framework therefore unifies bound-state spectroscopy, scattering, and decay in a single finite language:
exact unitary transport on a large position-side sector yields channel scattering;
projection to a quasi-bound block yields recurrent resonant response; and a
separately justified Markov/reservoir elimination of unresolved channels
yields irreversible decay.
This is the natural scattering completion of the model library developed in Section~\ref{p4:sec:part2-worked-models}.
\end{remark}




\paragraph{Finite-temperature kinetic QFT interface for baryogenesis}
\label{p4:sec:part2-thermal-baryogenesis-qft}

The finite mode/Fock bridge supplies the matrix and density-operator carrier
used by the later baryogenesis calculation.  The local relativistic
interaction, continuum two-body phase space, thermal plasma, kinetic
closure, and cosmological expansion are additional large-\(m_x\)
QFT/thermal interfaces; they are not consequences of a finite matrix algebra.
This subsection records those standard imported formulas with their
assumptions so that Paper~V-C can specialize them.  None of these interfaces
enlarges the protected compact momentum carrier \(m_p=1\).

Let \(L_\alpha\) be lepton doublet modes, let \(H\) be the Higgs-doublet mode, and
let \(N_i=N_i^c\) be gauge-neutral Majorana modes.  On the finite Fock carrier
selected by the previous subsections, the minimal local interaction block on the reciprocal-position/QFT interface has the
form
\begin{equation}
\label{p4:eq:type-I-seesaw-finite-interface}
\mathcal L_{N}
=
\frac12\overline{N_i}i\gamma^\mu\partial_\mu N_i
-
(Y_\nu)_{\alpha i}\,\overline{L_\alpha}\widetilde H N_i
-
\frac12\overline{N_i^c}M_{ij}N_j
+\mathrm{h.c.},
\qquad
\widetilde H=i\sigma_2H^* .
\end{equation}
Here \(Y_\nu\) and \(M\) are finite matrices acting on the family/multiplet blocks.
Thus the type-I seesaw interaction is, in the present framework, an ordinary
finite-mode Hamiltonian interaction plus the large-window/thermal interface
already licensed above.


\begin{theorem}[Conditional continuum/thermal type-I seesaw interface]
\label{p4:thm:minimal-finite-temperature-seesaw-interface}
\leavevmode\par\noindent
Within a finite Standard-Model carrier augmented by gauge-neutral Majorana
modes \(N_i\), and restricting to the usual renormalizable singlet-neutrino
Yukawa and Majorana mass interactions, the type-I block
\eqref{p4:eq:type-I-seesaw-finite-interface} is the finite QFT interface that
supplies a lepton-number-violating ingredient: the Majorana mass term is not
invariant under the continuous lepton-number phase.  That violation alone
does not generate a net \(B-L\) asymmetry.  On the standard decay-leptogenesis
branch, a nonzero source additionally requires a CP-odd flavor invariant, an
absorptive open-channel phase, departure from thermal equilibrium, and a
kinetic evolution domain.  In the unflavored realization a representative
necessary CP condition is
\begin{equation}
\label{p4:eq:seesaw-cp-odd-source-condition}
\Im\!\left[(Y_\nu^\dagger Y_\nu)_{ij}^{\,2}\right]\ne0
\qquad\text{for some }i\ne j,
\end{equation}
with the corresponding flavored invariants used on a flavored domain.  If all
relevant CP-odd invariants vanish, or if exact thermal detailed balance holds,
the net source vanishes.  This statement does not exclude other effective
\(B-L\)-violating operators or UV completions; it records the conditional
type-I interface used by the later sector-specific paper.

On the additional continuum relativistic two-body phase-space domain, at tree
level in the unbroken-phase/vacuum or Maxwell--Boltzmann approximation, with
final-state masses and finite-temperature statistical factors neglected, the
standard heavy Majorana decay width is
\begin{equation}
\label{p4:eq:heavy-majorana-width}
\Gamma_i
=
\frac{M_i}{8\pi}(Y_\nu^\dagger Y_\nu)_{ii},
\end{equation}
On the further Markovian/limiting-absorption approximation in which the
displayed single-width regulator is valid, a near-degenerate pair
\((i,j)\) has the conventional self-energy kernel
\begin{equation}
\label{p4:eq:resonant-self-energy-kernel}
\mathcal R_{ij}
=
\frac{(M_i^2-M_j^2)M_i\Gamma_j}
{(M_i^2-M_j^2)^2+M_i^2\Gamma_j^2},
\qquad
|\mathcal R_{ij}|\le\frac12 .
\end{equation}
The bound is saturated exactly when
\begin{equation}
\label{p4:eq:resonant-width-optimality}
M_i\Gamma_j=|M_i^2-M_j^2|.
\end{equation}
\end{theorem}


\begin{proof}
Under \(N_i\mapsto e^{i\alpha}N_i\), the Majorana bilinear transforms with
twice the lepton-number phase, so a nonzero Majorana mass breaks continuous
lepton number.  A particle--antiparticle rate difference is CP odd; in the
standard unflavored decay calculation its interference coefficient contains
\(\Im[(Y_\nu^\dagger Y_\nu)_{ij}^{\,2}]\), and it also requires the absorptive
part of an open-channel loop.  Hence lepton-number violation without a CP-odd
invariant or absorptive phase gives zero \(\varepsilon_i\).  In exact thermal
equilibrium, CPT/unitarity and detailed balance cancel source and inverse
processes, so departure from equilibrium and a kinetic domain are separately
required.

The finite Born rule supplies probabilities once the transition instrument
has been constructed, but it does not supply the factor \(1/(8\pi)\).
That factor is the standard continuum two-body Lorentz-invariant phase-space
integral, together with the Majorana conjugate-channel count, on the domain
stated in the theorem.  Likewise the nonzero \(\Gamma_j\) in
\eqref{p4:eq:resonant-self-energy-kernel} is not an exact pole width of a
finite Hermitian matrix; it is the boundary-value/Markov width supplied by
the stated open-channel interface.  Once that conventional regulator is
adopted, setting
\(X=M_i^2-M_j^2\) and \(Y=M_i\Gamma_j\), one has
\(\mathcal R_{ij}=XY/(X^2+Y^2)\).  The inequality
\(2|XY|\le X^2+Y^2\) gives the stated bound and the equality condition.
\end{proof}

Once the CP-odd source, out-of-equilibrium preparation, and plasma/cosmological
domain have been supplied, thermal propagation of an asymmetry is represented
as a finite open-system problem after the unobserved plasma modes are traced
out.  Write
\[
z=M/T,
\qquad
Y_X=n_X/s,
\]
where \(s\) is the entropy density. For a quasi-degenerate pair, \(M\) is the common reference mass; for well-separated heavy modes one should use \(z_i=M_i/T\) or rewrite the equations with a chosen reference mass and threshold factors. In the momentum-integrated Maxwell--Boltzmann
reduction the one-particle yields obey
\begin{equation}
\label{p4:eq:integrated-heavy-neutrino-boltzmann}
\frac{dY_{N_i}}{dz}
=
-D_i(z)\bigl(Y_{N_i}-Y_{N_i}^{\rm eq}\bigr),
\end{equation}
\begin{equation}
\label{p4:eq:integrated-BL-boltzmann}
\frac{dY_{B-L}}{dz}
=
\sum_i \varepsilon_iD_i(z)\bigl(Y_{N_i}-Y_{N_i}^{\rm eq}\bigr)
-
W(z)Y_{B-L},
\end{equation}
with
\begin{equation}
\label{p4:eq:D-W-kernels}
D_i(z)=\mathcal K_i z\frac{\mathrm K_1(z)}{\mathrm K_2(z)},
\qquad
\mathcal K_i=\frac{\Gamma_i}{H(T=M_i)},
\qquad
W_i^{\rm ID}(z)=\frac14\mathcal K_i z^3\mathrm K_1(z).
\end{equation}
Here \(\mathrm K_1,\mathrm K_2\) denote modified Bessel functions.  Equations
\eqref{p4:eq:integrated-heavy-neutrino-boltzmann}--\eqref{p4:eq:D-W-kernels}
are standard momentum-integrated kinetic interface equations under the
Maxwell--Boltzmann, kinetic-equilibrium, plasma-rate, entropy-conservation,
and expanding-background assumptions.  They are not derived from one fixed
finite GKLS reduction; the collision kernels and Hubble rate must be supplied
by the thermal continuum/reservoir and cosmological domains.

For a quasi-degenerate pair, the scalar equations may be replaced by a two-state
finite density matrix on the pair subspace:
\begin{equation}
\label{p4:eq:two-state-thermal-density-matrix}
\frac{d\rho_N}{dz}
=
-i[\Omega_N,\rho_N]
-
\frac12\{\Gamma_N,\rho_N-\rho_N^{\rm eq}\}
+\mathcal C_N,
\end{equation}
\begin{equation}
\label{p4:eq:flavored-asymmetry-density-matrix}
\frac{dY_{\Delta_\alpha}}{dz}
=
{\rm Tr}\!\left[\varepsilon_\alpha D(z)(\rho_N-\rho_N^{\rm eq})\right]
-
\sum_\beta W_{\alpha\beta}(z)Y_{\Delta_\beta}.
\end{equation}
This is the finite-dimensional reduced density-matrix carrier used by the
later resonant baryogenesis interface, but its thermal collision kernels are
continuum/plasma inputs.  Sector-specific geometric source amplitudes,
finite-trace efficiencies, and plasma refinements are not fixed by
Paper~IV.  Paper~V-C may derive or calibrate them only after its own
Standard-Model, neutral-sector, interaction, and thermal domains have been
stated.

\begin{remark}[What the finite-temperature interface does and does not fix]
\label{p4:rem:thermal-interface-status}
The present subsection supplies a finite heavy-neutral carrier and records
standard conditional continuum/thermal formulas: decay and resonance
kernels, a two-state kinetic density matrix, and the imported
sphaleron/photon conversion factors below.  It does not derive the plasma,
the Markov limit, a geometric CP source, or a sector-specific efficiency.
Those data and their validity belong to the later sector paper.
\end{remark}

\begin{lemma}[Uniform finite-trace efficiency]
\label{p4:lem:uniform-finite-trace-efficiency}
Let \(P_{\rm src}\) be a rank-\(r\) coherent source projector inside a finite thermal visible complement \(P_{\rm vis}\) of rank \(R\).  If the reduced thermal coupling is democratic on \(P_{\rm vis}\) and no channel weights survive the coarse-graining, then the finite trace efficiency is
\[
\kappa_{\rm tr}
=
\frac{\Tr(P_{\rm src})}{\Tr(P_{\rm vis})}
=
\frac{r}{R}.
\]
\end{lemma}

\begin{proof}
Democratic coupling means that, after coarse-graining, every orthonormal trace cell in \(P_{\rm vis}\) carries equal weight.  The reduced efficiency of the coherent source is therefore the normalized trace weight of its projector inside the visible complement.  Since \(\Tr(P_{\rm src})=r\) and \(\Tr(P_{\rm vis})=R\), the stated ratio follows.  The hypotheses explicitly exclude residual channel weights; if such weights survive, the result becomes model-dependent and is not a theorem of the uniform finite-trace interface.
\end{proof}


\begin{proposition}[Imported Standard-Model plasma conversion factors]
\label{p4:prop:sphaleron-and-photon-conversion}
\leavevmode\par\noindent
This proposition records standard thermal-Standard-Model plasma conversion factors for later interface use; it is not a Paper-IV derivation of the Standard-Model plasma.  Work in the Standard-Model plasma with three fermion families and one Higgs doublet.  Electroweak sphaleron equilibrium gives
\begin{equation}
\label{p4:eq:sphaleron-28-79}
Y_B
=
\frac{8N_f+4N_H}{22N_f+13N_H}Y_{B-L}
=
\frac{28}{79}Y_{B-L}.
\end{equation}
The present entropy-to-photon conversion is
\begin{equation}
\label{p4:eq:entropy-photon-704}
\mathcal C_{\gamma s}
:=
\frac{s_0}{n_{\gamma,0}}
=
\frac{\pi^4}{45\zeta(3)}g_{\ast S,0},
\qquad
g_{\ast S,0}=2+\frac78\,2\cdot3\left(\frac4{11}\right)=\frac{43}{11},
\end{equation}
so that
\begin{equation}
\label{p4:eq:etaB-conversion}
\eta_B
=
\frac{n_B-n_{\bar B}}{n_\gamma}
=
\mathcal C_{\gamma s}Y_B
\simeq 7.04\,Y_B.
\end{equation}
If the asymmetry variable is instead written as \(Y_{L-B}\), then
\(Y_{B-L}=-Y_{L-B}\); only the sign changes.  There is no separate standard
\(6.96\) conversion factor associated with the change from \(B-L\) to \(L-B\).
\end{proposition}


\begin{proof}
The sphaleron coefficient is the chemical-potential solution of the electroweak
plasma constraints for \(N_f\) fermion families and \(N_H\) Higgs doublets.
Here it is specialized to
\[
  (N_f,N_H)=(3,1).
\]
The entropy and photon densities are
\[
  s_0=\frac{2\pi^2}{45}g_{\ast S,0}T_{\gamma,0}^3,
  \qquad
  n_{\gamma,0}=\frac{2\zeta(3)}{\pi^2}T_{\gamma,0}^3 .
\]
Their ratio gives \eqref{p4:eq:entropy-photon-704}.  The identity
\(Y_{L-B}=-Y_{B-L}\) is definitional.
\end{proof}




\subsection{Release v4.96: boundary reduction versus thermal kinetics}
\label{p4:sec:v496-boundary-thermal-interface}

\begin{proposition}[Euclidean boundary depth is not thermal inverse temperature]
\label{p4:prop:v496-depth-not-beta}
For a positive finite Hamiltonian reconstructed from a Euclidean boundary
kernel, the transfer operator has the form
\[
 T_L=e^{-LH}.
\]
After trace normalization, $L\to\infty$ projects onto the lowest-energy
sector.  Therefore a boundary/gluing depth which tends to infinity cannot by
itself be identified with the inverse temperature of a simultaneously hot
physical plasma.  In particular,
\[
 \boxed{\ell_{\rm boundary}\ne\beta_{\rm thermal}}
\]
unless an independent theorem identifies the two parameters on a declared
thermal domain.
\end{proposition}

\begin{theorem}[Local quantum-kinetic H-theorem interface]
\label{p4:thm:v496-quantum-kinetic-h}
Assume a declared quasiparticle/scattering domain with microreversible
$2\leftrightarrow2$ transition kernels
$W_{12\leftrightarrow34}\ge0$.  On the strictly interior occupation branch
$f>0$ and $1+\epsilon f>0$ in every participating mode, let
\[
 X=f_1f_2(1+\epsilon_3f_3)(1+\epsilon_4f_4),\qquad
 Y=f_3f_4(1+\epsilon_1f_1)(1+\epsilon_2f_2),
\]
where $\epsilon=+1$ for bosons and $-1$ for fermions.  Then the quantum
kinetic entropy satisfies
\[
 \boxed{
 \dot S_{\rm coll}
 =\frac14\sum_{\rm reactions}\int d\Pi\,
 W\,(X-Y)\log\frac{X}{Y}\ge0.}
\]
For zero entropy production, every positive-rate reaction satisfies
$X=Y$ almost everywhere.  Equivalently
$\eta_a(p)=\log[(1+\epsilon_af_a(p))/f_a(p)]$ obeys
$\eta_1+\eta_2=\eta_3+\eta_4$, so $\eta$ is an additive collision
invariant.  \emph{If, in addition,} the full real space of such invariants on
the admitted positive-rate collision network is exactly the span of the
conserved energy--momentum coordinates and stated exact charges (with their
specified independent linear relations), and their coefficients admit the
physical inverse-temperature parametrization, then the stationary
occupations form the Bose/Fermi Gibbs family
\[
 \boxed{
 f_a^{\rm eq}(p)=
 \frac{1}{\exp[\beta(u\cdot p-\sum_A\mu_Aq_A(a))]-\epsilon_a}}
\]
with chemical potentials only for that complete invariant basis.  A connected
reaction graph alone does \emph{not} imply invariant-space completeness or
uniqueness: extra additive collision invariants yield more general stationary
families.  At zero-occupation endpoints one must additionally control
potentially frozen reactions before extending the equality characterization.
\end{theorem}

\begin{proof}
For $X,Y>0$, $(X-Y)\log(X/Y)\ge0$, with equality only for $X=Y$.
Nonnegative rates therefore force detailed balance almost everywhere on
all positive-rate reaction channels when the summed production vanishes.
Taking logarithms of $X=Y$ gives the displayed additive-invariant equation.
On the explicitly complete invariant space, its general solution has
$\eta_a(p)=\beta(u\cdot p-\sum_A\mu_Aq_A(a))$; solving
$e^\eta=(1+\epsilon f)/f$ yields the Gibbs formula.  Without this
classification the entropy inequality remains valid but the Gibbs
uniqueness assertion has no proof.
\end{proof}

\begin{remark}[Scope]
\label{p4:rem:v496-local-not-global-kms}
The theorem is a local kinetic statement on an admitted scattering domain.  It
does not assert global all-time KMS convergence of the complete universe and
does not create the scattering/quasiparticle domain from finite algebra alone.
\end{remark}
% END BOUNDARY-THERMAL-INTERFACE v4.96


% BEGIN INTERNAL-RESPONSE-POSITIVE-FREQUENCY v4.97


\subsection{Release v4.97: positive-frequency realization of internal response}
\label{p4:sec:v497-internal-response-positive-frequency}

\begin{proposition}[Positive-frequency polarization of a physical internal response]
\label{p4:prop:v497-internal-response-positive-frequency}
Let a separately constituted physical bosonic carrier have positive
stationary one-particle energy operator $H_1=\hbar\Omega$ with $\Omega>0$.
On its real Cauchy response space define
\[
 J(q,p)=(-\Omega^{-1}p,\Omega q),
 \qquad
 W(q,p)=\frac{1}{\sqrt{2\hbar}}
 \left(\Omega^{1/2}q+i\Omega^{-1/2}p\right).
\]
Then $WJ=iW$ and
\[
 \langle W(q,p),H_1W(q,p)\rangle
 =\frac12\bigl(\|p\|^2+\|\Omega q\|^2\bigr).
\]
Hence, once particle ownership is independently licensed, the one-particle
Hilbert realization is the positive-frequency polarization of the corresponding
real internal response.  The proposition does not assign particle semantics to
every positive spectral block.
\end{proposition}

\begin{corollary}[State inheritance does not imply boundary-current saturation]
\label{p4:cor:v497-state-not-saturation}
Let an occurring root state be transported unitarily/symplectically to a later
constituted carrier, so the later state is fixed by ancestry rather than by a
new selector.  This state inheritance alone does not imply
\[
 \frac1V\left\langle\partial_\eta H_{\rm ord}\right\rangle
 =\mathcal I
\]
for an independently generated structural boundary-current owner $\mathcal I$.
Such an equality is a common-owner/saturation theorem and must be proved from
the action, stress allocation and declared physical owner relation.
\end{corollary}
% END INTERNAL-RESPONSE-POSITIVE-FREQUENCY v4.97

% BEGIN CHARGE-PRESERVING-QUADRATIC-EVOLUTION v4.98


\subsection{Release v4.98: charge-preserving quadratic evolution and conditional CP pump}
\label{p4:sec:v498-charge-preserving-quadratic-evolution}

\begin{theorem}[No net charge production from a charge-preserving quadratic history]
\label{p4:thm:v498-no-charge-production}
Let $\mathbb H(t)$ be a strongly continuous finite/quasi-free quadratic
Hamiltonian on a Nambu carrier and let $\mathbb Q$ be a time-independent
self-adjoint additive charge.  If
\[
 [\mathbb Q,\mathbb H(t)]=0
 \qquad\text{for all }t\in[t_i,t_f],
\]
then the propagator $U(t_f,t_i)$ commutes with $\mathbb Q$.  Consequently an
initial state with vanishing charge expectation remains charge neutral:
\[
 \boxed{\operatorname{Tr}(\rho_i\mathbb Q)=0
 \Longrightarrow
 \operatorname{Tr}(U\rho_iU^\dagger\mathbb Q)=0.}
\]
In particular, time dependence, nonadiabaticity and Bogoliubov mixing do not
by themselves generate an asymmetry in an exactly conserved charge.
\end{theorem}

\begin{proof}
For $F(t)=U(t,t_i)^\dagger\mathbb Q U(t,t_i)$,
\[
 \dot F(t)=iU^\dagger[\mathbb H(t),\mathbb Q]U=0,
\]
so $F(t)=\mathbb Q$.  The expectation-value statement follows by cyclicity of
the trace.
\end{proof}

The covariant pairing trajectory diagnostic and its fixed-parallel-frame
specialization are stated in
Corollary~\ref{p4:cor:v498-phase-locked-pairing-invariant}; the charge-response
functional is supplied by Theorem~\ref{p4:thm:v4102-pair-charge-source}.

% END CHARGE-PRESERVING-QUADRATIC-EVOLUTION v4.98

% BEGIN FIXED-OWNER-PAIR-PHASE-SCOPE v4.103


\subsection{Release v4.103: fixed-owner real pair transport and CP scope}
\label{p4:sec:v4103-fixed-owner-pair-phase-scope}

\begin{proposition}[Real fixed-owner response generates no additional pair phase]
\label{p4:prop:v4103-real-pair-response-no-new-phase}
Let an already occurring finite pair datum $\Delta_0$ be transported on a
fixed owner by a real scalar response $\eta$, so that
\[
 \Delta_\partial=\eta\Delta_0.
\]
For the inherited metric-compatible connection,
\[
 D\Delta_\partial=(d\eta)\Delta_0+\eta D\Delta_0,
\]
and therefore
\[
 \boxed{
 \operatorname{Im}\langle\Delta_\partial,D\Delta_\partial\rangle_{\rm HS}
 =\eta^2\operatorname{Im}\langle\Delta_0,D\Delta_0\rangle_{\rm HS}.}
\]
Thus the real response contributes no independent covariant phase rate.  On a
parallel realization $D\Delta_0=0$, the self-trajectory diagnostic vanishes
identically.
\end{proposition}

\begin{proof}
Expand the metric-compatible pairing.  The term
$\eta(d\eta)\langle\Delta_0,\Delta_0\rangle_{\rm HS}$ is real, while the
remaining imaginary part is $\eta^2$ times the inherited diagnostic of
$\Delta_0$.  No new phase coordinate has been introduced.
\end{proof}

\begin{remark}[Relative-sector CP remains a different observable]
\label{p4:rem:v4103-relative-sector-cp}
The proposition closes only phase generation by a real fixed-owner radial
response.  It neither sets every CP-odd invariant to zero nor determines a
matter--antimatter asymmetry.  Relative-sector current invariants, the initial
state and the exact charge evolution remain separately typed.  In particular,
adding a scalar factor $e^{i\Theta}$ to the same frozen pair coefficient is a
new deformation unless an independent physical constructor supplies it.
\end{remark}
% END FIXED-OWNER-PAIR-PHASE-SCOPE v4.103



% BEGIN INLINE SOURCE: Paper4-ChargedResponseDiagnostics-Appendix-v1_0.tex
% Framework v2.168: integrated finite charged-source derivations.


\section{Complete charged-source dynamics and transport}
\label{p4b:sec:v238-b5}
\noindent\textit{Source and claim boundary.} Exact charge-sector construction, stationary support, finite pair insertion and time-dependent charge transport consume the complete source operator. A bilinear mass-only block must not be silently promoted to a fully interacting pole.

\subsection{Internal sectors, mixing, and threshold structure}
\label{p4:sec:part2-internal-thresholds}

The same finite quantum calculus that organizes ordinary textbook channels also organizes
distinguished internal sectors.
This is the framework point needed later when one passes from abstract quantum mechanics to concrete sector-specific applications:
once a family of mutually orthogonal projectors is present, transitions, threshold opening,
mixing, and metastability are all ordinary finite quantum phenomena inside the same algebra.




\paragraph{Sector compatibility, admissible observables, and factorized family spaces}
\label{p4:sec:part2-sector-compatibility}

Let
\[
\mathcal H=\bigoplus_{a=1}^r \mathcal H_a,
\qquad
P_aP_b=\delta_{ab}P_a,
\qquad
\sum_{a=1}^r P_a=\mathbf 1,
\]
be a distinguished orthogonal sector decomposition.
Once the projector family $\{P_a\}$ is given, the framework automatically provides three natural operator classes.
First, the \emph{sector-compatible} operators are those that preserve every sector; the physical observables are the self-adjoint elements among them,
\begin{equation}
\label{p4:eq:sector-compatible-observables}
\mathcal A_{\mathrm{sec}}
:=
\{A\in\End(\mathcal H): [A,P_a]=0\ \text{for all $a$}\}
=
\bigoplus_{a=1}^r \End(\mathcal H_a).
\end{equation}
Second, the diagonal algebra generated by the projectors,
\begin{equation}
\label{p4:eq:sector-center}
\mathcal Z_{\mathrm{sec}}:=\mathrm{span}\{P_1,\dots,P_r\},
\end{equation}
encodes coarse sector occupancy.
Third, the full operator algebra decomposes into block maps
\begin{equation}
\label{p4:eq:block-operator-decomp}
\End(\mathcal H)
=
\bigoplus_{a,b=1}^r \mathrm{Hom}(\mathcal H_b,\mathcal H_a)
=
\bigoplus_{a,b=1}^r P_a\End(\mathcal H)P_b,
\end{equation}
so off-diagonal blocks $P_aAP_b$ are the precise carriers of sector mixing.

\begin{theorem}[Local sector-compatible observables and dynamics]
\label{p4:thm:sector-compatible-observables}
For a distinguished orthogonal sector decomposition $\mathcal H=\bigoplus_a \mathcal H_a$ with projectors $\{P_a\}$:
\begin{enumerate}[label=(\roman*)]
\item A self-adjoint operator $A$ is sector compatible if and only if
\begin{equation}
\label{p4:eq:sector-compatible-form}
A=\sum_a P_aAP_a,
\qquad [A,P_a]=0\ \text{for all $a$},
\end{equation}
and then $A$ restricts to an ordinary observable on each sector $\mathcal H_a$.
\item A unitary $U$ preserves the sector decomposition if and only if
\begin{equation}
\label{p4:eq:sector-preserving-unitary}
UP_a=P_aU\ \text{for all $a$},
\end{equation}
in which case
\begin{equation}
\label{p4:eq:sector-preserving-unitary-sum}
U=\bigoplus_a U_a,
\qquad U_a\in U(\mathcal H_a).
\end{equation}
\item A Hamiltonian $H$ preserves all sectors if and only if
\begin{equation}
\label{p4:eq:sector-preserving-hamiltonian}
[H,P_a]=0\ \text{for all $a$},
\end{equation}
and then its evolution satisfies
\begin{equation}
\label{p4:eq:sector-preserving-flow}
e^{-itH}=\bigoplus_a e^{-itH_a},
\qquad H_a:=P_aHP_a.
\end{equation}
\item More generally, any self-adjoint block operator
\begin{equation}
\label{p4:eq:admissible-block-hamiltonian}
H=\sum_a H_{aa}+\sum_{a\neq b} H_{ab},
\qquad H_{ab}:=P_aHP_b,
\end{equation}
is a legitimate finite quantum Hamiltonian.
Its off-diagonal blocks $H_{ab}$ are exactly the admissible sector-mixing couplings.
\end{enumerate}
Thus, once a projector family is supplied by an imported geometric/spectral theorem or declared as theorem-domain data, the decomposition automatically determines both the maximal sector-preserving observable algebra and the precise quantum meaning of inter-sector mixing.  The finite algebra theorem interprets a sector decomposition; it does not select the decomposition as a primitive sector input.
\end{theorem}

\begin{proof}
The equivalence in (i) is immediate from
$A=(\sum_a P_a)A(\sum_b P_b)=\sum_{a,b}P_aAP_b$:
commutation with every $P_a$ forces all off-diagonal blocks to vanish, and conversely a block-diagonal operator plainly commutes with each projector.
Statement (ii) is the same argument for unitaries.
Statement (iii) follows because commutation with all $P_a$ is preserved by functional calculus, so $e^{-itH}$ is block diagonal whenever $H$ is.
Statement (iv) is just the exact block decomposition of the finite operator algebra together with self-adjointness $H_{ba}=H_{ab}^*$.
\end{proof}

\begin{corollary}[Sector observables versus mixing observables]
\label{p4:cor:sector-vs-mixing-observables}
Relative to a distinguished projector family $\{P_a\}$, the framework separates two conceptually different operator types:
\begin{equation}
\label{p4:eq:sector-vs-mixing-split}
A=A_{\mathrm{diag}}+A_{\mathrm{mix}},
\qquad
A_{\mathrm{diag}}:=\sum_a P_aAP_a,
\qquad
A_{\mathrm{mix}}:=\sum_{a\neq b} P_aAP_b.
\end{equation}
The diagonal part is the sector-compatible operator seen by coarse sector measurements,
while the off-diagonal part carries coherence and transitions between sectors.
No additional quantum axiom is required to interpret either part.
\end{corollary}

\begin{proof}
Insert the identity \(\sum_aP_a=\mathbf1\) on both sides of \(A\).  The terms with equal labels give \(A_{\mathrm{diag}}\), and the terms with unequal labels give \(A_{\mathrm{mix}}\).  Operators in the diagonal part commute with all sector projectors and are therefore visible to sector-coarse measurements; the off-diagonal part is exactly the remaining transition/coherence component.
\end{proof}


\begin{proposition}[Trace-neutral finite word selection]
\label{p4:prop:trace-neutral-finite-word-selection}
Let \(\mathcal A\subseteq\End(\mathcal H)\) be a finite matrix algebra equipped with a trace-orthogonal basis \(\{B_\lambda\}\) containing the identity basis element \(B_0=\mathbf 1\), and let
\[
\tau(A):=\frac{1}{\dim\mathcal H}\operatorname{tr}(A)
\]
be the normalized scalar readout.
If a finite word \(W\in\mathcal A\) decomposes as
\[
W=c_0\mathbf 1+\sum_{\lambda\ne0}c_\lambda B_\lambda,
\]
with \(\operatorname{tr}(B_\lambda)=0\) for \(\lambda\ne0\), then
\[
\tau(W)=c_0.
\]
Thus a scalar finite coefficient obtained by trace readout receives contributions only from the neutral identity component of the finite word.
In particular, for the rank-three Weyl basis of Proposition~\ref{p4:prop:rank-three-weyl-basis}, a product contributes to the trace only when its total \(\ZZ_3\times\ZZ_3\) Weyl character is trivial.
For the protected Clifford--Fierz basis of Paper~I-A, Proposition~\ref{p1:prop:math-protected-clifford-fierz}, only the scalar Clifford component survives the trace.
\end{proposition}


\begin{proof}
Trace-orthogonality to the identity means \(\operatorname{tr}(B_\lambda)=0\) for all non-identity basis elements.
Taking the normalized trace of the displayed expansion therefore gives \(\tau(W)=c_0\).
For the Weyl basis, \(X^aZ^b\) has zero trace unless \((a,b)=(0,0)\), and products reduce to a phase times the basis element whose indices are the summed Weyl characters.
For the Clifford--Fierz basis, all non-scalar grades are traceless, so only the scalar grade contributes.
\end{proof}

\begin{remark}[Later correction grammars]
\label{p4:rem:trace-neutral-selection-status}
\leavevmode\par\noindent
Proposition~\ref{p4:prop:trace-neutral-finite-word-selection} is only a finite
quantum-algebra statement.  It does not choose any numerical correction
coefficient for a later paper.  If a later sector-specific paper wants to
promote a sector-generated finite word to a scalar correction term, it must
supply the appropriate channel and then apply this trace-neutrality test there.
\end{remark}



\paragraph{Factorized family spaces and reduced internal states}
\label{p4:sec:part2-factorized-family-spaces}

A particularly important special case occurs when the distinguished sectors are all isomorphic and combine into a tensor factor.
Suppose
\begin{equation}
\label{p4:eq:family-factorization}
\mathcal H\cong \mathcal H_{\mathrm{ext}}\otimes \mathcal H_{\mathrm{fam}},
\qquad
\dim \mathcal H_{\mathrm{fam}}=n,
\end{equation}
with a chosen orthonormal basis $\{e_a\}_{a=1}^n$ of $\mathcal H_{\mathrm{fam}}$ and projectors
\begin{equation}
\label{p4:eq:family-projectors}
P_a=\mathbf 1_{\mathrm{ext}}\otimes |e_a\rangle\langle e_a|.
\end{equation}
Then
\begin{equation}
\label{p4:eq:family-sector-decomposition}
\mathcal H
=
\bigoplus_{a=1}^n \bigl(\mathcal H_{\mathrm{ext}}\otimes \mathbb C e_a\bigr)
\cong
\bigoplus_{a=1}^n \mathcal H_{\mathrm{ext}},
\end{equation}
so the family label is an ordinary finite quantum subsystem.

\begin{proposition}[Finite family factors]
\label{p4:prop:finite-family-factors}
Under the factorization \eqref{p4:eq:family-factorization}:
\begin{enumerate}[label=(\roman*)]
\item The reduced family state of any density matrix $\rho$ is
\begin{equation}
\label{p4:eq:reduced-family-state}
\rho_{\mathrm{fam}}:=\Tr_{\mathrm{ext}}(\rho)\in \End(\mathcal H_{\mathrm{fam}}),
\end{equation}
and the sector probabilities satisfy
\begin{equation}
\label{p4:eq:family-sector-probs}
p_a=\Tr(\rho P_a)=\langle e_a,\rho_{\mathrm{fam}}e_a\rangle.
\end{equation}
\item Every self-adjoint operator $B\in\End(\mathcal H_{\mathrm{fam}})$ defines an admissible internal observable
\begin{equation}
\label{p4:eq:family-observables}
\widehat B:=\mathbf 1_{\mathrm{ext}}\otimes B,
\end{equation}
and every unitary $U_f\in U(\mathcal H_{\mathrm{fam}})$ defines an admissible internal basis change or mixing action
\begin{equation}
\label{p4:eq:family-unitaries}
\widehat U_f:=\mathbf 1_{\mathrm{ext}}\otimes U_f.
\end{equation}
\item If a Hamiltonian splits as
\begin{equation}
\label{p4:eq:family-hamiltonian-split}
H=H_{\mathrm{ext}}\otimes \mathbf 1 + \mathbf 1\otimes H_{\mathrm{fam}} + H_{\mathrm{int}},
\end{equation}
then $H_{\mathrm{fam}}$ governs the autonomous internal mixing in the family subsystem, while $H_{\mathrm{int}}$ couples family and external degrees of freedom.
\end{enumerate}
Hence a distinguished rank-$n$ family space is already a complete finite quantum subsystem inside the framework.
\end{proposition}

\begin{proof}
These are standard properties of tensor-product quantum systems already established in Part~I, here specialized to a distinguished finite factor.
Equation \eqref{p4:eq:family-sector-probs} is the Born rule after partial trace.
The operators in \eqref{p4:eq:family-observables} and \eqref{p4:eq:family-unitaries} are admissible because the full algebra contains all tensor-product operators.
The Hamiltonian interpretation follows by direct inspection of \eqref{p4:eq:family-hamiltonian-split}.
\end{proof}

\begin{corollary}[Rank-three family sectors]
\label{p4:cor:rank-three-family-sectors}
If $\dim \mathcal H_{\mathrm{fam}}=3$, then the framework already supports the full finite quantum structure of a three-level internal subsystem:
projective measurements in any family basis, unitary relative diagonalization, phase-sensitive interference, and the rephasing-invariant mixing data discussed in Section~\ref{p4:sec:part2-threelevel-mixing}.
\end{corollary}

\begin{proof}
If \(\dim\mathcal H_{\mathrm{fam}}=3\), then the finite observable algebra on that factor is \(M_3(\CC)\).  Hence it contains all rank-one projective measurements, all unitary changes of family basis, all diagonalizations of self-adjoint family observables, and all phase-sensitive transition data.  The rephasing-invariant mixing quantities are precisely the quotient data of this finite three-level system described earlier.
\end{proof}




\paragraph{Distinguished sector decompositions and transition probabilities}
\label{p4:sec:part2-sector-transitions}

Let
\[
\mathcal H=\bigoplus_{a=1}^r \mathcal H_a,
\qquad
P_a:\mathcal H\to\mathcal H_a
\]
be a finite orthogonal decomposition into distinguished sectors.
Write
\[
H = H_{\mathrm{diag}} + H_{\mathrm{mix}},
\qquad
H_{\mathrm{diag}}:=\sum_a P_a H P_a,
\qquad
H_{\mathrm{mix}}:=\sum_{a\neq b} P_a H P_b.
\]
If a state \(\rho\) is prepared in sector \(a\), i.e.
\[
\rho=P_a\rho P_a,
\]
then the exact probability to be found in sector \(b\) after time \(t\) is
\begin{equation}
\label{p4:eq:sector-transition-prob}
p_{b\leftarrow a}(t)
=
\Tr\!\bigl(P_b\,U(t)\rho\,U(t)^*\bigr),
\qquad
U(t)=e^{-itH}.
\end{equation}

\begin{proposition}[Finite sector-transition law]
\label{p4:prop:finite-sector-transition}
For any distinguished sector decomposition, the family
\(\{P_a\}\) defines an exact projective measurement of sector occupancy, and
all sector transitions are generated solely by the off-diagonal block
\(H_{\mathrm{mix}}\).
If \(H_{\mathrm{mix}}=0\), then every sector is invariant under the dynamics.
If \(H_{\mathrm{mix}}\) is weak and \(\ket{\alpha}\in\mathcal H_a\),
\(\ket{\beta}\in\mathcal H_b\) are eigenvectors of \(H_{\mathrm{diag}}\)
with eigenvalues \(E_\alpha,E_\beta\), then the first-order transition amplitude is
\begin{equation}
\label{p4:eq:sector-transition-amplitude}
\mathcal A^{(1)}_{\beta\leftarrow\alpha}(t)
=
-\frac{i}{\hbar}\int_0^t e^{is(E_\beta-E_\alpha)/\hbar}
\bra{\beta}H_{\mathrm{mix}}\ket{\alpha}\,ds.
\end{equation}
\end{proposition}

\begin{proof}
Equation \eqref{p4:eq:sector-transition-prob} is the Born rule for the orthogonal family
\(\{P_a\}\).
If \(H_{\mathrm{mix}}=0\), then \(P_aH=HP_a\) for all \(a\), so each \(\mathcal H_a\) is invariant.
The first-order formula is the Dyson expansion in the interaction picture with
respect to \(H_{\mathrm{diag}}\).
\end{proof}

\begin{remark}[Sector-domain generality]
\label{p4:rem:sector-transition-meaning}
This proposition is deliberately general.
The sectors may encode external channels, internal multiplets, family blocks,
or any other distinguished finite decomposition selected by the later physics.
The quantum-mechanical legitimacy of such sectors does not depend on their detailed interpretation.
\end{remark}



\paragraph{Threshold opening as a spectral-overlap criterion}
\label{p4:sec:part2-threshold-opening}

Threshold behavior in the framework is not fundamentally a statement about an infinite continuum.
It is a statement about when the prepared spectral support begins to overlap the
spectral support of an outgoing channel family over the finite scattering window.

Let \(P_{\mathrm{in}}(I)\) be the spectral projector of the incoming free
Hamiltonian for an energy interval \(I\), and let \(P_{\mathfrak c}\) project
onto an outgoing channel \(\mathfrak c\).  Assume that the channel projector
reduces the free Hamiltonian,
\([P_{\mathfrak c},H_0]=0\), and define the channel free Hamiltonian on its
actual channel space by
\[
H_{\mathfrak c}
:=
\left.
P_{\mathfrak c}H_0P_{\mathfrak c}
\right|_{P_{\mathfrak c}\mathcal H}.
\]
This restriction is essential: treating the compression as an operator on
all of \(\mathcal H\) would introduce spurious zero eigenvalues on
\(P_{\mathfrak c}^{\perp}\mathcal H\).  Let
\begin{equation}
\label{p4:eq:channel-threshold}
E_{\mathfrak c}^{\mathrm{thr}}
:=\min\mathrm{spec}(H_{\mathfrak c}).
\end{equation}

\begin{proposition}[Finite threshold criterion]
\label{p4:prop:finite-threshold}
Let the incoming state be spectrally concentrated in a narrow window
\(I=[E-\delta,E+\delta]\) of the free Hamiltonian \(H_0\), and let
\(\mathfrak c\) be an outgoing channel with threshold
\(E_{\mathfrak c}^{\mathrm{thr}}\).
If
\begin{equation}
\label{p4:eq:no-threshold-overlap}
I\cap \mathrm{spec}(H_{\mathfrak c})=\varnothing,
\end{equation}
then every first-order off-shell amplitude is bounded and oscillatory, so its
time-normalized amplitude and probability rate tend to zero.  It vanishes
identically whenever \(H_{\mathrm{mix}}\) has no matrix elements connecting
the incoming support to the channel support.

If \(I\cap\mathrm{spec}(H_{\mathfrak c})\ne\varnothing\), the channel is
only \emph{kinematically open}.  A nonzero transition additionally requires
a nonvanishing coupling matrix element compatible with the conserved charges
and selection rules.  Thus spectral overlap is necessary for an on-shell
finite-window contribution, not sufficient for an allowed transition.  The
number \(E_{\mathfrak c}^{\mathrm{thr}}\) is the first possible finite
channel threshold; continuum-style opening is recovered only in large-window
regimes where channel spectra become densely sampled.
\end{proposition}

\begin{proof}
The first-order amplitude is a sum of oscillatory integrals of the form
\eqref{p4:eq:first-born-amplitude} or \eqref{p4:eq:sector-transition-amplitude}.
Without spectral overlap, the energy differences are bounded away from zero,
so the explicit integral in
\eqref{p4:eq:first-born-amplitude} is uniformly bounded and its
time-normalized amplitude and rate vanish.  If the corresponding matrix
elements are absent, the amplitude vanishes identically.  With overlap,
stationary phases are possible, but the amplitude still contains the
coupling matrix element; symmetry may force it to zero.
\end{proof}

\begin{remark}[Effective textbook threshold language]
\label{p4:rem:effective-threshold-language}
In a large reciprocal position-side regime where the channel spectra become densely sampled,
Proposition~\ref{p4:prop:finite-threshold} reproduces the continuum threshold picture of
successive channel opening at threshold energies.
The exact framework statement is the conjunction of spectral overlap and a
nonzero symmetry-allowed coupling matrix element.  Neither condition alone
assumes a literal infinite continuum.
\end{remark}



\paragraph{Unitary mixing on distinguished finite multiplets}
\label{p4:sec:part2-unitary-mixing-multiplets}

Suppose now that a distinguished \(n\)-dimensional sector \(\mathcal H_f\subset\mathcal H\)
carries states with the same externally resolved quantum numbers, so that the observable distinction lies only in an internal label.
Any self-adjoint effective Hamiltonian or positive observable on \(\mathcal H_f\) is then an \(n\times n\) matrix and hence diagonalizable by a unitary basis change.

\begin{theorem}[Finite multiplet mixing]
\label{p4:thm:finite-multiplet-mixing}
Let \(\mathcal H_f\) be a distinguished \(n\)-dimensional sector and let
\(Y_f\in\End(\mathcal H_f)\) be normal.  Then there exists \(U_f\in U(n)\) and
eigenvalues \(\lambda_1,\dots,\lambda_n\) such that
\begin{equation}
\label{p4:eq:multiplet-diagonalization}
Y_f = U_f \,\mathrm{diag}(\lambda_1,\dots,\lambda_n)\, U_f^*.
\end{equation}
If two such observables \(Y_f\) and \(Z_f\) have simple spectra and are
diagonalized by \(U_f\) and \(V_f\), then the relative mixing matrix
\begin{equation}
\label{p4:eq:relative-multiplet-mixing}
W_f := U_f^*V_f
\end{equation}
is unitary.  All basis-independent mixing information is encoded in the
rephasing invariants of \(W_f\).  In the degenerate case the invariant data are
instead taken modulo unitary rotations inside the degenerate eigenspaces, not
merely modulo phases.  For \(n=3\), this is exactly the finite unitary mixing
situation analyzed in Section~\ref{p4:sec:part2-threelevel-mixing}.
\end{theorem}

\begin{proof}
Finite-dimensional spectral theory gives the diagonalizations.
The relative matrix \(W_f\) is unitary because it is a product of unitaries.
For simple spectra, rephasing of the chosen eigenbases multiplies rows and columns by phases without changing the underlying relative diagonalization data.  If a spectrum is degenerate, the same argument is replaced by quotienting by the full unitary freedom in the degenerate eigenspaces.
\end{proof}



\paragraph{Unstable internal blocks and effective widths}
\label{p4:sec:part2-unstable-blocks}

Combining the preceding subsections with the resonance analysis of Section~\ref{p4:sec:part2-feshbach-resonances},
one obtains the exact finite block template together with a conditional
open-system/continuum width interface for unstable internal sectors.
A distinguished block \(Q\mathcal H\) may carry an internal level structure or mixing matrix,
while couplings \(PHQ\) to channel blocks shift those levels.  A decay width
appears only after a stated reservoir, limiting-absorption, or Markov
coarse-graining domain has supplied irreversible channel loss.
Thus
\begin{equation}
\label{p4:eq:unstable-block-template}
\text{internal sector}
\;+\;
\text{channel coupling}
\;\Longrightarrow\;
\text{resonance shift} + \text{effective width}.
\end{equation}
This is the framework form of unstable-species phenomenology:
exactly unitary on the full finite algebra, resonant on the enlarged channel sector,
and irreversible only after channel elimination.



\subsection{Release v4.101: generated finite pairing-cell handoff}
\label{p4:sec:v4101-generated-pair-cell-handoff}

\begin{corollary}[Generated finite pairing member on an occurring quadratic owner]
\label{p4:cor:v4101-generated-pair-member-handoff}
Let $\mathcal K$ be a constituted finite CAR carrier and let an occurring
$\mathsf{DYN}$ owner use the predeclared quadratic Hamiltonian coordinate of
Eq.~\eqref{p4:eq:many-body-hamiltonian}.  Suppose an upstream rooted
constructor independently generates a specific fermionic pairing coefficient
\[
 \Delta_\star\in\Lambda^2\mathcal K^*,
 \qquad \Delta_\star^{\mathsf T}=-\Delta_\star,
 \qquad \Delta_\star\ne0,
\]
and a licensed owner-preserving support map places that generated member in
the pairing coordinate of the same $\mathsf{DYN}$ owner.  Then Paper~II-A,
Corollary~\ref{p2a:cor:v3109-finite-pair-cell-specialization}, gives
\[
 \boxed{
 H_{\rm pair}(\Delta_\star)
 =\frac12\bigl(
  \mathfrak a^\dagger\Delta_\star\mathfrak a^\dagger
  +\mathfrak a\Delta_\star^\dagger\mathfrak a
 \bigr)}
\]
as a co-realized required constituent of that owner, with no duplicate
occurrence or relative normalization token.

The conclusion is member-specific.  It does not assert that every available
$\Delta$ occurs, and it does not create a $\mathsf{DYN}$ owner from a
structural pairing coefficient.
\end{corollary}

\begin{proof}
The pairing coordinate and its Hermitian conjugate are already part of the
finite quadratic Hamiltonian schema.  The upstream theorem supplies a
specific nonzero generated member and the hypothesis supplies the typed
owner-preserving support map.  The complete pair cell therefore satisfies all
antecedents of Paper~II-A's generated-member handoff.  For fermionic CAR the
antisymmetric coefficient is the physical coefficient of the pair-creation
operator; if it is nonzero, the pair operator is nonzero on the finite Fock
carrier.  Generated-owner normalization inheritance excludes an additional
multiplicative factor absent from the frozen constructor.
\end{proof}

\begin{remark}[Exact-charge scope]
\label{p4:rem:v4101-exact-charge-scope}
The exact-charge no-asymmetry theorem remains unchanged on a quadratic owner
that commutes with the charge.  A generated pairing member carrying nonzero
charge changes that owner and therefore lies outside the commuting-owner
hypothesis.  A nonzero pairing block alone still does not supply a CP-odd
trajectory; Corollary~\ref{p4:cor:v498-phase-locked-pairing-invariant} retains
its separate phase-ownership requirement.
\end{remark}
% END GENERATED-FINITE-PAIR-CELL-HANDOFF v4.101


% BEGIN FINITE-PAIR-CHARGE-RESPONSE v4.102


\subsection{Release v4.102: generated pair-charge response and finite propagation}
\label{p4:sec:v4102-pair-charge-response}

The generated pairing handoff above supplies a physical pair member only on
its stated occurring-owner domain.  The following consequences retain that
domain, its kinetic metric, its metric connection, its clock, and its initial
state.  They introduce no independently normalized charge-production law.

\begin{theorem}[Exact charge source of an owned finite pair]
\label{p4:thm:v4102-pair-charge-source}
Let an occurring finite CAR Hamiltonian, in a parallel orthonormal frame for
its retained metric connection, be
\[
 H(t)=H_0(t)+\frac12\bigl(A(t)+A(t)^\dagger\bigr),
 \qquad A(t)=\mathfrak a^\dagger\Delta(t)\mathfrak a^\dagger,
 \qquad \Delta^{\mathsf T}=-\Delta.
\]
Let $Q=Q^\dagger$ be a covariantly constant dimensionless additive charge.
Assume $[Q,H_0(t)]=0$ and $[Q,A(t)]=q_{\rm pair}A(t)$ for all $t$,
with $q_{\rm pair}\in\mathbb R$.
Here $q_{\rm pair}$ is the charge of the pair-creation operator; it is the
negative of the compensating source/spurion charge.  Then every state
propagated by this same Hamiltonian obeys
\begin{equation}
\label{p4:eq:v4102-exact-pair-charge-source}
 \boxed{
 \frac{d}{dt}\langle Q\rangle_t
 =\frac{i}{\hbar}\langle[H,Q]\rangle_t
 =\frac{q_{\rm pair}}{\hbar}\operatorname{Im}\langle A(t)\rangle_t.}
\end{equation}
On the formal charge orbit
\[
 H_\theta=H_0+\frac12
 \bigl(e^{i\theta}A+e^{-i\theta}A^\dagger\bigr),
\]
evolution under this same $H_\theta$ obeys
\begin{equation}
\label{p4:eq:v4102-pair-phase-ward}
 \boxed{
 \frac{d}{dt}\langle Q\rangle_t
 =-\frac{q_{\rm pair}}{\hbar}
   \langle\partial_\theta H_\theta\rangle_t.}
\end{equation}
The orbit derivative is an algebraic derivative of the already retained
charge representation; it does not select or actualize a phase history.
\end{theorem}

\begin{proof}
Taking adjoints gives $[Q,A^\dagger]=-q_{\rm pair}A^\dagger$.
Hence $[H,Q]=q_{\rm pair}(A^\dagger-A)/2$.  The finite von Neumann equation
and cyclicity of the trace give the first identity.  Differentiating
$H_\theta$ gives $i(e^{i\theta}A-e^{-i\theta}A^\dagger)/2$, yielding the
second identity with the stated sign.  No kinetic or state approximation is
used.  In particular the exact accumulated charge is
\[
 \langle Q\rangle_f-\langle Q\rangle_i
 =\frac{q_{\rm pair}}{\hbar}
  \int_{t_i}^{t_f}\operatorname{Im}\langle A(t)\rangle_t\,dt.
\]
The source is a same-Hamiltonian commutator descendant, so Paper~II-A
generated co-consumption retains its coefficient and its conjugate jointly.
\end{proof}

\begin{theorem}[Dependent finite CAR charge propagation]
\label{p4:thm:v4102-finite-car-charge-propagation}
On the quadratic domain of
Theorem~\ref{p4:thm:curved-background-finite-bridge}, fix the occurring clocked
Hamiltonian, its metric connection, and an initial state.  In a parallel
orthonormal frame put
\[
 \Psi=\binom{\mathfrak a}{\mathfrak a^\dagger},\qquad
 \mathbb H=
 \begin{pmatrix}h&\Delta\\\Delta^\dagger&-h^{\mathsf T}\end{pmatrix},
 \qquad
 \mathcal R_{ij}=\langle\Psi_j^\dagger\Psi_i\rangle,
\]
where $h=h^\dagger$ and $\Delta^{\mathsf T}=-\Delta$.  The induced Nambu
connection gives the exact covariance equation and its unique solution
\begin{equation}
\label{p4:eq:v4102-covariance-propagation}
 \boxed{
 i\hbar D_t\mathcal R=[\mathbb H,\mathcal R],\qquad
 \mathcal R(t)=\mathcal U(t,t_i)\mathcal R_i\mathcal U(t,t_i)^\dagger,}
\end{equation}
where $i\hbar D_t\mathcal U=\mathbb H\mathcal U$ and
$\mathcal U(t_i,t_i)=I$.  For a covariantly constant one-particle charge
$q_1=q_1^\dagger$, let
$Q=d\Gamma(q_1)$ and $\mathbb Q=\operatorname{diag}(q_1,-q_1^{\mathsf T})$.
Then
\begin{equation}
\label{p4:eq:v4102-covariance-charge-readout}
 \boxed{
 \langle Q\rangle_t
 =\frac12\operatorname{Tr}(\mathbb Q\mathcal R(t))
  +\frac12\operatorname{Tr}q_1,\qquad
 \frac{d}{dt}\langle Q\rangle_t
 =\frac1{2i\hbar}
   \operatorname{Tr}([\mathbb Q,\mathbb H]\mathcal R(t)).}
\end{equation}
All traces are on the displayed finite one-particle or doubled Nambu carrier.
No further efficiency or density-matrix selector occurs after these owned
data have been fixed.
\end{theorem}

\begin{proof}
The CAR give $i\hbar D_t\Psi=\mathbb H\Psi$ for the Heisenberg mode column.
Differentiating $\langle\Psi_j^\dagger\Psi_i\rangle$ gives the commutator
equation.  Norm-continuous finite coefficients have a unique propagator, and
differentiation verifies its displayed covariance solution.  Finally the CAR
normal-ordering identity is
$Q=\Psi^\dagger\mathbb Q\Psi/2+\operatorname{Tr}q_1/2$.
Taking its expectation and differentiating proves the readout.  The theorem
holds for the covariance of any initial state; quasi-free initial data make
that covariance sufficient for all higher correlations as well.
\end{proof}

\begin{remark}[Exact domain and retained state data]
\label{p4:rem:v4102-charge-propagation-scope}
The covariance closure is quadratic.  With quartic interactions one retains
the exact finite many-body equation $i\hbar D_t\rho=[H,\rho]$ and
Theorem~\ref{p4:thm:v4102-pair-charge-source} when its charge hypotheses hold;
the covariance equation generally couples to higher correlations.  Neither
form supplies an unowned initial state, physical clock/history, bath
reduction, Markovian collision law, washout kernel, or entropy normalization.
An efficiency computed from such a completed evolution is an output
functional of those data, not a second normalization of the pair owner.
\end{remark}

\begin{corollary}[Covariant pairing trajectory diagnostic]
\label{p4:cor:v498-phase-locked-pairing-invariant}
Let a licensed pair history $\Delta$ carry its inherited kinetic metric and
induced metric connection $D_u$, where $u$ is its declared history parameter.
Using the inherited Hilbert--Schmidt pairing, the frame-independent
self-trajectory diagnostic is
\begin{equation}
\label{p4:eq:v4102-covariant-pair-diagnostic}
 \boxed{J_u:=\operatorname{Im}
 \langle\Delta,D_u\Delta\rangle_{\rm HS}.}
\end{equation}
On a connection-preserved pair line write
$\Delta(u)=a(u)e^{i\Theta(u)}\Delta_*(u)$ with $a\ge0$,
$\|\Delta_*\|_{\rm HS}$ constant, and
$D_u\Delta_*=-i\mathcal A_u\Delta_*$.  Then
\begin{equation}
\label{p4:eq:v4102-relative-pair-phase}
 \boxed{J_u=\|\Delta_*\|_{\rm HS}^2a^2
             (\Theta'-\mathcal A_u).}
\end{equation}
In a retained parallel frame $\mathcal A_u=0$, this reduces to
$\operatorname{Im}\operatorname{Tr}(\Delta^\dagger\partial_u\Delta)
=\|\Delta_*\|_{\rm HS}^2a^2\Theta'$.
In particular a purely radial history $\Delta=a\Delta_*$ with
$D_u\Delta_*=0$ has $J_u=0$.
\end{corollary}

\begin{proof}
Metric-compatible frame changes transport both $\Delta$ and $D_u\Delta$,
so their inherited inner product is invariant.  On the indicated line,
$D_u\Delta=e^{i\Theta}[a'+ia(\Theta'-\mathcal A_u)]\Delta_*$;
the imaginary part gives the formula.  The ordinary derivative expression
alone is invariant under constant frame changes, not arbitrary varying
ones: $\Delta\mapsto e^{i\chi(u)}\Delta$ adds
$\|\Delta\|_{\rm HS}^2\chi'$ to it.  The induced connection supplies the
compensating term.
\end{proof}

\begin{remark}[Phase diagnostic versus physical CP yield]
\label{p4:rem:v4102-relative-phase-scope}
The diagnostic retains the actual pair norm and kinetic metric; a family
coefficient norm cannot replace it without evaluating its support insertion.
A nonzero $J_u$ is not a complete CP invariant or a baryon-yield theorem,
and $J_u=0$ excludes only this self-trajectory phase diagnostic.  The full
action, relative sectors, initial state, and charge propagation still matter.
A structural phase and a pair coefficient acquire a physical relative phase
only through their licensed common-owner connection; the displayed formulas
do not create that binding.
\end{remark}
% END FINITE-PAIR-CHARGE-RESPONSE v4.102


% BEGIN INLINE SOURCE: Paper4-InheritedPairCharge-Main-v1_0.tex
% Framework v2.167: exact pair-insertion response, arbitrary inherited state.


\subsection{Pair-insertion charge response on the inherited state}
\label{p4:sec:v2167-inherited-pair-charge}
This is a specialization of the exact finite pair-charge source theorem.
Use a licensed canonical finite CAR chart, its complete state, the total
charge $Q=-N_{\rm lep}$, a numerical antisymmetric pair coefficient
$\Gamma(t)$, and number-preserving quadratic normal evolution $V$.
All inherited constraints and intermediate projectors must be retained
where the physical word consumes them. The native cyclic application uses
the already inherited $\omega_{16}$; it does not prepare another vacuum.

\begin{lemma}[Exact CAR pair commutator]
\label{p4:lem:v2167-pair-commutator}
For antisymmetric matrices $D,E$ and
$\mathsf P(D)=\frac12\sum_{ij}D_{ij}a_i^\dagger a_j^\dagger$,
\[
 [\mathsf P(D),\mathsf P(E)^\dagger]
 =a^\dagger DE^\dagger a-\tfrac12\operatorname{Tr}(DE^\dagger)I.
\]
\end{lemma}
\begin{proof}
Expand the two quadratic products and use the canonical
anticommutators. Their quartic terms cancel; the remaining bilinear and
central terms are the displayed contractions. This is an operator
identity and uses no Wick factorization.
\end{proof}

\begin{theorem}[One- and two-pair coefficients for an arbitrary state]
\label{p4:thm:v2167-inherited-pair-coefficients}
In the normal interaction picture put
\[
 D_I(t)=V(t,t_0)^\dagger\Gamma(t)\overline{V(t,t_0)},\qquad
 \mathcal B=\hbar^{-1}\int D_I(t)\,dt,
\]
\[
 C_{ij}=\omega(a_j^\dagger a_i),\qquad
 F_{ij}=\omega(a_i^\dagger a_j^\dagger).
\]
The coefficients with one and two pair insertions are
\[
 \boxed{q_1=-4\operatorname{Im}\omega(\mathsf P(\mathcal B))
 =-2\operatorname{Im}\sum_{ij}\mathcal B_{ij}F_{ij},}
 \qquad
 \boxed{q_2=\operatorname{Tr}[(2C-I)\mathcal B\mathcal B^\dagger].}
\]
No Gaussian-state hypothesis is needed. In particular
$|q_2|\le\|\mathcal B\|_{\rm HS}^2$.
\end{theorem}
\begin{proof}
With $H_{P,I}=\mathsf P(D_I)+\mathsf P(D_I)^\dagger$,
$[Q,\mathsf P]=-2\mathsf P$ gives the first Dyson coefficient.
The second nested commutator, using the preceding lemma, gives
\[
 \frac4{\hbar^2}\int_{t>s}
 \operatorname{Re}\operatorname{Tr}
 [D_I(s)D_I(t)^\dagger(C-I/2)]\,ds\,dt.
\]
Since $C$ is Hermitian the real integrand is symmetric in $s,t$.
The ordered triangle is half the square, and the two integrated pair
factors yield the stated $q_2$. The CAR bound $0\le C\le I$ gives its norm
bound. Only the exact pair commutator was used, so non-Gaussian inherited
states are included.
\end{proof}
An empty comparison has $C=F=0$, $q_1=0$ and
$q_2=-\|\mathcal B\|_{\rm HS}^2$. At $C=I/2$ the second coefficient
vanishes, while anomalous moments can still enter the first. These are
response identities, not state-selection rules.
One and two count pair insertions with the complete normal propagation
retained. A formal parameter gives
$\Delta Q=\lambda q_1+\lambda^2q_2+\cdots$; it does not assert the onset
order of the complete native theory or control truncation at its fixed
generated pair normalization.

\begin{corollary}[Exact returned pair kernel]
\label{p4:cor:v2167-returned-pair-kernel}
For an actual canonical pair insertion
$\Gamma(t)=J_R(t)\Delta_R(t)J_R(t)^{\mathsf T}$,
\[
 \begin{aligned}
 &\operatorname{Tr}[\Gamma(t)^\dagger V(t,s)\Gamma(s)V(t,s)^{\mathsf T}]\\
 &\quad=\operatorname{Tr}[\Delta_R(t)^\dagger V_R(t,s)
             \Delta_R(s)V_R(t,s)^{\mathsf T}],
 \end{aligned}
\]
where $V_R(t,s)=J_R(t)^\dagger V(t,s)J_R(s)$. No unitarity of $V_R$ is
required.
\end{corollary}
\begin{proof}
Substitute both pair insertions and use cyclicity of trace.
\end{proof}
For a normal owner containing quartic or quantum boson--fermion
interactions, retain its actual many-body interaction picture. The exact
first coefficients are instead evaluated as
\[
 q_1=\frac i\hbar\int\omega([H_{P,I}(t),Q])\,dt,
 \qquad
 q_2=-\frac1{\hbar^2}\int_{t>s}
 \omega([H_{P,I}(s),[H_{P,I}(t),Q]])\,ds\,dt.
\]
The closed matrix formula above is their quadratic-normal specialization,
not a covariance closure of the full interacting theory.

% END INLINE SOURCE: Paper4-InheritedPairCharge-Main-v1_0.tex




% BEGIN INLINE SOURCE: Paper4-StationaryChargedResponse-Main-v1_0.tex
% Framework v2.168: integrated finite charged-source derivations.


\subsection{Complete stationary charged-source dynamics and exact fidelity tests}
\label{p4:sec:v2168-stationary-charged-response}
Let $M=\diag(\varepsilon_e,\varepsilon_\mu,\varepsilon_\tau)>0$ retain the fixed terminal coefficients generated by V-B. Use the complete finite physical $H=H^\dagger$, its actual
odd source ports, and an admitted stationary state/constraint chart.
Every consumed interaction, central energy, kinetic and reduction contact
is part of this $H$. Finite mode rank alone does not make a bosonic Fock
domain finite; the occupation domain must itself be supplied by the
licensed constructor. A nonstationary history uses the complete two-time
calculus instead of the stationary specialization below.


\paragraph{A positive source Hilbert space for the complete dynamics}
Let $H=H^\dagger$ act on the admitted finite carrier and let
$\rho\succeq0$, $\Tr\rho=1$, $[H,\rho]=0$ on this source chart. No
faithfulness or purity of $\rho$ is assumed. On the complete operator algebra,
with the retained odd source ports, define
\[
 (X,Y)_\rho:=\Tr\rho\{X^\dagger,Y\},\qquad
 \mathcal L X:=[H,X].
\]
The energy-valued generator $\mathcal L$ retains the inherited clock:
time evolution uses $e^{-it\mathcal L/\hbar}$.

\begin{proposition}[State quotient and self-adjoint source dynamics]
\label{p4:prop:v2168-stationary-source-quotient}
The pairing is positive semidefinite. Its null space is invariant under
$\mathcal L$, and the induced operator on the quotient
$\mathscr O_\rho$ is finite and self-adjoint. For the native normalized CAR
sources,
\[
 \{q_r^\dagger,q_s\}=\delta_{rs}I,
 \qquad (q_r,q_s)_\rho=\delta_{rs}.
\]
Thus the source embedding
$\mathcal Q:\mathbb C^3\to\mathscr O_\rho$, $\mathcal Qe_r=[q_r]$, is isometric.
It is a readout of the supplied source ports, not a physical family selector.
\end{proposition}
\begin{proof}
The norm is
\[
 \|X\|_\rho^2
 =\|X\rho^{1/2}\|_{\rm HS}^2+
   \|X^\dagger\rho^{1/2}\|_{\rm HS}^2.
\]
Both terms are nonnegative. Since $[H,\rho^{1/2}]=0$, their simultaneous
vanishing implies the same for $[H,X]$ and its adjoint. The null space is
therefore invariant. Trace cyclicity and $[H,\rho]=0$ give
$(X,\mathcal LY)_\rho=(\mathcal LX,Y)_\rho$. The induced operator on the
finite quotient is self-adjoint. The CAR identity gives the source Gram
without selecting a state or inserting a new source normalization.
\end{proof}

If an actual constrained source chart has a different source Gram, that Gram
and its inherited normalization must be transported first. The displayed CAR
condition is a source-domain antecedent, not a claim that an arbitrary
constraint projection preserves CAR.



\paragraph{The exact finite two-point function and its memory block}
The stationary fermionic anticommutator response is
\[
 \boxed{G(z)=\mathcal Q^\dagger(zI-\mathcal L)^{-1}\mathcal Q.}
\]
For $\operatorname{Im}z>0$ its entries are equivalently
\[
 G_{rs}(z)=-\frac{i}{\hbar}\int_0^\infty e^{izt/\hbar}
 \Tr\rho\{q_r^\dagger,e^{-it\mathcal L/\hbar}q_s\}\,dt.
\]
This is a fixed-source stationary specialization, not a replacement for
moving-frame, energy-dependent-source or nonstationary contacts.

One may calculate on the source-generated cyclic subspace
$\operatorname{span}\{\mathcal L^n\mathcal Qx:n\ge0\}$. It is reducing
because $\mathcal L$ is finite and self-adjoint. Its orthogonal complement
has zero source response and cannot re-enter that response. This is an exact
calculation restriction of the complete Hamiltonian, not deletion of a
physical interaction or a new independent carrier.

With $\mathbb P=\mathcal Q\mathcal Q^\dagger$, decompose the source cyclic
space into $\operatorname{ran}\mathcal Q\oplus\operatorname{ran}(I-\mathbb P)$.
Its actual blocks are
\[
 A=\mathcal Q^\dagger\mathcal L\mathcal Q,
 \quad B=(I-\mathbb P)\mathcal L\mathcal Q,
 \quad C=(I-\mathbb P)\mathcal L(I-\mathbb P).
\]
Finite block elimination gives, whenever the displayed inverses exist,
\[
 \boxed{G(z)=\bigl[zI-A-\Sigma(z)\bigr]^{-1},\qquad
 \Sigma(z)=B^\dagger(zI-C)^{-1}B.}
\]
The self-energy is the exact memory of the supplied $H,\rho,q_r$. It is not
a free counterterm or a phenomenological input. At a singular eliminated
complement use a larger chart, as in III-F's finite pole theorem.



\paragraph{An exact dynamical replacement for the free source-lift test}
Put
\[
 \mu_n:=\mathcal Q^\dagger\mathcal L^n\mathcal Q,
 \qquad D:=M^2,\qquad
 \mathcal E_r:=[H,[H,q_r]]-\varepsilon_r^2q_r.
\]
Let $\Xi=\mathcal L^2\mathcal Q-\mathcal QD$. Its $r$th column is the
state-quotient class of $\mathcal E_r$.

\begin{theorem}[Complete finite dynamical fidelity certificate]
\label{p4:thm:v2168-dynamical-fidelity}
On the stated stationary source domain the following are equivalent:
\begin{enumerate}
\item The whole tagged response is
\[
 G(z)=\diag_r\frac{z}{z^2-\varepsilon_r^2}.
\]
\item $\mu_1=0$ and $\Xi=0$ on $\mathscr O_\rho$.
\item $\mu_1=0$, $\mu_2=D$, and $\mu_4=D^2$.
\item $A=0$, $B^\dagger B=D$, and $CB=0$ in the exact memory block.
\end{enumerate}
In particular the dynamically generated partner ports
$p_r=[H,q_r]/\varepsilon_r$ obey
\[
 \boxed{\mathcal L[q_r]=\varepsilon_r[p_r],\qquad
        \mathcal L[p_r]=\varepsilon_r[q_r]}
\]
when the certificate holds. Their classes and the source classes are
orthonormal. No particular right family frame is chosen.
\end{theorem}
\begin{proof}
For large $z$, $G(z)=\sum_{n\ge0}\mu_n z^{-n-1}$. The response in (1)
therefore implies the moments in (3). Self-adjointness gives the exact
positive identity
\[
 \boxed{\Xi^\dagger\Xi
   =\mu_4-\mu_2D-D\mu_2+D^2\succeq0.}
\]
The moments in (3) make this matrix zero, hence $\Xi=0$.
Conversely $\Xi=0$ gives $\mathcal L^2\mathcal Q=\mathcal QD$ and
$\mu_2=D$. Polynomial multiplication then gives
\[
 (zI-\mathcal L)^{-1}\mathcal Q
   =(z\mathcal Q+\mathcal L\mathcal Q)(z^2I-D)^{-1}.
\]
Compression with $\mathcal Q^\dagger$ and $\mu_1=0$ proves (1).

In source/complement blocks, $A=0$ gives
$\mathcal L\mathcal Q=(0,B)^T$ and
$\mathcal L^2\mathcal Q=(B^\dagger B,CB)^T$.
Thus $\Xi=0$ is precisely $B^\dagger B=D$ and $CB=0$, proving (4).
The partner classes are $\mathcal L\mathcal Q M^{-1}$; their source cross
Gram is zero and their own Gram is $M^{-1}DM^{-1}=I$. The double
commutator supplies their second evolution equation.
\end{proof}

\begin{corollary}[Positive leakage test; no cancellation of hidden channels]
\label{p4:cor:v2168-positive-source-leakage}
If $\mu_1=0$ and $\mu_2=D$, then
\[
 \boxed{\mu_4-D^2=(CB)^\dagger(CB)\succeq0.}
\]
More generally the scalar test
\[
 \boxed{\Tr_{\rm role}(\Xi^\dagger\Xi)
 =\sum_r\Tr\rho\{\mathcal E_r^\dagger,\mathcal E_r\}=0}
\]
together with $\mu_1=0$ certifies the whole response. At a pure stationary
state $\rho=|\Omega\rangle\langle\Omega|$, this requires exactly
$\mathcal E_r\Omega=0$ and $\mathcal E_r^\dagger\Omega=0$ for every $r$.
At a faithful state it requires $\mathcal E_r=0$ as operators.
\end{corollary}
\begin{proof}
Under the first two moment conditions, the block calculation in the theorem
has first component $B^\dagger B-D=0$ and second component $CB$. The general
scalar identity is the trace of the positive Gram. A positive finite Gram
has zero trace exactly when every column is null. Apply the two Hilbert--
Schmidt norm formula of Proposition~\ref{p4:prop:v2168-stationary-source-quotient} to each
residual.
\end{proof}

The fourth moment is essential: the first two moments can agree while $CB$
is nonzero. The source-dynamics diagnostic appendix gives an exact finite witness. This does not require
all interactions to vanish; it requires that the complete dynamics have no
additional response channel on the stated source/state domain.



\paragraph{Closure on a generated stationary family without choosing a member}
Let $\mathfrak S_H=\{\rho\succeq0:\Tr\rho=1,[H,\rho]=0\}$ be IV's complete
stationary density family on the admitted carrier. Let
\[
 H=\sum_\lambda E_\lambda P_\lambda,
 \qquad \mathfrak P_H(X)=\sum_\lambda P_\lambda XP_\lambda,
 \qquad K_{rs}=\{q_r^\dagger,[H,q_s]\}.
\]

\begin{theorem}[Common whole-response certificate on the stationary family]
\label{p4:thm:v2168-stationary-family-fidelity}
The target whole response holds for every $\rho\in\mathfrak S_H$ if and only if
\[
 \boxed{\mathcal E_r=0\ \text{for all }r,\qquad
        \mathfrak P_H(K_{rs})=0\ \text{for all }r,s.}
\]
If these are proved from the actual word, Paper II-A common-family readout
closure retains the response without selecting a stationary member.
\end{theorem}
\begin{proof}
If the response holds throughout the family, it holds for the faithful
comparison state $I/\dim\mathcal H$, so the positive residual test implies
$\mathcal E_r=0$. The first moment vanishes for every stationary density.
Within each $H$ eigenspace every density matrix is stationary. Pure-state
expectations and polarization therefore force each block
$P_\lambda K_{rs}P_\lambda$ to vanish. Conversely these block identities
give $\Tr\rho K_{rs}=0$ throughout the stationary family, and the operator
residual identities give $\Xi=0$ for each state. The previous theorem applies.
\end{proof}

The faithful trace state is used only to prove the equivalence; it is not
selected as the physical vacuum. A physical stationary subfamily may satisfy
the weaker state-quotient test without these whole-family operator identities.
For an actual nonstationary history use the full two-time kernel, including
its source/frame jets; this theorem does not supply stationarity.



\paragraph{Free reduction and the exact interaction terms that must be evaluated}
On the admitted constant quadratic clock comparison let
\[
 h(R)=\begin{pmatrix}0&R\\R^\dagger&0\end{pmatrix},\qquad
 J_L=\binom{U}{0},\qquad
 H_0=d\Gamma(h(R)),\qquad q_r=a^\dagger(J_Le_r).
\]
Here $h(R)$ is the existing finite clock comparison symbol; it is not a new
primitive rest-momentum preparation. This second-quantized formula realizes
the comparison algebraically; its six transfer coordinates are not asserted
to be the complete physical field carrier. Exact CAR gives
$[H_0,a^\dagger(j)]=a^\dagger(h(R)j)$, and therefore, on every stationary
state of this quadratic system,
\[
 \mu_n=J_L^\dagger h(R)^nJ_L,
 \qquad \mu_1=0,\qquad \mu_2=U^\dagger RR^\dagger U.
\]
With $\Delta=RR^\dagger-UDU^\dagger$,
\[
 \mathcal E_r^{(0)}=a^\dagger\!\left(\binom{\Delta Ue_r}{0}\right),
 \qquad
 \boxed{\Xi^\dagger\Xi=U^\dagger\Delta^2U,
        \Tr_{\rm role}(\Xi^\dagger\Xi)=\|\Delta\|_{\rm HS}^2.}
\]
Thus the native dynamic test is exactly the the free-return criterion Gram test on this branch.
The explicit V-B transfer lift has $\Delta=0$, so all stationary states of that
constructed quadratic comparison have the target response. This statement
does not identify its $R$ with $R_{\rm act}$.

For an actual action-owned decomposition $H=H_0+V$, retain the native
generated quadratic and interaction words with their inherited kinetic/source
transport. Define
$\mathcal E_r^{(0)}=[H_0,[H_0,q_r]]-\varepsilon_r^2q_r$.
On the preceding reducing quadratic clock chart it is the displayed Gram
defect; other consumed free sectors stay in $H_0$. Direct expansion gives
\[
 \boxed{\begin{aligned}
 \mathcal E_r={}&\mathcal E_r^{(0)}
   +[H_0,[V,q_r]]+[V,[H_0,q_r]]+[V,[V,q_r]],\\
 K_{rs}={}&\{q_r^\dagger,[H_0,q_s]\}
            +\{q_r^\dagger,[V,q_s]\}.
 \end{aligned}}
\]
These are the exact next expressions to evaluate. Terms are not discarded by
grade, perturbative order or an assumed zero self-energy. The total residual
can vanish on a state domain even if its free part does not; no equality to
one particular geometric right frame is then required. The actual generated $V$ must prove
that cancellation. It cannot be chosen to enforce the displayed terminal
values.


\paragraph{A physical pole needs less than the whole free response}
The target of the preceding certificate is the entire tagged free response,
including its residue $1/2$ and absence of additional poles. A physical
interacting mass comparison can instead consume only one visible pole; it
need not satisfy the stronger certificate.

Let $\Pi_\omega$ be the spectral projectors of the induced self-adjoint
$\mathcal L$. Then
\[
 \boxed{G(z)=\sum_\omega\frac{Z_\omega}{z-\omega},\qquad
 Z_\omega=\mathcal Q^\dagger\Pi_\omega\mathcal Q\succeq0,
 \qquad \sum_\omega Z_\omega=I.}
\]
Its possible frequencies are the Bohr differences of the complete $H$.
The exact finite particle-tag condition at $+\varepsilon_r$ is
\[
 \boxed{\|\Pi_{\varepsilon_r}[q_r]\|_\rho^2>0.}
\]
If the stated target requires both signs, test both. This is equivalent to a
visible solution of $[H,O]=\varepsilon_r O$ in the state quotient. The operator
space and the source overlap are fixed by the actual action; no spectral
mass is used to redefine the terminal coefficient.

On an admitted memory chart with $\varepsilon_r\notin\spec(C)$, let
$S(z)=zI-A-B^\dagger(zI-C)^{-1}B$. For a simple source-channel zero,
$\ker S(\varepsilon_r)=\mathbb C u$ with $u^\dagger u=1$,
\[
 a=u^\dagger S'(\varepsilon_r)u
   =1+\|(\varepsilon_r I-C)^{-1}Bu\|^2>0,
 \qquad
 \boxed{\operatorname*{Res}_{z=\varepsilon_r}G(z)=uu^\dagger/a.}
\]
The tagged residue is $|e_r^\dagger u|^2/a$. This is the concrete native
Hamiltonian specialization of III-F's finite pole/source theorem; the
derivative norm is retained. A singular complement requires a larger chart,
and an energy-dependent source keeps its full transported source jet.

This permits a dynamic finite pole theorem even when the free source Gram
does not equal $UDU^\dagger$, provided the actual memory has the required
zero and visible residue. It supplies no justification for omitting that
memory or for equating finite poles with charged infrared or continuum poles.


% END INLINE SOURCE: Paper4-StationaryChargedResponse-Main-v1_0.tex


% BEGIN INLINE SOURCE: Paper4-ChargedMemorySpectrum-Main-v1_0.tex
% Framework v2.168: integrated finite charged-source derivations.


\subsection{Complete-action support stationarity and native Slater leakage}
\label{p4:sec:v2169-state-certificate}
\paragraph{Result and actual domain.}
The stationary-support input of Theorem~\ref{p4:v2169:b10:thm:support} can be tested directly on the
existing state constructor. The complete normal word supplies all leakage
coefficients; no additional preparation is needed. The native weak mark's
leakage on a Slater support is evaluated explicitly, including all external
slot cross terms. This matters because III-F's free clock construction fills
negative-energy modes: its ground is generally not the empty fermion unit.
The existing fresh-unit theorem is reused only within its own branch.
No support is chosen to obtain a terminal pole.



\paragraph{An exact normal-word test on the occurring support}
On a licensed finite tensor CAR chart let $\mathcal F=\Lambda\mathcal K$
and let $\mathcal B$ be the actual finite bosonic carrier. Relative to an
already generated fermionic Slater vector, use the induced particle--hole CAR
$b_i$: $b_i=a_i$ on unoccupied modes and $b_i=a_i^\dagger$ on occupied modes.
Their unit vector $\Omega$ is that same Slater vector. This is a reversible
presentation of the actual state, not a new Bogoliubov preparation or a
mass-generating diagonalization. For the empty branch simply $b_i=a_i$.

Choose the ordered exterior basis $|I\rangle=b_I^\dagger\Omega$ and write the
complete finite Hamiltonian uniquely as
\begin{equation}
 H=\sum_{I,J}b_I^\dagger b_J\otimes B_{IJ},\qquad
 H=H^\dagger,\qquad H_B^{(0)}=B_{\varnothing\varnothing}.
 \label{p4:eq:v2169-b11-l1}
\end{equation}
Here $b_J$ is the adjoint of $b_J^\dagger$. The $4^{\dim\mathcal K}$ normal
CAR monomials form a basis of the fermion matrix algebra; all action,
occupation and contact terms are included in their actual coefficients.
On a projected chart not preserving this tensor CAR representation, the
complete finite support block is used instead of asserting \eqref{p4:eq:v2169-b11-l1}.

\begin{theorem}[Action-generated stationary-support certificate]
\label{p4:v2169:b11:thm:normalstate}
Let the actual bosonic density be $\sigma$, with exact support $P_B$, and
$P=|\Omega\rangle\langle\Omega|\otimes P_B$. Then
\begin{equation}
 \boxed{[H,P]=0\ \Longleftrightarrow\
 [H_B^{(0)},P_B]=0,\quad B_{I\varnothing}P_B=0\ (I\ne\varnothing).}
 \label{p4:eq:v2169-b11-l2}
\end{equation}
For the actual density $\rho=|\Omega\rangle\langle\Omega|\otimes\sigma$,
\begin{equation}
 \boxed{[H,\rho]=0\ \Longleftrightarrow\
 [H_B^{(0)},\sigma]=0,\quad B_{I\varnothing}P_B=0\ (I\ne\varnothing).}
 \label{p4:eq:v2169-b11-l3}
\end{equation}
The creation leakage is the generated positive operator
\begin{equation}
 \boxed{\mathcal E_{\rm cr}
 =\sum_{I\ne\varnothing}P_BB_{I\varnothing}^\dagger
                  B_{I\varnothing}P_B.}
 \label{p4:eq:v2169-b11-l4}
\end{equation}
It vanishes precisely when all creation tests in \eqref{p4:eq:v2169-b11-l2}--\eqref{p4:eq:v2169-b11-l3} hold.
For a pure bosonic ray $b$, the complete energy eigenstate condition is
$H_B^{(0)}b=E_*b$ and $B_{I\varnothing}b=0$ for every nonempty $I$.
\end{theorem}
\begin{proof}
An annihilation word with $J\ne\varnothing$ kills $\Omega$, so
\[
 H(\Omega\otimes b)=\Omega\otimes H_B^{(0)}b+
            \sum_{I\ne\varnothing}|I\rangle\otimes B_{I\varnothing}b.
\]
The distinct exterior occupation vectors are orthonormal. Thus $(I-P)HP=0$
is exactly the listed creation tests and invariance of $P_B$ under $H_B^{(0)}$.
Self-adjointness supplies the reverse off-diagonal block. Compression of
$[H,\rho]$ gives $[H_B^{(0)},\sigma]$, while its other blocks give
$B_{I\varnothing}\sigma=0$. Strict positivity of $\sigma$ on its exact
support makes this equivalent to $B_{I\varnothing}P_B=0$.
Finally each summand of \eqref{p4:eq:v2169-b11-l4} is a square, so its zero expectation on every
support vector is equivalent to each coefficient killing the support.
\end{proof}

The coefficients in \eqref{p4:eq:v2169-b11-l4} are summed \emph{after} assembling the complete
action. If $B_{I\varnothing}=\sum_\alpha B^\alpha_{I\varnothing}$, its square
retains every $\alpha,\beta$ cross term. Positivity of one mark alone does
not exclude an actual cancellation with another mark. These are evaluated
state equations, not a requirement to resupply an already occurring state.




% BEGIN FRAMEWORK v2.239: GIBBS/SUPPORT SCOPE
\begin{remark}[Stationary Gibbs support is not Slater leakage cancellation]
\label{p4:rem:v239-gibbs-support-scope}
For any finite self-adjoint $H$ and $0<\beta<\infty$,
$\rho_\beta=e^{-\beta H}/\Tr(e^{-\beta H})$ commutes with $H$,
but its exact support is the identity of the full finite carrier.
Consequently the identity $[H,P_{\rho_\beta}]=0$ does not prove
$B_{I\varnothing}P_B=0$ for a separately specified free-coordinate
Slater support. For the two-state witness
$H=\left(\begin{smallmatrix}0&g\\\bar g&\Delta\end{smallmatrix}\right)$
with $g\ne0$ and $P_0=|0\rangle\langle0|$, one has
$[H,\rho_\beta]=0$ but
$\|(I-P_0)HP_0\|_{\rm HS}^2=|g|^2>0$.
A \emph{ground-energy spectral projection} $P_*$ instead satisfies
$HP_*=E_*P_*$ and is an admissible single-energy support once the
complete action and that physical state have been licensed; it need
not agree with the free Slater support. Finite-temperature Gibbs
stationarity cannot be substituted into the adjacent-degree formula
requiring $HP=E_*P$.
\end{remark}
% END FRAMEWORK v2.239: GIBBS/SUPPORT SCOPE

\paragraph{The native weak mark on a supplied Slater state}
Let $F_e,F_\nu$ be the occupied one-particle projections of an actual
charge-separated Slater vector on its full charged and neutral slot spaces.
An individual left-current compression need not itself be a projection;
these $F$ are the complete occupied subspace projections. Let
$T_\alpha:\mathcal K_e\to\mathcal K_\nu$ be the existing external/current
map, including $\overline C_{rj}$ and $K_\alpha$ in the native vertex. Put
\begin{equation}
 V_W=\pi\sum_\alpha
   \bigl(a_\nu^\dagger T_\alpha a_e\otimes X_{\alpha,-}
        +a_e^\dagger T_\alpha^\dagger a_\nu\otimes X_{\alpha,+}\bigr),
 \qquad X_{\alpha,+}=X_{\alpha,-}^\dagger.
 \label{p4:eq:v2169-b11-l5}
\end{equation}
Define the action-determined particle--hole maps
\[
 M_\alpha=(I-F_\nu)T_\alpha F_e,\qquad
 N_\alpha=(I-F_e)T_\alpha^\dagger F_\nu .
\]

\begin{theorem}[Complete native weak particle--hole leakage]
\label{p4:v2169:b11:thm:weakslater}
The weak mark has no component along this Slater vector and its generated
leakage on $P_B$ is
\begin{equation}
 \boxed{\mathcal E_W=\pi^2\sum_{\alpha\beta}
 \left[
 \Tr(M_\alpha^\dagger M_\beta)P_BX_{\alpha,+}X_{\beta,-}P_B
 +\Tr(N_\alpha^\dagger N_\beta)P_BX_{\alpha,-}X_{\beta,+}P_B
 \right]\succeq0.}
 \label{p4:eq:v2169-b11-l6}
\end{equation}
It vanishes exactly when $V_W$ kills that full Slater--boson support.
No Wick approximation, bosonic CCR or product assumption within the
bosonic density is used.
\end{theorem}
\begin{proof}
For a Slater vector, $a^\dagger(u)a(v)\Omega$ is zero if $v$ is unoccupied
or $u$ is an already occupied distinct direction. Its nontrivial component
replaces an occupied vector by an unoccupied vector. Inner products of such
replacements are the Hilbert--Schmidt products of $(I-F)TF$.
In \eqref{p4:eq:v2169-b11-l5}, charged and neutral one-particle spaces are distinct, so
$\Tr(FT_\alpha)=0$. The forward replacements change $(N_e,N_\nu)$ by
$(-1,+1)$ and the reverse replacements by $(+1,-1)$; the two output
families are orthogonal. Pair their actual bosonic coefficient words in their
displayed order. This proves \eqref{p4:eq:v2169-b11-l6}, and also realizes it as the Gram of
$V_W(\Omega\otimes P_B)$, proving positivity and the zero criterion.
\end{proof}

For a common three-family constituent $T=C^\dagger$, $X_{\alpha,-}=X$,
write
\[
 \eta_\to=\|(I-F_\nu)C^\dagger F_e\|_{\mathrm{HS}}^2,\qquad
 \eta_\leftarrow=\|(I-F_e)CF_\nu\|_{\mathrm{HS}}^2.
\]
Equation \eqref{p4:eq:v2169-b11-l6} is now evaluated as
\begin{equation}
 \boxed{\mathcal E_W=\pi^2[
       \eta_\to P_BX^\dagger XP_B+
       \eta_\leftarrow P_BXX^\dagger P_B].}
 \label{p4:eq:v2169-b11-l7}
\end{equation}
If both bosonic coefficient expectations are strictly positive in the
actual faithful-on-$P_B$ density, this vanishes iff
\begin{equation}
 \boxed{F_\nu=C^\dagger F_eC.} \label{p4:eq:v2169-b11-l8}
\end{equation}
Indeed the two zero maps give both inclusions of occupied subspaces.
The empty and completely filled constituents obey \eqref{p4:eq:v2169-b11-l8}, as do current-matched
occupied subspaces. This does not assert that the actual free Dirac
ground has the matched occupancies. For a one-sided null bosonic port,
only its remaining map is required; \eqref{p4:eq:v2169-b11-l6}, not \eqref{p4:eq:v2169-b11-l8}, is the exact criterion.



\paragraph{An evaluated leakage on the filled chiral clock branch}
There is a more direct calculation for the filled-negative-mode construction
already used by III-F. Suppose the \emph{actual} invertible charged source
has the admitted chiral one-body symbol
\[
 h_R=\begin{pmatrix}0&R\\R^\dagger&0\end{pmatrix}.
\]
Its negative spectral projection is
\begin{equation}
 F_e=\tfrac12\begin{pmatrix}I&-S\\-S^\dagger&I\end{pmatrix},
 \qquad S=R(R^\dagger R)^{-1/2}.
 \label{p4:eq:v2169-b11-l9}
\end{equation}
In particular the occupied left-current compression is exactly $I_3/2$.
This follows from the actual off-diagonal symbol and invertibility; it does
not consume the unproved terminal-return certificate
$RR^\dagger=UM^2U^\dagger$.

Let $E_e,E_\nu$ be the occurring left-current isometries and let the native
common-slot weak map be $T=E_\nu C^\dagger E_e^\dagger$. For the actual
neutral occupied projection put $F_{\nu,L}=E_\nu^\dagger F_\nu E_\nu$.
This compression is not assumed to be a projection. The Slater leakage
coefficients are now explicitly
\begin{equation}
 \boxed{\eta_\to=\tfrac12\Tr(I_3-F_{\nu,L}),\qquad
        \eta_\leftarrow=\tfrac12\Tr F_{\nu,L}.} \label{p4:eq:v2169-b11-l10}
\end{equation}
For an actual bosonic density $\sigma$, write
$\kappa_-=\Tr\sigma X^\dagger X$ and $\kappa_+=\Tr\sigma XX^\dagger$. Then
\begin{equation}
 \boxed{\Tr(\sigma\mathcal E_W)
  =\frac{\pi^2}{2}
    [(3-\Tr F_{\nu,L})\kappa_-+\Tr F_{\nu,L}\kappa_+].} \label{p4:eq:v2169-b11-l11}
\end{equation}
On the already licensed free coordinate unit of Theorem~\ref{p5c:v2169:b13:thm:vacuumreturn},
$\kappa_-=\kappa_+=\kappa>0$, so
\begin{equation}
 \boxed{\Tr(\sigma\mathcal E_W)=\tfrac32\pi^2\kappa>0.} \label{p4:eq:v2169-b11-l12}
\end{equation}
This is independent of the neutral occupied orientation and does not require
a neutral terminal-pole or Dirac-mass dictionary.
\begin{proof}
Square $h_R$ and use its polar unitary to obtain \eqref{p4:eq:v2169-b11-l9} by its sign function.
Since $E_e^\dagger F_eE_e=I_3/2$, cyclicity in the squared particle--hole
maps gives
$\Tr(TF_eT^\dagger(I-F_\nu))=\Tr(I_3-F_{\nu,L})/2$ and
$\Tr(T(I-F_e)T^\dagger F_\nu)=\Tr F_{\nu,L}/2$.
Unitarity of the already generated $C$ is the only current identity used.
Insert these two values into \eqref{p4:eq:v2169-b11-l7} and the finite coordinate vacuum covariance.
\end{proof}

Thus this supplied free filled Slater--coordinate unit is not stationary
under the weak mark. To use it as a stationary state of the complete action,
the complete creation coefficients must cancel the leakage in \eqref{p4:eq:v2169-b11-l4}, or the
existing interacting state constructor must supply its actual dressed support.
The marked calculation alone does not prohibit that full-action cancellation.



\paragraph{What existing state constructors now supply}
The fresh-branch unit theorem \ref{p4:thm:v2114-neutral-fock-unit} supplies
$\Omega$ only on its stated neutral-origin domain. A free Slater ground
can instead come from the negative spectral subspace of the \emph{actual}
one-body action on a licensed free-vacuum branch. III-F explicitly fills the
negative modes of its comparison $h_*$; that comparison cannot substitute
for the actual source exporter. The declared full interacting Gibbs vacuum
is stationary by its existing IV theorem and may be non-Slater. In that
case Theorem~\ref{p4:v2169:b10:thm:support} directly uses its full support; no Slater certificate is needed.

For an actual Slater/unit state, the first new application is to evaluate
\eqref{p4:eq:v2169-b11-l2}--\eqref{p4:eq:v2169-b11-l4} from the complete action coefficients and its existing bosonic
state, including pairing and pure-creation contacts. On the weak mark that
evaluation is reduced to \eqref{p4:eq:v2169-b11-l6}. Every bilinear normal $a^\dagger a$ term and
normal quartic $a^\dagger a^\dagger aa$ term kills the empty fermion unit;
a pairing term can supply nonzero $B_{I\varnothing}$. Its absence cannot be
deduced from absence of a preparation token. No ground claim follows solely
from stationarity: the energies on the complementary sectors must be compared
if this is the state constructor's claimed output.

% END INLINE SOURCE: Paper4-ActionStateCertificate-Main-v1_0.tex


% BEGIN INLINE SOURCE: Paper4-AdjacentDegreeReduction-Main-v1_0.tex
% Framework v2.169: Batches 7--13, integrated at the owning source.


\subsection{Balanced action grading and adjacent physical source degrees}
\label{p4:sec:v2169-adjacent-degrees}
\paragraph{Result and required action identity.}
The complete charged source polynomial can be calculated on adjacent
fermion-number sectors when the actual complete action preserves that number.
On an occurring empty-fermion support this reduces to the exact one-fermion
Hamiltonian. The forced density column of Theorem~\ref{p5c:v2169:b8:thm:weaksecond} then collapses on that
support, which closes that specific extra-column obligation on this domain.
A filled Dirac ground instead uses its actual adjacent degrees. Pairing or a
number-changing contact requires the full generated target cyclic space.



\paragraph{The complete word already supplies the number certificate}
In the original source CAR presentation $a_i$, the exact normal-word grading is
\begin{equation}
 \boxed{[N_F,H]=
       \sum_{I,J}(|I|-|J|)a_I^\dagger a_J\otimes B_{IJ}.} \label{p4:eq:v2169-b12-n0}
\end{equation}
Normal-monomial independence therefore gives $[N_F,H]=0$ precisely when every
assembled coefficient of unequal creation and annihilation degree is zero.
If the actual complete grammar contains only normal bilinears and balanced
quartics/higher words, this certificate is already proved by its typing.
The native weak and gauge bilinears obey it. Only an actually inherited
unbalanced coefficient requires a further cancellation or larger target space.
The bosonic coefficients commute with this \emph{fermionic} grading; no
bosonic occupation conservation is being assumed.
Particle--hole reexpression in Theorem~\ref{p4:v2169:b11:thm:normalstate} also transports $N_F$; its relative
unit number is not a replacement for the original source grading.
To prove \eqref{p4:eq:v2169-b12-n0}, commute $N_F$ through each creation and annihilation factor
and add the resulting integer coefficients. This is an action equation,
not an additional conservation postulate.



\paragraph{Exact adjacent-degree reduction}
Let the complete licensed finite Hamiltonian satisfy
\[
 [H,N_F]=0,\qquad HP=E_*P,\qquad N_FP=n_*P,
\]
where $P$ is the actual stationary single-energy support. Let $q_r$ be the
actual raw creation tag, satisfying $[N_F,q_r]=q_r$. Physical source
normalization and any equivariant constraint projections are already included;
no CAR is asserted on a projection which does not preserve it. Define
\begin{equation}
 A_+=q_rP,\qquad A_-=q_r^\dagger P,\qquad
 T_+=H_{n_*+1}-E_*I,\qquad T_-=H_{n_*-1}-E_*I .
 \label{p4:eq:v2169-b12-n1}
\end{equation}
Absent sectors or zero maps give the polynomial one.

\begin{theorem}[The complete physical source uses only adjacent degrees]
\label{p4:v2169:b12:thm:adjacentnumber}
For every density $\rho$ supported in $P$,
\begin{equation}
 \boxed{G_{rr,\rho}(z)=\Tr_P\rho[
   A_+^\dagger(z-T_+)^{-1}A_+
  +A_-^\dagger(z+T_-)^{-1}A_-].}
 \label{p4:eq:v2169-b12-n2}
\end{equation}
For exact support $P$, the source polynomial is the least common multiple
of the matrix-source cyclic polynomial of $T_+^kA_+$ and the sign-reversed
polynomial of $T_-^kA_-$. Their degrees are at most the two actual target
dimensions. Further licensed charge and internal conserved-sector restrictions
can be intersected with these sectors.
\end{theorem}
\begin{proof}
The commutator with $N_F$ puts $q_rP$ and $q_r^\dagger P$ in the indicated
degrees, and $[H,N_F]=0$ keeps each complete cyclic word there.
Because $HP=E_*P$,
$(\operatorname{ad}_H)^kq_rP=(H-E_*)^kq_rP$.
The adjoint source produces the removal frequencies with the opposite sign.
The energy resolution and positive spectral residue matrices give \eqref{p4:eq:v2169-b12-n2}.
The finite matrix-source closures and their least common multiple are
Theorem~\ref{p4:v2169:b10:thm:support}'s construction applied to the additional conserved grading.
All action terms are evaluated before restricting to the actual reducing
sector; no term is discarded by its label.
\end{proof}

On a licensed full ground support $P_0$, number conservation also gives
$[N_F,P_0]=0$. Its generated decomposition is $P_0=\sum_nP_{0,n}$.
Calculate both adjacent legs for each nonzero $P_{0,n}$ and take the least
common multiple of their source polynomials. This is exactly the full
ground-support polynomial: the comparison density $P_0/\Tr P_0$ gives
strictly positive weights to every nonzero summand, so positive residues
cannot cancel. Every density with exact support $P_0$ has those same roots
by Theorem~\ref{p4:v2169:b10:thm:support}. Thus an already licensed degenerate Gibbs ground requires
neither a unique vector nor selection of one fermion-number member.



\paragraph{The full one-fermion Hamiltonian on the empty branch}
On an actual tensor CAR chart preserving $N_F$, write the complete normal word
\begin{equation}
 H=I\otimes H_B+\sum_{ij}a_i^\dagger a_j\otimes K_{ij}
       +\sum_{\substack{|I|=|J|\ge2}}a_I^\dagger a_J\otimes B_{IJ}.
 \label{p4:eq:v2169-b12-n3}
\end{equation}
All $K_{ij}$ are outputs of the complete action, including kinetic, scalar,
gauge and normal-ordering contacts; $K_{ij}^\dagger=K_{ji}$.
Put $P=|0_F\rangle\langle0_F|\otimes P_B$ with
$H_BP_B=E_*P_B$. Define
\begin{equation}
 \boxed{\mathsf H_1=I_{\mathcal K}\otimes H_B+
                 \sum_{ij}|i\rangle\langle j|\otimes K_{ij},\qquad
 \mathsf A_r=j_r\otimes P_B.} \label{p4:eq:v2169-b12-n4}
\end{equation}

\begin{corollary}[Exact full-word one-fermion construction]
\label{p4:v2169:b12:cor:onefermion}
The removal map is zero, and the physical source polynomial is exactly the
matrix-source polynomial of
$(\mathsf H_1-E_*)^k\mathsf A_r$. Every normally ordered higher-fermion term in
\eqref{p4:eq:v2169-b12-n3} vanishes on both degree zero and degree one, by its actual annihilation
degree. The closure degree is at most $\dim\mathcal K\dim\mathcal B$, or the
smaller actually conserved charge-sector dimension. For exact support $P$,
its root at $\varepsilon_r$ already includes tagged visibility.
\end{corollary}
\begin{proof}
An annihilation word of degree at least two kills every vector of fermion
degree at most one. A bilinear $a_i^\dagger a_j$ acts on degree one as
$|i\rangle\langle j|$. The same complete action therefore restricts to \eqref{p4:eq:v2169-b12-n4}.
The unit gives $q_r^\dagger P=0$, so \eqref{p4:eq:v2169-b12-n2} has only its addition leg.
\end{proof}

A sufficient local replacement for global number conservation is proof that
the actual cyclic target space generated by $q_rP$ lies in degree one and
is invariant under the complete $H$. Merely killing pair creation on the
initial vacuum support does not prove this: a later bosonic return can activate
that coefficient. In that case the graded target Krylov word must retain the
higher sectors. The normal-word test of Theorem~\ref{p4:v2169:b11:thm:normalstate} diagnoses this difference.



\paragraph{The previously forced density column now has an exact state quotient}
For the common native weak constituent of Theorem~\ref{p5c:v2169:b8:thm:weaksecond} define
$q_r=e_r^\dagger$, $F_{rs}=n_r^\dagger n_s-\delta_{rs}/2$,
$Z_r=\sum_s e_s^\dagger F_{rs}$,
$K=[X^\dagger,X]$, and $D=(X^\dagger X+XX^\dagger)/2$.
For any bosonic support on the empty fermion branch,
\begin{equation}
 \boxed{Z_rP=-\tfrac12q_rP,\qquad Z_r^\dagger P=0.} \label{p4:eq:v2169-b12-n5}
\end{equation}
The exact weak second column is consequently
\begin{equation}
 \boxed{\operatorname{ad}_{V_W}^2(q_r)P
      =\pi^2q_r X^\dagger XP,\qquad
       (\operatorname{ad}_{V_W}^2(q_r))^\dagger P=0.} \label{p4:eq:v2169-b12-n6}
\end{equation}
These identities hold for every finite $X$, without CCR or a chosen
occupation cap.
\begin{proof}
Every $n_r^\dagger n_s$ kills the fermion unit, leaving only its centered
constant. In $Z_r^\dagger$, the charged annihilator also kills the unit.
Insert \eqref{p4:eq:v2169-b12-n5} into Theorem~\ref{p5c:v2169:b8:thm:weaksecond}'s exact
$\pi^2(q_rD-Z_rK)$ and use $D+K/2=X^\dagger X$.
The weak vertex kills the fermion unit, so the adjoint of each iterated
source commutator kills this support as well.
\end{proof}

Thus full trace independence of $Z_rK$ and physical quotient dependence are
compatible. The density direction creates no independent higher-fermion
obligation on this empty support. Its \emph{bosonic} return $X^\dagger X$
still remains and can create new occupation columns.
The first two marked physical moments are, for unit-normalized $j_r$,
\begin{equation}
 \boxed{\mu_2^W=\pi^2\Tr(\sigma X^\dagger X)I_3,\qquad
 \mu_4^W=\pi^4\Tr(\sigma (X^\dagger X)^2)I_3.} \label{p4:eq:v2169-b12-n7}
\end{equation}
Their defect is the actual variance of $X^\dagger X$, not Theorem~\ref{p5c:v2169:b8:thm:weaksecond}'s
full-trace density norm. These marked moments do not substitute for moments
of the complete $H$.



\paragraph{First unproved application and existing generators}
III-A's complete action, IV's actual quantization and V-C's existing vertex
give all coefficients in \eqref{p4:eq:v2169-b12-n3}--\eqref{p4:eq:v2169-b12-n4}. Theorem~\ref{p4:v2169:b11:thm:normalstate} supplies the exact support
stationarity test for the occurring state. The next application is:
\begin{enumerate}
\item reuse $[H,N_F]=0$ when the complete normal grammar is already
balanced, by \eqref{p4:eq:v2169-b12-n0}; if the actual complete word contains
unbalanced coefficients, prove their assembled cancellation or certify the
smaller complete cyclic invariant space before this reduction;
\item use the actual state's $n_*,P$ to assemble the appropriate complete
adjacent blocks; the filled negative-energy free ground must not be replaced
by $n_*=0$;
\item calculate their source polynomial and prove its fixed terminal root.
\end{enumerate}
No additional occupation-state selector or rest preparation enters this list.
The particle/rest dictionary and each consumed physical limit remain separate.
The primitive $m_p=1$, current incidence and terminal coefficients are retained.

% END INLINE SOURCE: Paper4-AdjacentDegreeReduction-Main-v1_0.tex



% v2.171: conversation proofs, with actual history obligations retained.

% BEGIN INLINE SOURCE: Paper4-CompleteChargeTests-Main-v1_0.tex
% Framework v2.171: integrated conversation proof on its stated domain.


\subsection{Complete particle and conjugate charge transport}
\label{p4:v2171:sec:fullchain}
The particle carrier is the internal boundary response already constituted
from $M_8$, rather than a separately installed collection of continuum
particles. Its gauge and charged mass slots, K8--K16 owner and root-state
ancestry are reused. The common-action kinetic constructor of III-F,
\texttt{p3f:thm:v1148-finite-metric-constructor}, has the exact form
\begin{equation}
 \widehat H=d\Gamma(G^{-1/2}KG^{-1/2})
 -i\hbar d\Gamma(\gamma)+\widehat V_{\rm inherited,can}
 +\varepsilon_0I,
 \qquad G=J^\dagger G_pJ.
 \label{p4:v2171:eq:fullham}
\end{equation}
The displayed connection term belongs to the coordinate Hamiltonian.
In that coordinate convention $D_t$ in the charge balance below denotes the
ordinary time derivative of coordinate operators. Equivalently, use a
covariant derivative and remove its connection contribution from $H$.
The state and charge equations use one convention consistently; the
connection is never counted twice.

Here $J$ is the actual field embedding, $\gamma$ its complete metric-compatible
connection, and $\widehat V_{\rm inherited,can}$ contains every pulled-back
interaction and stationary contact. Writing this constructor is not an
evaluation of $J$, the physical weak overlaps, or the Hamiltonian along the
cyclic history. Terminal masses or family traces cannot be substituted for
these missing contractions.

For a licensed physical conjugate-sector realization, the additive charges
on a canonical particle/antiparticle occupation presentation are
\begin{equation}
 Q_B=\tfrac13(N_q-N_{\bar q}),\qquad
 Q_L=N_\ell-N_{\bar\ell},\qquad Q_D=Q_B-Q_L.
 \label{p4:v2171:eq:signedcharges}
\end{equation}
All spin, color, family and occurring external slots are included in these
counts. The opposite charges follow from V-A's conjugate representation.
Their physical particle/antiparticle interpretation consumes its stated
PCT domain; no extra finite trace copy is inserted. The finite algebraic
counting statement is distinct from this analytic interpretation.
An evolving retained frontier restricts the already evolved parent state,
rather than preparing a new vacuum or deleting outflow from a charge balance.
For the closed whole spatial slice there is no omitted spatial flux; a
smaller material region retains its actual boundary flux separately.

\begin{theorem}[Exact many-body charge history]
\label{p4:v2171:thm:fullcharge}
On an actually licensed finite self-adjoint Hamiltonian branch, with the
inherited state $\rho_i$, let $U$ solve its full covariant Schr\"odinger
equation and $\rho_t=U(t,t_i)\rho_iU(t,t_i)^\dagger$.
For every physical charge $Q(t)$ define
\begin{equation}
 \mathcal J_Q=D_tQ+\frac{i}{\hbar}[H,Q].
 \label{p4:v2171:eq:chargeoperator}
\end{equation}
Then, retaining the moving-frame derivative,
\begin{equation}
 q_Q(t_f)=q_Q(t_i)+\int_{t_i}^{t_f}
                 \operatorname{Tr}(\rho_t\mathcal J_Q(t))\,dt.
 \label{p4:v2171:eq:fullintegral}
\end{equation}
This is an exact complete-state calculation. Any washout, charge
redistribution or spectator effect represented in this owner is already
inside $H$ and $\rho_t$; no additional multiplicative efficiency follows.
For a fully inherited non-Hamiltonian reduced domain, its actual generator
replaces the commutator in \eqref{p4:v2171:eq:chargeoperator}.
\end{theorem}
\begin{proof}
Differentiate $\operatorname{Tr}(\rho_tQ)$ using
$D_t\rho=-i[H,\rho]/\hbar$ and cyclicity. Integration proves the result.
No closure of the covariance or Wick factorization is invoked.
\end{proof}

If $D_tQ=0$ and a complete charge grading of the inherited interaction is
$H=H_c+\sum_r(V_r+V_r^\dagger)$ with $[Q,V_r]=rV_r$, then
\begin{equation}
 \dot q_Q=\frac{2}{\hbar}\sum_r r\,
                      \Im\operatorname{Tr}(\rho_tV_r).
 \label{p4:v2171:eq:gradedsource}
\end{equation}
The owned pair has $V_{-2}=A_{BL2}/2$, giving precisely
$\dot q_D=-2\Im\omega_t(A_{BL2})/\hbar$. If other charged contacts
actually occur, their complete graded sources are retained rather than
assigned to a charge-conserving normal block by notation.



\subsection{Complete time-dependent charge calculation in a finite test space}
Fix an already licensed finite physical carrier of dimension $d$ on a
constant-rank compact history interval and a parallel orthonormal frame.
The actual Hamiltonian includes the induced connection and every normal,
pair, quartic, bosonic, constraint and contact term consumed by the claim.
It may be written
\[
 H(t)=\sum_{\alpha=1}^{r}f_\alpha(t)H_\alpha,
 \qquad H_\alpha=H_\alpha^\dagger,\quad f_\alpha(t)\in\mathbb R,
\]
by taking a real basis of the span of its actual Hermitian values.
The functions are the existing history coefficients, not chosen profiles.
Assume they are integrable on the compact interval. Fix the inherited state
$\rho_i$ and the dimensionless covariantly constant charge $Q$.

\begin{theorem}[Least complete charge-test flow]\label{p4:v2171:thm:flow}
Let $W_0=\operatorname{span}_{\mathbb R}\{I,Q\}$, enlarged by actual
consumed source-jet tests if the target uses them. Set
\[
 W_{n+1}=W_n+\sum_\alpha i[H_\alpha,W_n].
\]
This stabilizes after at most $d^2-\dim W_0$ strict increases at the unique
least real Hermitian operator space $W_*$ invariant under every owned
generator and containing the tests. Choose a fixed Hermitian basis $E_j$ and
define real matrices $A^\alpha$ by
\[
 \frac{i}{\hbar}[H_\alpha,E_j]
       =\sum_k A^\alpha_{jk}E_k.
\]
Then the complete expectation vector obeys exactly
\begin{equation}\label{p4:v2171:eq:testflow}
 x_j(t)=\Tr(\rho(t)E_j),\qquad
 \dot x(t)=\sum_\alpha f_\alpha(t)A^\alpha x(t),\qquad
 x_j(t_i)=\Tr(\rho_iE_j).
\end{equation}
This is the full finite charge history; it uses no Gaussian assumption,
Wick factorization, covariance truncation or truncation in pair strength.
\end{theorem}
\begin{proof}
Hermitian matrices have real dimension $d^2$. Every strict increase adds
a dimension; once stable, all generators preserve the space. Any other
invariant space containing the seeds contains every $W_n$ by induction,
proving leastness. The exact finite equation
$i\hbar\dot\rho=[H,\rho]$ and cyclicity of trace give
\[
 \dot x_j=\frac{i}{\hbar}\Tr\rho[H,E_j],
\]
which is \eqref{p4:v2171:eq:testflow}. Its time-ordered exponential exists on the
integrable finite interval. Thus the reduced expectation calculation agrees
with full many-body unitary propagation by uniqueness. This is a linear-test
specialization of II-A's already licensed least-invariant-test construction;
it neither selects $H$ nor replaces the actual inherited state.
\end{proof}

On the quadratic CAR branch, the existing Nambu covariance equation can be
used as an equivalent shorter calculation. When quartic normal contacts are
present, Theorem~\ref{p4:v2171:thm:flow} retains their actual operator action.
If an actual reduction or kinetic limit is consumed, its complete generator
and inherited correlation/assignment domain must first be certified. A
larger finite calculation carrier is only the existing operator realization,
not a new physical field or a primitive occupation cutoff.



\subsection{The finite coordinate boundary fixes the same-charge commutator}
Paper IV, \texttt{p4:thm:gp1-bosonic-occupation-boundary}, licenses a complete
initial occupation sublevel when its generated-boundary condition is supplied
or derived; GP1 finiteness alone is not that condition. V-C,
\texttt{p5c:eq:v2169-b13-o1}, uses the actual finite positive-frequency
coordinate branch. On that branch retain two independently occurring charged
coordinate constituents among the existing $d_B\ge4$ oscillator modes:
\[
 X_i=\chi_iQ_N(a_{i,-}^\dagger+a_{i,+})Q_N,\quad
 X_j=\chi_jQ_N(a_{j,-}^\dagger+a_{j,+})Q_N,\quad
 \kappa_k=|\chi_k|^2>0.
\]
$Q_N$ is the actual total-occupation projection through the actual $N\ge1$.
It is not replaced by a product of independently truncated mode spaces.
All other admitted modes remain in $d_B$. No mode count, cap or coordinate
coefficient is selected here. If two such ports have not occurred in the
consumed chart, this formula does not create them.

\begin{theorem}[Exact coordinate commutator and complete-cap trace norm]
\label{p4:v2171:thm:capcommutator}
Let $P_N=Q_N-Q_{N-1}$. For $i\ne j$ on this complete-sublevel chart,
\begin{equation}
 \boxed{[X_i,X_j]=\chi_i\chi_jP_N
 (a_{i,-}^\dagger a_{j,+}-a_{j,-}^\dagger a_{i,+})P_N.}
 \label{p4:v2171:eq:capcommutator}
\end{equation}
The normalized trace on the complete bosonic carrier gives
\begin{equation}
 \boxed{\|[X_i,X_j]\|_\tau^2
 =\frac{2N}{d_B+1}\kappa_i\kappa_j>0.}
 \label{p4:v2171:eq:capcommutatornorm}
\end{equation}
\end{theorem}
\begin{proof}
Below degree $N$, the two coordinate multiplication orders agree. At degree
$N$, creation is excluded while annihilation followed by creation is retained.
Distinct-mode creation--creation and annihilation--annihilation terms cancel.
The two mixed terms give \eqref{p4:v2171:eq:capcommutator}.

In its squared norm the cross term changes four distinct occupations and
has zero trace. The two diagonal terms have equal sums. With $u$ the $i,-$
mode and $v$ the $j,+$ mode, one such sum is
\[
 \sum_{n_1+\cdots+n_{d_B}=N}(n_u+1)n_v
 =\binom{N+d_B-1}{d_B+1}+\binom{N+d_B-1}{d_B}
 =\binom{N+d_B}{d_B+1}.
\]
The complete sublevel dimension is $\binom{N+d_B}{d_B}$. Twice this sum
divided by the dimension is $2N/(d_B+1)$. Restore $\kappa_i\kappa_j$ to obtain the norm.
V-C applies this norm to its single-row marked coefficient on that chart.
\end{proof}
The commutator acts on the actual top occupation degree and annihilates the
Fock unit. The physical response consumes its inherited occupation support.
For any actual positive state $\omega$,
\[
 0\le\omega([X_i,X_j]^\dagger[X_i,X_j])
 \le4N^2\kappa_i\kappa_j\,\omega(P_N).
\]
This uses \eqref{p4:v2171:eq:capcommutator} and
$\|a_{i,-}^\dagger a_{j,+}\|\le N$ on the top degree.
The mathematical trace does not set $\omega(P_N)$ or the physical charge.

% END INLINE SOURCE: Paper4-CompleteOccupationCommutator-Main-v1_0.tex





\subsection{Carrier-four reheating source on compact $S^3$ (framework v2.219)}
\label{p4:sec:v2219-s3-reheating}
\begin{proposition}[Discrete closed-space excitation from a coherent UV source]
\label{p4:prop:v2219-uv-bogoliubov}
For a physically selected scalaron $\phi(t)$ and the full positive
Schur-reduced particle kernel $K_P(\phi,\eta)$, the normalized spatial
scalar harmonics on a round $S^3$ obey
$-\Delta_{S^3}Y_\ell=\ell(\ell+2)Y_\ell$ and the physical volume is
$2\pi^2a^3$.  A resolved positive oscillator frequency $\Omega_j(t)$
has the standard instantaneous mixing law
\[
 \dot a_j=-i\Omega_ja_j+\frac{\dot\Omega_j}{2\Omega_j}a_j^\dagger,
\]
with suitable matrix generalization when eigenvectors rotate.
The transition amplitudes and occupations depend on
$\partial_\phi K_P$ and the actual time-dependent finite modes.
The squared carrier energy identity $C_2=E_{32}^2$ is not a
frequency resonance $m_{\rm UV}=2\Omega_j$.
\end{proposition}
\begin{remark}[Conformal photons and compact-mode selection]
\leavevmode\par\noindent
For constant-normalization conformal Maxwell theory on a closed FLRW
background the transverse conformal modes have
$A_n''+n^2A_n=0$, $n\ge2$, and no expansion-only Bogoliubov
production.  Photon excitation by the UV scalaron requires an
actual nonconformal gauge, trace-anomaly or internal-boundary vertex.
The particle and scalaron stresses must obey a single conserved ledger;
no independent $K16$ energy deposit may be added to the work already
attributed to these modes.
\end{remark}


% BEGIN FRAMEWORK v2.228: CHARGED ORIGIN AND CHARACTER TESTS


\subsection{Propagation of the certified charged mass-only block}
\label{p4:sec:v228-charged-bilinear}
\begin{corollary}[Protected bilinear response on its certified source domain]
\label{p4:cor:v228-charged-bilinear}
For an actual canonically normalized number-conserving bilinear CAR chart,
assume the charged hierarchy-arrow origin certificate and the mass-only
condition of Corollary~\ref{p3f:cor:v228-derived-electron-pair}.
The normalized left electron anticommutator resolvent is $z/(z^2-e^2)$.
The scalar character test or its finite operator-valued version may certify
this condition after its actual word-class identification.
\end{corollary}
\begin{proof}
Theorem~\ref{p4:thm:v2168-bilinear-source-commutators} identifies the
canonical bilinear source dynamics with its one-particle coefficient.
The reducing pair gives the displayed resolvent. The CAR normalization
does not require a Gaussian preparation.
\end{proof}
This result is confined to the stated bilinear block. An interacting
source-visible pole consumes the full action, constraints, stationary
quotient and visibility test. The source handoff does not discharge the
hierarchy-arrow origin certificate of
Definition~\ref{p5:def:v228-hierarchy-origin}.


% BEGIN FRAMEWORK v2.240: EXACT SOURCE-POLE AND REPRESENTATION HANDOFF
\begin{proposition}[Derived positive/negative Dirac source splitting]
\label{p4b:prop:v240-protected-chiral-source-splitting}
Let the \emph{actual certified} chiral mass block have a reducing
pair $(J_L,J_R)$ with $TJ_R=eJ_L$ and $T^\dagger J_L=eJ_R$.
In the doubled chiral space define
$J_\pm=2^{-1/2}(J_L,\pm J_R)^T$.
Then $h_TJ_\pm=\pm eJ_\pm$ and the original left source is
$2^{-1/2}(J_++J_-)$. Hence its marked free resolvent is exactly
\[
 \boxed{G_L^{(0)}(z)=\frac12\left(\frac1{z-e}
               +\frac1{z+e}\right)=\frac{z}{z^2-e^2}.}
\]
Replacing this marked source by a sole positive-energy source in
an interacting pole calculation requires the corresponding generated
source projection; it is not a harmless change of representation.
\end{proposition}
\begin{proof}
Apply the chiral block directly to $J_\pm$ and take the
orthogonal resolvent matrix element of the left source.
\end{proof}

\begin{theorem}[Complete finite one-source Schur pole and positive residues]
\label{p4b:thm:v240-complete-one-source-stieltjes}
Let an actual finite stationary source-sector Hamiltonian $H=H^\dagger$
and a normalized source line $j$ have a block representation
\[
 H=\begin{pmatrix}a&b^\dagger\\b&C_R\end{pmatrix},
 \qquad a\in\mathbb R,\quad C_R=C_R^\dagger.
\]
Let $\lambda_1<\cdots<\lambda_m$ be the distinct eigenvalues of
$C_R$ with strictly positive source weights
$\beta_\alpha=\|P_{\lambda_\alpha}(C_R)b\|^2$; other complement
modes are source-invisible. Then
\[
 \boxed{\langle j,(zI-H)^{-1}j\rangle
 =\left[z-a-\sum_{\alpha=1}^m
           \frac{\beta_\alpha}{z-\lambda_\alpha}\right]^{-1}.}
\]
There are exactly $m+1$ simple real source-visible poles, one in
each interval cut out by the $m$ coupled complementary eigenvalues.
For each pole $z_p$ its scalar source residue is
\[
 \boxed{w_p=\left[1+\sum_{\alpha=1}^m
       \frac{\beta_\alpha}{(z_p-\lambda_\alpha)^2}\right]^{-1}>0,
       \qquad \sum_{p=1}^{m+1}w_p=1.}
\]
The conclusion includes all inherited interactions already present
in $C_R$ and $b$; it does not evaluate either from a symbolic action.
\end{theorem}
\begin{proof}
The exact $1\oplus\dim C_R$ Schur complement gives the
resolvent. Its real denominator $f(z)$ has
$f'(z)=1+\sum_\alpha\beta_\alpha/(z-\lambda_\alpha)^2>0$
between poles, with the endpoint signs ensuring one zero in every
interval, including both exterior intervals. The residue equals
$1/f'(z_p)$. The large-$z$ expansion is $z^{-1}+O(z^{-2})$;
comparison with its partial fractions proves the sum rule.
\end{proof}
% END FRAMEWORK v2.240: EXACT SOURCE-POLE AND REPRESENTATION HANDOFF

% END FRAMEWORK v2.228: CHARGED ORIGIN AND CHARACTER TESTS

% Formal appendix delimiter retired in the Paper IV split.
% BEGIN INLINE SOURCE: Paper4-SourceCofactorChecks-Appendix-v1_0.tex
% Framework v2.191: supplementary finite-matrix diagnostics only.


\section{Source-cyclic resolvents and physically visible poles}
\label{p4b:sec:v238-b6}
\noindent\textit{Source and claim boundary.} Finite commutator closures, source-cyclic quotients and tagged polynomials locate poles of an actual source-state response. A determinant zero without an occupied source-visible residue is not a measured particle.

\subsection{Generated memory columns, charge-sector poles and state-family visibility}
\label{p4:sec:v2168-charged-memory-spectrum}
The following constructions act on the actual finite source chart.
The bilinear formulas are constituents of the complete word; all remaining
columns and cross returns stay in the response. The spectral construction
requires the actual charge-equivariant physical reduction and does not
assume CAR for a projected source.


\paragraph{Actual vertex words generate the memory columns}
Consider first the inherited number-preserving bilinear constituent
\[
\begin{aligned}
 H_1 &= d\Gamma(h)\otimes I+I\otimes H_B
       +\sum_\alpha d\Gamma(T_\alpha)\otimes X_\alpha,\\
 h &= h^\dagger,\qquad T_\alpha=T_\alpha^\dagger,\\
 H_B &= H_B^\dagger,\qquad X_\alpha=X_\alpha^\dagger.
\end{aligned}
\]
Real bosonic ports are used, including real components of complex fields.
Their finite occupation compression is retained. Let
$q_r=a^\dagger(j_r)\otimes I$ with $J_c=(j_1,\ldots,j_k)$ and
$J_c^\dagger J_c=I$. Fermionic CAR are used only on a chart that actually
preserves them. This section evaluates this constituent; if the complete
Hamiltonian contains pairing, higher fermion terms or reduction contacts,
they are included as $V_{\rm rem}$ below.

\begin{theorem}[First and second commutators with the inherited bilinear word]
\label{p4:thm:v2168-bilinear-source-commutators}
The source columns generated by $H_1$ are exactly
\begin{align*}
 [H_1,q_r]={}&a^\dagger(hj_r)\otimes I
       +\sum_\alpha a^\dagger(T_\alpha j_r)\otimes X_\alpha,\\
 [H_1,[H_1,q_r]]={}&a^\dagger(h^2j_r)\otimes I
       +\sum_\alpha a^\dagger((hT_\alpha+T_\alpha h)j_r)
                   \otimes X_\alpha\\
 &+\sum_\alpha a^\dagger(T_\alpha j_r)\otimes[H_B,X_\alpha]\\
 &+\sum_{\alpha\beta}a^\dagger(T_\beta T_\alpha j_r)
                   \otimes X_\beta X_\alpha\\
 &+\sum_{\alpha\beta}a^\dagger(T_\alpha j_r)d\Gamma(T_\beta)
                   \otimes[X_\beta,X_\alpha].
 \tag{C1}
\end{align*}
In particular the last line vanishes on a commuting-coordinate chart.
No product state, perturbative expansion or Wick factorization is used.
\end{theorem}
\begin{proof}
Exact CAR give
$[d\Gamma(T),a^\dagger(v)]=a^\dagger(Tv)$. For a bosonic coefficient $F$,
\[
 [d\Gamma(T_\beta)\otimes X_\beta,
             a^\dagger(v)\otimes F]
 =a^\dagger(T_\beta v)\otimes X_\beta F
  +a^\dagger(v)d\Gamma(T_\beta)\otimes[X_\beta,F].
\]
Apply this identity to each term of the first commutator and add the
$d\Gamma(h)$ and $H_B$ contributions. This proves every term and its order.
\end{proof}

For the complete $H=H_1+V_{\rm rem}$ define
$R_r=[V_{\rm rem},q_r]$. The columns actually consumed are
\[
 [H,q_r]=[H_1,q_r]+R_r,
\]
\[
 [H,[H,q_r]]=[H_1,[H_1,q_r]]+[H_1,R_r]
       +[V_{\rm rem},[H_1,q_r]]+[V_{\rm rem},R_r].
\]
Every remainder column and cross return stays in the source Gram. It is
not licensed to vanish because the gauge coefficients have been computed.
$H_B$ retains its native scalar potential and all gauge/constraint terms;
the third line of (C1) is a generated bosonic momentum/force contribution,
not an unspecified new interaction.



\paragraph{An exact first memory return on an entangled state}
Write $\langle O\rangle=\Tr(\rho O)$ on the finite carrier. The coefficient-pairing
identities in this section hold for any density matrix. If $\rho$ is
stationary for the full $H_1$, they are the moments of its self-adjoint
source Liouvillian $\mathcal L=\operatorname{ad}_{H_1}$ in the anticommutator quotient
of the stationary source construction.

\begin{proposition}[Coefficient pairing including the finite boundary correction]
\label{p4:prop:v2168-coefficient-pairing}
For $O=a^\dagger(v)\otimes F$ and $P=a^\dagger(w)\otimes G$,
\[
 \boxed{
 (O,P)_\rho
 =\langle v,w\rangle\langle I\otimes F^\dagger G\rangle
  +\left\langle a^\dagger(w)a(v)\otimes[G,F^\dagger]\right\rangle.}
\]
Consequently, put $v_{0r}=hj_r$, $F_0=I$ and
$v_{\alpha r}=T_\alpha j_r$, $F_\alpha=X_\alpha$. Then
\[
 \boxed{
 \mu_{2,rs}=\sum_{ab}\langle v_{ar},v_{bs}\rangle
                         \langle F_a^\dagger F_b\rangle
 +\sum_{ab}\left\langle a^\dagger(v_{bs})a(v_{ar})
                         \otimes[F_b,F_a^\dagger]\right\rangle.}
 \tag{C2}
\]
The second sum is the exact occupation/coefficient-commutator correction.
It includes correlations in an entangled $\rho$.
\end{proposition}
\begin{proof}
Use $a(v)a^\dagger(w)=\langle v,w\rangle I-a^\dagger(w)a(v)$ in
$\{O^\dagger,P\}$. The two products have bosonic orders $F^\dagger G$
and $GF^\dagger$. Their difference gives the correction. Sum the
identity over the generated first-commutator columns.
\end{proof}

Let $m_\alpha=\langle X_\alpha\rangle$,
$\delta X_\alpha=X_\alpha-m_\alpha I$,
$\bar h=h+\sum_\alpha m_\alpha T_\alpha$ and
$C_{\alpha\beta}=\langle\delta X_\alpha\delta X_\beta\rangle$.
The matrix $C$ is a positive Hermitian covariance, as follows by applying
the state to $Y^\dagger Y$ for $Y=\sum z_\alpha\delta X_\alpha$.

\begin{theorem}[Generated first memory Gram on a commuting-coordinate chart]
\label{p4:thm:v2168-commuting-port-memory-gram}
Assume $[X_\alpha,X_\beta]=0$ and $[H_1,\rho]=0$. No independence of
fermions and bosons is required. Then
\[
 \boxed{\mu_1=J_c^\dagger\bar hJ_c,\qquad
 \mu_2=J_c^\dagger\bar h^2J_c
       +\sum_{\alpha\beta}C_{\alpha\beta}
                   J_c^\dagger T_\alpha T_\beta J_c.}
\]
For the exact source/complement blocks $A,B,C_{\rm mem}$ of the stationary source construction,
\[
 \boxed{
 A=J_c^\dagger\bar hJ_c,\qquad
 B^\dagger B
 =J_c^\dagger\bar h(I-J_cJ_c^\dagger)\bar hJ_c
  +\sum_{\alpha\beta}C_{\alpha\beta}
                   J_c^\dagger T_\alpha T_\beta J_c.}
 \tag{C3}
\]
Both displayed terms in $B^\dagger B$ are positive semidefinite.
\end{theorem}
\begin{proof}
The CAR anticommutator of $q_r^\dagger$ with the first commutator is
$j_r^\dagger hj_sI+\sum j_r^\dagger T_\alpha j_sX_\alpha$.
This proves $\mu_1$. Center the bosonic ports in (C2); every coefficient
commutator vanishes and every linear centered mean is zero. This gives
the stated $\mu_2$. In the stationary quotient
$B^\dagger B=\mu_2-A^2$, which proves (C3). The first term is
the Gram of $(I-J_cJ_c^\dagger)\bar hJ_c$.
The second is a Gram with the positive covariance $C$; since the centered
coordinates commute, $C$ is real symmetric and can be factored over the
real coefficient space. Each factor produces a square of a Hermitian
linear combination of the $T_\alpha$, compressed by $J_c$.
\end{proof}

Thus a centered electromagnetic vertex need not disappear from the memory
when $\langle X_\alpha\rangle=0$. Its first return uses the covariance
generated by the same state and dynamics. On a factor-preserving mode chart
with $T_{Q,\alpha}=-eK_\alpha\otimes I_G$ and
$J_c=j_{\rm ext}\otimes U$, its own covariance return is explicitly
\[
 \boxed{
 \mathcal B_Q
 =\frac{3\pi^2}{4}
   \sum_{\alpha\beta}C_{\alpha\beta}
    \langle j_{\rm ext},K_\alpha K_\beta j_{\rm ext}\rangle I_3.}
\]
This scalar family return is conditional on that actual factor-preserving
chart and commuting ports. It is not a proof that the complete mass source,
propagator or memory is family diagonal. Cross returns with the $Z_0$,
charged weak and scalar vertices remain in the total sum.
If finite occupation compression makes the bosonic ports noncommuting,
equation (C2) and the last line of (C1) are used instead of deleting them.
These formulas evaluate first memory returns; a pole depends on the entire
complement dynamics $C_{\rm mem}$, not only on $B^\dagger B$.



\paragraph{The native charge reduction is exact on its declared chart}
Let $\mathcal Q_{\rm tot}$ be the total charge represented on the actual
finite field carrier, including the charged bosonic, neutral, conjugate and
constraint sectors consumed by the action. It is not just the charge of the
three observed source ports.

\begin{proposition}[Charge conservation under equivariant quantization]
\label{p4:prop:v2168-equivariant-charge-reduction}
Suppose every inherited Hamiltonian monomial has net charge zero and the
licensed occupation and constraint projections commute with
$\mathcal Q_{\rm tot}$. Then
\[
 [H,\mathcal Q_{\rm tot}]=0.
\]
If the retained charged source obeys
$[\mathcal Q_{\rm tot},q_r]=-q_r$, its cyclic operator space lies in
\[
 \mathfrak B_{-1}
 =\bigoplus_c\operatorname{Hom}(\mathcal H_c,\mathcal H_{c-1}),
 \qquad \mathcal H_c=\operatorname{Ran}P_c(\mathcal Q_{\rm tot}).
\]
This is a reducing space for $\operatorname{ad}_H$.
\end{proposition}
\begin{proof}
Commutation with charge is a derivation. The charge of a monomial is the
sum of its slot charges, so every neutral monomial commutes with charge,
independently of ordering. An equivariant compression preserves this
identity and the charge of each compressed source. The Jacobi identity
gives $[\mathcal Q_{\rm tot},[H,X]]=-[H,X]$ for $X\in\mathfrak B_{-1}$.
Self-adjointness of $\operatorname{ad}_H$ for the Hilbert--Schmidt pairing makes this
invariant finite space reducing.
\end{proof}

This verifies the finite charge identity once the actual projection's
equivariance has been certified. A classical gauge stabilizer alone does
not certify an arbitrary quantum constraint reduction. Charged $W$ and
scalar ports have compensating slot charges; they remain in $H$ even though
the source response changes total charge by one unit. Charge conservation
does not set any energy difference equal to a terminal coefficient.



\paragraph{An explicit finite charged matrix and complete spectral readout}
On the charge-reduced physical chart put $H_c=P_cHP_c$. For a matrix
$X_c:\mathcal H_c\to\mathcal H_{c-1}$,
\[
 \mathcal L_{-1}X_c=H_{c-1}X_c-X_cH_c.
\]
In column vectorization this is the explicit block matrix
\[
 \boxed{L_c=I_{d_c}\otimes H_{c-1}-H_c^T\otimes I_{d_{c-1}},
 \qquad L_{-1}=\bigoplus_cL_c.}
 \tag{C4}
\]
It is Hermitian in the Hilbert--Schmidt metric. Set
$v_r=\bigoplus_c\operatorname{vec}(P_{c-1}q_rP_c)$ and
$V=(v_1,\ldots,v_k)$. These are coordinates of the actual source operators,
including any inherited constraint compression. This construction requires
no CAR assumption about a projected source. Any dressing or charged boundary
domain consumed by the physical claim is part of this actual $q_r$ and its
carrier; a bare local source cannot replace it after a constraint quotient.

For a stationary density $\rho$ one may first pinch in total charge:
$\rho^Q=\sum_cP_c\rho P_c$. It is again stationary and positive. Every
charged-source anticommutator resolvent has charge zero, so its expectation
is exactly unchanged by this pinching. It is a response-equivalent
descendant, not a selected physical state.
Write $\rho_c=P_c\rho^QP_c$. The source metric on (C4) is
\[
 \boxed{W_{\rho,c}=\rho_c^T\otimes I_{d_{c-1}}
                         +I_{d_c}\otimes\rho_{c-1}.}
\]
It commutes with $L_c$. With $W_\rho=\bigoplus_cW_{\rho,c}$,
\[
 \boxed{G_\rho(z)=V^\dagger W_\rho(zI-L_{-1})^{-1}V.}
 \tag{C5}
\]
When $\rho$ is not faithful, the null space of $W_\rho$ is quotiented;
equation (C5) still holds. The source Gram is
$N_\rho=V^\dagger W_\rho V$. Exact CAR imply $N_\rho=I$; otherwise this
is the actual inherited source Gram. A positive inherited normalization
transports this Gram, residues and the source tag together. It cannot create
a vanished pole in that transported tag.

\begin{theorem}[Positive residues and a state-uniform pole-location test]
\label{p4:thm:v2168-faithful-stationary-poles}
Let $\Pi_\omega$ be the spectral projector of $L_{-1}$ at $\omega$.
Then the complete finite source response is
\[
 \boxed{G_\rho(z)=\sum_\omega\frac{Z_{\omega,\rho}}{z-\omega},
 \quad Z_{\omega,\rho}=V^\dagger W_\rho\Pi_\omega V\succeq0,
 \quad\sum_\omega Z_{\omega,\rho}=N_\rho.}
\]
For every faithful stationary $\rho$, a pole is visible in tag $r$ exactly
when $\Pi_\omega v_r\ne0$. Thus its location is common to the entire
faithful stationary family. At a proposed terminal coordinate
$\varepsilon_r$, this is equivalent to the concrete rank test
\[
 \boxed{
 \rank[\,\varepsilon_rI-L_{-1}\ \ v_r\,]
       >\rank(\varepsilon_rI-L_{-1}).}
 \tag{C6}
\]
\end{theorem}
\begin{proof}
For two charge-changing matrices $X,Y$, direct vectorization of
$\Tr\rho\{X^\dagger,Y\}$ gives the displayed $W_\rho$.
Since $[H_c,\rho_c]=0$, $W_\rho$ commutes with $L_{-1}$.
The retarded anticommutator resolvent is therefore (C5), and its spectral
expansion gives the residue matrices. The product
$W_\rho\Pi_\omega$ is positive because its two positive Hermitian factors
commute. If $\rho$ is faithful, $W_\rho$ is positive definite, so the
$r$th residue is positive exactly when $\Pi_\omega v_r\ne0$.

Put $A_r=\varepsilon_rI-L_{-1}$. Hermiticity gives
$\operatorname{Ran}A_r=(\ker A_r)^\perp$.
Consequently $v_r$ has a nonzero component in $\ker A_r$ exactly when
$v_r\notin\operatorname{Ran}A_r$, which is precisely (C6).
This proves the target-location test without fitting a Hamiltonian entry.
\end{proof}

Equivalently, simultaneous energy/charge eigenvectors give
\[
 (Z_{\omega,\rho})_{rs}
 =\sum_{\substack{E_b-E_a=\omega\\c_b=c_a-1}}
    (p_a+p_b)\overline{(q_r)_{ba}}(q_s)_{ba}.
 \tag{C7}
\]
Here the basis is also chosen to diagonalize $\rho^Q$ inside the joint
degeneracies. The frequencies are Bohr differences of the full $H$, not
automatically particle masses. Equations (C4)--(C7) are an executable
finite construction, and contain the same full memory as the Schur route.
Only charge-forbidden operator blocks are removed.



\paragraph{Which state input is discharged, and which remains}
Finite Gibbs densities $e^{-\beta H}/\Tr e^{-\beta H}$ are faithful for
every finite real $\beta$ on a simple physical algebra. Thus no temperature
value is required to determine the common finite pole locations within
this faithful family. Direct-sum states have the same conclusion when
their central weights are all strictly positive. Residues depend on the
weights and on $\beta$. None of these facts identifies a finite-temperature
pole as a physical vacuum mass.

For a narrower physical state family, define the actual eigenoperator
\[
 q_r^{(\omega)}=\Pi_\omega q_r,
 \qquad
 \mathcal A_{r,\omega}
 =\{(q_r^{(\omega)})^\dagger,q_r^{(\omega)}\}\succeq0.
\]
The notation $\Pi_\omega q_r$ means the spectral action of $\operatorname{ad}_H$ on
operators, not a field-carrier projector. Then
\[
 [H,\mathcal A_{r,\omega}]=0,qquad
 [\mathcal Q_{\rm tot},\mathcal A_{r,\omega}]=0,qquad
 (Z_{\omega,\rho})_{rr}=\Tr\rho\mathcal A_{r,\omega}.
\]
The last equality uses stationarity: cross terms between different
Liouvillian frequencies have zero expectation.

\begin{corollary}[Exact visibility test for an entire stationary support family]
\label{p4:cor:v2168-support-family-visibility}
Let $P$ be an already generated projector commuting with $H$, and consider
all stationary densities supported in $\operatorname{Ran}P$. The smallest
tagged residue in this family is
\[
 \boxed{\min_\rho\Tr\rho\mathcal A_{r,\omega}
 =\lambda_{\min}\!\left(P\mathcal A_{r,\omega}P
                               |_{\operatorname{Ran}P}\right).}
\]
Hence the pole is visible in every member precisely when this eigenvalue
is strictly positive. For the ground-support family one may use the
Hamiltonian's generated minimum-energy projector $P_0$; no unique vector
in a degenerate ground space is required for this test.
\end{corollary}
\begin{proof}
Every normalized density on the support gives an expectation bounded
below by the smallest eigenvalue. The compressed positive operator
commutes with $H|_{\operatorname{Ran}P}$, so an eigenvector attaining
this minimum can simultaneously be chosen as an energy eigenvector.
Its pure density is stationary and attains the bound.
\end{proof}

The ground projector here is a mathematical descendant of a specified
physical $H$. Its use as the actual vacuum support must still be licensed
by the physical state/history branch. For an actual state subfamily with
additional constraints, minimizing $\Tr\rho\mathcal A_{r,\omega}$ over
that supplied family is the exact visibility test. A theorem on faithful
states cannot silently be extended to a pure or nonfaithful vacuum.

Paper II-A's co-consumed family-readout principle can now apply to the
common pole-location output on the faithful family, or to a narrower
stationary support family after the corollary's positivity test.
It transports this proved family property. It does not assert that
the common location is $\varepsilon_r$ before (C6) has been proved for
the actual charged Hamiltonian.



\paragraph{The generated root residual also bounds the next pole}
For this fixed finite source chart define
$\kappa_i(\varepsilon)=\lambda_i(S(\varepsilon))$ in increasing order.
These are residual values calculated after the actual action and source
are fixed. They do not select the source or replace a terminal energy.

\begin{theorem}[Regular-gap root certification and displacement bound]
\label{p4:thm:v2168-regular-gap-root-bound}
Let $I$ be a real interval disjoint from $\spec(C_{\rm mem})$. For
$z_2>z_1$ in $I$,
\[
 S(z_2)-S(z_1)\succeq(z_2-z_1)I.
\]
Each ordered eigenvalue of $S$ is strictly increasing and crosses zero
at most once in $I$. For $\varepsilon\in I$, if
$\kappa_i=\lambda_i(S(\varepsilon))\ne0$ and the segment between
$\varepsilon$ and $\varepsilon-\kappa_i$ lies in $I$, that eigenvalue
has exactly one zero $\omega_i$ on that segment, with
\[
 \boxed{|\omega_i-\varepsilon|\le|\kappa_i|.}
\]
If $\kappa_i=0$, the target is an exact root. Source visibility is then
checked on its null channel; root certification alone does not supply it.
\end{theorem}
\begin{proof}
Integrate $S'=I+B^\dagger(zI-C_{\rm mem})^{-2}B\succeq I$.
The min--max characterization of ordered Hermitian eigenvalues gives
strict increase by at least $z_2-z_1$. If $\kappa_i>0$, the eigenvalue at
$\varepsilon-\kappa_i$ is at most zero. If $\kappa_i<0$, its value at
$\varepsilon-\kappa_i$ is at least zero. Continuity and strict increase
give the unique zero and the bound.
\end{proof}

For an actual smooth action/source parameter $\theta$ on the same fixed-rank
chart, a simple root satisfies
\[
 \boxed{\partial_\theta\omega
 =-\frac{u^\dagger(\partial_\theta S)u}{u^\dagger S'u},}
\]
where, at fixed spectral argument,
\[
 \partial_\theta S
 =-\partial_\theta A-(\partial_\theta B)^\dagger RB
                   -B^\dagger R(\partial_\theta B)
                   -B^\dagger R(\partial_\theta C_{\rm mem})RB.
\]
All source-state, kinetic, embedding, constraint and interaction jets
consumed by that parameter are included in these block derivatives.
The marked formula may be differentiated with $\partial_\theta C$ if
the actual current presentation moves. It does not introduce a variable
weak strength to fit a root. No native weak-coupling truncation is assumed.


% END INLINE SOURCE: Paper4-ChargedMemorySpectrum-Main-v1_0.tex



% v2.169: generated action/support constructors and consuming scope.

% BEGIN INLINE SOURCE: Paper4-ActionGeneratedPolePolynomial-Main-v1_0.tex
% Framework v2.169: Batches 7--13, integrated at the owning source.


\subsection{Terminating action-generated tagged pole polynomials}
\label{p4:sec:v2169-action-pole-polynomial}
\paragraph{Result and scope.}
On the already admitted complete finite charged action chart, the repeated
commutators of each actual named source generate a terminating polynomial
calculation. Its roots are exactly that source's visible finite frequencies
throughout the faithful stationary family. A terminal-pole comparison becomes
one scalar equality, $p_r(\varepsilon_r)=0$. The same commutators explicitly
generate the frequency component and its positive physical-state visibility
effect. No separately supplied self-energy, spectral projector, temperature,
right frame or free-response fidelity identity is needed for this calculation.
It remains necessary to evaluate the polynomial of the actual complete action;
its existence does not prove this terminal equality. The finite chart, physical
particle/rest dictionary and consumed limiting claims retain their existing
domains. A05 retains its physical identification and limit status.



\paragraph{Inputs already owned by the framework}
The actual source occurrence and synthesis remain closed in V-A:
\[
 R_{\rm act}=\sum_r\varepsilon_r\iota_r,\qquad
 \iota_r=\Phi_{\rm HG}(A_r^{\rm hier}).
\]
The fixed upstream energies are
\[
 \varepsilon_e=E_{40}(1+\sqrt5)^{-1}(24/25)^{-3/2},\quad
 \varepsilon_\mu=(E_{36}/2)(24/25),\quad
 \varepsilon_\tau=E_{36}(2\sqrt6+\sqrt{12})(24/25).
\]
They are neither defined nor changed by the polynomial below. The inherited
current frame is $U=(g_e,g_\mu,c)$, the neutral order is inverse/direct/Gram,
and $C=U^\dagger(e_2,e_1,e_0)$ is the already generated unitary current.
The primitive momentum resolution remains $m_p=1$; the construction neither
inserts a primitive rest state nor adds a position lattice.

Use the complete finite physical $H=H^\dagger$ and its actual source
operators $q_r$ on the licensed occupation/constraint carrier $\mathcal H$.
Every quadratic, gauge, scalar, pairing, higher-fermion and reduction term
consumed by the response remains in $H$. A finite mode count alone does not
provide a bosonic occupation boundary. The existing A05.1 chart certification
is consumed before a claim about the physical target, where not already proved.
The calculation below itself requires no CAR for a compressed source.

If the licensed charge reduction gives $[H,\mathcal Q_{\rm tot}]=0$ and
$[\mathcal Q_{\rm tot},q_r]=-q_r$, IV already supplies
\[
 \mathfrak B_{-1}=\bigoplus_c\operatorname{Hom}(\mathcal H_c,\mathcal H_{c-1}),
 \qquad D_{-1}=\sum_c d_c d_{c-1},\quad d_c=\dim\mathcal H_c.
\]
Use this reducing operator space. Otherwise the same finite algebraic theorem
acts on the full operator algebra, with bound $(\dim\mathcal H)^2$; that
fallback does not certify an unproved physical charge reduction.



\paragraph{A terminating commutator calculation, before state selection}
Put $\mathcal L=\operatorname{ad}_H$. On the operator algebra define the mathematical
comparison pairing
\[
 (X,Y)_*=\frac{2}{d}\Tr(X^\dagger Y),\qquad d=\dim\mathcal H.
\]
It equals IV's anticommutator pairing for $\rho_*=I/d$, and is positive
definite. This trace is a computation on the constituted finite algebra;
it is not a selected physical state, thermal parameter or vacuum premise.
Trace cyclicity gives $(X,\mathcal LY)_*=(\mathcal LX,Y)_*$.

For a nonzero raw physical tag $q_r$, form
\[
 K_{r,n}=\mathcal L^n q_r,\qquad
 \gamma_{r,n}=(q_r,\mathcal L^nq_r)_*,\qquad n\geq0.
\]
The $\gamma_{r,n}$ are real, with their inherited energy powers. No new
energy normalization is made. If the physical Hamiltonian has a marked
complete word $H=\sum_\alpha H_\alpha$, each column is explicitly
\begin{equation}
 K_{r,n}=\sum_{\alpha_1,\ldots,\alpha_n}
 \operatorname{ad}_{H_{\alpha_n}}\cdots\operatorname{ad}_{H_{\alpha_1}}q_r.
 \label{p4:eq:v2169-b7-p1}
\end{equation}
Thus all mixed interaction and reduction returns are generated before a
moment is read out. This is II-A's preterminal relation rule applied to the
existing complete action. It introduces no covariance or memory function.

\begin{theorem}[First exact source closure and its moment certificate]
\label{p4:v2169:b7:thm:closure}
There is a first $s_r\geq1$ such that
$K_{r,0},\ldots,K_{r,s_r-1}$ are linearly independent and
\begin{equation}
 K_{r,s_r}+\sum_{j=0}^{s_r-1}a_{r,j}K_{r,j}=0.
 \label{p4:eq:v2169-b7-p2}
\end{equation}
On the charge chart $s_r\leq D_{-1}$. For an independent prefix of length
$s$ define
\[
 \mathsf H_{r,s}=(\gamma_{r,i+j})_{i,j=0}^{s-1},\qquad
 b_{r,s}=(\gamma_{r,s+i})_{i=0}^{s-1}.
\]
Then $\mathsf H_{r,s}>0$, the uniquely generated tentative coefficients are
$a=-\mathsf H_{r,s}^{-1}b_{r,s}$, and the exact residual norm is
\begin{equation}
 \delta_{r,s}=\gamma_{r,2s}
       -b_{r,s}^\dagger\mathsf H_{r,s}^{-1}b_{r,s}\geq0.
 \label{p4:eq:v2169-b7-p3}
\end{equation}
Closure \eqref{p4:eq:v2169-b7-p2} holds exactly at the first $\delta_{r,s}=0$. Every previous
residual is strictly positive. A nonzero residual cannot be dropped to
declare an exact pole calculation.
\end{theorem}
\begin{proof}
All iterates lie in the finite reducing operator space. The first linear
dependence gives \eqref{p4:eq:v2169-b7-p2} with a unique monic leading coefficient. Its independent
prefix has the positive Gram $\mathsf H_{r,s}$, since
$(K_{r,i},K_{r,j})_* =\gamma_{r,i+j}$. Orthogonal projection of $K_{r,s}$
onto this prefix gives $a=-\mathsf H^{-1}b$ and squared residual \eqref{p4:eq:v2169-b7-p3}.
The pairing is faithful, so zero residual is the operator identity \eqref{p4:eq:v2169-b7-p2},
rather than a state-null approximation. Once \eqref{p4:eq:v2169-b7-p2} holds, applying
$\mathcal L$ proves that the prefix span is invariant. Self-adjointness
makes it reducing. The dimension bound follows.
\end{proof}

Define the monic source polynomial
\begin{equation}
 \boxed{p_r(z)=z^{s_r}+\sum_{j=0}^{s_r-1}a_{r,j}z^j.}
 \label{p4:eq:v2169-b7-p4}
\end{equation}
Its coefficients are outputs of the complete source word \eqref{p4:eq:v2169-b7-p1}--\eqref{p4:eq:v2169-b7-p3}.
If an actual constraint quotient makes the raw $q_r$ the zero operator,
set $p_r=1$ and all tagged residues to zero; there is no visible pole in
that tag. No source normalization repairs this case.



\paragraph{The minimal finite terminal-pole obligation}
\begin{theorem}[Exact scalar source-pole test on the faithful family]
\label{p4:v2169:b7:thm:poles}
The nonconstant $p_r$ has real simple roots. They are exactly the frequencies
$\omega$ for which the actual operator component
$\Pi_\omega q_r$ is nonzero. Consequently, for every faithful stationary
density on this finite chart,
\begin{equation}
 \boxed{\text{tag }r\text{ has a visible finite pole at }\varepsilon_r
       \quad\Longleftrightarrow\quad p_r(\varepsilon_r)=0.}
 \label{p4:eq:v2169-b7-p5}
\end{equation}
The condition includes visibility in the named tag, not merely an eigenvalue
of the ambient Hamiltonian or a zero of its characteristic determinant.
It remains valid at a singular eliminated memory block. It is equivalent
to IV's full charged rank test on its certified charge chart.
\end{theorem}
\begin{proof}
By finite self-adjointness write $q_r=\sum_\omega q_r^{(\omega)}$ with
orthogonal frequency components. An annihilating polynomial must vanish at
every $\omega$ having $q_r^{(\omega)}\ne0$. Conversely the product of the
distinct corresponding linear factors annihilates $q_r$. The Vandermonde
matrix makes their frequency components a basis of its scalar cyclic span.
Thus this product is exactly the first monic polynomial \eqref{p4:eq:v2169-b7-p4}; its roots are
simple even if the ambient frequency eigenspace has higher multiplicity.
For any faithful stationary density the source residue is positive precisely
when that operator component is nonzero, by IV's positive pairing theorem.
This proves \eqref{p4:eq:v2169-b7-p5}. The rank test is the same nonzero projection condition,
and no complementary inverse was used here.
\end{proof}

An equivalent exact moment-pencil test is particularly concrete. After
closure define $\mathsf K_{r,s}=(\gamma_{r,i+j+1})_{i,j=0}^{s-1}$. Then
\begin{equation}
 \boxed{p_r(z)=\frac{\det(z\mathsf H_{r,s}-\mathsf K_{r,s})}
                        {\det\mathsf H_{r,s}},\qquad s=s_r.}
 \label{p4:eq:v2169-b7-p6}
\end{equation}
Indeed, in the cyclic basis the matrix of $\mathcal L$ is the companion
matrix $T_r$ with last column $-a$ and subdiagonal ones. Its exact Gram
identity is $\mathsf K=\mathsf H T_r$, so \eqref{p4:eq:v2169-b7-p6} follows. A shorter unclosed
moment pencil supplies a projected calculation and is not this theorem.

This discharges a separate mathematical existence/provision problem: all
required memory and pole data can be generated in finitely many exact
commutator operations from the complete finite action. It does not discharge
the evaluation $p_r(\varepsilon_r)=0$ for the framework's actual physical
action. II-A permits retaining a proved common location on the faithful
family; it does not turn this still unevaluated equality into an identity.



\paragraph{The frequency mode and physical visibility are generated too}
For a certified root $\varepsilon$ of $p_r$, define the divided polynomial
$\widetilde p_{r,\varepsilon}(z)=p_r(z)/(z-\varepsilon)$. Then
\begin{equation}
 \boxed{O_{r,\varepsilon}
    =\frac{\widetilde p_{r,\varepsilon}(\mathcal L)q_r}
           {p_r'(\varepsilon)},\qquad
       [H,O_{r,\varepsilon}]=\varepsilon O_{r,\varepsilon}.}
 \label{p4:eq:v2169-b7-p7}
\end{equation}
The denominator is nonzero because the source polynomial is square-free.
Evaluation on its distinct frequency components proves
$O_{r,\varepsilon}=\Pi_\varepsilon q_r$. Thus the original action word
generates the needed mode without a separately supplied projector.

For an actual stationary physical density or licensed support family put
\begin{equation}
 \boxed{\mathcal A_{r,\varepsilon}
       =O_{r,\varepsilon}^\dagger O_{r,\varepsilon}
        +O_{r,\varepsilon}O_{r,\varepsilon}^\dagger\succeq0,\qquad
       Z_{r,\varepsilon,\rho}=\Tr\rho\mathcal A_{r,\varepsilon}.}
 \label{p4:eq:v2169-b7-p8}
\end{equation}
Both products commute with $H$; on the charge chart they also have charge
zero. Formula \eqref{p4:eq:v2169-b7-p8} is the exact raw-tag residue. A physical nonfaithful
state sees the root iff this expectation is positive. Equivalently, its
absence means both $O\rho^{1/2}=0$ and $O^\dagger\rho^{1/2}=0$.
For every stationary density on an already licensed reducing support $P$, a
uniform positive lower bound is
\begin{equation}
 \inf_\rho Z_{r,\varepsilon,\rho}
       =\lambda_{\min}(P\mathcal A_{r,\varepsilon}P|_{P\mathcal H}).
 \label{p4:eq:v2169-b7-p9}
\end{equation}
The minimum is attained by an eigenvector that can also be taken in an
$H$ eigenspace, since the compressed effect and $H$ commute. This assumes
the full stationary support family; a narrower actually generated family
uses its own expectation domain. A ground support is used only where the
physical claim has licensed that support.

For completeness the whole tagged response is generated from the same
polynomial. Write $p_r(z)=\sum_{k=0}^{s_r}a_kz^k$, $a_{s_r}=1$, and use the
actual stationary moments $\mu_{r,j,\rho}=(q_r,\mathcal L^jq_r)_\rho$. Then
\begin{equation}
 G_{rr,\rho}(z)=\frac{n_{r,\rho}(z)}{p_r(z)},\quad
 n_{r,\rho}(z)=\sum_{k=1}^{s_r}a_k
                 \sum_{j=0}^{k-1}z^{k-1-j}\mu_{r,j,\rho}.
 \label{p4:eq:v2169-b7-p10}
\end{equation}
This follows by dividing $p_r(z)-p_r(\mathcal L)$ by $z-\mathcal L$ and
using $p_r(\mathcal L)q_r=0$. At a root its residue is
$n_{r,\rho}(\varepsilon)/p_r'(\varepsilon)$, equal to \eqref{p4:eq:v2169-b7-p8}. In a nonfaithful
state a numerator cancellation is precisely a vanished physical visibility
effect; the faithful-family denominator itself is unchanged.



\paragraph{Physical tags survive inherited source normalization}
Suppose the actual operator source span has an invertible positive Gram $N$, and
$\mathcal Q_{\rm raw}=\mathcal Q N^{1/2}$ is its inherited normalization.
The physical tag $r$ is the same operator
\[
 q_r=\mathcal Q_{\rm raw}e_r=\mathcal Q t_r,\qquad t_r=N^{1/2}e_r.
\]
Run \eqref{p4:eq:v2169-b7-p1}--\eqref{p4:eq:v2169-b7-p10} on this $q_r$, or on precisely its transported expression.
It has the same polynomial and effect. Replacing it by $\mathcal Q e_r$
after a mixing normalization can change the scalar source polynomial.
An invertible mixing preserves the joint source cyclic space and its union
of visible roots, but not necessarily the roots visible in each named tag.
A same-owner unitary presentation change of $H,q_r$ preserves all operator
relations, polynomial coefficients and transported visibility effects.
These are comparison invariances, not permissions to choose a physical
kinetic block or a new source. The propagation kinetic Gram remains distinct.
If a nonfaithful state-null quotient is taken first, it can remove invisible
frequency components. Its normalization preserves that quotient's response,
but need not preserve the original faithful-family polynomial. Form \eqref{p4:eq:v2169-b7-p4}
before that state quotient, then apply \eqref{p4:eq:v2169-b7-p8} to the actual physical state.

% END INLINE SOURCE: Paper4-ActionGeneratedPolePolynomial-Main-v1_0.tex


% BEGIN INLINE SOURCE: Paper4-PhysicalSupportPolynomial-Main-v1_1.tex
% Framework v2.169: Batches 7--13, integrated at the owning source.


\subsection{Actual stationary support and source-visible polynomials}
\label{p4:sec:v2169-support-polynomial}
\paragraph{Result.}
For a physical stationary state, its support and the complete action fix all
tagged finite pole locations; the density weights inside that support are
unneeded for this particular output. The source polynomial can be generated
directly on that physical support, even when the state is nonfaithful on the
full carrier. Any root on that exact support already includes physical
visibility. Testing whether a fixed terminal coefficient is one of those roots
is a later comparison, not an input to the polynomial construction. This
removes a separate overlap obligation after that support-root test has been
proved. On a licensed single-energy or ground support, a further exact
construction uses only adjacent charge-sector excitation matrices.
Neither theorem selects a support or a physical vacuum from mathematical
availability. The terminal coefficients and $m_p=1$ remain fixed.



\paragraph{Existing physical support, rather than a chosen stationary member}
Let the complete finite physical Hamiltonian be $H=H^\dagger$ and the actual
stationary state have support $P$. Stationarity gives $[H,P]=0$. The support
is an output of the occurring preparation/history or of an already licensed
state constructor; it is not supplied by choosing a desired pole.
Use the mathematical support comparison density
\[
 \rho_P=P/k,\qquad k=\operatorname{rank} P,
\]
only to calculate the existing support's response. It is not a replacement
physical state. Its anticommutator pairing is
\begin{equation}
 (X,Y)_P=\frac1k\Tr P\{X^\dagger,Y\},\qquad
 \|X\|_P^2=\frac1k
       \left(\|XP\|_{\rm HS}^2+\|X^\dagger P\|_{\rm HS}^2\right).
 \label{p4:eq:v2169-b10-s1}
\end{equation}

\begin{theorem}[Pole locations depend only on actual stationary support]
\label{p4:v2169:b10:thm:support}
Every stationary density $\rho$ with exactly support $P$ has the same
tagged finite source pole locations. For any operator $X$,
\begin{equation}
 k\lambda_{\min}(\rho|_P)\|X\|_P^2
 \leq\|X\|_\rho^2
 \leq k\lambda_{\max}(\rho|_P)\|X\|_P^2.
 \label{p4:eq:v2169-b10-s2}
\end{equation}
Thus all these densities have exactly the same invariant source-null space
and the same scalar source polynomial $p_{r,P}$ of the induced
$\mathcal L=\operatorname{ad}_H$. In particular
\begin{equation}
 \boxed{\text{the actual physical tag }r\text{ has a finite pole at }
 \varepsilon_r\quad\Longleftrightarrow\quad
                         p_{r,P}(\varepsilon_r)=0.}
 \label{p4:eq:v2169-b10-s3}
\end{equation}
No additional physical visibility test is required after \eqref{p4:eq:v2169-b10-s3} on that same
certified actual support. Residues still depend on $\rho$.
\end{theorem}
\begin{proof}
On $P\mathcal H$, $\rho$ is positive definite. Apply its lower and upper
eigenvalue bounds to the positive compression
$P(X^\dagger X+XX^\dagger)P$ to get \eqref{p4:eq:v2169-b10-s2}.
Hence all pairings have the common null space $XP=X^\dagger P=0$.
Since $P$ commutes with $H$, that null space is invariant under $\operatorname{ad}_H$.
The induced generator is finite self-adjoint by IV's source-quotient proof.
Theorem~\ref{p4:v2169:b7:thm:closure}'s cyclic construction on this quotient has simple real roots and
includes a root iff its frequency component is not null. The strictly
positive factors in \eqref{p4:eq:v2169-b10-s2} make this criterion common to every state with
exactly that support. This proves \eqref{p4:eq:v2169-b10-s3}.
\end{proof}

Evaluation is forward: construct $p_{r,P}$ from $(H,q_r,P)$, determine its
source-visible root set, and only then apply
\eqref{p4:eq:v2169-b10-s3} when the claim specifically identifies
$\varepsilon_r$ with a physical pole.  A nonzero
$p_{r,P}(\varepsilon_r)$ is therefore a failed terminal-to-pole equality on
that branch, not a reason to choose a different support, delete an action term,
or fit a replacement terminal coefficient.

The distinction between exact support and a containing support matters.
Knowing only $\rho=P\rho P$ permits smaller supports with different visible
roots. For a family allowing all such states, the full positive effect test
or the actual smaller support is still needed. An already proved physical
root is not retested by a second independent visibility assumption.



\paragraph{The quotient polynomial is generated by the complete action}
Use the actual raw source $q_r$, retaining its physical tag through any
kinetic normalization. Set
\[
 \mu^P_{r,n}=(q_r,\mathcal L^nq_r)_P,
 \quad \mathsf H^P_{r,s}=(\mu^P_{r,i+j}),
 \quad b^P_{r,s}=(\mu^P_{r,s+i}).
\]
For the first independent quotient prefix, solve
$a=-(\mathsf H^P_{r,s})^{-1}b^P_{r,s}$. Its exact closure test is
\begin{equation}
 \boxed{\delta^P_{r,s}=\mu^P_{r,2s}
       -(b^P_{r,s})^\dagger(\mathsf H^P_{r,s})^{-1}b^P_{r,s}=0.}
 \label{p4:eq:v2169-b10-s4}
\end{equation}
The residual $p_{r,P}(\mathcal L)q_r$ may be a nonzero operator on the
whole carrier. What \eqref{p4:eq:v2169-b10-s4} proves exactly is
\begin{equation}
 p_{r,P}(\mathcal L)q_rP=0,\qquad
 (p_{r,P}(\mathcal L)q_r)^\dagger P=0.
 \label{p4:eq:v2169-b10-s5}
\end{equation}
That is the actual source quotient identity. There is no discarded
interaction: the full $H$ generates every commutator before the null test.
If $q_r$ is already null, use $p_{r,P}=1$, which has no visible root.
An unclosed projected moment pencil cannot assert \eqref{p4:eq:v2169-b10-s3}.

At a certified root the divided polynomial supplies a representative
$O^P_{r,\varepsilon}$ whose commutator equation holds in the quotient.
Its compressed positive effect is
\begin{equation}
 \boxed{\mathcal M^P_{r,\varepsilon}
     =P\{(O^P_{r,\varepsilon})^\dagger,O^P_{r,\varepsilon}\}P,\qquad
 Z_{r,\varepsilon,\rho}=\Tr\rho\mathcal M^P_{r,\varepsilon}.}
 \label{p4:eq:v2169-b10-s6}
\end{equation}
The representative need not be a full-carrier eigenoperator. Its action on
$P$ and that of its adjoint equal those of the actual frequency component,
which suffices for \eqref{p4:eq:v2169-b10-s6}. A state-null residual cannot be promoted to a
global operator identity. If the exact support is later changed by a real
process, recompute or transport its quotient with that process.



\paragraph{A smaller charge-response dimension is also generated}
On an equivariant actual charge chart, IV permits pinching the state in
charge without changing a charged-source anticommutator response. Its
pinched support is generated by $P$ alone:
\begin{equation}
 \overline P=\operatorname{supp}\left(\sum_cP_cPP_c\right).
 \label{p4:eq:v2169-b10-s7}
\end{equation}
Indeed \eqref{p4:eq:v2169-b10-s2}'s positive bounds show that $P_c\rho P_c$ and $P_cPP_c$ have
the same support. This pinched state is a response-equivalent descendant,
not a physical state selection. The polynomial of a definite-charge source
has the same roots, and hence the same monic square-free polynomial, on
$P$ and on $\overline P$.

Write $d_c=\dim\mathcal H_c$ and
$k_c=\operatorname{rank}(P_c\overline P)$. The exact response quotient dimension of
the full charge-minus-one operator space is
\begin{equation}
 \boxed{D_P=\sum_c
 [k_{c-1}d_c+d_{c-1}k_c-k_{c-1}k_c].}
 \label{p4:eq:v2169-b10-s8}
\end{equation}
In each block an operator is null precisely if both its rows on the retained
target support and its columns on the retained input support vanish.
The null block therefore has dimension
$(d_{c-1}-k_{c-1})(d_c-k_c)$. Subtract from $d_{c-1}d_c$ to get \eqref{p4:eq:v2169-b10-s8}.
The scalar cyclic closure degree is at most this generated dimension.
Without licensed charge equivariance the full operator quotient has
dimension $k(2d-k)$; no unproved charge projection is required for that
fallback. These are calculation bounds, not added field multiplicities.



\paragraph{A single-energy support gives an adjacent-sector construction}
Suppose the actual licensed support has fixed energy and charge,
\[
 HP=E_*P,\qquad Q_{\rm tot}P=c_*P,
\]
and the raw source has charge $-1$. Define the actual maps
\[
 Q_+=q_rP:P\mathcal H\to\mathcal H_{c_*-1},\qquad
 Q_-=q_r^\dagger P:P\mathcal H\to\mathcal H_{c_*+1},
\]
and $T_\pm=H_{c_*\mp1}-E_*I$ on their respective carriers.
Every density with support contained in $P$ is stationary. Its exact response is
\begin{equation}
 \boxed{G_{rr,\rho}(z)=\Tr_P\rho
 \left[Q_+^\dagger(z-T_+)^{-1}Q_+
             +Q_-^\dagger(z+T_-)^{-1}Q_-\right].}
 \label{p4:eq:v2169-b10-s9}
\end{equation}
To prove this, insert the energy resolution of the complete $H$ on the two
vectors $q_rP$ and $q_r^\dagger P$. Their source frequencies are
$E_b-E_*$ and $E_*-E_a$, respectively, which gives the two resolvents with
the displayed signs. This also retains the removal channel when $P$ is only
an energy eigen-support rather than the global ground.

Compute the two exact cyclic polynomials from
$T_+^nQ_+$ and $T_-^nQ_-$ with the positive matrix Hilbert--Schmidt pairing.
Their degrees are at most $d_{c_*-1}$ and $d_{c_*+1}$, respectively, because
each acts only on the target energy space. If their monic polynomials are
$p_+$ and $p_-$, the physical support polynomial is
\begin{equation}
 \boxed{p_{r,P}(z)=\operatorname{lcm}
       \left(p_+(z),(-1)^{\deg p_-}p_-(-z)\right).}
 \label{p4:eq:v2169-b10-s10}
\end{equation}
Zero source maps contribute the polynomial one. Positive residue matrices
prevent cancellations between the two legs at a shared root.

For a positive target $\varepsilon$, the compressed visibility matrix is
explicitly generated without the full operator Liouvillian:
\begin{equation}
 \boxed{\mathcal M_{r,\varepsilon}
   =Q_+^\dagger\Pi_{+\varepsilon}(T_+)Q_+
          +Q_-^\dagger\Pi_{-\varepsilon}(T_-)Q_-.}
 \label{p4:eq:v2169-b10-s11}
\end{equation}
The cyclic divided polynomials produce each projector on its source image;
no new physical spectral projector is provided independently.



\paragraph{Ground-family and branch-origin reuse, at their actual domains}
If this $P$ is an already licensed global-ground support, $T_\pm\succeq0$.
For $\varepsilon>0$, the second term of \eqref{p4:eq:v2169-b10-s11} then vanishes.
On this branch the remaining finite particle-energy comparison is simply
$p_+(\varepsilon_r)=0$ with the raw tag $Q_+=q_rP$; physical source visibility
is built into that polynomial for densities with exact support $P$.
The finite symmetric Gibbs vacuum of
\ref{p4:thm:v164-finite-gibbs-vacuum-no-thermalization} has exact support
$P_0$ and is therefore one already proved reuse route on its declared vacuum
branch. Its mathematical construction does not establish that a generic
charged physical claim uses that branch.

For a family allowing every density supported in a degenerate ground $P$,
the same positive pole is visible throughout the family iff
\begin{equation}
 \boxed{\mathcal M_{r,\varepsilon}>0\ \text{on }P\mathcal H.}
 \label{p4:eq:v2169-b10-s12}
\end{equation}
Its smallest eigenvalue is the uniform lower bound. For a unique ground,
nonzero $\mathcal M$ already implies \eqref{p4:eq:v2169-b10-s12}, so no extra state member or
overlap datum is needed. For higher rank, a nonzero but singular $\mathcal M$
does not license a common pole for every pure ground orientation.

The neutral fresh-branch theorem
\ref{p4:thm:v2114-neutral-fock-unit} supplies the Fock unit when its
no-inherited-state/no-preparation hypotheses actually hold. If that constructor
is the state origin for the claimed charged chart, reuse its state instead of
requesting a new occurrence. Its next consumed stationary condition is the
complete $H\Omega=E_*\Omega$. The Fock unit is not automatically the complete
interacting ground, and the previously closed neutral-root QROOT theorem is
not reopened or promoted to a different charged physical domain.

% BEGIN C01 V2174 FULL GROUND BOUND

% BEGIN INLINE SOURCE: Paper4-FullGroundPolynomialBound-Main-v1_0.tex
% v2.174: counting consequence of the already closed exact-support theorem.


\paragraph{A finite bound on the full-ground tagged calculation}
\begin{corollary}[Full-ground polynomial degree and moment bound]
\label{p4:cor:v2174-full-ground-degree-bound}
Let the specified complete finite $H$ act on dimension $D$, with full ground
$P_0$ of rank $d_0$ and $m$ distinct excited energies $E_j>E_0$.
Set $\lambda_j=E_j-E_0>0$,
\[
 J_+=\{j:\Pi_jq_eP_0\ne0\},\qquad
 J_-=\{j:\Pi_jq_e^\dagger P_0\ne0\},\qquad
 \delta=\mathbf1_{P_0q_eP_0\ne0}.
\]
On this exact stationary support the existing tagged polynomial is
\begin{equation}
 p_{e,P_0}(z)=z^\delta
       \prod_{j\in J_+}(z-\lambda_j)
       \prod_{j\in J_-}(z+\lambda_j),\qquad
 \deg p_{e,P_0}\le2m+\delta\le2(D-d_0)+1.
 \label{p4:eq:v2174-full-ground-degree-bound}
\end{equation}
There are at most $D-d_0$ positive tagged roots. The exact moment
construction requires orders only through $2\deg p_{e,P_0}$, hence at most
$4(D-d_0)+2$.
\end{corollary}
\begin{proof}
The spectral resolution in
Theorem~\ref{p4:v2169:b10:thm:support} and
\eqref{p4:eq:v2169-b10-s9} gives only zero and the signed excitation gaps.
The two source legs have positive squared matrix weights, so a gap occurs
exactly under the displayed nonzero test. Distinct positive gaps are
different, and a positive gap cannot coincide with a negative one.
At zero the adjoint ground compression is nonzero exactly when the ground
compression is nonzero. This proves the square-free product, including
the polynomial one for a null tag. Each distinct excited energy requires
at least one vector outside the ground, so $m\le D-d_0$.
The residual norm test~\eqref{p4:eq:v2169-b10-s4} uses moments through
order twice the closure degree.
\end{proof}
Every nonzero conserved-fermion-degree ground summand and both adjacent
source legs remain in this calculation. Its bound applies after the actual
finite $H,q_e,P_0$ have been supplied; it neither selects $D$ nor bounds
the search for their missing constructor. Exact execution also consumes an
effective coefficient domain with decidable linear dependence and real
spectral ordering. No terminal coefficient is imposed as a root.

% END INLINE SOURCE: Paper4-FullGroundPolynomialBound-Main-v1_0.tex

% END C01 V2174 FULL GROUND BOUND

% END INLINE SOURCE: Paper4-PhysicalSupportPolynomial-Main-v1_1.tex


% BEGIN INLINE SOURCE: Paper4-SourceCyclicDeterminacy-Main-v1_0.tex
% Framework v2.186: exact source-cyclic determination and scalar-centre scope.


\subsection{Determination on the actual source-cyclic quotient}
\label{p4:sec:v2186-source-cyclic-determinacy}

The complete action and actual stationary support determine a tagged pole
through their source-cyclic quotient.  A complete finite representation is
one sufficient way to evaluate that quotient.  The following criterion also
retains a common pole polynomial when different independently licensed
representations have the same source-cyclic action.  All action terms that
can enter this action or change the consumed support remain included.

Fix the admitted finite stationary chart of
Proposition~\ref{p4:prop:v2168-stationary-source-quotient}, with complete
$H=H^\dagger$, actual raw source $q$ and exact stationary support $P$.
Use the comparison pairing of
Theorem~\ref{p4:v2169:b10:thm:support}, and put
\[
 \mathcal N_P=\{X:XP=X^\dagger P=0\},\qquad
 \mathscr O_P=\End(\mathcal H)/\mathcal N_P,\qquad
 \mathcal L[X]=[[H,X]],\qquad v=[q].
\]
Here $[[H,X]]$ denotes the quotient class of the complete commutator;
the inherited energy and clock normalization is retained.  The null space
is $\mathcal L$-invariant and the induced operator is self-adjoint.

\begin{lemma}[Exact cyclic action and absence of re-entry]
\label{p4:lem:v2186-exact-source-cyclic-action}
Let
\[
 \mathscr C_{q,P}
 =\operatorname{span}\{\mathcal L^n v:n\ge0\}\subseteq\mathscr O_P
\]
and let $\Pi_{q,P}$ be its orthogonal projector.  Then
\[
 [\Pi_{q,P},\mathcal L]=0,\qquad
 \Pi_{q,P}\mathcal L(I-\Pi_{q,P})=0.
\]
For every nonreal $z$ the exact tagged response is
\[
 (v,(zI-\mathcal L)^{-1}v)_P
 =\bigl(v,(zI_{\mathscr C_{q,P}}
             -\mathcal L|_{\mathscr C_{q,P}})^{-1}v\bigr)_P.
\]
Its first monic annihilating polynomial is the same on this cyclic space
as on the full support quotient.  The orthogonal complement has zero
response in this tag and cannot return to it.
\end{lemma}
\begin{proof}
The span is finite and invariant under $\mathcal L$ by its definition.
For $x\in\mathscr C_{q,P}^{\perp}$ and $y\in\mathscr C_{q,P}$,
self-adjointness gives
$(y,\mathcal Lx)_P=(\mathcal Ly,x)_P=0$.
Thus the complement is also invariant and the projector commutes with
$\mathcal L$.  The resolvent is block diagonal on the two reducing spaces,
while $v$ belongs to the first; this proves the response identity.
For any polynomial $a$, the vector $a(\mathcal L)v$ already lies in the
cyclic space, so its vanishing is the same condition in both spaces.
\end{proof}

This is a restriction of the \emph{actual complete} source action, not
permission to delete a term whose iterated returns or support change have
not been evaluated.  An independently derived cyclic restriction must
agree with this action, including constraint, occupation-boundary and mixed
interaction returns.  The lemma does not infer that agreement from a
short projected moment pencil.

\begin{theorem}[Common cyclic-intertwiner polynomial criterion]
\label{p4:thm:v2186-common-source-cyclic-polynomial}
Let $\theta$ range over a nonempty independently licensed family of complete
finite stationary source charts, with actual raw tags $q_\theta$, exact
supports $P_\theta$ and the same inherited energy/clock unit.  Let
\[
 v_\theta=[q_\theta]\ne0,\qquad
 T_\theta=\mathcal L_\theta|_{\mathscr C_\theta},\qquad
 \mathscr C_\theta=\mathscr C_{q_\theta,P_\theta}.
\]
Write $p_\theta$ for the first monic polynomial satisfying
$p_\theta(T_\theta)v_\theta=0$.  The following are equivalent:
\begin{enumerate}
\item The polynomial $p_\theta=p$ is common throughout the family.
\item There is one finite cyclic pair $(T,v)$ and, for every $\theta$,
a linear isomorphism $U_\theta:\mathscr C\to\mathscr C_\theta$ with
\[
 \boxed{U_\theta v=v_\theta,\qquad
               T_\theta U_\theta=U_\theta T,\qquad
 \mathscr C=\operatorname{span}\{T^nv:n\ge0\}.}
\]
\end{enumerate}
When either condition is proved, $p$ has simple real roots, and those roots
are the common visible tagged finite pole locations on the exact supports.
No isomorphism of the total action carriers, unique total Hamiltonian or
unique total occupation member is required by this criterion.
\end{theorem}
\begin{proof}
Under (2), for every polynomial $a$,
\[
 a(T_\theta)v_\theta=U_\theta a(T)v.
\]
Injectivity makes the annihilating polynomial ideals equal, proving (1).
Conversely, choose one $\theta_0$ and take
$(T,v)=(T_{\theta_0},v_{\theta_0})$.  Define
\[
 U_\theta\bigl(a(T)v\bigr):=a(T_\theta)v_\theta.
\]
For a finite cyclic pair, its first monic annihilator $p$ generates the
annihilating ideal: division $a=bp+r$ with $\deg r<\deg p$ reduces the
claim to $r(T)v=0$, and minimality makes that remainder zero.  Equality
of the first monic polynomials therefore makes the displayed map
well-defined and injective.  Cyclicity makes it surjective; multiplication
by the variable gives the intertwining equation.  This proves (2).
Each $T_\theta$ is a restriction of the finite self-adjoint quotient
operator.  Its scalar cyclic annihilator is the product of the distinct
real eigenvalue factors with nonzero components of $v_\theta$.
The actual-support theorem identifies precisely those components with
visible tagged poles.  Lemma~\ref{p4:lem:v2186-exact-source-cyclic-action}
shows that no complementary response was omitted.
\end{proof}

The converse constructs a comparison map from an already proved common
polynomial; it does not prove commonness by defining the physical family
to have a desired response.  To use the forward criterion, the actual
source/support provenance and the complete intertwining identities must
be established independently.  A null tag has $p=1$ and zero cyclic space;
that trivial case is handled separately and supplies no visible pole.

\begin{corollary}[Residues require the transported physical pairing]
\label{p4:cor:v2186-source-cyclic-residue-scope}
The common-polynomial criterion fixes pole locations.  If, in addition,
each $U_\theta$ is an isometry for the \emph{actual} stationary source
pairings and one common pairing on $\mathscr C$, then the complete scalar
tagged responses and their residues are common:
\[
 (v_\theta,(zI-T_\theta)^{-1}v_\theta)_{\rho_\theta}
 =(v,(zI-T)^{-1}v)_*.
\]
Without this additional condition, the common roots need not have common
residues or common density weights.
\end{corollary}
\begin{proof}
The intertwining equation gives
$(zI-T_\theta)^{-1}U_\theta=U_\theta(zI-T)^{-1}$.
Apply the isometry and transported-vector identities.  The meromorphic
functions are equal, hence so are their residues.  The actual-support
theorem only identifies the invariant null space when the density weights
vary; it does not equate those pairings or residues.
\end{proof}

\begin{corollary}[Scalar-centre invisibility for a fixed stationary pole]
\label{p4:cor:v2186-scalar-center-pole-invisibility}
On a fixed admitted stationary source chart, put $H'=H+cI$ with real
$c$, retaining the same raw source, physical density and support.  Then
\[
 \boxed{\operatorname{ad}_{H'}=\operatorname{ad}_H,\qquad
          p'_{q,P}=p_{q,P},\qquad G'_{q,\rho}=G_{q,\rho}.}
\]
Moreover $P_0(H')=P_0(H)$, with full ground degeneracy retained.  A
normalized finite Gibbs state at a fixed inverse temperature is also
unchanged.  Thus a unique scalar central value is not a separate input
for this fixed pole-location or fixed-state residue readout.
\end{corollary}
\begin{proof}
$[cI,X]=0$ for every $X$, so the quotient, cyclic action and response
resolvent are unchanged.  The eigenvalues of $H$ all increase by $c$,
without changing their projectors, order or ground multiplicity.
In the normalized Gibbs density the common factor $e^{-\beta c}$
cancels between numerator and partition function.
\end{proof}

This conclusion is restricted to that fixed chart.  A metric/source family
$H(g)+c(g)I$ still has central first and second derivatives in its
energy, stress and metric response; absolute ground energies and partition
functions also retain the scalar value.  Sector-dependent central terms,
terms changing the source/support chart, and couplings that can re-enter
the source action are retained on their own consumed domains.

\paragraph{Charged-source application.}
For a native pole claim, determine the actual complete source-cyclic action
and exact support, or prove their pole polynomial common throughout the
independently licensed family by the theorem above.  Constructing one
complete finite member is a sufficient implementation of the first route.
Neither route has been evaluated here for the actual electron: its external,
dual-field and complete remainder returns remain to determine.  Native
terminal coefficients are unchanged; comparison with them, physical
rest/particle identification and every limit actually consumed retain
their existing domains.  In particular this criterion is not itself a
closure of C01/A05.

% END INLINE SOURCE: Paper4-SourceCyclicDeterminacy-Main-v1_0.tex


% BEGIN INLINE SOURCE: Paper4-SourceMomentPencil-Main-v1_0.tex
% Framework v2.188: source-cyclic finite moment pencil; exact-support and
% complete-Hamiltonian scope, no chosen mass and no projected-moment shortcut.


\subsection{Additional receiver closures and interacting pole comparisons}
\label{p4:app:v2169-polynomial-comparisons}


\paragraph{An additional closed weak-receiver chart and its native polynomial}
\label{p4:v2169:b7:app:weakchart}
This appendix is an exact diagnostic of a further dynamical closure, not a
claim that the actual complete interacting action has only these nine
cyclic directions. The extra hypothesis is that its source/partner/receiver
space closes with the complete self-adjoint block
\begin{equation}
 \mathcal T=\begin{pmatrix}
 0&M&\lambda C\\ M&0&0\\ \lambda C^\dagger&0&N_\nu
 \end{pmatrix},\quad
 M=\diag(\varepsilon_e,\varepsilon_\mu,\varepsilon_\tau),\quad
 N_\nu=\diag(\nu_{\rm inv},\nu_{\rm dir},\nu_{\rm Gram}).
 \label{p4:eq:v2169-b7-d1}
\end{equation}
Here $\lambda$ is the actual normalized common weak receiver scale, when
such a chart has been proved. With the trace ports of \eqref{p5c:eq:v2169-b7-p12} and $\kappa>0$,
the receiver frame is $f_j=Y_j/\sqrt\kappa$ and
$\lambda=\pi\sqrt\kappa$; it is not an independently chosen coupling.
The $\nu_j$ are receiver eigenvalues generated
by its action, not terminal neutral coefficients installed as denominators.
Proving that no additional source-visible commutator columns exist is part
of the extra closure; neither \eqref{p5c:eq:v2169-b7-p12} nor the current matrix proves it.

Nevertheless its complete determinant can be calculated without deleting
the coherent current. Put $d_r=z^2-\varepsilon_r^2$, $n_j=z-\nu_j$, and
$t=\lambda^2z$. The generated exact current weights are
\[
 \mathcal P=|C|^2=\begin{pmatrix}
 1/4&1/4&1/2\\5/12&5/12&1/6\\1/3&1/3&1/3
 \end{pmatrix}.
\]
Block elimination and complementary unitary minors give the polynomial
identity, valid even at singular intermediate denominators,
\begin{equation}
 \boxed{\det(zI-\mathcal T)
 =\prod_rd_r\prod_jn_j
 -t\sum_{rj}\mathcal P_{rj}\prod_{s\ne r}d_s\prod_{k\ne j}n_k
 +t^2\sum_{rj}\mathcal P_{rj}d_rn_j-t^3.}
 \label{p4:eq:v2169-b7-d2}
\end{equation}
For a proof, away from intermediate zeros eliminate the right partners.
The remaining block is $\left(\begin{smallmatrix}\diag(d_r)&-z\lambda C\\
-\lambda C^\dagger&\diag(n_j)\end{smallmatrix}\right)$.
In its determinant, the coefficients with one or two returns are the
squared $1\times1$ and $2\times2$ minors of $C$. For a $3\times3$ unitary,
the latter equal the squared entries of their complementary rows/columns;
the full determinant has modulus one. This gives \eqref{p4:eq:v2169-b7-d2} on a dense open set,
and polynomial equality extends it everywhere.

For this additional chart the determinant is invariant under exchanging
$\nu_{\rm inv}$ and $\nu_{\rm dir}$, because those columns of $\mathcal P$
coincide. This conclusion concerns this complete block's pole set. Its mixed
current response still retains the inverse/direct distinction; it is not a
general symmetry of the full actual memory. A zero of \eqref{p4:eq:v2169-b7-d2} still needs the
named-tag source test, supplied automatically by its scalar $p_r$.

If the receiver is actually common, $N_\nu=\nu I$, a unitary receiver
presentation gives one three-dimensional block per charged role. Its exact
scalar source response is
\begin{equation}
 G_{rr}(z)=\frac{z(z-\nu)}{
 (z-\nu)(z^2-\varepsilon_r^2)-\lambda^2z}.
 \label{p4:eq:v2169-b7-d3}
\end{equation}
For $\lambda\ne0$ and $\varepsilon_r>0$, the displayed denominator at
$z=\varepsilon_r$ equals $-\lambda^2\varepsilon_r\ne0$. Thus this
closed action chart shifts the terminal pole. When $\nu=0$, the scalar
response cancels the ambient dark zero and its minimal polynomial is
$z^2-\varepsilon_r^2-\lambda^2$. A characteristic determinant alone would
incorrectly count that dark frequency. These comparisons explain why the
actual polynomial evaluation remains load-bearing; they choose no new
physical action, receiver mass or counterterm.



\paragraph{A nonzero interaction can preserve a pole without free-response fidelity}
This exact finite CAR example is a diagnostic, not a proposed framework
Hamiltonian or a VEV-based mass mechanism. Take the already supplied constant
word
\[
 H=\varepsilon n_1+b n_2+g n_1n_2,\quad q=a_1^\dagger,
 \quad\varepsilon>0,\ b>0,
\]
and suppose $\varepsilon+b+g>0$, so the empty state is the unique ground
state. Exact CAR gives $[H,q]=(\varepsilon+g n_2)q$. For $g\ne0$,
\begin{equation}
 p_q(z)=(z-\varepsilon)(z-\varepsilon-g),\quad
 O_{q,\varepsilon}=(1-n_2)q,\quad
 \mathcal A_{q,\varepsilon}=1-n_2.
 \label{p4:eq:v2169-b7-d4}
\end{equation}
The ground state sees the original terminal pole with residue one;
the faithful trace response sees both frequencies with weights $1/2$.
If $g=-\varepsilon$, the second frequency is zero and the same calculation
still applies. Thus a pole-only target can close dynamically even when the
whole free-response certificate fails. The conclusion follows from this
specified action word, not from GP2 alone or a general promise of interaction
cancellation.

% END INLINE SOURCE: Paper4-ChargedPolynomialDiagnostics-Appendix-v1_0.tex




% BEGIN GEOMETRY-FIRST CONSEQUENCES v2.201
\setlength{\emergencystretch}{2em}


\section{Measurement histories, boundary work and retained excitation}
\label{p4b:sec:v238-b7}
\noindent\textit{Source and claim boundary.} Weak measurements, causal quantum histories, Gaussian boundary amplitudes and finite work/occupation tests retain their necessary detector, state, background and kinetic assumptions.

\paragraph{Resolution Change, Decoherence, Control, and Conditional Reeb Interfaces}
\label{p4:part:part3}

\paragraph{Orientation.}
This part collects the resolution-change layer: effective classicality,
entropy inequalities on specified coarse-graining domains, continuous
monitoring, decoherence, inference, feedback, metrology, controllability,
protected codes, and a conditional Reeb phase-transfer interface.  None of
these results reopen the finite state/effect/operator closure proved in the
first part.



\subsection[Why effective classicality and thermodynamic asymmetry can appear]
{Effective classicality and thermodynamic asymmetry}
\label{p4:sec:part3-classicality-entropy}

\paragraph{Purpose.}
The preceding parts derived the closed quantum calculus and organized its representation, finite-QFT, measurement, scattering, and internal-sector presentations.
This resolution-change part asks under which additional effective domains
irreversible-looking behavior can arise even though the resolved microscopic
sector is quantum and finite.  Resolution change, subsystem reduction, and
restricted algebra access can support such a description, but do not by
themselves define an oriented irreversible time law.




\subsection{Weak instruments and the diffusive monitoring limit}
\label{p4:sec:part3-weak-monitoring}

Let $L$ be the measured system operator and let $\eta\in[0,1]$ denote detection efficiency.
A short weak-measurement step of duration $dt$ may be represented by Kraus operators whose statistics depend on a real readout increment $dY_t$.
In the diffusive limit one obtains the standard conditional stochastic master equation on the finite system algebra:
\begin{equation}
\label{p4:eq:sme}
d\rho_t = -i[H,\rho_t]dt + \mathcal D[L](\rho_t)dt + \sqrt{\eta}\,\mathcal H[L](\rho_t)\,dW_t,
\end{equation}
where
\begin{equation}
\label{p4:eq:d-super}
\mathcal D[L](\rho)=L\rho L^* - \frac12\{L^*L,\rho\},
\end{equation}
\begin{equation}
\label{p4:eq:h-super}
\mathcal H[L](\rho)=L\rho+\rho L^* - \Tr\!\bigl((L+L^*)\rho\bigr)\rho,
\end{equation}
and $W_t$ is a Wiener process.
The measured current is
\begin{equation}
\label{p4:eq:meas-current}
dY_t = \sqrt{\eta}\,\Tr\!\bigl((L+L^*)\rho_t\bigr)dt + dW_t.
\end{equation}


\begin{theorem}[Finite weak-instrument diffusive limit]
\label{p4:thm:weak-monitoring-diffusive}
Fix \(H=H^*\), \(L\in\mathcal A_m^{\mathrm{env}}\), and \(\eta\in[0,1]\).
For each short time step \(dt\), define a two-outcome observed instrument by
the CP maps
\[
\mathcal I_{\pm,dt}(\rho)
:=
M_\pm(dt)\rho M_\pm(dt)^*
+
N_\pm(dt)\rho N_\pm(dt)^*,
\]
where
\[
M_\pm(dt)
=
\frac1{\sqrt2}
\Bigl(
\mathbf 1-(iH+\tfrac12L^*L)dt \pm \sqrt{\eta}\,L\,\sqrt{dt}
\Bigr)
+O(dt^{3/2}),
\]
\[
N_\pm(dt)
=
\sqrt{\frac{1-\eta}{2}}\,L\,\sqrt{dt}+O(dt).
\]
Here \(N_\pm\) distributes the unobserved inefficiency channel over each
observed sign rather than treating it as a third recorded outcome.  Choose
the higher-order remainders uniformly on the compact finite-dimensional state
space so that, for all sufficiently small \(dt\),
\[
\sum_{\sigma=\pm}
\bigl(M_\sigma(dt)^*M_\sigma(dt)
+N_\sigma(dt)^*N_\sigma(dt)\bigr)
=\mathbf 1
\]
exactly, without changing the displayed coefficients through order \(dt\).
For a state \(\rho_t\), set
\[
p_\sigma=\Tr\!\bigl(\mathcal I_{\sigma,dt}(\rho_t)\bigr),
\qquad
\rho_{t+dt}^{(\sigma)}
=\frac{\mathcal I_{\sigma,dt}(\rho_t)}{p_\sigma}.
\]
Then the unconditional one-step map satisfies
\begin{equation}
\label{p4:eq:weak-monitor-uncond}
\sum_{\sigma=\pm} p_\sigma\,\rho_{t+dt}^{(\sigma)}
=
\rho_t
+
\Bigl(-i[H,\rho_t]+\mathcal D[L](\rho_t)\Bigr)dt
+
o(dt).
\end{equation}
Let the observed increment take the values \(dY_t=\sigma\sqrt{dt}\).
Then
\[
\mathbb E[dY_t\mid \rho_t]
=
\sqrt{\eta}\,\Tr\!\bigl((L+L^*)\rho_t\bigr)dt+o(dt),
\qquad
dY_t^2=dt+o(dt),
\]
and define the centered innovation
\[
dW_t:=dY_t-\sqrt{\eta}\,\Tr\!\bigl((L+L^*)\rho_t\bigr)dt.
\]
Then the normalized posterior state obeys
\begin{equation}
\label{p4:eq:weak-monitor-step}
\rho_{t+dt}-\rho_t
=
\Bigl(-i[H,\rho_t]+\mathcal D[L](\rho_t)\Bigr)dt
+
\sqrt{\eta}\,\mathcal H[L](\rho_t)\,dW_t
+
o(dt).
\end{equation}
Then the piecewise-constant posterior-state interpolation converges weakly,
as \(dt\downarrow0\), to the unique solution of the stochastic master
equation \eqref{p4:eq:sme}; averaging over records recovers the unconditional
GKLS evolution.
\end{theorem}


\begin{proof}
\ProofPlan{We obtain the diffusive weak-instrument limit by expanding the unconditional channel, expanding the conditioned monitored updates, separating hidden inefficiency from the observed record, and passing to the It\^o innovation equation and unconditional average.}
\begin{proofsteps}

\item \textbf{Exact finite-step normalization.}
Before expanding, collect the four displayed leading Kraus operators as
\(\{K_a(dt)\}_a=\{M_+,M_-,N_+,N_-\}\) and put
\[
R_{dt}:=\sum_a K_a(dt)^*K_a(dt)
=\mathbf1+O(dt^{3/2}).
\]
For sufficiently small \(dt\), \(R_{dt}>0\).  Replacing every \(K_a(dt)\)
by \(K_a(dt)R_{dt}^{-1/2}\) gives exact completeness, while
\(R_{dt}^{-1/2}=\mathbf1+O(dt^{3/2})\) leaves all coefficients through order
\(dt\) unchanged.  Thus the finite-step maps used below are an exact
two-outcome instrument, not merely trace preserving to leading order.

\item \textbf{Unconditional first-order channel expansion.}
Expand the unconditional channel
\[
\mathcal E_{dt}(\rho)
:=\sum_{\sigma=\pm}\mathcal I_{\sigma,dt}(\rho)
\]
to first order in \(dt\).
The \(M_\pm\)-terms contribute
\[
\rho-i[H,\rho]dt-\frac12\{L^*L,\rho\}dt+\eta\,L\rho L^*dt+o(dt),
\]
while the two \(N_\pm\)-terms together contribute
\[
(1-\eta)L\rho L^*dt+o(dt).
\]
Adding both observed CP outcomes gives
\[
\mathcal E_{dt}(\rho)
=
\rho+\Bigl(-i[H,\rho]+L\rho L^*-\tfrac12\{L^*L,\rho\}\Bigr)dt+o(dt),
\]
which is exactly \eqref{p4:eq:weak-monitor-uncond}.


\item \textbf{Conditioned monitored update.}
For the conditioned update, the monitored probabilities satisfy
\[
p_\pm
=
\frac12\Bigl(1\pm \sqrt{\eta}\,\Tr((L+L^*)\rho_t)\sqrt{dt}\Bigr)+O(dt),
\]
and direct expansion of
\(\mathcal I_{\pm,dt}(\rho_t)/p_\pm\) gives
\[
\begin{aligned}
\rho_{t+dt}^{(\pm)}
={}&
\rho_t
+\Bigl(-i[H,\rho_t]+\mathcal D[L](\rho_t)\Bigr)dt\\
&+\Bigl(\eta m_t^2\rho_t
-\eta m_t(L\rho_t+\rho_tL^*)\Bigr)dt\\
&\mathrel{\pm}
\sqrt{\eta}\,\mathcal H[L](\rho_t)\sqrt{dt}
+o(dt).
\end{aligned}
\]
Here \(m_t=\Tr((L+L^*)\rho_t)\).

\item \textbf{Hidden inefficiency and innovation process.}
The \(N_\pm\)-Kraus operators are unobserved within each sign outcome, so
their contribution is marginalized in the conditional state and supplies
the missing \((1-\eta)L\rho L^*dt\) drift without creating a recorded jump.
The sign probabilities give
\(\mathbb E[dY_t\mid\rho_t]=\sqrt{\eta}\,m_tdt+o(dt)\) and
\(dY_t^2=dt+o(dt)\).  Hence
\[
dW_t=dY_t-\sqrt\eta\,m_tdt
\]
has conditional mean zero and quadratic variation \(dt+o(dt)\).  Since
\[
-\eta m_t\mathcal H[L](\rho_t)
=\eta m_t^2\rho_t-\eta m_t(L\rho_t+\rho_tL^*),
\]
the additional deterministic term in the normalized sign update is exactly
the It\^o centering term contained in
\(\sqrt\eta\,\mathcal H[L](\rho_t)dW_t\).

\item \textbf{It\^o limit and unconditional average.}
Substituting this innovation into the expansion above produces
\eqref{p4:eq:weak-monitor-step}.  The state space is compact and the drift
and diffusion coefficients are polynomial, hence globally Lipschitz on that
state space.  The uniform remainder bound, conditional first-moment
convergence, conditional quadratic-variation convergence, and vanishing
higher conditional moments therefore meet the standard martingale
diffusion-approximation criterion.  The interpolated posterior processes
converge weakly to the unique It\^o solution of \eqref{p4:eq:sme}.
Finally, taking conditional expectation removes the martingale term because \(\mathbb E[dW_t\mid \rho_t]=0\), so one recovers the unconditional GKLS equation.
\end{proofsteps}

\end{proof}

Averaging \eqref{p4:eq:sme} over measurement records removes the stochastic term and recovers the unconditional GKLS evolution.
Continuous monitoring is an unraveling of the completely positive reduced dynamics derived in Part~I.



\subsection{Filtering, weak values, and disturbance}
\label{p4:sec:part3-filtering}

Equation \eqref{p4:eq:sme} is the finite-system quantum filter: the conditioned state $\rho_t$ is the optimal state assignment given the measurement record $Y_{[0,t]}$.

\begin{proposition}[Weak-value response and model-defined noise--disturbance quantities]
\label{p4:prop:weak-value-ozawa-domain}
Let a finite probe with initial unit vector \(\xi\) couple to a system
observable \(A=A^*\) through
\[
U_g=\exp(-igA\otimes G),
\]
where \(G=G^*\) is bounded.  For preselection in \(\ket{\psi}\) and
postselection on \(\ket{\phi}\) with
\(\braket{\phi}{\psi}\ne0\), define
\begin{equation}
\label{p4:eq:weak-value}
A_w = \frac{\bra{\phi}A\ket{\psi}}{\braket{\phi}{\psi}}.
\end{equation}
Then the unnormalized postselected probe vector satisfies
\begin{equation}
\label{p4:eq:weak-value-response}
\bra{\phi}U_g\ket{\psi}\xi
=
\braket{\phi}{\psi}
\bigl(\Id-igA_wG\bigr)\xi+O(g^2)
\end{equation}
in norm.  Thus \(A_w\) is the first-order response coefficient of this
specified coupling.  Identifying its real or imaginary part with a particular
pointer shift additionally requires the corresponding centering, commutator,
and covariance conditions on the chosen pointer observable and state.

For the noise--disturbance statement, let
\((\mathcal K,\sigma,U,M)\) be a supplied indirect measurement model:
\(\mathcal K\) is the probe carrier, \(\sigma\) its state, \(U\) the
system--probe interaction, and \(M=M^*\) the probe pointer.  Define
\[
N(A):=U^*(\Id\otimes M)U-A\otimes\Id,
\qquad
D(B):=U^*(B\otimes\Id)U-B\otimes\Id,
\]
\[
\epsilon_\rho(A)^2
:=\Tr\!\bigl((\rho\otimes\sigma)N(A)^2\bigr),
\qquad
\eta_\rho(B)^2
:=\Tr\!\bigl((\rho\otimes\sigma)D(B)^2\bigr).
\]
Then the Ozawa inequality \cite{Ozawa2003} is
\begin{equation}
\label{p4:eq:ozawa}
\epsilon_\rho(A)\eta_\rho(B)
+ \epsilon_\rho(A)\sigma_\rho(B)
+ \sigma_\rho(A)\eta_\rho(B)
\ge \frac12\bigl|\Tr(\rho[A,B])\bigr|,
\end{equation}
where \(\sigma_\rho(A)\) and \(\sigma_\rho(B)\) are the input-state standard
deviations.  The abstract effect or instrument alone does not define these
root-mean-square quantities; the measurement model does.
\end{proposition}

\begin{proof}
The finite Taylor expansion
\[
U_g=\Id-igA\otimes G+O(g^2)
\]
gives \eqref{p4:eq:weak-value-response} after applying
\(\bra{\phi}\) on the system factor and factoring out
\(\braket{\phi}{\psi}\).

For the second claim, the output pointer
\(M_{\rm out}=U^*(\Id\otimes M)U\) commutes with the output observable
\(B_{\rm out}=U^*(B\otimes\Id)U\).  Substituting
\(M_{\rm out}=A_{\rm in}+N(A)\) and
\(B_{\rm out}=B_{\rm in}+D(B)\) into
\([M_{\rm out},B_{\rm out}]=0\), taking expectation in
\(\rho\otimes\sigma\), and applying the triangle and
Cauchy--Schwarz/Robertson bounds to the three error commutators gives
\eqref{p4:eq:ozawa}.
\end{proof}



\subsection{Pointer projectors and decoherence rates}
\label{p4:sec:part3-pointer}

On a pure-dephasing interface domain, suppose the supplied coupling observables $A_{\alpha}=A_{\alpha}^*$ commute, admit common eigenvectors $\ket{i}$, and that the supplied system generator $H_S$ is diagonal in the same pointer basis (equivalently, it commutes with the pointer projectors) with
\begin{equation}
\label{p4:eq:pointer-eigs}
A_{\alpha}\ket{i}=a_{\alpha,i}\ket{i}.
\end{equation}
Then the diagonal projectors
\begin{equation}
\label{p4:eq:pointer-projectors}
P_i = \ket{i}\bra{i}
\end{equation}
form the pointer family identified by the supplied dephasing interaction.  They are not constitutive projectors of the closed carrier unless fixed by an upstream theorem.


\begin{proposition}[Pure dephasing on the common eigenbasis]
\label{p4:prop:pure-dephasing-rates}
Let the reduced generator have the pure-dephasing form \eqref{p4:eq:pure-dephasing}; let the commuting self-adjoint operators \(A_\alpha\) satisfy \eqref{p4:eq:pointer-eigs}; and let \(H_S\) be diagonal in the same common eigenbasis.
Then each matrix unit \(\ket{i}\!\bra{j}\) is an eigenoperator of the generator:
\begin{equation}
\label{p4:eq:dephasing-eigenoperator}
\mathcal L(\ket{i}\!\bra{j})
=
\Bigl(
-i\omega_{ij}
-\frac12\sum_\alpha \gamma_\alpha(a_{\alpha,i}-a_{\alpha,j})^2
\Bigr)\ket{i}\!\bra{j},
\end{equation}
where \(\omega_{ij}\) is the Bohr frequency coming from the Hamiltonian part.
Consequently
\begin{equation}
\label{p4:eq:decoh-rate}
\rho_{ij}(t)=e^{-i\omega_{ij}t}\exp\!\Bigl[-\frac12\sum_{\alpha}\gamma_{\alpha}(a_{\alpha,i}-a_{\alpha,j})^2 t\Bigr]\rho_{ij}(0),
\end{equation}
and the decoherence time between sectors \(i\) and \(j\) is
\begin{equation}
\label{p4:eq:decoh-time}
\tau^{\mathrm{dec}}_{ij}
=
\frac{2}{\sum_{\alpha}\gamma_{\alpha}(a_{\alpha,i}-a_{\alpha,j})^2}
\end{equation}
whenever the denominator is nonzero.
\end{proposition}


\begin{proof}
Because the \(A_\alpha\) are diagonal in the basis \(\{\ket{i}\}\), one has
\[
[A_\alpha,\ket{i}\!\bra{j}]
=
(a_{\alpha,i}-a_{\alpha,j})\ket{i}\!\bra{j}.
\]
Commuting once more gives
\[
[A_\alpha,[A_\alpha,\ket{i}\!\bra{j}]]
=
(a_{\alpha,i}-a_{\alpha,j})^2\ket{i}\!\bra{j}.
\]
If the Hamiltonian is also written in this basis, its commutator acts by
\[
-i[H_S,\ket{i}\!\bra{j}]=-i\omega_{ij}\ket{i}\!\bra{j}.
\]
Substituting these identities into \eqref{p4:eq:pure-dephasing} yields \eqref{p4:eq:dephasing-eigenoperator}.
Since the generator is diagonal on the operator basis \(\{\ket{i}\!\bra{j}\}\), exponentiating it immediately gives \eqref{p4:eq:decoh-rate}.
The decoherence time is the reciprocal of the damping rate, which is exactly \eqref{p4:eq:decoh-time}.
\end{proof}

This is the explicit quantitative version of the classicality mechanism already described conceptually earlier: the environment suppresses off-diagonal coherences in the basis singled out by the couplings it actually reads.



\subsection{Decoherence as effective superselection}
\label{p4:sec:part3-einselection}

If an accessible laboratory algebra is approximately block diagonal with respect to a pointer family $\{P_a\}$, then once
\begin{equation}
\label{p4:eq:decoh-separation}
t \gg \max_{a\neq b}\tau_{ab}^{\mathrm{dec}},
\end{equation}
all operationally relevant expectation values are reproduced by the dephased state
\begin{equation}
\label{p4:eq:dephased-state}
\rho_{\mathrm{ptr}} = \sum_a P_a\rho P_a.
\end{equation}


\begin{corollary}[Operational superselection from pointer dephasing]
\label{p4:cor:pointer-operational}
Let
\[
\mathcal A_{\mathrm{acc}}
:=
\{X\in\mathcal A_m^{\mathrm{env}}:[X,P_a]=0\ \text{for all }a\}
\]
be the block-diagonal accessible algebra determined by the pointer family.
Then
\begin{equation}
\label{p4:eq:pointer-operational-exact}
\Tr(\rho X)=\Tr(\rho_{\mathrm{ptr}}X)
\qquad
\text{for every }X\in\mathcal A_{\mathrm{acc}}.
\end{equation}
More generally, for an arbitrary observable \(X\),
\begin{equation}
\label{p4:eq:pointer-operational-bound}
\bigl|\Tr((\rho-\rho_{\mathrm{ptr}})X)\bigr|
\le
\sum_{a\neq b}\|P_a\rho P_b\|_1\,\|P_b X P_a\|.
\end{equation}
Hence exponential damping of the off-diagonal blocks implies operational convergence to the dephased state for every observable whose off-diagonal part is small in norm.
\end{corollary}


\begin{proof}
\ProofPlan{We prove operational superselection from pointer dephasing by comparing block-diagonal accessible observables with the dephased state, isolating the off-diagonal error, and bounding it by trace duality.}
\begin{proofsteps}

\item \textbf{Accessible block-diagonal expectation.}
Write \(X=\sum_{a,b}P_aXP_b\) and \(\rho=\sum_{a,b}P_a\rho P_b\).
If \(X\in\mathcal A_{\mathrm{acc}}\), then \(P_aXP_b=0\) for \(a\neq b\), so
\[
\Tr(\rho X)
=
\sum_a \Tr(P_a\rho P_a X)
=
\Tr\!\Bigl(\sum_a P_a\rho P_a\,X\Bigr)
=
\Tr(\rho_{\mathrm{ptr}}X),
\]
which proves \eqref{p4:eq:pointer-operational-exact}.

\item \textbf{Off-diagonal error contribution.}
For general \(X\), only the off-diagonal blocks of \(\rho-\rho_{\mathrm{ptr}}\) contribute, hence
\[
\Tr((\rho-\rho_{\mathrm{ptr}})X)
=
\sum_{a\neq b}\Tr(P_a\rho P_b\,P_bXP_a).
\]

\item \textbf{Trace-duality bound and operational convergence.}
Applying trace duality,
\(
|\Tr(YZ)|\le \|Y\|_1\|Z\|,
\)
term by term gives \eqref{p4:eq:pointer-operational-bound}.
Therefore once the off-diagonal blocks are suppressed by decoherence, every block-diagonal accessible observable is reproduced exactly by \(\rho_{\mathrm{ptr}}\), and every nearly block-diagonal observable is reproduced up to the same small error.
\end{proofsteps}

\end{proof}

Classical alternatives therefore arise as environment-stabilized sectors of the same quantum state space.
In the framework, this is not a separate postulate of collapse; it is an effective superselection structure generated by open-system decoherence and restricted algebra access.



\subsection{Gaussian physical work and causal source jets}
\label{p4:sec:v2203-gaussian-work}
The functional source calculus of
Theorem~\ref{p4:thm:v4106-functional-hilbert-response} has an exact
finite-mode Gaussian specialization.  It consumes an already licensed,
constraint-reduced common action, its complete physical source family and
the supplied preparation.  A finite number of oscillator modes still has
an infinite-dimensional Fock space.  The statements below use a common
invariant quadratic-operator domain, differentiable source histories and
finite Gaussian moments on the stated time interval; they do not borrow
finite-Hilbert-space operator-norm bounds.  Centres, metric-dependent
frames, projectors, reductions and endpoint terms retain their existing
owners.  In particular, a perturbation propagation generator does not by
itself specify an additional physical source-energy allocation.

Let \(z\) be a vector of \(2n\) real canonical operators with
\([z_i,z_j]=i\hbar J_{ij}\), where \(J\) is the standard symplectic
matrix.  Put \(m=\langle z\rangle\) and
\(\Gamma_{ij}=\langle\{z_i-m_i,z_j-m_j\}\rangle/2\).
All quadratic products below are Weyl ordered, and their coefficient
matrices are real symmetric.

\begin{theorem}[Physical Gaussian work on a regular canonical atlas]
\label{p4:thm:v2203-gaussian-physical-work}
On a closed admitted quadratic interface with generator
\(H=z^TGz/2+cI\), write \(A=JG\) and let \(S(t,s)\) be its symplectic
propagator.  Then
\[
 \dot\Gamma=A\Gamma+\Gamma A^T,\qquad
 \Gamma(t)=S(t,t_i)\Gamma_iS(t,t_i)^T,
 \qquad \det\Gamma(t)=\det\Gamma_i.
\]
For an already derived physical insertion
\(\mathcal O_a=z^TB_az/2+l_a^Tz+c_aI\),
\[
 \boxed{\langle\mathcal O_a\rangle
 =\tfrac12\operatorname{tr}(B_a\Gamma)
  +\tfrac12m^TB_am+l_a^Tm+c_a.}
\]
Under a homogeneous symplectic chart change \(z'=Tz\),
\[
 \Gamma'=T\Gamma T^T,\quad m'=Tm,\quad
 B_a'=T^{-T}B_aT^{-1},\quad l_a'=T^{-T}l_a,
\]
and the physical expectation is unchanged.  If the supplied initial
Gaussian state is pure on the full \(n\)-mode interface, then
\(\det\Gamma=(\hbar^2/4)^n\).  For one independent canonical pair
this is \(\hbar^2/4\); the same equality is not asserted for a marginal
pair that becomes entangled with other retained modes.
\end{theorem}
\begin{proof}
The linear Heisenberg equation is \(\dot z=Az\).  Its centred second
moment gives the covariance equation.  Symplecticity gives
\(\det S=1\), while a pure Gaussian covariance has all symplectic
eigenvalues \(\hbar/2\).  Weyl ordering gives the expectation formula,
and trace cyclicity proves its chart invariance.
\end{proof}
Consequently a physical work insertion whose transformed coefficients and
centre remain bounded on a regular chart has finite Gaussian expectation
there.  Covariance purity is an uncertainty statement, not an Euler
charge, particle occupation, action determinant or energy normalization.
Homogeneous source density and pressure still consume the actual
functional lapse and radius insertions; the contraction does not
construct those insertions or a new source owner.

\paragraph{Functional source contacts and time-dependent charts}
\begin{proposition}[Gaussian functional work insertion]
\label{p4:prop:v2203-gaussian-functional-work}
Let the complete quadratic generator family in a fixed independent
canonical frame be
\(H(t,x,x^{(1)},\ldots,x^{(s)})=z^TGz/2+cI\).
Along the supplied history define
\[
 \mathcal D_tB=\dot B+A^TB+BA,\qquad A=JG.
\]
The complete on-equation action insertion
\(\mathcal O_a=\sum_r(-\mathfrak d_t)^rH_{,x^{a(r)}}\) has coefficients
\[
 \boxed{B_a=\sum_{r=0}^s(-\mathcal D_t)^rG_{,x^{a(r)}},\qquad
 c_a=\sum_{r=0}^s(-\partial_t)^rc_{,x^{a(r)}}.}
\]
The inherited endpoint operator is retained.  Full force variation uses
\(\delta\mathcal O_a[v]+i[K[v],\mathcal O_a]/\hbar\), together with
its actual initial/preparation tangent, as in
Theorem~\ref{p4:thm:v4106-functional-hilbert-response}.
\end{proposition}
\begin{proof}
The Heisenberg derivative of \(z^TBz/2\) is the quadratic with
coefficient \(\mathcal D_tB\).  Quadratic commutators close exactly.
Repeated integration by parts in the complete source variation gives
the formula and the existing endpoint operator.  A central coefficient
has only its ordinary total derivative.
\end{proof}
All parameter derivatives are taken in the actual common-action family.
If a moving frame, constraint elimination or metric-dependent projector
has already changed that family, its derivative and endpoint contacts
must be included before this proposition is applied.  A frozen derivative
of a propagation coefficient does not replace a functional Hilbert
variation.

\begin{proposition}[Canonical connection versus physical work]
\label{p4:prop:v2203-time-dependent-chart-work}
For a differentiable time-dependent symplectic chart \(z'=T(t)z\),
the quadratic evolution matrices obey
\[
 A'=\dot TT^{-1}+TAT^{-1},\qquad
 G'=T^{-T}GT^{-1}-J\dot TT^{-1}.
\]
The final term is a canonical connection of the time generator.  An
already derived physical observable instead transforms by the
congruence of Theorem~\ref{p4:thm:v2203-gaussian-physical-work}.
For the one-pair shear \(\pi=p-Kq\), the action identity is
\[
 p\dot q-H(q,p)
 =\pi\dot q-H'(q,\pi)+\frac{d}{dt}\frac{Kq^2}{2},
 \quad H'=H(q,\pi+Kq)+\frac{\dot Kq^2}{2}.
\]
Thus the continued time generator is not automatically the physical
Hilbert energy.  Source-dependent charts retain their generating
endpoint and preparation-frame variations.
\end{proposition}
\begin{proof}
Differentiate \(z'=Tz\) and use symplecticity to convert \(A'\) to
\(JG'\).  The shear identity follows by differentiating \(Kq^2/2\).
The common-action variation includes that total derivative and its
source jets, whereas a separately varied quadratic coefficient need not.
\end{proof}

\paragraph{Mean, causal response and connected noise}
For the same physical preparation and complete source word, use the
already licensed two-history source functional
\[
 \mathcal I[x_+,x_-]
 =\operatorname{Tr}(U[x_+]\rho_iU[x_-]^\dagger),\qquad
 \Gamma_s=-i\hbar\log\mathcal I,
 \quad x_c=(x_++x_-)/2,\quad x_d=x_+-x_-.
\]
The logarithm is continued near equal histories.  At \(x_d=0\), the
complete mean is \(-\delta\Gamma_s/\delta x_d^a=F_a\), with
\(F_a=\langle\mathcal O_a\rangle\).  The identity
\(\Gamma_s[x,x]=0\) does not set this difference-direction derivative
to zero.  For \(H=NH_0(a)\), with \(V=V_\gamma a^3\), it gives
\[
 -\left.\frac{\delta\Gamma_s}{\delta N_d}\right|_{x_d=0}=V\rho_s,
 \qquad
 \left.\frac{\delta\Gamma_s}{\delta(\log a)_d}\right|_{x_d=0}=3NVp_s.
\]
Derivative-dependent words use the complete functional insertions and
endpoints rather than these instantaneous derivatives.

\begin{theorem}[The first two causal source jets]
\label{p4:thm:v2203-causal-mean-response-noise}
Fix the physical preparation, and use compactly supported history
variations, or retain their actual endpoint operators.  Let
\(\chi^R_{ab}\) be the complete direct contact plus retarded work
response of \(F_a\), and put
\[
 \Pi^R_{ab}(t,s)=-\frac{i}{\hbar}\Theta(t-s)
  \langle[\mathcal O_a^H(t),\mathcal O_b^H(s)]\rangle,
 \quad
 N_{ab}(t,s)=\tfrac12
  \langle\{\Delta\mathcal O_a^H(t),\Delta\mathcal O_b^H(s)\}\rangle.
\]
For common-history perturbation \(v_c\) and difference history \(d\),
the first two jets are
\[
 \boxed{\begin{aligned}
 \Gamma_s={}&-\int d^aF_a
 -\iint d^a(t)\chi^R_{ab}(t,s)v_c^b(s)\,dt\,ds\\
 &+\frac{i}{2\hbar}\iint d^a(t)N_{ab}(t,s)d^b(s)\,dt\,ds+O(3).
 \end{aligned}}
\]
The real smeared noise form is nonnegative.  These are mean, response
and fluctuation jets of one functional, not separately additive means.
\end{theorem}
\begin{proof}
Unitary first variation gives the difference-direction mean with the
displayed sign.  Its common-history derivative is the complete contact
and causal commutator response.  The second connected cumulant in the
difference history gives the imaginary term.  For real \(d\), its
noise form is the variance of
\(\int d^a(t)\Delta\mathcal O_a^H(t)\,dt\), hence is nonnegative.
\end{proof}
A central energy contributes to the mean and its direct derivatives,
but cancels from connected variance and commutator memory.  Independent
preparation, endpoint and source-state variations retain their existing
additional jets.

\begin{proposition}[Exact Gaussian quadratic response and noise]
\label{p4:prop:v2203-gaussian-connected-noise}
On a centred Gaussian quadratic interface, let
\(\mathcal O_a=z^TB_az/2+c_aI\).  Define the ordered covariance
\[
 W(t,s)=\langle z(t)z(s)^T\rangle
 =S(t,s)[\Gamma(s)+i\hbar J/2].
\]
The connected ordered correlator is
\[
 \boxed{C_{ab}(t,s)
 =\tfrac12\operatorname{tr}
 [B_a(t)W(t,s)B_b(s)W(t,s)^T].}
\]
Consequently \(N_{ab}=\operatorname{Re}C_{ab}\) and
\(\Pi^R_{ab}=2\Theta(t-s)\operatorname{Im}C_{ab}/\hbar\).
Equivalently, with
\(\widetilde B_a=S(t,s)^TB_a(t)S(t,s)\),
\[
 \Pi^R_{ab}=\frac{\Theta(t-s)}2\operatorname{tr}
 [(\widetilde B_aJB_b-B_bJ\widetilde B_a)\Gamma(s)].
\]
This is the sign for a perturbation \(\delta H=+v^b\mathcal O_b\);
the memory sign reverses for its negative.
At equal times,
\[
 \boxed{\operatorname{Var}(\mathcal O_B)
 =\tfrac12\operatorname{tr}(B\Gamma B\Gamma)
  +\frac{\hbar^2}{8}\operatorname{tr}(BJBJ).}
\]
\end{proposition}
\begin{proof}
Wick contraction of the two Weyl quadratics gives two equal ordered
cross-contractions.  The final factor is an ordinary transpose, not a
complex conjugate: both cross-contractions use \(W(t,s)\).
The reversed ordered expectation is the complex conjugate, giving the
noise and commutator formulas.  Equivalently,
\([\mathcal O_B,\mathcal O_C]\) is \(i\hbar\) times the quadratic
with matrix \(BJC-CJB\).  Substitution of
\(W(t,t)=\Gamma+i\hbar J/2\) gives the variance formula.
\end{proof}
For a stable oscillator ground, the quantum \(J\) term cancels the
covariance-only term when \(B\) is its Hamiltonian matrix: its energy
variance is zero.  For \(q^2/2\), the variance is instead
\(\Gamma_{qq}^2/2>0\).  On its admitted quadratic root Hamiltonian,
the fresh-root state of
Theorem~\ref{p8:thm:v2114-fresh-root-neutral-qroot} therefore has zero
covered vacuum-subtracted excitation mean and Hamiltonian-energy noise
at birth, while field quadratics may fluctuate and respond.  This does
not set every metric-work variance to zero, erase the REN0 vacuum
baseline, or supply a state for scalar directions excluded from that
root interface.

\begin{corollary}[Linear forcing has no first-order even-work response]
\label{p4:cor:v2203-gaussian-linear-even-parity}
Let unperturbed quadratic evolution transport a zero-mean Gaussian.
Suppose a source enters only as \(\delta H=-y(t)l(t)^Tz\), and the
target work insertion is even quadratic.  If there is no explicit
first-order quadratic or central work contact and no independently
changed initial covariance, then
\[
 D_y\langle\mathcal O_{\rm even}\rangle\big|_{y=0}=0.
\]
\end{corollary}
\begin{proof}
A linear force changes the mean at first order and leaves the centred
covariance unchanged.  The even expectation changes through
\(m^TBm/2\), which begins at second order.  Equivalently, its
commutator with a linear source is linear and has zero unperturbed mean.
Purely linear direct contacts also have zero expectation.
\end{proof}
An actual quadratic source--mode--mode vertex, nonzero coherent mean,
or independently owned state tangent lies outside these hypotheses.
The corollary restricts a linear route; it does not prohibit nonlinear
source creation.

\paragraph{Quantum Schur contraction and one-time allocation}
\begin{proposition}[Gaussian covariance and determinant count the same work]
\label{p4:prop:v2203-quantum-schur-nonduplication}
On an actually licensed fixed-rank positive Euclidean Gaussian word,
let
\[
 S=S_0+\tfrac12y^TCy+j^Ty,\quad C>0,\qquad
 W_E=-\hbar\log\left(\mu\int dy\,e^{-S/\hbar}\right).
\]
Its mean and covariance are \(m=-C^{-1}j\), \(\Sigma=\hbar C^{-1}\),
and its complete variation is
\[
 \boxed{\delta W_E=\delta S_0+\delta j^Tm
 +\tfrac12m^T\delta C m
 +\tfrac12\operatorname{tr}(\delta C\,\Sigma)
 -\hbar\delta\log\mu.}
\]
The covariance contraction is the derivative of
\(\hbar\log\det C/2\); it is not an additional copy of that term.
\end{proposition}
\begin{proof}
Completing the square gives
\(W_E=S_0-j^TC^{-1}j/2+\hbar\log\det C/2-
\hbar\log\mu\), up to a fixed constant.  Its derivative is exactly
the expectation of the parent variation minus the measure jet.
\end{proof}
A classical stationary Schur reduction has only the mean-field part;
it does not manufacture the determinant contribution.  At \(j=0\)
its induced classical first variation vanishes, while an admitted
quadratic quantum contraction can remain nonzero.  This Euclidean
identity does not replace the actual propagated Lorentzian covariance
by \(\hbar C^{-1}\): contour, preparation, measure, endpoint and
constituent matching remain necessary.  Any determinant already
allocated to gravity or a retained source is counted once.  More
generally, the complete stationary envelope reproduces the parent mean
on the eliminated solution, rather than adding a second source mean.

\paragraph{A conditional Gaussian cold tail without finite Fock bounds}
\begin{theorem}[Finite-mode Gaussian cold survival]
\label{p4:thm:v2203-gaussian-cold-tail}
Let an actually owned instantaneous proper-time source have
\(H_0(a)=z^TG(a)z/2+c(a)I\) in a fixed canonical frame, with finite
initial centred Gaussian covariance, and let this same complete
generator evolve the state.  Assume no omitted exchange interaction,
\(a(t)\to\infty\), and
\[
 G(t)\to G_\infty>0,\qquad
 \int_{t_i}^\infty\|G(t)-G_\infty\|\,dt<\infty,
\]
\[
 c(t)\to c_\infty,\qquad
 \|a\partial_aG\|\to0,\qquad a\partial_ac\to0.
\]
Then the interaction-picture covariance has a finite limit and
\[
 \boxed{V\rho\to E_\infty
 =\tfrac12\operatorname{tr}(G_\infty\Gamma_{I,\infty})+c_\infty,
 \qquad Vp\to0.}
\]
If the physically allocated \(E_\infty>0\) and \(V=V_\gamma a^3\),
then \(a^3\rho\to E_\infty/V_\gamma\) and \(p/\rho\to0\).
\end{theorem}
\begin{proof}
Put \(S_\infty=e^{JG_\infty(t-t_i)}\).  Positivity of
\(G_\infty\) makes this flow conjugate to an orthogonal flow, so
it and its inverse are uniformly bounded.  For the actual symplectic
flow \(S\), \(R=S_\infty^{-1}S\) obeys
\[
 \dot R=S_\infty^{-1}J(G-G_\infty)S_\infty R.
\]
Its coefficient is integrable.  Gronwall bounds \(R\) and its inverse;
their derivatives are integrable, so they converge to an invertible
symplectic \(R_\infty\).  Thus
\(\Gamma_{I,\infty}=R_\infty\Gamma_iR_\infty^T\) exists.
Since \(S_\infty^TG_\infty S_\infty=G_\infty\), bounded physical
covariance and \(G\to G_\infty\) give the energy limit.  The complete
instantaneous source work is
\[
 3Vp=-\tfrac12\operatorname{tr}(aG_{,a}\Gamma)-ac_{,a},
\]
whose limit is zero by the hypotheses.
\end{proof}
The pointwise limit is explicit: integrability alone permits narrowing
late spikes.  Nonzero means add their transported quadratic energy and
obey the same finite-matrix argument.  The physical covariance itself
may continue rotating; only its interaction-picture covariance must
converge.  No bounded-Fock-operator claim or WKB approximation is used.
The centre, preparation and vacuum subtraction retain their inherited
allocations; a late instantaneous reset would change the physical word.
A local FLRW interpretation additionally consumes the actual
homogeneous, isotropic and no-flux source certificate.  The theorem is
not automatically a result for a constrained scalaron propagation
generator, nor does a finite-time source-free replay certify its
positive asymptotic energy matrix, source allocation or history limits.


% BEGIN FINITE SOURCE AND CONTROLLED RETENTION v2.209
% Framework v2.209: finite source loading, compact response and controlled retention.


\subsection{Finite source work and controlled retained excitation}
\label{p4:sec:v2209-finite-retained-source}
These results consume the actual complete finite physical carrier, its
admitted Hamiltonian history, inherited density and source-energy allocation.
They neither select a state nor add a bosonic constituent to a finite source.
Production, subsequent capture, retention and cold stress are separate tests.
An interaction-picture time generator retains its connection; physical energy
and metric work retain the functional insertions and centres already owned by
the action. The Gaussian interface of
Theorem~\ref{p4:thm:v2203-gaussian-cold-tail} remains a distinct domain.



\paragraph{Passive work and its finite capacity}
\begin{theorem}[Same-endpoint passive work and unitary-orbit capacity]
\label{p4:thm:v2209-passive-work-capacity}
Let the fixed complete endpoint Hamiltonian have ordered eigenvalues
\(\epsilon_1\le\cdots\le\epsilon_N\). Let its initial passive density be
\(\rho_i=\sum_a w_a|a\rangle\langle a|\), where
\(w_1\ge\cdots\ge w_N\ge0\) and \(\sum_a w_a=1\).
For any physical closed unitary history with final state
\(\rho_f=U\rho_iU^\dagger\), set \(w'_a=\langle a|\rho_f|a\rangle\).
Comparison using this same endpoint Hamiltonian gives
\begin{equation}
\label{p4:eq:v2209-passive-work}
 \Delta E=\sum_{k=1}^{N-1}(\epsilon_{k+1}-\epsilon_k)
 \left(\sum_{a=1}^kw_a-\sum_{a=1}^kw'_a\right)\ge0.
\end{equation}
It is strictly positive iff some strictly positive gap has a strictly
depleted cumulative low-energy population. On the full unitary orbit,
\begin{equation}
\label{p4:eq:v2209-passive-capacity}
 \Delta E_{\max}=\sum_{a=1}^N\epsilon_a(w_{N+1-a}-w_a).
\end{equation}
If physical conservation or superselection restricts the available
unitaries, this is an upper bound; attaining it requires that the reverse
permutation be physically reachable on the actual source domain.
\end{theorem}
\begin{proof}
The matrix \(B_{ab}=|U_{ab}|^2\) is doubly stochastic and \(w'=Bw\).
Consequently \(\sum_{a\le k}w'_a\le\sum_{a\le k}w_a\) for each \(k\).
Summation by parts, using equal total populations, gives
\eqref{p4:eq:v2209-passive-work} and its strictness criterion.
The doubly stochastic convex decomposition into permutations and the
rearrangement inequality bound the final energy by
\(\sum_a\epsilon_aw_{N+1-a}\). A reversing permutation is itself unitary,
so it attains the bound on the unrestricted orbit.
\end{proof}
The endpoint comparison can be used for a cyclic drive. For unequal endpoint
Hamiltonians it is a diagnostic against the fixed reference, rather than the
whole physical work; endpoint work, reference and binding contributions must
be kept. This theorem proves positivity when its population test is passed;
it does not imply that every changing source history absorbs energy.



\paragraph{A sharp restriction for a monotone free sign pair}
\begin{theorem}[Mass-independent monotone free-pair loading]
\label{p4:thm:v2217-monotone-free-loading}
On a complete free sign-pair block let
$h(t)=(\lambda/a(t))\sigma_1+m\sigma_3$, with fixed $m>0$, $\lambda$,
locally absolutely continuous nondecreasing $a>0$, and initial instantaneous
lower eigenstate. If $N=\log(a_f/a_i)\ge0$, its final instantaneous upper
probability obeys
\[
 P_f\le\frac{1-\operatorname{sech}(N/2)}2\le\frac{N^2}{16}.
\]
This is uniform in $m,\lambda$ and the diagonal dynamical phases.
For an already loaded pure pair with probability $P_*$, any subsequent
monotone massive tail with $a\to\infty$ satisfies
\[
 P_\infty\ge\sin^2\!\left(\max\{0,
       \arcsin\sqrt{P_*}-\tfrac12\theta_*\}\right),\qquad
 \theta_* =\arctan(|\lambda|/(ma_*)).
\]
\end{theorem}
\begin{proof}
Diagonalize with $\theta=\arctan(|\lambda|/(ma))$. The instantaneous-frame
off-diagonal connection has magnitude $\hbar|\dot\theta|/2$.
The two-component Schr\"odinger equation gives
$|\dot P|\le|\dot\theta|\sqrt{P(1-P)}$, hence the change of
$\arcsin\sqrt P$ is at most $\frac12\int|\dot\theta|dt$.
Starting at zero, this bounds $P_f$ by
$\sin^2[(\theta_i-\theta_f)/2]$. For $x=|\lambda|/(ma_i)$,
the angle difference is $\arctan x-\arctan(xe^{-N})$.
Its derivative in $x$ vanishes at $x=e^{N/2}$ (for $N>0$), is positive
before and negative after; the cosine of its maximum is
$\operatorname{sech}(N/2)$. Also
$1-\operatorname{sech}y=\int_0^y\operatorname{sech}v\tanh v\,dv
\le y^2/2$, proving the second bound. The reverse angle inequality from
$P_*$, with remaining variation $\theta_*$, proves the retention bound.
The same projector-variation estimate makes $P_\infty$ exist. A sudden
monotone change approaches the optimized upper bound. Zero momentum or a
massless constant-sign block has no rotating eigenspace and no such loading.
\end{proof}
This bound concerns the stated free drive and preparation. A changing mass,
interaction, inherited earlier population or nonmonotone history requires
its own complete source. Finite-band energy also retains its actual
degeneracies, gaps and central allocation.



\paragraph{An oscillator tangent of an existing sign pair}
\begin{proposition}[Compact Bloch source and saturating force]
\label{p4:prop:v2209-compact-bloch-source}
On an actually owned two-level source \(H_*=E\tau_3\), \(E>0\), let a
passive initial state have Bloch length \(0<r\le1\). In units \(\hbar=1\),
its unitary orbit admits coordinates
\[
 0\le J\le r,\qquad r_z=-r+2J,\qquad
 r_x=2\sqrt{J(r-J)}\cos\psi,\qquad E_{\rm exc}=2EJ.
\]
Near the passive pole the canonical chart
\(q=\sqrt{2J}\cos\psi\), \(p=-\sqrt{2J}\sin\psi\) obeys
\(q^2+p^2<2r\). The outer circle is collapsed to the opposite pole in the
global spherical orbit; it is not an additional continuum of source states.
An actual transverse insertion \(d\tau_1\) has expectation
\[
 d\sqrt2\,q\sqrt{r-(q^2+p^2)/2}.
\]
Its small-amplitude response has frequency \(\Omega=2E\) and, after the
unit-mass oscillator normalization, squared coupling \(g^2=4rEd^2\).
The exact source capacity is \(0\le E_{\rm exc}\le2rE\).
\end{proposition}
\begin{proof}
The Bloch sphere of radius \(r\) gives
\(r_x^2+r_y^2=r^2-r_z^2=4J(r-J)\). Its spin Poisson structure makes
\((\psi,J)\) canonical in the displayed units, hence \(\{q,p\}=1\).
Subtracting the passive-pole energy \(-rE\) gives \(2EJ\).
Substitution gives the saturating force. The tangent Hamiltonian is
\(E(q^2+p^2)\); setting \(x=q/\sqrt\Omega\),
\(\pi=\sqrt\Omega p\) gives the stated frequency and coupling.
\end{proof}
For \(r=0\) the orbit is a point and this response tangent is absent.
The local oscillator is a response presentation of the existing source,
with its inherited clock and energy. It is not an unlimited extra CCR
species. Pair or Bogoliubov descriptions require the already admitted
finite source or regular owned positive-frequency interface.

\begin{remark}[Kernel matching does not duplicate source energy]
\label{p4:rem:v2209-noise-energy-witness}
For \(\rho=(I-r\tau_3)/2\) and \(F=d\tau_1\), the exact symmetrized
connected noise is \(N(t)=d^2\cos(2Et)\). A vacuum auxiliary oscillator
matching the retarded pair kernel instead has noise \(rd^2\cos(2Et)\).
For \(0<r<1\), matching both requires
\(n_{\rm eff}=(1-r)/(2r)\). Its auxiliary excitation energy would be
\(E(1-r)/r\), while the actual pair energy above its pure ground is
\(E(1-r)\). Moreover its exact equal-time fourth cumulant is
\(\langle F^4\rangle-3\langle F^2\rangle^2=-2d^4\), whereas the Gaussian
auxiliary value is zero. These follow from \(\tau_1^2=I\) and two-level
Heisenberg rotation. Matching a retarded susceptibility alone therefore
does not establish a state, higher correlators or an additional physical mean.
\end{remark}



\paragraph{A quantitative retention certificate}
\begin{theorem}[Controlled leakage of a physical excitation sector]
\label{p4:thm:v2209-controlled-excitation-survival}
On the complete finite carrier after a loading time \(t_\ell\), let
\(P_C\) be a fixed physical excitation projector and let
\(H(t)=H_0(t)+H_{\rm mix}(t)\), with \([H_0(t),P_C]=0\).
Use locally integrable self-adjoint generators and the actual inherited
density. Set \(p_C(t)=\operatorname{Tr}(\rho(t)P_C)\) and
\[
 I(t)=\hbar^{-1}\int_{t_\ell}^t\|H_{\rm mix}(s)\|\,ds.
\]
Then
\begin{equation}
\label{p4:eq:v2209-survival-bound}
 p_C(t)\ge\left[\max\{0,\sqrt{p_C(t_\ell)}-I(t)\}\right]^2.
\end{equation}
If \(I_{\rm tail}=I(\infty)<\infty\), the interaction-picture propagator
converges in norm and \(p_C(t)\) has a limit satisfying the same bound.
The physical density need not converge. For an actually allocated excitation
observable \(H_{\rm exc}\succeq\Delta_CP_C\), \(\Delta_C>0\),
\[
 \liminf_{t\to\infty}\langle H_{\rm exc}\rangle_t
 \ge\Delta_C[\max\{0,\sqrt{p_C(t_\ell)}-I_{\rm tail}\}]^2.
\]
Writing a limiting energy additionally requires its existence. The same
conclusion applies to physical insertions approaching \(H_{\rm exc}\)
in operator norm. A moving projector must first be transported to a fixed
physical frame and its transport connection included in the generator.
\end{theorem}
\begin{proof}
Purify the initial density and let \(U,U_0\) act on its physical factor.
Duhamel's identity and unitarity give \(\|U-U_0\|\le I(t)\).
Since \(U_0\) preserves \(P_C\), the reverse triangle inequality for
the projected purification gives
\(\sqrt{p_C(t)}\ge\sqrt{p_C(t_\ell)}-I(t)\).
Nonnegativity, squaring, monotonicity of \(I\), and the operator gap
inequality prove the bounds. In the \(U_0\) interaction picture the
perturbation has integrable norm, making its unitary propagator Cauchy.
Since \(U_0\) commutes with \(P_C\), this proves the population limit;
it does not assert convergence of other rotating observables.
Norm convergence of an insertion changes
its normalized expectation by a vanishing amount.
\end{proof}
This is retention conditional on computed loading. It does not establish
the initial population, a physical positive gap, the support surviving a
terminal restriction, or its finite-observation pressure. Within-block
terms may be moved into \(H_0\) to sharpen the leakage bound; counting the
entire perturbation in \(H_{\rm mix}\) remains a valid larger estimate.

\begin{corollary}[The existing finite-source late tail]
\label{p4:cor:v2209-native-tail-retention}
On the complete fixed-band source of
Theorem~\ref{p3f:thm:v2203-local-source-cold-tail}, take
\(H=H_m+a^{-1}\widehat D+c(a)I\) and a physical \(P_C\) commuting
with \(H_m\). Choosing \(H_0=H_m+c(a)I\) gives
\[
 I_{\rm tail}\le\frac{\|\widehat D\|}{\hbar}
 \int_{t_\ell}^{\infty}\frac{dt}{a(t)}
 \le\frac{\|\widehat D\|}{\hbar H_{\min}a_\ell}
\]
on a tail satisfying \(a(t)\ge a_\ell e^{H_{\min}(t-t_\ell)}\),
\(H_{\min}>0\). Thus a loading amplitude
\(\sqrt{p_C(t_\ell)}>I_{\rm tail}\), together with the physical gap
inequality, gives a strictly positive retained gap-energy lower bound.
Under the centre and physical work hypotheses of that imported theorem,
\(V\rho\to E_\infty\) and \(Vp\to0\); cold survival additionally needs
the physically allocated \(E_\infty>0\). The finite-mode Gaussian analogue
uses \ref{p4:thm:v2203-gaussian-cold-tail}, rather than finite-Fock norms.
\end{corollary}
\begin{proof}
The displayed norm estimate follows directly from the perturbation.
Its integrability makes the \(H_m\) interaction propagator converge
in operator norm, as in the imported finite-source theorem; the energy
limit and scale-work limit require its stated centre conditions.
Apply \eqref{p4:eq:v2209-survival-bound} to this same complete source.
\end{proof}

\begin{proposition}[Inherited mixedness after genuine retained allocation]
\label{p4:prop:v2209-inherited-mixedness-gap}
For a complete actually retained source with unique final ground projector
\(P_0\), let its physically allocated excitation operator satisfy
\(H_{\rm exc}\succeq\Delta(I-P_0)\), \(\Delta>0\). Then
\[
 \operatorname{Tr}(\rho H_{\rm exc})
 \ge\Delta[1-\operatorname{Tr}(\rho P_0)]
 \ge\Delta[1-\lambda_{\max}(\rho)].
\]
\end{proposition}
\begin{proof}
Take the expectation of the operator gap inequality and bound the
rank-one ground probability by the largest eigenvalue of \(\rho\).
\end{proof}
The density is the actual inherited retained state, including its restriction
history. It is not reset to a fresh neutral root during descent. A mixed
marginal of an interacting ground does not establish additional cosmic
energy: omitted interaction, binding or reference terms can compensate
the bare marginal energy. Local FLRW stress still requires the complete
homogeneous, isotropic, no-flux physical incidence and its work insertions.

% END FINITE SOURCE AND CONTROLLED RETENTION v2.209

% BEGIN LITERATURE-CHECKED QUANTUM CERTIFICATES v2.210


\subsection{Exact massive conformal--shape source response}
\label{p4:sec:v2212-conformal-shape-source}

% BEGIN EXACT MASSIVE CONFORMAL--SHAPE SOURCE KERNEL v2.212
\begin{corollary}[Massive conformal--shape causal source kernel]
\label{p4:cor:v2212-massive-conformal-shape-kernel}
Use an actually admitted invariant free Dirac mass member of
Proposition~\ref{p4:prop:spherical-free-dirac-source}, with constant
Jordan mass $m$, the inherited canonical spinor pairing and a regular
$F>0$, $f_{RR}>0$ scalaron chart on the existing certified scalar-action
and common-action overlap domain (Paper~III-F,
Proposition~\ref{p3f:prop:v1164-scalar-overlap-condition}).
Set $g_E=Fg_J$ and $q=\log F$.
Use Einstein proper time and hold the Einstein proper cell volume $V_E$
fixed in the shape variation $\alpha=\log(A_E/B_E)$:
$A_E\mapsto A_Ee^{2\alpha/3}$,
$B_E\mapsto B_Ee^{-\alpha/3}$.
In units $c_{\rm op}=1$, the noncentral one-body member is
\begin{equation}
\label{p4:eq:v2212-conformal-shape-hamiltonian}
 h_E=u\Gamma_z+v\Gamma_\Omega+\mu\beta,\qquad
 u=\hbar k/A_E,\quad v=\hbar\nu/B_E,\quad
 \mu=me^{-q/2},\quad E^2=u^2+v^2+\mu^2.
\end{equation}
Its actual vertices and bare mixed contact are
\begin{equation}
\label{p4:eq:v2212-conformal-shape-vertices}
 \boxed{O_q=h_{E,q}=-\frac\mu2\beta,\qquad
 O_\alpha=h_{E,\alpha}=-\frac{2u}{3}\Gamma_z
                  +\frac v3\Gamma_\Omega,\qquad
 h_{E,q\alpha}=0.}
\end{equation}
For the complete retained many-body energy
$\widehat H_E=d\Gamma(h_E)+\varepsilon_{0,E}I$, define
$\widehat O_a=d\Gamma(O_a)+\varepsilon_{0,E,a}I$.
The physical averaged Hilbert ports are
\begin{equation}
\label{p4:eq:v2212-conformal-shape-ports}
 \boxed{T_E=\frac2{V_E}\operatorname{Tr}(\sigma\widehat O_q),
 \qquad \Delta_E:=p_{z,E}-p_{\perp,E}
 =-\frac3{2V_E}\operatorname{Tr}(\sigma\widehat O_\alpha).}
\end{equation}
Thus the complete mixed contact retains
$\varepsilon_{0,E,q\alpha}$ and every additional actually consumed mass,
embedding or source contact. Vanishing of the displayed bare one-body
contact does not discard them.

For an evolving member put
$a_q=(0,0,-\mu/2)$ and
$a_\alpha=(-2u/3,v/3,0)$ in the existing two-orientation notation.
The noncentral mixed causal kernel is explicitly
\begin{equation}
\label{p4:eq:v2212-conformal-shape-retarded}
 \boxed{\Pi^R_{q\alpha}(t,s)=\Theta(t-s)
 \sum_{\eta=\pm1}\frac{2\eta}{\hbar}
 [(R_\eta(t,s)^Ta_q(t))\times a_\alpha(s)]
                       \cdot s_\eta(s).}
\end{equation}
The actual source preparation and its finite evolution determine
$s_\eta$; these vertices introduce no independent response coefficient.

For a frozen stationary evaluation suppose the actual retained band
contains a matched axial pair $u,-u$, with equal $v,\mu$ and equal
pole-diagonal preparation in the same Clifford orientation. On each
two-level component write its one-body expectation matrix as
$C=(pI+pfh_E/E)/2$, where $p$ is its actual trace and $pf$ its signed
pole-population difference. Keep those preparation data fixed in
variation. Let $\overline\Pi^R$ denote the mean of the two actual packet
kernels. Then
\begin{equation}
\label{p4:eq:v2212-conformal-shape-stationary}
 \boxed{\begin{aligned}
 \overline\Pi^R_{q\alpha}(t)
 &=\Theta(t)\frac{pf\mu^2(v^2-2u^2)}{3\hbar E^2}
                         \sin(2Et/\hbar),\\
 \overline\Pi^R_{q\alpha}(\omega)
 &=\frac{2pf\mu^2(v^2-2u^2)/(3E)}
              {4E^2-\hbar^2(\omega+i0)^2}.
 \end{aligned}}
\end{equation}
The Fourier convention is $\int_0^\infty e^{i\omega t}(\cdot)\,dt$.
In particular,
\begin{equation}
\label{p4:eq:v2212-conformal-shape-static}
 \boxed{\overline\Pi^R_{q\alpha}(0)
 =\frac{pf\mu^2(v^2-2u^2)}{6E^3}
 =\partial_q\partial_\alpha(pfE).}
\end{equation}
The complete axial pair contributes twice this mean. Every other
already retained Clifford component and spherical multiplicity is summed
once in the original trace; no additional degeneracy or efficiency is
inserted. The massless kinetic constituent has $O_q=0$ and hence zero
mixed source response in this physical clock.
\end{corollary}
\begin{proof}
Canonical conformal spinor transport gives $\mu=m/\sqrt F$.
For the massless kinetic term $h_J=\sqrt Fh_E$, whereas
$N_J=N_E/\sqrt F$, so its coordinate Hamiltonian is
$N_Jh_J=N_Eh_E$. Its apparent homogeneous conformal rescaling cancels
against the physical clock. The inherited fixed-volume shape exponents
and differentiation of $\mu$ give
Eq.~\eqref{p4:eq:v2212-conformal-shape-vertices}.
The source action has the Hamiltonian term
$-\int N_E\operatorname{Tr}(\sigma\widehat H_E)dt$.
At fixed Einstein metric, $\delta_qg_J=-g_J\delta q$;
at fixed volume its spatial shape variation gives
$\delta_\alpha S_s=(2/3)\int N_EV_E\Delta_E\delta\alpha\,dt$.
The Hilbert variation therefore yields the two signs in
Eq.~\eqref{p4:eq:v2212-conformal-shape-ports}.
Inserting the two explicit vectors in
Proposition~\ref{p4:prop:spherical-dirac-retarded-kernel} proves
Eq.~\eqref{p4:eq:v2212-conformal-shape-retarded}.

For the stationary evaluation let $n=d/E$ and write perpendicular
components relative to $n$. Direct contraction gives
\[
 a_{q,\perp}\cdot a_{\alpha,\perp}
 =\frac{\mu^2(v^2-2u^2)}{6E^2},\qquad
 n\cdot(a_q\times a_\alpha)=\frac{\mu uv}{2E}.
\]
The constant-$d$ rotation in the preceding proposition gives
$2pf/\hbar$ times the first contraction multiplied by
$\sin(2Et/\hbar)$, plus
$2\eta pf/\hbar$ times the second contraction multiplied by
$\cos(2Et/\hbar)$.
Under the stated matched axial pairing the second term is odd in $u$
and cancels, while the first is even. This proves the time formula with
the stated mean normalization. Its retarded Fourier transform gives
the frequency formula. Finally
$E_q=-\mu^2/(2E)$ and
$E_\alpha=(v^2-2u^2)/(3E)$ give
Eq.~\eqref{p4:eq:v2212-conformal-shape-static} without a population
reset or a new thermal assumption.
\end{proof}

\begin{remark}[Response status and remaining physical execution]
\label{p4:rem:v2212-conformal-shape-scope}
This is an evaluated causal mixed source kernel of the already inherited
free massive member. Its identification with the native scalaron consumes
the existing action and normalization certificates, rather than proving
them from this finite Hamiltonian. Its sign depends on its actual population and
geometry, and its pole is the inherited $2E/\hbar$. The general evolving
formula retains the supplied state and two-time rotation; the stationary
formula requires its explicitly matched preparation. Central allocation,
interactions, derivative vertices, mixing mass sectors and moving-carrier
contacts keep their earlier complete-source treatment.
Together with the constrained local metric block in Paper~VI-A,
Theorem~\ref{p8:thm:v2212-constrained-shape-hessian}, this supplies
explicit entries for the existing same-action response construction on
their common admitted domain; a shared source and metric history must
still be evolved. The result neither identifies this physical source kernel with an
Euclidean heat Hessian nor constructs the complete nonlocal gravitational
action. Actual state formation, admissible inhomogeneous source embedding,
coupled scalaron--metric propagation and a complete-curvature barrier on
a matched collapse history remain physical execution tasks. The earlier
finite response, work, localization and source-domain statements are
unchanged.
\end{remark}
% END EXACT MASSIVE CONFORMAL--SHAPE SOURCE KERNEL v2.212

% BEGIN SOURCE-POLARIZATION-AND-NATIVE-UV-LEDGER v2.216


\section{Quantum readout exports and physical scope}
\label{p4b:sec:v238-b8}
\noindent\textit{Source and claim boundary.} Collects source-visibility, response-equivalence, frequency and downstream application conditions. Paper~IV-A supplies native finite quantum structure; this paper tests actual process/readout certificates; Paper~IV-C contains detailed auxiliary verifications.

\subsection{Quasi-local many-body and thermodynamic bridge}
\label{p4:sec:part2-thermodynamic-bridge}

The finite framework also admits the standard algebraic route to many-body thermodynamic language.


\begin{theorem}[Quasi-local Gibbs accumulation and limiting-KMS criterion]
\label{p4:thm:quasilocal-thermodynamic-bridge}
On a supplied thermodynamic interface domain, let
\(\Lambda_1\subset\Lambda_2\subset\cdots\) be an increasing family of finite
subsystems with local algebras
\[
\mathcal A_{\Lambda_n}=\bigotimes_{x\in\Lambda_n}\mathcal A_x
\]
together with specified compatible injective unital embeddings.  The sites,
tensor products, embeddings, boundary data, and finite Hamiltonians
constitute the thermodynamic net; they are not the Paper-II Peter--Weyl
truncation ladder.  Then the union of local algebras admits the
\(C^*\)-completion
\[
\mathcal A_{\mathrm{ql}}:=\overline{\bigcup_n \mathcal A_{\Lambda_n}},
\]
the quasi-local algebra.
For each inverse temperature \(\beta>0\), the finite-volume Gibbs states
\[
\omega_{n,\beta}(A)=\frac{\Tr(e^{-\beta H_{\Lambda_n}}A)}{\Tr(e^{-\beta H_{\Lambda_n}})}
\]
admit state extensions to \(\mathcal A_{\mathrm{ql}}\), and every chosen
sequence of such extensions has weak-* accumulation points.  Any such
limit is a finite-volume Gibbs accumulation state.

For a KMS conclusion, additionally suppose that the finite local dynamics
\(\alpha_t^{(n)}=\Ad(e^{itH_{\Lambda_n}})\) converges in norm on the local
algebra, uniformly for \(t\) in compact real intervals, to a strongly
continuous quasi-local automorphism group \(\alpha_t\).  Suppose there is a
dense \(\alpha\)-entire \(^*\)-subalgebra
\(\mathcal D\subset\mathcal A_{\mathrm{ql}}\) such that, for every
\(A,B\in\mathcal D\), compatible finite-volume approximants
\(A_n,B_n\in\mathcal A_{\Lambda_n}\) satisfy
\[
A_n\longrightarrow A,\qquad B_n\longrightarrow B
\]
in norm and
\[
\alpha_z^{(n)}(B_n)\longrightarrow\alpha_z(B)
\]
in norm, uniformly for \(z\) in compact subsets of the closed strip
\(0\le\Im z\le\beta\).  Then every Gibbs accumulation state obtained along
the same convergent subnet satisfies, for \(A,B\in\mathcal D\),
\[
\omega_\beta\!\left(A\alpha_{t+i\beta}(B)\right)
=
\omega_\beta\!\left(\alpha_t(B)A\right),
\qquad t\in\mathbb R,
\]
and therefore extends by the usual density criterion to a
\(\beta\)-KMS state for \(\alpha\).  Real-time convergence alone does not
justify this imaginary-time passage.
\end{theorem}


\begin{proof}
Each supplied injective unital embedding is an isometric
\(^*\)-homomorphism.  Hence the directed union
\(\bigcup_n\mathcal A_{\Lambda_n}\) is a normed \(^*\)-algebra, and its
completion is the quasi-local algebra \(\mathcal A_{\mathrm{ql}}\).

For every \(n\), \(\omega_{n,\beta}\) is a state on the
\(C^*\)-subalgebra
\(\mathcal A_{\Lambda_n}\subset\mathcal A_{\mathrm{ql}}\).
The state-extension theorem gives at least one extension
\(\widetilde\omega_{n,\beta}\) to
\(\mathcal A_{\mathrm{ql}}\).  The state space of a unital
\(C^*\)-algebra is weak-* compact, so any chosen sequence of extensions has
a convergent subnet.  Its limit \(\omega_\beta\) is a state, and on each
fixed local observable the corresponding finite-volume expectation values
converge along that subnet.

The finite Gibbs states satisfy
Theorem~\ref{p4:thm:finite-thermal-kms-kinetic-interface}.  Along the same
subnet, weak-* convergence of the states and the assumed norm convergence
of the analytic-time evolutions allow passage to both sides of that
identity:
\[
\omega_\beta\!\left(A\alpha_{t+i\beta}(B)\right)
=
\omega_\beta\!\left(\alpha_t(B)A\right).
\]
The dense analytic-subalgebra criterion then gives the KMS property.
Without the limiting automorphism group and closed-strip convergence, the
exact conclusion stops at weak-* Gibbs accumulation.
\end{proof}

\begin{remark}[What is exact here and what remains a limiting criterion]
\label{p4:rem:thermo-conditionality}
The quasi-local completion and weak-* compactness statements in Theorem~\ref{p4:thm:quasilocal-thermodynamic-bridge} are exact algebraic consequences of the compatible finite local system.
The local subsystem net, site factorization, embeddings, boundary
conditions, interaction family, and finite-volume Gibbs sequence are
themselves supplied thermodynamic-interface data.  Paper~II's spectral
symbol ladder does not construct that net.  What remains a further limiting
criterion is the quasi-local Heisenberg dynamics together with the analytic
closed-strip convergence required for the KMS identity.
\end{remark}

Thermodynamic, bath, and plasma language is therefore licensed only after
the relevant subsystem net, interaction, states, and dynamics have been
supplied.  Paper~IV proves the finite and quasi-local consequences on that
domain; it does not derive the thermodynamic net from the Paper-II spectral
ladder.

\begin{remark}[Local kinetic equilibrium is weaker than exact global KMS convergence]
\label{p4:rem:v490-local-kinetic-versus-kms}
A cosmological gauge plasma may license a local temperature and radiation
equation of state on a supplied effective-kinetic domain when its microscopic
collision/splitting rates drive the relevant one-particle distributions to
local kinetic/chemical equilibrium on timescales short compared with the
background evolution.  This statement is not the same as, and does not require,
an exact GKLS limit or a global thermodynamic-limit KMS theorem for the full
closed cosmological state.  Conversely, Paper~IV's finite Gibbs, reservoir,
GKLS, and quasi-local KMS results remain available whenever that stronger
analytic claim is explicitly requested.  Exact conserved charges may survive
as chemical potentials or central-sector weights in either description.
\end{remark}



\subsection{Atlas status of thermal and scattering interfaces}
\label{p4:sec:v236-state-scattering-atlas}
\begin{proposition}[Equilibrium versus nonequilibrium thermal branches]
\label{p4:prop:v236-thermal-atlas}
For a supplied finite self-adjoint Hamiltonian \(H\) and inverse temperature
\(\beta\) on a simple algebra, the normalized Gibbs/KMS density
\[
 \rho_\beta=Z_\beta^{-1}e^{-\beta H}
\]
is uniquely determined, as in
Theorem~\ref{p4:thm:finite-simple-kms-uniqueness}.  On a direct sum,
equilibrium fixes the normalized Gibbs density within each block and leaves
the central probabilities free.  In particular, on the native
\(M_{11}\oplus M_4\) coordinate algebra the full family is
\[
 \boxed{\rho_{\beta,p}
 =p\,\frac{e^{-\beta H_g}}{Z_g}
  \oplus(1-p)\,\frac{e^{-\beta H_s}}{Z_s},
 \qquad 0\le p\le1,\qquad
 Z_j=\Tr(e^{-\beta H_j}).}
\]
Faithfulness requires \(0<p<1\).  The globally normalized exponential is
the particular member \(p=Z_g/(Z_g+Z_s)\), selected only when the ensemble
or an actual exchange process supplies that central weighting; the
blockwise KMS identities alone do not select it.  Cosmological reheating,
freeze-out and baryogenesis branches generally consume a nonequilibrium
initial member and kinetic history; these remain actual state/history data
unless a later kinetic theorem generates them.
\end{proposition}
\begin{proposition}[Finite scattering, pole and decay hierarchy]
\label{p4:prop:v236-scattering-atlas}
The exact finite-window channel map and finite Feshbach/Schur resonance are
finite-theory outputs.  A large-window asymptotic scattering operator requires
the strong-Cauchy/limit certificate already stated in this paper.  A nonreal
resonance pole or nonzero width additionally requires a continuum/limiting-
absorption continuation, and a time-homogeneous exponential decay law requires
a declared irreversible reservoir/Markov branch.  None of these analytic
certificates introduces new interaction vertices or particle representations.
\end{proposition}
% END v2.36 STATE/SCATTERING ATLAS



% BEGIN INLINE SOURCE: Paper4-FiniteRootStateTransport-v1_1.tex
% Source: G30bRootVacuum v1.10 plus v2.114 QROOT neutral-root closure; retained finite-state transport and fresh-branch source-response results.


\subsection{Finite root-state restrictions and conditional absorption}
\label{p4:sec:v2109-finite-root-state-transport}
The positive ground-support theorem already distinguishes zero excitation
from a vacuum-energy offset.  Here the finite root candidates acquire
explicit restriction maps and a propagation error bound.  These construct states on a specified quantum branch.  The generic branch still requires a state-occurrence word; the native fresh cyclic root below supplies that word by the neutral Fock-unit theorem.  Parent VACCAP remains a separate claim.

\begin{theorem}[Finite free root candidates and compatible state restrictions]
\label{p4:thm:v2109-finite-root-restrictions}
Let \(\mathcal K\) be a finite one-particle space, \(h>0\), and
\(\mathcal F_N(\mathcal K)=\bigoplus_{n=0}^N\operatorname{Sym}^n\mathcal K\)
with a complete total-occupation cap \(N\ge1\).  The positive free energy
\(E_N=d\Gamma(h)|_{\mathcal F_N}\) has unique ground projector \(P_0\)
and gap \(\min\operatorname{spec}h\).  The compressed canonical
coordinates have the usual ground quadratic covariance, because their
vacuum--one-particle matrix elements are retained.
For \(\mathcal L=\mathcal K\oplus\mathcal Q\), the natural occupation
isometry
\[
 V^N:\mathcal F_N(\mathcal L)\longrightarrow
 \mathcal F_N(\mathcal K)\otimes\mathcal F_N(\mathcal Q)
\]
defines the CPTP restriction
\[
 \mathcal R^N_{\mathcal K\leftarrow\mathcal L}(\rho)
 =\operatorname{Tr}_{\mathcal Q}(V^N\rho V^{N\dagger}).
\]
For nested mode spaces these maps compose exactly, and they send the larger
vacuum projector to the smaller one.  Their dual maps are unital completely
positive, not generally multiplicative.
\end{theorem}
\begin{proof}
Every occupied eigenvector has energy equal to a sum of positive
one-particle energies; the vacuum alone has zero energy, and the smallest
one-particle vector is admitted when \(N\ge1\).  Each quadratic vacuum
moment passes only through the zero- or one-particle subspace, which proves
the covariance statement without asserting finite exact CCR.
The isometry followed by partial trace is CPTP.  In occupation coordinates,
\[
 (\mathcal R^N\rho)_{\mathbf n,\mathbf m}
 =\sum_{\substack{\mathbf r:\ |\mathbf n|+|\mathbf r|\le N\,;
 |\mathbf m|+|\mathbf r|\le N}}
 \rho_{(\mathbf n,\mathbf r),(\mathbf m,\mathbf r)}.
\]
Nested restriction sums over the same complementary tuples with the same
total-occupation condition; associativity of finite sums gives the cocycle.
The vacuum factors under \(V^N\).  The dual is
\(A\mapsto V^{N\dagger}(A\otimes I)V^N\); compression need not preserve
products.  Thus this is a projective family of state channels, not an
inclusion system of full matrix algebras.
\end{proof}

\begin{theorem}[Neutral branch-origin Fock-unit actualization]
\label{p4:thm:v2114-neutral-fock-unit}
Let a constituted branch origin generate one or more already licensed
Bose/Fermi mode carriers \(\mathcal K_\alpha\) and their Fock constructors,
with degree-zero units
\[
 \operatorname{Sym}^0\mathcal K_\alpha=\mathbb C
 \quad\hbox{or}\quad
 \Lambda^0\mathcal K_\alpha=\mathbb C .
\]
Assume that, in the state role under discussion, the occurring branch word
contains no inherited pre-branch state and no independently generated
state/preparation token, creation word, coherent displacement, Bogoliubov
squeeze, thermal weight, mixing channel, or boundary-state/process insertion.
Then the occurring output of the neutral Fock constructor is the tensor
product of the degree-zero units,
\[
 \boxed{\Omega_{\rm br}=\widehat\bigotimes_\alpha\Omega_\alpha,
 \qquad \rho_{\rm br}=|\Omega_{\rm br}\rangle
                    \langle\Omega_{\rm br}|.}
\]
In particular every generated occupation number obeys
\[
 \boxed{N_{\alpha j}\Omega_{\rm br}=0.}
\]
This theorem does not say that the one-particle carriers are absent.  It says
that an available mode is not an occupied mode: a nonzero occupation or a
nonvacuum normal state requires a nontrivial state/preparation/process word
(or an explicit independent occurrence token) beyond the neutral constructor.
\end{theorem}

\begin{proof}
The Fock functor contains a canonical degree-zero summand and hence a canonical
unit map \(1\mapsto\Omega_\alpha\).  Every positive occupation basis vector
contains at least one creation operation.  Coherent, squeezed, thermal and
mixed alternatives likewise require their corresponding nontrivial
preparation/process output or an independently occurring state member.
Paper~II-A GP2 no-insertion and rooted generated-image co-realization permit
only outputs of the licensed occurring word; they do not promote unused
mathematical members of the ambient state space to physical occurrences.
Under the stated branch-origin hypotheses all nonneutral preparation words
are absent, while the neutral unit word is part of the already licensed Fock
constructor.  Hence the state-role output is the product unit.  The number
operator statement is the defining degree-zero property.
\end{proof}

\begin{theorem}[Positive excitation-family ground-source compression]
\label{p4:thm:v2114-positive-family-ground-source}
Let \(x\) range over an open set of admitted geometric/source configurations
and let \(\mathcal E(x)\succeq0\) be a \(C^1\) family of positive excitation
operators on a fixed finite Hermitian carrier, or on a smoothly identified
finite carrier bundle with its metric-compatible connection.  Put
\[
 P=\mathbf1_{\{0\}}(\mathcal E(x_*)).
\]
For every admitted two-sided variation \(v\),
\[
 \boxed{P(D_v\mathcal E)(x_*)P=0.}
\]
If the already generated claim-complete quantum source interface at \(x_*\)
has covered components
\[
 \widehat{\mathcal T}_\alpha
 =c_\alpha\mathcal E+\sum_i c_{\alpha i}D_i\mathcal E,
\]
then every ground-supported density operator \(\sigma=P\sigma P\) satisfies
\[
 \boxed{\operatorname{Tr}(\sigma\widehat{\mathcal T}_\alpha)=0}
\]
for every covered source component \(\alpha\).  The statement requires no
simple ground and no constant kernel rank in a neighbourhood; it does require
that the displayed source family is the complete off-shell source interface
for the claim being made.
\end{theorem}

\begin{proof}
For \(\xi\in\operatorname{Ran}P\), the function
\(s\mapsto\langle\xi,\mathcal E(x_*+sv)\xi\rangle\) is nonnegative on both
sides of zero and vanishes at zero.  Its derivative therefore vanishes.
Polarization gives the compressed derivative identity.  In a moving unitary
trivialization, the connection contributes a commutator
\([A_v,\mathcal E]\), whose ground compression vanishes because
\(\mathcal EP=P\mathcal E=0\).  The source statement follows by linearity
and \(\sigma=P\sigma P\).
\end{proof}

\begin{theorem}[Finite evolution controlled by the boundary flux]
\label{p4:thm:v2109-finite-evolution-error}
On a supplied finite-mode regular quadratic branch, let \(\psi(t)\) be
the full evolution of the vacuum, \(P_N\) the fixed total-occupation
projection, \(Q_N=I-P_N\), and \(\psi_N(t)\) the evolution generated by
\(P_NH(t)P_N\) from the same vacuum.  On intervals with the indicated
Duhamel integral finite, put
\[
 \tau_N(t)=\|Q_N\psi(t)\|^2,\qquad
 \eta_N(t)=\frac1\hbar\int_0^t\|P_NH(s)Q_N\psi(s)\|\,ds.
\]
Then
\[
 \|\psi(t)-\psi_N(t)\|^2\le\tau_N(t)+\eta_N(t)^2.
\]
In the convention
\(H_{\rm pair}=\hbar\sum_{ij}(B_{ij}a_i^\dagger a_j^\dagger+
\overline B_{ij}a_ja_i)/2\), with \(B\) symmetric, one may bound
\[
 \eta_N(t)\le\frac12\int_0^t\|B(s)\|_{\rm HS}
 \sqrt{N(N+1)p_{N+1}(s)+(N+1)(N+2)p_{N+2}(s)}\,ds,
\]
where \(p_n\) is the full state's total-number probability.
\end{theorem}
\begin{proof}
The difference \(P_N\psi-\psi_N\) is the Duhamel integral of
\(-iP_NHQ_N\psi/\hbar\) against the finite unitary propagator.
The omitted component \(Q_N\psi\) is orthogonal to that retained
difference.  Taking norms proves the first inequality.
Only pair annihilation takes \(Q_N\) into \(P_N\), and only occupations
\(N+1,N+2\) contribute.  Their output sectors are orthogonal.  The
identity \(\sum_{ij}\|a_ja_i\xi_n\|^2=n(n-1)\|\xi_n\|^2\), followed
by Cauchy--Schwarz over \(i,j\), gives the displayed bound.
\end{proof}
For a Bogoliubov-evolved vacuum let \(n_j\) be the eigenvalues of
\(\beta\beta^\dagger\) and \(\ell_j=n_j/(1+n_j)\).  Passive mode
rotations preserve total number; reducing to single-mode squeezed vacua
and summing their binomial series gives
\[
 M_t(z)=\langle z^{\mathsf N}\rangle
 =\prod_j\sqrt{\frac{1-\ell_j(t)}{1-z^2\ell_j(t)}},\qquad
 z>1,\quad z^2\max_j\ell_j(t)<1.
\]
Consequently \(p_n\le M_t(z)z^{-n}\) and
\(\tau_N\le M_t(z)z^{-(N+1)}\), yielding a computable history bound.
Trace distance is at most the norm bound above and is contractive under
restriction.  Bounded-observable errors are at most twice the observable
norm times this bound.  Fixed-degree unbounded moments require the
corresponding weighted-tail estimates, not a bare trace-distance argument.

\begin{theorem}[One unnormalized Euclidean weight characterizes zero excitation]
\label{p4:thm:v2109-euclidean-weight-support}
On a finite carrier let \(E\succeq0\), \(\rho\succeq0\),
\(\operatorname{Tr}\rho=1\), and \(\tau>0\).  Then
\[
 \operatorname{Tr}(e^{-\tau E/2}\rho e^{-\tau E/2})=1
 \quad\Longleftrightarrow\quad
 \operatorname{supp}\rho\subseteq\ker E.
\]
If \(E\ge\Delta(I-P_0)\), \(\Delta>0\), then
\[
 \operatorname{Tr}\rho(I-P_0)\le
 \frac{1-\operatorname{Tr}(\rho e^{-\tau E})}{1-e^{-\tau\Delta}}.
\]
\end{theorem}
\begin{proof}
The positive operator \(I-e^{-\tau E}\) has precisely \(\ker E\) as
kernel.  Its zero expectation forces the state into that kernel by the
existing positive-support theorem.  The converse is immediate.  The gap
bound follows from
\(I-e^{-\tau E}\ge(1-e^{-\tau\Delta})(I-P_0)\).
\end{proof}
An occurring cap must independently supply this unnormalized weight
condition before it becomes a physical selection theorem.  Mathematical
availability of the heat map does not supply its incidence.  The distinct
normalized-stationarity and Gaussian-construction branches are set out in
Appendix~\ref{p4:app:v2109-gaussian-boundary}.

% END INLINE SOURCE: Paper4-FiniteRootStateTransport-v1_1.tex




\subsection{Release v1.64: action poles, finite vacuum, and kinetic boundary}
\label{p4:sec:v164-action-poles-vacuum}

\begin{theorem}[Action-Hessian pole owner]
\label{p4:thm:v164-action-hessian-pole-owner}
For an occurring action and state root, the quadratic response is
\(K=\delta^2S\).  On a renormalized field block the physical two-point
operator is
\[
 G_R(z)=\bigl(K_R(z)-\Sigma_R(z)\bigr)^{-1}.
\]
An interacting mass is an isolated zero of the inverse two-point operator
with nonzero resolvent residue.  A finite lookup-row energy, Gram norm, or
representation count is not a pole unless a generated intertwiner identifies
it with this resolvent datum.  No such identification follows from equality
of numbers or common global ancestry.
\end{theorem}

\begin{theorem}[Finite Gibbs vacuum and no closed thermalization]
\label{p4:thm:v164-finite-gibbs-vacuum-no-thermalization}
Let \(H\) be a constituted finite Hamiltonian with ground projector \(P_0\)
of rank \(d_0\).  Then
\[
 \rho_\beta=\frac{e^{-\beta H}}{\operatorname{Tr}e^{-\beta H}}
 \longrightarrow \rho_0=\frac{P_0}{d_0}
 \qquad(\beta\to\infty).
\]
If the transfer is positivity improving, Perron--Frobenius gives \(d_0=1\).
Thus the symmetric finite vacuum on the declared vacuum branch is derived by
functional calculus.  By contrast, a closed finite unitary evolution cannot
converge irreversibly to a Gibbs state: unitary conjugation preserves trace
distance and finite spectra recur.  Thermal convergence therefore requires a
generated mixing, repeated-probe, large-bath, or kinetic limit; only the KMS
target is fixed by the supplied \((H,\beta)\).
\end{theorem}




\subsection{Parent meaning of the degree-two comparison layer}
\label{p4:sec:v187-cartan-square-comparison}
The symmetric-pair Cartan embedding of the completed $M_8$ parent obeys
\[
 c(k)=(k,k^{-1}),\qquad r(c(k))=k^2.
\]
Thus the primitive factor carrier remains degree one, while the first
intrinsic parent relation/composition layer is quadratic.  On $U(1)$ this
produces the $\pm2$ characters from primitive $\pm1$ characters; on $SU(2)$
products of fundamental $j=\frac12$ matrix coefficients generate the
$j=0\oplus1$ sector.  This explains why the first relation-closed comparison
algebra reaches degree two.

This argument is explanatory, not a replacement for the independent sharp
bounds already proved in the comparison audit:
\[
\begin{aligned}
 m_{\rm cmp}&\ge2 &&\text{from the clock-operation requirement},\\
 m_{\rm cmp}&\le2 &&\text{from exact quartic }2I\text{ evaluation}.
\end{aligned}
\]
Hence
\[
 \boxed{\begin{gathered}
 m_p=1\ \text{(primitive momentum carrier)},\\
 m_{\rm cmp}=2\ \text{(unique first relation/action comparison band)}.
 \end{gathered}}
\]



\subsection{Same-action cosmological quadratic-kernel import}
\label{p4:sec:v189-same-action-cosmological-kernels}
Paper~VI-A now supplies the scalar and tensor cosmological sectors as explicit
quadratic Hamiltonian blocks descended from the same action, rather than as a
phenomenological transfer matrix.  For a resolved mode write
\[
 Y_A=(q_A,\Pi_A)^T,\qquad A\in\{S,T\}.
\]
In a fixed symplectic trivialization its classical evolution has the form
\[
 \dot Y_A=J\mathbb H_A(t)Y_A,
 \qquad
 \mathcal T_A(t_+,t_-)=\mathcal P\exp\!\int_{t_-}^{t_+}J\mathbb H_A(t)\,dt.
\]
On a moving finite-rank mode bundle, the metric/symplectic connection already
required by the curved-background bridge is included in the covariant
generator before this path-ordered exponential is formed.  Consequently the
classical $K=16$ transfer is generated once the Paper-VI-A JUNC member,
background history and mode trivialization are fixed; it is not an
independent physical coefficient.

The Bogoliubov transformation still requires the positive-energy/complex
structure, state and implementability hypotheses of this paper.  Therefore a
classical transfer does not by itself determine occupation numbers or a
numerical reheating efficiency.


% BEGIN FERMION-FACTOR-MULTIPLICATION v4.91
\begin{corollary}[Factor-local fermion source multiplication]
\label{p4:cor:v491-fermion-factor-source-multiplication}
Let the relation-free matter carrier factor as $F=W\otimes Q\otimes G$.
If $\Sigma:W_R\to W_L$ is an odd chiral factor and $R\in\End(G)$ is an even
family factor generated on independent tensor slots, then the existing
relation-free multiplication map gives
\[
 (\Sigma\otimes I_Q\otimes I_G)(I_W\otimes I_Q\otimes R)
 =\boxed{\Sigma\otimes I_Q\otimes R}.
\]
No additional cross-factor coefficient is generated.  An additive expression
$\Sigma\otimes I+I\otimes R$ is not a pure odd source because it mixes odd and
even grading.
\end{corollary}
% END FERMION-FACTOR-MULTIPLICATION v4.91


% BEGIN EXACT-GRADED-EAA v4.100


\subsection{Transported presentation and the actual source-band obligation}
\label{p4:app:v4118-transported-embedding-diagnostics}
This appendix diagnoses the distinction needed to apply
Theorem~\ref{p4:thm:v4118-transported-embedding-covariance}. It selects
no native state, interaction, carrier or physical cutoff. The generic
matrix example below is not an admissible physical countermodel or a
refutation of the native parent action's diffeomorphism symmetry.

\paragraph{Actual reference-band domain.}
The current free source in
Section~\ref{p3f:sec:v1148-free-metric-dirac} uses the actual reference
component band $P_m$. Its consumed Ward theorem
\ref{p8:thm:v2149-guarded-free-metric-ward} supplies the free first-response
guard $m\ge K+2n$ and its stated support, force, state and contact
conditions. Its parent multiplication words are compressed only at the
outside; separately compressed products require the actual guard.
The generic interacting finite constructor consumes $J$ and its jets
rather than choosing a source-dependent physical carrier.

Re-expressing that same band as $T_\varphi P_mT_\varphi^{-1}$ proves
presentation covariance. Identifying it with the carrier consumed by
the actual metric-source family requires the complete family diagram
of Corollary~\ref{p4:cor:v4118-source-bound-spatial-ward}, or direct
evaluation of the actual full Ward defect and its first variation.
A native constructor identity can establish the diagram before any
component matrix is printed. If an independently retained reference-band
or background mark varies in the complete presentation, its actual work
belongs to the external-force term of
Theorem~\ref{p4:thm:v4106-canonical-stress-ward}. Total presentation
covariance alone does not discard that term in a metric-only derivative.
This is a typed same-owner/source-map obligation, not an added physical
symmetry premise and not a proof that the actual native identity fails.

\paragraph{Exact frozen-compression comparison.}
Appendix~\ref{p4:app:v4106-cutoff-ward-defect} gives, for a frozen
many-body $P$ and $Q=I-P$, the exact leakage identity
\[
 [PHP,PGP]-P[H,G]P=PGQHP-PHQGP.
\]
For an elementary comparison set $\hbar=1$ and use the finite parent
\[
 H=\begin{pmatrix}1&0&1\\0&2&0\\1&0&3\end{pmatrix},\qquad
 G=\begin{pmatrix}0&0&0\\0&0&1\\0&1&0\end{pmatrix},\qquad
 P=\operatorname{diag}(1,1,0).
\]
The full family generated by $T_\epsilon=e^{-i\epsilon G}$ has
$\delta H=-i[G,H]$ and zero parent defect. On $\operatorname{Ran}P$,
$PGP=0$, while its frozen compressed defect is exactly
\[
 \mathfrak D_F=i(PHQGP-PGQHP)
       =i\begin{pmatrix}0&1\\-1&0\end{pmatrix}\ne0.
\]
It is Hermitian and noncentral. Transporting the actual embedding
$J'=T_\epsilon J$ cancels the presentation change exactly by
Theorem~\ref{p4:thm:v4118-transported-embedding-covariance}; leaving the
embedding frozen does not. The example establishes only this logical
distinction. The actual native interacting defect must be evaluated
with its own generated tensors, constraints, projectors and contacts.

\paragraph{Optional spectral-band comparison.}
Appendix~\ref{p10:app:v2186-covariant-spectral-band} supplies a conditional
Dirac spectral-band covariance construction on its additional global
domain and comparison radius. As stated there, it is not consumed by
the native positive proof chain. Its Riesz/Kato jets compute a licensed
cluster rather than selecting that cluster. Neither that comparison
radius nor an arbitrarily covariant moving projector can replace the
actual admitted reference band. An already licensed native covariant
band can use the positive theorem directly on its own complete diagram.

\paragraph{Scope retained by the positive chain.}
Complete passive-clock pullback of the existing finite word proves the
homogeneous time identity and complete source work of
Corollary~\ref{p4:cor:v4118-homogeneous-source-work}. It keeps the actual
proper-time calibration and pulls back existing marks. It does not
license a new calibrated physical rate. Its spatially constant tangent
does not imply a spatially varying normal-time Ward identity. Likewise,
spatial presentation covariance does not construct a different foliation,
positive-frequency split, constraint reduction or bosonic occupation
word. Those application claims retain their actual canonical lift and
source-contact domains.

% END INLINE SOURCE: Paper4-TransportedEmbeddingDiagnostics-Appendix-v1_0.tex


% BEGIN INLINE SOURCE: Paper4-GaussianBoundaryAppendix-v1_0.tex
% Source: G30bRootVacuum v1.9, Sections 37--41; optional amplitude branch.


\subsection{Geometry-first consequences and physical scope (framework v2.201)}
\label{p4:sec:v2201-geometry-first-consequences}
\paragraph{Geometry ancestry.}
Protected finite carrier and positive physical pairing $\to$ occurring source/instrument maps $\to$ finite frequencies and shared spectral residues.
\paragraph{Status.}
These are consequences on the domains stated below.
A mathematical certificate or an executed restriction retains every
physical source, state, clock and readout hypothesis of its inputs.



\paragraph{D04: Born probabilities to frequencies require a correlation bound}
\label{p4:sec:v2201-derived-d04}
\noindent\textit{Status: proved result on stated domain.}

For an occurring sequential binary readout $X_i\in\{0,1\}$ define
$\bar X_N=N^{-1}\sum_iX_i$. On a stationary instrument history with
mean $p$ and lag covariance $c_k$,
\[
 \operatorname{Var}(\bar X_N)=\frac{p(1-p)}N
 +\frac2{N^2}\sum_{k=1}^{N-1}(N-k)c_k.
\]
If the actual instrument transfer gives $|c_k|\le A\zeta^k$,
$0\le\zeta<1$, then
\[
 \operatorname{Var}(\bar X_N)\le
 \frac{p(1-p)+2A\zeta/(1-\zeta)}N.
\]
\begin{proof}
Expand the variance of the sum and group ordered pairs by their
lag. Bound $N-k\le N$ and sum the geometric series. Joint probabilities are
computed from the successive CP instrument maps, not from carrier counts.
\end{proof}
Chebyshev's inequality now gives an explicit finite sampling-error bound.
For a persistent projective readout, all $X_i$ can coincide and the variance
stays $p(1-p)$: the repeated count itself does not create independence.
The remaining native task is to derive the concrete apparatus/instrument
and its correlation bound from its action. \par\noindent\textit{Inherited inputs:} \ref{p4:thm:finite-choi-kraus-instruments}, \ref{p4:prop:effect-instrument-nonselection}.



\paragraph{A8: shared simple poles give cross-observable consistency tests}
\label{p4:sec:v2201-derived-a8}
\noindent\textit{Status: proved corollary on stated domain.}

On a stationary finite Hermitian source chart with its positive physical
pairing, let
\(G_{ab}(z)=j_a^\dagger(zI-H)^{-1}j_b\). At a nondegenerate visible
eigenvalue \(\lambda\), put \(v_a=j_a^\dagger u_\lambda\). Then
\[
\boxed{\operatorname{Res}_{z=\lambda}G_{ab}=v_a\overline v_b,\qquad
\operatorname{rank}\operatorname{Res}G\le1.}
\]
Every two-by-two minor vanishes, and
\(|R_{ab}|^2=R_{aa}R_{bb}\). A pole of multiplicity \(r\) instead has
residue rank at most \(r\).

\textbf{Proof.} The Hermitian spectral resolution contributes
\(u_\lambda u_\lambda^\dagger/(z-\lambda)\). Apply the same source maps
on both sides. Analytic contact terms do not change its residue.
\par\noindent\textit{Inherited inputs:} \ref{p4:thm:v2188-moment-pencil}, \ref{p3f:prop:v2167-complete-neutral-return}.

This can test whether several readouts really arise from one common mode:
mass/current probes, or future stationary/frozen-window metric response
probes. It does not assert that K48 has one physical propagating mode,
that a full retarded cosmological response is rank one, or that all contacts
and multiple poles can be discarded. The actual source maps and analytic
domain must be shared before applying the test to dynamics or lensing.

% END GEOMETRY-FIRST CONSEQUENCES v2.201



% Framework v2.203: Gaussian physical work, causal noise, and cold-tail domains.
% This fragment adds no source owner, state choice, or action coefficient.


\subsection{Literature-checked finite readout and effective-response certificates}
\label{p4:sec:v2210-lit-quantum-certificates}
The following results specialize already owned finite instruments, source
blocks and detector channels. They introduce no preparation, apparatus,
reservoir or physical action member. In particular the finite state/effect
and Born evaluation theorems remain closed; an actual experiment consumes
its occurring instrument rather than another probability postulate.



\paragraph{An event-sector certificate for the D04 frequency bound}
\begin{theorem}[Exact instrument correlations and finite event-sector mixing]
\label{p4:thm:v2210-lit-instrument-mixing}
Let $\{\mathcal I_i\}$ be an occurring finite CP instrument on its
constituted matrix algebra, let $\Phi=\sum_i\mathcal I_i$ be trace
preserving, and let the actual stationary preparation obey
$\Phi(\sigma)=\sigma$. For one binary recorded event write
$\mathcal I_1$ for its outcome operation and put
\[
 p=\operatorname{Tr}\mathcal I_1(\sigma),\qquad
 v=\mathcal I_1(\sigma)-p\sigma,\qquad
 P(X)=\operatorname{Tr}(X)\sigma,\qquad R=\Phi-P.
\]
Define the finite event-cyclic space
$\mathcal V_1=\operatorname{span}_{\mathbb C}\{R^nv:n\ge0\}$.
The stationary lag covariance of the actual successive instrument record is
\begin{equation}
 c_k=\operatorname{Tr}\mathcal I_1\Phi^{k-1}\mathcal I_1(\sigma)-p^2
 =\operatorname{Tr}\mathcal I_1 R^{k-1}v,\qquad k\ge1.
 \label{p4:eq:v2210-lit-event-covariance}
\end{equation}
If $\mathcal V_1=0$, all these covariances vanish. Otherwise, if the
computed spectral radius $r_1=r(R|_{\mathcal V_1})<1$, then for every
$r_1<\zeta<1$ there is a finite, computable $C_\zeta$ such that
\[
 |c_k|\le C_\zeta\|v\|_1\zeta^{k-1},\qquad
 \operatorname{Var}(\bar X_N)
 \le\frac{p(1-p)+2C_\zeta\|v\|_1/(1-\zeta)}N.
\]
Only the event-cyclic restriction is consumed; uniqueness of the whole
stationary state or mixing of every unused central sector is unnecessary.
\end{theorem}
\begin{proof}
The successive CP operations give the joint event probability in
\eqref{p4:eq:v2210-lit-event-covariance}. Stationarity and trace
preservation imply $P\Phi=\Phi P=P$, and $\operatorname{Tr}v=0$.
Thus $\Phi^{k-1}v=R^{k-1}v$, including $k=1$. A CP trace-nonincreasing
operation has dual event effect $0\le E_1\le I$, so
$|\operatorname{Tr}\mathcal I_1(X)|=|\operatorname{Tr}(E_1X)|\le\|X\|_1$.
Finite Jordan or Schur analysis of $R|_{\mathcal V_1}$ gives
$\|R^n|_{\mathcal V_1}\|_{1\to1}\le C_\zeta\zeta^n$: the finite
polynomial factors of Jordan blocks are absorbed by any $\zeta>r_1$.
Apply this to $v$, insert the result into the exact D04 variance formula
of Section~\ref{p4:sec:v2201-derived-d04}, bound $N-k\le N$, and sum
the geometric series.
\end{proof}
The finite positive-map and primitivity results
\cite{LitQFTEvansHoeghKrohn1978,LitQFTSanzWielandt2010} provide useful
full-channel sufficient tests. The smaller event-sector calculation above
can succeed even when such a full-channel test fails. The remaining Q01/Q02
execution is the actual action-induced instrument, its preparation and the
displayed spectral/correlation certificate; finite carrier counts do not
prove $r_1<1$. Chebyshev's inequality gives the corresponding finite
sampling-error bound without adding independence of the recorded events.

\begin{corollary}[Finite event resolvent and computable contraction certificate]
\label{p4:cor:v2211-finite-event-resolvent-certificate}
On the stationary instrument domain of
Theorem~\ref{p4:thm:v2210-lit-instrument-mixing}, suppose
\(\mathcal V_1\ne0\) and put \(A=R|_{\mathcal V_1}\) with \(r(A)<1\).
In any orthonormal Hilbert--Schmidt basis of this actual event-cyclic
space, use \(v\) for the coordinates of the already defined event vector
and the row functional
\(\ell(X)=\operatorname{Tr}\mathcal I_1(X)\).  Then the exact finite
variance is
\begin{equation}
 \boxed{
 \operatorname{Var}(\bar X_N)=\frac{p(1-p)}N+
 \frac2{N^2}\ell\bigl[(N-1)I-NA+A^N\bigr](I-A)^{-2}v.}
 \label{p4:eq:v2211-event-resolvent-variance}
\end{equation}
The positive matrix \(P\) defined by the finite Lyapunov equation
\[
 P-A^\dagger PA=I
\]
exists uniquely.  With
\[
 \zeta=\sqrt{1-\lambda_{\max}(P)^{-1}}<1,\qquad
 B=\sqrt{\ell P^{-1}\ell^\dagger}\sqrt{v^\dagger Pv},\qquad
 C=p(1-p)+\frac{2B}{1-\zeta},
\]
one has
\begin{equation}
 |c_k|\le B\zeta^{k-1},\qquad
 \operatorname{Var}(\bar X_N)\le\frac CN,\qquad
 \Pr\{|\bar X_N-p|>\varepsilon\}
 \le\min\!\left(1,\frac C{N\varepsilon^2}\right).
 \label{p4:eq:v2211-computable-frequency-error}
\end{equation}
For \(\zeta=0\), the power at \(k=1\) is interpreted as one.  If
\(\mathcal V_1=0\), all lag covariances vanish and the same statements
hold with \(B=0,\ C=p(1-p)\).  All matrices and constants are descendants
of the actual instrument and preparation; \(P\) is a calculation
certificate, not a physical state or an added dynamics.
\end{corollary}
\begin{proof}
The covariance formula of the theorem gives \(c_k=\ell A^{k-1}v\).
Multiplying the finite geometric series by \((I-A)^2\) gives
\[
 \sum_{k=1}^{N-1}(N-k)A^{k-1}
 =[(N-1)I-NA+A^N](I-A)^{-2}.
\]
Substitution in the exact variance identity proves
\eqref{p4:eq:v2211-event-resolvent-variance}.  Since \(r(A)<1\),
finite Jordan calculus makes
\[
 P=\sum_{n=0}^{\infty}(A^\dagger)^nA^n
\]
norm convergent; it is at least \(I\) and solves the Lyapunov equation.
If two solutions differ by \(D=A^\dagger DA\), iteration and \(A^n\to0\)
give \(D=0\).  Moreover
\[
 \|Ax\|_P^2=\|x\|_P^2-\|x\|^2
 \le(1-\lambda_{\max}(P)^{-1})\|x\|_P^2.
\]
Dual Cauchy--Schwarz therefore gives
\(|\ell A^{k-1}v|\le B\zeta^{k-1}\).
Bounding \(N-k\le N\) in the variance identity and summing the geometric
series proves \(C/N\); Chebyshev gives the probability bound.  When
\(\zeta=0\), the contraction identity forces \(A=0\), so only \(c_1\)
can be nonzero.  The zero event-space case follows directly from the
earlier theorem.
\end{proof}

\begin{corollary}[Almost-sure frequencies from the actual finite event sector]
\label{p4:cor:v2211-event-frequency-strong-law}
On the domain of the preceding certificate, the successive stationary
instrument record obeys
\[
 \boxed{\bar X_N\longrightarrow p
       =\operatorname{Tr}\mathcal I_1(\sigma)
       =\operatorname{Tr}(\sigma E_1)
       \quad\hbox{almost surely}.}
\]
Thus a computed \(r(R|_{\mathcal V_1})<1\) closes the non-i.i.d.\
frequency law for that event on the actual instrument domain.  It does
not assert convergence for every individual outcome sequence or construct
a missing instrument/preparation.
\end{corollary}
\begin{proof}
Successive CP instrument maps and normalization of their sum give
consistent finite cylinder probabilities for the outcome records.
They define the corresponding probability measure on the countable
record space.  For any \(\varepsilon>0\), the preceding bound gives
\[
 \sum_{m=1}^{\infty}
 \Pr\{|\bar X_{m^2}-p|>\varepsilon\}
 \le\frac C{\varepsilon^2}\sum_{m=1}^{\infty}m^{-2}<\infty.
\]
The first Borel--Cantelli lemma, which requires no independence, implies
convergence along \(m^2\), simultaneously for the countable choices
\(\varepsilon=1/j\).  If \(m^2\le N<(m+1)^2\), the binary bound
\(0\le X_i\le1\) gives
\[
 |\bar X_N-\bar X_{m^2}|
 \le 2\,\frac{N-m^2}{N}
 \le2\,\frac{2m+1}{m^2}\longrightarrow0.
\]
Interpolation proves the full almost-sure limit.  The dual event effect
\(E_1=\mathcal I_1^\sharp(I)\) gives the final Born pairing.  The
argument promotes the finite event-sector variance certificate to an
almost-sure statement on the actual record measure, without a continuous
projective uniquely ergodic map.
\end{proof}

\paragraph{The sharp stationary event-frequency criterion}
\begin{theorem}[Conserved event obstruction and exact nonmixing frequency law]
\label{p4:thm:v255-sharp-event-frequency}
Retain the occurring finite instrument, its actual stationary preparation
$\sigma$, and the event notation $\Phi,p,v,R,\mathcal V_1,\ell$ of
Theorem~\ref{p4:thm:v2210-lit-instrument-mixing}, where
$\ell(X)=\operatorname{Tr}\mathcal I_1(X)$. No spectral-radius
hypothesis is imposed. If $\mathcal V_1\ne0$, set
\[
 A=R|_{\mathcal V_1}.
\]
The eigenvalue $1$ of $A$, if present, is semisimple. Let $Q$ be its
spectral projector, with $Q=0$ if $1$ is absent, and define the finite
calculation operators and scalar
\begin{equation}
 S=(I-A+Q)^{-1}(I-Q),\qquad
 d_1=\ell(Qv),\qquad t_1=\|S^2v\|_1.
 \label{p4:eq:v255-event-group-inverse}
\end{equation}
Here the trace norm is that of the represented matrix in the actual event
space. Then $d_1$ is real and
\begin{equation}
 0\le d_1\le p(1-p),\qquad
 \lim_{N\to\infty}\operatorname{Var}(\bar X_N)=d_1.
 \label{p4:eq:v255-event-variance-floor}
\end{equation}
For every $N\ge1$ the exact variance is
\begin{equation}
 \boxed{
 \operatorname{Var}(\bar X_N)
 =d_1+\frac{p(1-p)-d_1}{N}
 +\frac2{N^2}\ell\!\left(
       [(N-1)I-NA+A^N]S^2v\right).}
 \label{p4:eq:v255-exact-event-variance}
\end{equation}
In particular,
\begin{equation}
 \left|\operatorname{Var}(\bar X_N)-d_1\right|
 \le\frac{p(1-p)-d_1+4t_1}{N}.
 \label{p4:eq:v255-event-variance-budget}
\end{equation}
The sharp deterministic Born-frequency criterion on this stationary
instrument domain is
\begin{equation}
 \boxed{
 \bar X_N\longrightarrow p\ \hbox{almost surely}
 \quad\Longleftrightarrow\quad d_1=0.}
 \label{p4:eq:v255-sharp-born-frequency-iff}
\end{equation}
When $d_1=0$, the fully finite constant
$C_1=p(1-p)+4t_1$ gives
\begin{equation}
 \operatorname{Var}(\bar X_N)\le\frac{C_1}{N},\qquad
 \Pr\{|\bar X_N-p|>\varepsilon\}
 \le\min\!\left(1,\frac{C_1}{N\varepsilon^2}\right).
 \label{p4:eq:v255-nonmixing-frequency-error}
\end{equation}
If $\mathcal V_1=0$, use $d_1=t_1=0$ and omit the empty-space matrices;
the same scalar conclusions hold. Nontrivial unit-circle eigenvalues
of $A$ are allowed. Every object in this certificate is calculated from
the actual instrument and preparation; $Q$ and $S$ add no physical
state, dynamics, or selector.
\end{theorem}
\begin{proof}
First, $v$ is Hermitian and traceless. Since $R$ preserves Hermiticity
and trace-zero matrices, $\mathcal V_1$ is closed under adjoints and
consists of traceless matrices. For $X\in\mathcal V_1$ one therefore has
\[
 A^nX=\Phi^n(X),\qquad n\ge0.
\]
A positive trace-preserving map contracts the trace norm of a Hermitian
matrix: write $X=X_+-X_-$ with $X_\pm\ge0$, use positivity of the two
images, and bound
\[
 \|\Phi^n(X)\|_1
 \le\operatorname{Tr}\Phi^n(X_+)
      +\operatorname{Tr}\Phi^n(X_-)=\|X\|_1.
\]
For a general matrix write $X=H+iK$ with $H,K$ Hermitian. The inequalities
$\|H\|_1,\|K\|_1\le\|X\|_1$ give
$\|A^nX\|_1\le2\|X\|_1$ on $\mathcal V_1$.
Thus $A$ is power bounded. Finite Jordan calculus consequently gives
$\operatorname{spec}A\subset\{|z|\le1\}$ and no nontrivial Jordan
block at any unit-circle eigenvalue. In particular $1$ is semisimple.
Its projector can be calculated either by finite spectral algebra or by
\[
 Q=\lim_{M\to\infty}\frac1M\sum_{n=0}^{M-1}A^n,
 \qquad
 \mathcal V_1=\ker(I-A)\oplus\operatorname{ran}(I-A).
\]
The projector preserves Hermiticity, commutes with $A$, and obeys
$AQ=Q$. The displayed direct sum proves invertibility of $I-A+Q$ and
\[
 SQ=QS=0,\qquad (I-A)S=S(I-A)=I-Q.
\]
These identities also show that $S$ preserves Hermiticity: $I-A+Q$ is
an invertible complex-linear operator commuting with the adjoint
involution, and its inverse has the same property. In particular
$S^2v$ is Hermitian and traceless.

On $\ker(I-A)$ a triangular sum has value $N(N-1)/2$, while on
$\operatorname{ran}(I-A)$ the ordinary finite geometric-series identity
applies. Together they give, for every $N\ge1$,
\[
 \sum_{k=1}^{N-1}(N-k)A^{k-1}
 =\frac{N(N-1)}2Q+
       [(N-1)I-NA+A^N]S^2.
\]
For $N=1$ the sum is empty and both sides vanish. Substitute this
identity and the already proved covariance
$c_k=\ell(A^{k-1}v)$ into the exact stationary variance formula. This
proves \eqref{p4:eq:v255-exact-event-variance}.

The event effect is $0\le E_1\le I$, so
$|\ell(Y)|=|\operatorname{Tr}(E_1Y)|\le\|Y\|_1$. Apply the Hermitian
trace-norm contraction above to $Y=S^2v$. For all $n\ge0$,
\[
 |\ell(A^nS^2v)|\le t_1.
\]
The absolute value of the last term in
\eqref{p4:eq:v255-exact-event-variance} is therefore at most
\[
 \frac2{N^2}\bigl((N-1)+N+1\bigr)t_1=\frac{4t_1}{N}.
\]
It follows that the variance tends to the real number $d_1$.
Nonnegativity of variance gives $d_1\ge0$. Jensen's inequality applied
to the scalar record gives
\[
 \mathbb E\left|\frac1N\sum_{i=1}^N(X_i-p)\right|^2
 \le\frac1N\sum_{i=1}^N\mathbb E|X_i-p|^2=p(1-p),
\]
so $d_1\le p(1-p)$. These facts prove
\eqref{p4:eq:v255-event-variance-floor} and
\eqref{p4:eq:v255-event-variance-budget}.

If $d_1=0$, the variance budget yields \eqref{p4:eq:v255-nonmixing-frequency-error}
by Chebyshev's inequality. Successive instrument maps define consistent
finite record probabilities and hence their probability measure. For each
$\varepsilon>0$,
\[
 \sum_{m=1}^\infty
 \Pr\{|\bar X_{m^2}-p|>\varepsilon\}
 \le\frac{C_1}{\varepsilon^2}\sum_{m=1}^\infty m^{-2}<\infty.
\]
The first Borel--Cantelli lemma proves convergence along $m^2$,
simultaneously for $\varepsilon=1/j$, $j\ge1$. For
$m^2\le N<(m+1)^2$, binary boundedness gives
\[
 |\bar X_N-\bar X_{m^2}|
 \le2\frac{N-m^2}{N}\longrightarrow0,
\]
so convergence holds along the whole sequence. Conversely, if
$\bar X_N\to p$ almost surely, bounded convergence applied to
$|\bar X_N-p|^2\le1$ implies
$\operatorname{Var}(\bar X_N)\to0$. The variance limit already proved
then forces $d_1=0$. This proves the equivalence. If
$\mathcal V_1=0$, the prior covariance formula gives $c_k=0$ directly,
and the same scalar variance and almost-sure arguments apply.
\end{proof}

\begin{proposition}[Alternating records and the exact persistent obstruction]
\label{p4:prop:v255-periodic-persistent-records}
On $\mathbb C^2$, put $P_j=|j\rangle\langle j|$ and consider the
alternating instrument
\[
 \mathcal I_0(X)=|1\rangle\langle0|X|0\rangle\langle1|,
 \qquad
 \mathcal I_1(X)=|0\rangle\langle1|X|1\rangle\langle0|,
 \qquad \sigma=I/2.
\]
Its binary record alternates, $p=1/2$, and its nonzero event-cyclic
space has $A=-I$. Thus $r(A)=1$, while $Q=0$ and $d_1=0$. Exactly,
\[
 c_k=\frac{(-1)^k}{4},\qquad
 \operatorname{Var}(\bar X_N)
 =\frac{1-(-1)^N}{8N^2}.
\]
Consequently its frequencies converge to the Born value despite the
absence of event-sector mixing.
For the persistent instrument
$\mathcal J_j(X)=P_jXP_j$ and stationary preparation
$\sigma=(1-p)P_0+pP_1$, $0<p<1$, instead
\[
 A=I,\qquad Q=I,\qquad d_1=p(1-p),\qquad S=0,
 \qquad\operatorname{Var}(\bar X_N)=p(1-p).
\]
It therefore fails the deterministic Born-frequency law on exactly the
obstruction identified by Theorem~\ref{p4:thm:v255-sharp-event-frequency}.
\end{proposition}
\begin{proof}
For the alternating instrument,
$\mathcal I_1(\sigma)=P_0/2$ and
$v=(P_0-P_1)/4$. The total channel interchanges $P_0,P_1$, so
$Av=-v$, while $\ell(v)=\operatorname{Tr}(P_1v)=-1/4$.
This gives the stated covariance, and substitution in the exact variance
formula gives the parity expression. Equivalently, every even record
has mean $1/2$, and an odd record has either of the two means
$1/2\pm1/(2N)$, each with probability $1/2$.
For the persistent instrument,
$v=p(1-p)(P_1-P_0)$ is fixed by the total channel and
$\ell(v)=p(1-p)$. The projector and group-inverse formulas give the
claimed values. Every record is constant after its first outcome,
which agrees with the same nonvanishing variance and the sharp criterion.
\end{proof}

\noindent\textit{Claim boundary.}
The earlier strict-contraction certificate remains a sufficient special
case: $r(A)<1$ implies $Q=0$. The present calculation removes that
unnecessary hypothesis and replaces it by the exact scalar condition
$d_1=0$ for the recorded event, allowing periodic modes and invisible
conserved directions. It is a complete deterministic frequency test for
an occurring finite stationary instrument. It does not construct a
missing apparatus, select its preparation, or force $d_1=0$ for every
physical instrument. Stationarity is retained throughout. The
persistent example proves why no unconditional universal frequency
assertion follows from finite carrier geometry alone.


\paragraph{Arbitrary preparations and the finite asymptotic frequency observable}
\begin{theorem}[A nonstationary finite instrument has a computable frequency law]
\label{p4:thm:v258-nonstationary-frequency-law}
Let $\{\mathcal I_i\}$ be an occurring, time-homogeneous finite CP
instrument on its constituted matrix algebra, and put
$\Phi=\sum_i\mathcal I_i$. Assume $\Phi$ is trace preserving. Its
initial preparation $\rho$ is any density matrix; stationarity,
primitivity and mixing are not assumed. For one recorded event use
$\mathcal I_1$ for its outcome operation, and write
\[
 E=\mathcal I_1^\sharp(I),\qquad
 \mathcal C=\lim_{M\to\infty}\frac1M\sum_{n=0}^{M-1}\Phi^n,
 \qquad
 \mathcal G=(I-\Phi+\mathcal C)^{-1}(I-\mathcal C),
 \qquad \mathcal K=\mathcal C\mathcal I_1\mathcal C.
\]
The finite projector $\mathcal C$ and inverse $\mathcal G$ exist.
$\mathcal C$ is CPTP and $\mathcal K$ is CP and trace nonincreasing.
On the actual successive-instrument record measure,
\begin{equation}
 \boxed{\bar X_N\longrightarrow Z\quad\hbox{almost surely},\qquad
 0\le Z\le1,\qquad
 \mathbb E_\rho Z^r=\operatorname{Tr}(\mathcal K^r\rho)
 \quad(r\ge0).}
 \label{p4:eq:v258-all-frequency-moments}
\end{equation}
Here $r$ is an integer and $\mathcal K^0=I$. Thus the full limiting event-frequency law is
calculated from the actual instrument and arbitrary preparation, with no
stationary-state selection. In particular, put
\begin{equation}
 m_1=\operatorname{Tr}(\mathcal K\rho),\qquad
 m_2=\operatorname{Tr}(\mathcal K^2\rho).
 \label{p4:eq:v258-frequency-first-two}
\end{equation}
For any specified $c\in[0,1]$,
\begin{equation}
 \boxed{\bar X_N\longrightarrow c\ \hbox{almost surely}
 \quad\Longleftrightarrow\quad m_2-2cm_1+c^2=0.}
 \label{p4:eq:v258-deterministic-frequency-iff}
\end{equation}
In particular the frequency is deterministic with value $m_1$ exactly
when $m_2-m_1^2=0$.

There is also a finite raw-moment error certificate. In the trace norm
on Hermitian matrices let $g\ge\|\mathcal G\|_{1\to1,\mathrm{Herm}}$
be any certified finite bound. For fixed integer $r\ge1$ and $N\ge r$,
\begin{equation}
 \left|\mathbb E_\rho\bar X_N^r-
                 \operatorname{Tr}(\mathcal K^r\rho)\right|
 \le\frac{r(r-1)+2r^2g}{N}.
 \label{p4:eq:v258-frequency-moment-budget}
\end{equation}
The constants are calculation certificates, not added dynamics.
\end{theorem}
\begin{proof}
A positive trace-preserving map contracts the trace norm of Hermitian
matrices: decompose $Y=Y_+-Y_-$ with $Y_\pm\ge0$, apply positivity,
and use preservation of their traces. Complexification gives
$\|\Phi^nY\|_1\le2\|Y\|_1$ for all matrices, uniformly in $n$.
Finite Jordan calculus therefore confines the spectrum to the closed
unit disk and makes every unit-circle eigenvalue semisimple. In
particular the Ces\`aro limit $\mathcal C$ exists, is the spectral
projector at $1$, and satisfies
\[
 \Phi\mathcal C=\mathcal C\Phi=\mathcal C,\qquad
 \mathcal C^2=\mathcal C.
\]
It is CPTP as a norm limit of finite CPTP averages. The direct sum of the
$1$-eigenspace and the range of $I-\Phi$ proves invertibility of
$I-\Phi+\mathcal C$. The resulting group inverse obeys
\begin{equation}
 (I-\Phi)\mathcal G=\mathcal G(I-\Phi)=I-\mathcal C,
 \qquad
 \sum_{n=0}^{M-1}(\Phi^n-\mathcal C)
 =(I-\Phi^M)\mathcal G.
 \label{p4:eq:v258-bounded-channel-partial-sums}
\end{equation}
In particular the sum on the right has induced Hermitian trace norm at
most $2g$. Complete positivity and the trace assertions for $\mathcal K$
follow by composition.

First prove existence of the almost-sure limit. Let
$T=\Phi^\sharp$, $e=\mathcal C^\sharp E$, and
$h=\mathcal G^\sharp E$. All these observable insertions are Hermitian,
$0\le e\le I$, and
\[
 Te=e,\qquad E=e+(I-T)h.
\]
Let $\rho_n$ denote the conditional state after the first $n$ recorded
outcomes. The instrument cylinder probabilities and their normalized CP
updates give
\[
 \mathbb E[\rho_n\mid\mathcal F_{n-1}]=\Phi(\rho_{n-1}),
 \qquad
 \mathbb E[X_n\mid\mathcal F_{n-1}]
       =\operatorname{Tr}(E\rho_{n-1}).
\]
Assignments on null-probability branches are immaterial. It follows that
$z_n=\operatorname{Tr}(e\rho_n)$ is a bounded martingale. The bounded
martingale convergence theorem gives $z_n\to Z$ almost surely, with
$Z\in[0,1]$. Define the real martingale differences
\[
 \xi_n=X_n-\operatorname{Tr}(E\rho_{n-1})
       +\operatorname{Tr}(h\rho_n)
       -\operatorname{Tr}(Th\rho_{n-1}).
\]
They have conditional mean zero and
$|\xi_n|\le1+2\|h\|$. With $M_N=\sum_{n=1}^N\xi_n$, the Poisson
identity for $h$ gives the exact relation
\begin{equation}
 \sum_{n=1}^N(X_n-z_{n-1})
 =M_N+\operatorname{Tr}(h\rho_0)-\operatorname{Tr}(h\rho_N).
 \label{p4:eq:v258-frequency-martingale-decomposition}
\end{equation}
Orthogonality of martingale increments gives
$\mathbb E M_N^2\le(1+2\|h\|)^2N$. Chebyshev and the first
Borel--Cantelli lemma prove $M_{m^2}/m^2\to0$ almost surely; bounded
increments interpolate between successive squares to give
$M_N/N\to0$. The boundary term divided by $N$ also vanishes.
Ces\`aro convergence of the bounded convergent sequence $z_n$ now proves
$\bar X_N\to Z$ almost surely.

Next compute the law through all its moments. For ordered distinct
event times $1\le n_1<\cdots<n_r\le N$, set
$u_0=n_1-1$ and $u_j=n_{j+1}-n_j-1$ for $1\le j<r$. Their joint
probability is
\[
 \operatorname{Tr}\!\left(
 \mathcal I_1\Phi^{u_{r-1}}\mathcal I_1\cdots
        \mathcal I_1\Phi^{u_0}\rho\right),
 \qquad u_j\ge0,\quad\sum_{j=0}^{r-1}u_j\le N-r.
\]
Replace each of the $r$ powers of $\Phi$ by $\mathcal C$, telescoping
one factor at a time. In each difference, fix all other gap indices and
sum the single factor $\Phi^{u_j}-\mathcal C$ over its allowed initial
integer interval. Equation~\eqref{p4:eq:v258-bounded-channel-partial-sums}
bounds that summed factor by $2g$ on Hermitian matrices. Every other
map in the telescoping term is $\Phi^u$, $\mathcal C$, or
$\mathcal I_1$, and contracts the Hermitian trace norm. The final trace
has norm at most one. There are
$\binom{N-1}{r-1}$ allowed choices of the other $r-1$ gaps. Thus the
sum over all ordered distinct times differs from
\[
 \binom Nr\operatorname{Tr}(\mathcal K^r\rho)
\]
by at most $2rg\binom{N-1}{r-1}$. The equality of the all-$\mathcal C$
term with the displayed trace uses trace preservation of the leftmost
$\mathcal C$.

Expand $\bar X_N^r$ into $N^r$ ordered index tuples. Distinct tuples
are the preceding sum multiplied by $r!$. At most
$\binom r2N^{r-1}$ tuples have a repeated index, by the union bound
over equal pairs. Their products lie in $[0,1]$. The same counting
bound gives
$|1-r!\binom Nr/N^r|\le\binom r2/N$.
Since $0\le\operatorname{Tr}(\mathcal K^r\rho)\le1$ and
$r!\binom{N-1}{r-1}\le rN^{r-1}$, these estimates prove
\eqref{p4:eq:v258-frequency-moment-budget}. Bounded convergence from
$\bar X_N\to Z$ proves \eqref{p4:eq:v258-all-frequency-moments}.
Finally,
$\mathbb E(Z-c)^2=m_2-2cm_1+c^2$ vanishes exactly when $Z=c$ almost
surely, proving \eqref{p4:eq:v258-deterministic-frequency-iff}.
\end{proof}

\begin{corollary}[An exact finite POVM for the asymptotic frequency]
\label{p4:cor:v258-asymptotic-frequency-povm}
Represent $\mathcal K$ in any real basis of the Hermitian matrix space,
and let $p_{\mathcal K}$ be its real characteristic polynomial. Let
$\lambda_1,\ldots,\lambda_s$ be all its distinct real roots in $[0,1]$.
There is at least one. Define their Lagrange polynomials and effects by
\begin{equation}
 L_j(x)=\prod_{k\ne j}\frac{x-\lambda_k}{\lambda_j-\lambda_k},
 \qquad
 F_j=L_j(\mathcal K)^\sharp(I).
 \label{p4:eq:v258-frequency-povm}
\end{equation}
Then
\[
 F_j\ge0,\qquad \sum_jF_j=I,\qquad
 \Phi^\sharp(F_j)=F_j,
 \qquad
 \boxed{\Pr_\rho\{Z=\lambda_j\}=\operatorname{Tr}(\rho F_j).}
\]
Zero effects may be discarded. Thus the limiting event frequency has a
finite atomic law, with rates and weights obtained from the actual
finite instrument. The conserved mean effect is
$e=\sum_j\lambda_jF_j=\mathcal C^\sharp E$. The effects need not be
orthogonal projectors on transient initial supports.
\end{corollary}
\begin{proof}
Cayley--Hamilton and the moment formula give, for every integer $a\ge0$,
\[
 \mathbb E_\rho[Z^ap_{\mathcal K}(Z)]
 =\operatorname{Tr}\!\left(\mathcal K^a
                              p_{\mathcal K}(\mathcal K)\rho\right)=0.
\]
Take the real linear combination given by the coefficients of
$p_{\mathcal K}$ to obtain
$\mathbb E_\rho p_{\mathcal K}(Z)^2=0$. Hence $Z$ lies among the listed
real roots almost surely. Its law has total mass one, so the list is
nonempty. Lagrange interpolation on that finite support gives
\[
 \Pr_\rho\{Z=\lambda_j\}
 =\mathbb E_\rho L_j(Z)
 =\operatorname{Tr}(L_j(\mathcal K)\rho)
 =\operatorname{Tr}(\rho F_j).
\]
This holds for every initial density matrix, proving positivity of
$F_j$; the same equality summed over $j$ proves $\sum_jF_j=I$.
The operator $I$ is fixed by $\Phi^\sharp$, and every positive power
$(\mathcal K^\sharp)^aI$, $a\ge1$, lies in the range of
$\mathcal C^\sharp$ because
$\mathcal K^\sharp=\mathcal C^\sharp\mathcal I_1^\sharp
\mathcal C^\sharp$. Thus every $F_j$ is fixed. Taking the first moment
of the finite law proves the mean-effect identity.
\end{proof}

\begin{corollary}[Stationary recovery and the initial-event comparison]
\label{p4:cor:v258-stationary-frequency-recovery}
If the actual preparation is stationary, $\Phi(\sigma)=\sigma$,
then $m_1=p=\operatorname{Tr}\mathcal I_1(\sigma)$ and
\[
 m_2-p^2
 =\operatorname{Tr}\mathcal I_1\mathcal C
              (\mathcal I_1(\sigma)-p\sigma).
\]
This is exactly the conserved event obstruction of
Theorem~\ref{p4:thm:v255-sharp-event-frequency}. The new theorem therefore
retains that sharp criterion while removing stationarity from the
existence and computation of the limiting frequency law.
For arbitrary initial $\rho$, write
$p_0=\operatorname{Tr}(\rho E)$. Frequencies converge to this initial
one-shot Born value precisely when
$m_2-2p_0m_1+p_0^2=0$; it is not automatic.
\end{corollary}
\begin{proof}
Stationarity gives $\mathcal C\sigma=\sigma$ and the displayed first
two moments. On the traceless event-cyclic space the full-channel
projector $\mathcal C$ restricts to the event projector at eigenvalue
one, giving the earlier obstruction. The arbitrary-preparation statement
is the deterministic criterion with $c=p_0$.
\end{proof}


\begin{proposition}[Periodic, persistent, transient and absorption diagnostics]
\label{p4:prop:v258-frequency-law-diagnostics}
The following explicit finite instruments give four exact comparison cases of
Theorem~\ref{p4:thm:v258-nonstationary-frequency-law}.
\begin{enumerate}
\item The alternating qubit instrument of
Proposition~\ref{p4:prop:v255-periodic-persistent-records}, with either
basis-state initial preparation, has the law $Z=1/2$ almost surely,
although its initial one-shot event probability is either zero or one.
\item The persistent sharp instrument $\mathcal I_j(X)=P_jXP_j$ and
$\rho=(1-q)P_0+qP_1$ has law
$(1-q)\delta_0+q\delta_1$, not a deterministic rate $q$ for $0<q<1$.
\item For the amplitude-damping instrument
\[
 M_0=P_0+\sqrt{1-\gamma}\,P_1,\qquad
 M_1=\sqrt\gamma\,|0\rangle\langle1|,\qquad
 \mathcal I_j(X)=M_jXM_j^*,\qquad 0<\gamma\le1,
\]
initial preparation $\rho=P_1$ has initial event probability $\gamma$,
while $Z=0$ almost surely. Every record contains at most one event.
\item A three-state measure-and-prepare instrument on basis $A,B,U$
can be specified by nonnegative transition matrices
\[
 P=\begin{pmatrix}1&0&q\\0&1&1-q\\0&0&0\end{pmatrix},\qquad
 J_1=\begin{pmatrix}1&0&q\\0&0&0\\0&0&0\end{pmatrix},\qquad
 J_0=P-J_1,
 \qquad 0<q<1,
\]
using
$\mathcal I_j(X)=\sum_{a,b}(J_j)_{ab}
 \langle b|X|b\rangle|a\rangle\langle a|$.
These are CP outcome maps with Kraus operators
$\sqrt{(J_j)_{ab}}|a\rangle\langle b|$ and trace-preserving sum.
For $\rho=|U\rangle\langle U|$, the initial state is nonstationary,
yet the frequency law is exactly
$(1-q)\delta_0+q\delta_1$. The harmonic frequency effects are
\[
 F_1=\operatorname{diag}(1,0,q),\qquad
 F_0=\operatorname{diag}(0,1,1-q).
\]
In particular the transient entry of the mean effect is $q$, while its
second frequency moment is $q$, rather than $q^2$.
\end{enumerate}
\end{proposition}
\begin{proof}
For the alternating instrument, on diagonal states
$\mathcal C$ averages the two coordinates, and
$\operatorname{Tr}(\mathcal K^r\rho)=2^{-r}$ for every $r\ge1$.
The moment theorem gives the asserted point mass. For the persistent
instrument $\mathcal K$ is its event projection, so all its positive
moments equal $q$, giving the displayed Bernoulli law.
For amplitude damping, the Ces\`aro state projection is
$\mathcal C(X)=\operatorname{Tr}(X)P_0$. Since
$\mathcal I_1(P_0)=0$, $\mathcal K=0$ and every positive frequency
moment is zero. The event leaves state $P_0$, where it cannot recur,
which also proves the pathwise assertion.
For the last instrument $P^2=P$ on diagonal states, so its full-channel
Ces\`aro projection restricts to $P$. Direct multiplication gives
$PJ_1P=J_1$ and $J_1^2=J_1$. Thus all positive frequency moments of the
initial $U$ state equal $q$. Interpolation on rates $0,1$ gives the two
stated effects, which satisfy $\Phi^\sharp(F_j)=F_j$. They encode the
transient state's actual absorption probabilities without selecting
one absorbing state in advance.
\end{proof}

\noindent\textit{Claim boundary.}
The result is unconditional for the repeated, time-homogeneous finite
instrument calculus once its actual instrument and preparation occur.
It adds no stationarity, mixing, reservoir, uniquely ergodic hit map, or
pure-state selector. It computes a frequency observable; it does not
choose the apparatus, force a deterministic rate, identify that rate
with the initial one-shot probability, or supply a time-dependent
instrument law that has not been constituted.


\begin{proposition}[Two exact nonmixing diagnostics]
\label{p4:prop:v2210-lit-frequency-obstructions}
For a persistent sharp nondemolition instrument
$\mathcal I_1(X)=P_1XP_1$ on a stationary state commuting with $P_1$,
all successive binary records coincide and
$c_k=p(1-p)$, $\operatorname{Var}(\bar X_N)=p(1-p)$.
Moreover, for $d\ge2$, a fixed unitary acting on $\mathbb P(\mathbb C^d)$
cannot be uniquely ergodic with Fubini--Study measure as its unique invariant
probability measure.
\end{proposition}
\begin{proof}
After either sharp outcome the subsequent identical nondemolition tests
repeat it, giving the joint probability $p$ and the stated covariance.
A finite unitary has an eigenvector; its projective eigenray is fixed and
supports an invariant point measure distinct from Fubini--Study volume.
\end{proof}
Corollary~\ref{p4:cor:unique-ergodic-born-frequencies} proves the stronger
topological obstruction for every continuous self-map on nontrivial
projective space, so that former Fubini--Study unique-ergodic route is
retired.  The actual finite instrument route now has the computable
resolvent/Lyapunov sufficient certificate and the sharp nonmixing criterion of
Theorem~\ref{p4:thm:v255-sharp-event-frequency}.
Open or apparatus-driven trajectory results
\cite{LitQFTBenoistTrajectories2019} are possible implementations, but their
irreducibility/purification conditions and invariant law must be checked on
the actual instrument; that law is not automatically Fubini--Study.



\paragraph{What a finite-window Markov comparison certifies}
\begin{proposition}[State and semigroup-defect error budgets]
\label{p4:prop:v2210-lit-markov-window}
Let $\Psi_t$ be the exact reduced CPTP maps of a supplied native finite
system/environment process, and let $\Xi_t$ be a supplied CPTP semigroup
comparison on the same carrier. If
$\|\Psi_t-\Xi_t\|_{1\to1}\le\epsilon(t)$ on a declared time window,
then every normalized preparation and event effect obey
\[
 |\operatorname{Tr}E(\Psi_t-\Xi_t)(\rho)|\le\epsilon(t),\qquad
 \|\Psi_{t+s}-\Psi_t\Psi_s\|_{1\to1}
 \le\epsilon(t+s)+\epsilon(t)+\epsilon(s)
\]
whenever $t,s,t+s$ belong to that window.
\end{proposition}
\begin{proof}
Use $\|E\|\le1$, $\|\rho\|_1=1$ and trace duality. For the second
bound insert $\Xi_{t+s}=\Xi_t\Xi_s$ and telescope the product. CPTP
maps have trace-norm operator norm one, by duality with unital CP maps.
\end{proof}
This certifies a state/readout approximation, not arbitrary multi-time
instrument statistics or an all-time irreversible finite-bath law.
The existing finite-bath rigidity and reservoir theorem domains remain in
force. Davies' weak-coupling and the repeated-fresh-probe construction
\cite{Davies1974,LitQFTAttalPautrat2006} can supply comparison maps only
after their actual correlation, preparation and scaling hypotheses are
proved. An equilibrium KMS or Gibbs label is not a cosmological kinetic
history.



\subsection{Ground-source compression, dynamical response and the native UV ledger}
\label{p4:sec:v2216-source-polarization-uv}

The positive excitation-family compression theorem concerns a covered
source mean at its ground. It does not set the time-dependent metric
polarization to zero. The following exact calculation makes this distinction
on the same finite carrier, before taking a local or continuum limit.

\begin{theorem}[Zero excitation mean and nonzero dynamical polarization]
\label{p4:thm:v2216-zero-excitation-dynamic-response}
Use units $\hbar=1$. On an actual positive-gap Pauli pair let
\[
 q=\log a,\qquad h(q)=\ell\sigma_3+x\sigma_1,\qquad
 x=am,\qquad \epsilon=(\ell^2+x^2)^{1/2}>0,
 \qquad H_{\rm exc}=h+\epsilon I.
\]
Here $m$ and $\ell$ are fixed under this variation. The mathematical
excitation operator $H_{\rm exc}\succeq0$ has identically zero ground
eigenvalue. Its ground projection $P$ obeys
$P(\partial_qH_{\rm exc})P=0$. Nevertheless its exact stationary-ground
source-action response, including the derivatives of the same $\epsilon I$,
is
\begin{equation}
\label{p4:eq:v2216-excitation-retarded-kernel}
 \boxed{K_{\rm exc}^{\rm action}(\omega)
 =\frac{b^2\omega^2}
 {\epsilon[4\epsilon^2-(\omega+i0)^2]},\qquad
 b^2=\frac{x^2\ell^2}{\epsilon^2}.}
\end{equation}
The numerator is understood as the retarded boundary value of
$(\omega+i0)^2$. In particular $K_{\rm exc}^{\rm action}(0)=0$, whereas
its subgap $\omega^4$ coefficient is
\[
 \boxed{\Sigma_4=\frac{b^2}{16\epsilon^5}
 =\frac{a^2m^2\ell^2}{16(\ell^2+a^2m^2)^{7/2}}\ge0.}
\]
It is strictly positive when $m\ell\ne0$. More generally, on an actual
stationary ground with positive excitation gaps $\Delta_r$ and complete
noncentral insertion $B$, the source-memory contribution satisfies
\[
 K_{\rm mem}^{\rm action}(\omega)
 =2\sum_r\frac{\Delta_r|\langle r|B|0\rangle|^2}
 {\Delta_r^2-(\omega+i0)^2},\qquad
 \Sigma_4=2\sum_r\frac{|\langle r|B|0\rangle|^2}{\Delta_r^5}.
\]
Independent tensor-ground sectors add their nonnegative weights. A central
instantaneous allocation has no transition matrix elements and does not
cancel this coefficient.
\end{theorem}
\begin{proof}
The ground expectation of $\partial_qh=x\sigma_1$ is
$-x^2/\epsilon$, while $\partial_q\epsilon=x^2/\epsilon$.
Its off-diagonal squared matrix element is $b^2$. Differentiating twice
gives
\[
 \langle0|\partial_q^2H_{\rm exc}|0\rangle
 =-\frac{x^2}{\epsilon}
   +\frac{2x^2}{\epsilon}-\frac{x^4}{\epsilon^3}
 =\frac{b^2}{\epsilon}.
\]
The Lorentzian source action has negative Hamiltonian contact, so the full
kernel is
\[
 -\frac{b^2}{\epsilon}
 +\frac{4\epsilon b^2}{4\epsilon^2-(\omega+i0)^2}.
\]
Combining the fractions proves
Eq.~\eqref{p4:eq:v2216-excitation-retarded-kernel}; expanding below the
gap gives the stated $\Sigma_4$. An energy resolution in the retarded
force commutator proves the general spectral formula, with the source-action
sign as displayed. On a tensor ground, insertions in distinct factors excite
orthogonal subspaces, even when gaps coincide, so the weights add. A central
matrix commutes with every state and contributes only its instantaneous
metric contacts.
\end{proof}

\begin{remark}[QROOT scope, fixed preparation and absolute action data]
\label{p4:rem:v2216-qroot-dynamic-scope}
Theorem~\ref{p4:thm:v2114-positive-family-ground-source} and the present
theorem are compatible: compression of first derivatives does not compress
all excited intermediate states in a second response. QROOT selects the
native primordial scalar/tensor excitation ground once at the fresh root;
the resulting state is then transported by its actual generator. It neither
prepares every auxiliary Dirac/matter band nor resets evolving sources to
their instantaneous ground. The replacement $h\mapsto h+\epsilon I$ in
the theorem is a stated mathematical excitation normalization, not a new
absolute gravitational vacuum prescription. Central shifts leave excitation
gaps, ground eigenspaces and transition weights unchanged, while their metric
means and contact derivatives change. QROOT alone consequently cannot
determine an absolute cosmological or Einstein coefficient. No independent
native cyclic subtraction is added here. Keeping matter explicit is
compatible with an affine terminal gravitational constituent; integrating
it out retains its actual nonlocal polarization and does not by itself
create another independent gravitational Legendre coordinate.
\end{remark}

\begin{theorem}[Native parity-even graded matter UV coefficients]
\label{p4:thm:v2216-graded-matter-uv}
Consider the already constituted SSQFT smooth-parent quadratic owner on its
common short-proper-time domain: unbroken gauge symmetry, zero fermion
background, the generated four-real-coordinate Higgs doublet, the generated
$48$ internal Weyl coordinates, and $12$ gauge generators. Include the
actual background-gauge quotient Jacobian. At vanishing background gauge
curvature the relevant operators are
\[
 \Delta_H=-\nabla^2+R/6-2\Lambda_g^2,\qquad
 (\Delta_1)_\mu{}^\nu=-\nabla^2\delta_\mu{}^\nu+R_\mu{}^\nu,
 \qquad \Delta_0=-\nabla^2.
\]
For the Weyl block use its least self-adjoint Dirac-square completion,
$\Delta_D=-\nabla^2+R/4$ in its purely geometric part. The parity-even
Gaussian/Berezin weights on bundle ranks $(4,48,12,192)$ are respectively
$(1/2,1/2,-1,-1/4)$. After removing the common $(4\pi)^{-2}$ factor,
the weighted local coefficients are
\begin{equation}
\label{p4:eq:v2216-graded-matter-table}
\begin{array}{c|r|r|r}
 \text{block}&\text{rank}&a_0^{\rm weighted}&a_{2,\rm geom}^{\rm weighted}
 \\\hline
 H&4&2&0\\
 A_\mu&48&24&-2R\\
 \text{quotient Jacobian}&12&-12&-2R\\
 \text{Weyl via Dirac square}&192&-48&4R\\\hline
 \text{sum}&&-34&0
\end{array}
\end{equation}
The actual Higgs mass slot additionally contributes $4\Lambda_g^2$ to
weighted $a_2$. Thus, when no further mass slots occur,
\[
 \boxed{\mathcal K(s)\sim(4\pi s)^{-2}\int\sqrt g\,
 [-34+4\Lambda_g^2s+O(s^2)].}
\]
The geometric Einstein cancellation does not cancel the volume coefficient.
These are interior local smooth-parent UV identities, rather than identities of
arbitrary finite-window supertraces or a computed full physical vacuum.
\end{theorem}
\begin{proof}
Paper~V-A's component composition gives
$48=6+18+3+3+9+9$ internal Weyl coordinates; its coordinate-Higgs
theorem gives four real scalars, and the gauge algebra has dimension
$8+3+1=12$. Its least Clifford-stable self-adjoint carrier has rank
$4\cdot48=192$. The complementary analytic chirality is not an extra
physical family, and no independent PCT/conjugate carrier is added.
The real Gaussian exponent is $1/2$ and the paired quotient Berezin
exponent is $-1$. The parity-even Weyl exponent is $-1/2$ on the
rank-two chiral square, equivalently $-1/4$ on its rank-four Dirac
completion. These are the determinant powers of the existing bridge, not
new adjustable loop multiplicities.

For $\Delta=-(\nabla^2+E)$, the local densities are
$\operatorname{tr}1$ and $\operatorname{tr}(E+R/6)$.
The vector trace gives $-12R+48R/6=-4R$, before its $1/2$
weight. The quotient scalar gives $12R/6$, with weight $-1$.
The Dirac square gives $192(-R/4+R/6)=-16R$, with weight
$-1/4$. The Higgs conformal term cancels its geometric $R/6$,
leaving $4(2\Lambda_g^2)/2=4\Lambda_g^2$.
The rank sum is $2+24-12-48=-34$, proving every entry.
Stabilizer zero-mode exclusions do not change these first two local
four-dimensional heat coefficients.
\end{proof}

\begin{remark}[Actual UV owner and physical effective-action frontier]
\label{p4:rem:v2216-graded-matter-scope}
The native flat-space Higgs mass squared is $-2\Lambda_g^2$. Its
unregulated infrared Gaussian is not positive; a sufficiently positive EAA
regulator can define a high-scale block, while a stable physical member
requires the actual completion and preparation. Further occurring masses
and thresholds retain their own traces and finite parts. Metric/scalaron
integration, its gauge and ghost complexes, the actual field measure and
Jacobian, global chiral phase, separately allocated sources, boundary terms,
and interacting mixed blocks have not been evaluated by this matter ledger.

For a separately stipulated common shell
$s\in[\Lambda_{\rm UV}^{-2},\mu^{-2}]$ and
$\Gamma^{(1)}=-\int ds\,\mathcal K(s)/s$, the mass-independent volume
term alone is
\[
 \Delta V_{\rm UV}^{(1)}
 =\frac{17}{16\pi^2}(\Lambda_{\rm UV}^4-\mu^4),\qquad
 \Delta\kappa_{\rm geom}^{(1)}=0.
\]
If one additionally identifies $\Lambda_{\rm UV}=\Lambda_g$, neglects
the lower endpoint and all mass/higher terms, and licenses the absolute
measure, combining this illustration with the bare common-heat coefficient
gives $V_{\rm diagnostic}/\kappa_J=(41/2)\Lambda_g^2$.
This is only a common-cutoff diagnostic: the native Higgs mass is itself
of order that cutoff, so that omission is not controlled there. Neither
this ratio nor the rank ledger is a stationary-vacuum prediction, an
evaluation of the complete $\Gamma$, or a license for a new subtraction.
The exact finite flow of Theorem~\ref{p4:thm:v4100-exact-graded-eaa}
still requires its convergent measure and nonsingular Legendre domain.
\end{remark}

\begin{proposition}[Finite-heat integrability and compact-factor test]
\label{p4:prop:v2216-finite-heat-integrability}
Let $q\in\mathbb R^d$, $d>0$, be an actually unrestricted additive
bosonic owner with Lebesgue measure, and let $D(q)=D(q)^*$ act on one
fixed rank-$N$ carrier. The literal positive heat action satisfies
\[
 0\le S_{\rm heat}(q)=\operatorname{Tr}e^{-D(q)^2/\Lambda_g^2}\le N,
 \qquad \int_{\mathbb R^d}e^{-S_{\rm heat}(q)}dq=\infty.
\]
A fixed finite real linear combination of such traces, or a fixed positive
localizer, remains bounded and fails the same absolute-integrability test.
For a Gaussian regulator $\rho|q|^2/2$, $\rho>0$,
\[
 e^{-N}(2\pi/\rho)^{d/2}\le Z_\rho\le(2\pi/\rho)^{d/2};
\]
the unregulated integral is therefore not obtained merely by writing its
finite EAA flow.

In contrast, let $\Omega$ have a supplied compact, normalized,
metric-independent measure, let $c$ be continuous and strictly positive,
and suppose the complete
retained collective action factorizes as $S_E(q,g)=c(q)H[g]$.
Then $\Gamma(H)=-\log\int_\Omega e^{-Hc(q)}d\mu(q)$ is finite and
\[
 \Gamma'(H)=\langle c\rangle_H>0,\qquad
 \Gamma''(H)=-\operatorname{Var}_H(c)\le0,\qquad
 \boxed{\delta_g\Gamma=\langle c\rangle_H\,\delta_gH.}
\]
The metric stationary set is exactly that of $H$ on the stated domain.
For $H>0$, $\Gamma(H)\ge H\Gamma'(H)$; for an unnormalized measure
replace the left side by $\Gamma(H)-\Gamma(0)$.
\end{proposition}
\begin{proof}
Every eigenvalue of the positive heat operator lies in $[0,1]$.
Thus $e^{-S_{\rm heat}}\ge e^{-N}$ on an infinite-volume domain;
comparison with the Gaussian integral proves the regulated bounds.
For a positive localizer $L$, its trace is bounded by
$\operatorname{Tr}L$; finite fixed linear combinations have finite
uniform upper and lower bounds. On the compact factorized domain,
differentiation under the finite integral gives the mean and variance,
and the chain rule gives the metric identity. Since the mean is strictly
positive, its stationary set is unchanged. Finally the derivative is
decreasing, so integrating it from $0$ to $H$ proves the inequality.
\end{proof}

\begin{remark}[Domain restriction is part of the quantum claim]
\label{p4:rem:v2216-finite-heat-domain}
The obstruction assumes an unrestricted infinite-volume additive owner;
it does not impose that domain on every quantum claim. In particular,
the existing positive graviton chart
$\mathscr F=\kappa_J-h^2/12>0$ gives a bounded constant-Higgs
coefficient ball. Such an actual domain can evade the integrability
obstruction, with its inherited Jordan/Einstein measure, Jacobian and
boundary retained. Compact gauge Haar factors also remain distinct from
additive vector factors. Conversely a bounded constant-field chart does
not establish convergence of every metric or interacting field integral.
The compact factor theorem requires exact factorization and a
metric-independent measure: spacetime-dependent modes, mixed interactions
and metric-dependent measures contribute their own variations. A quartic
coercive truncation is a different action whose stability and domain must
be certified on their own. None of these tests computes the full physical
effective action from the bare heat trace alone.
\end{remark}
% END SOURCE-POLARIZATION-AND-NATIVE-UV-LEDGER v2.216

% BEGIN FRAMEWORK v2.226: KINEMATIC CLOSURES


% Positive finite-source closure after v2.246. Proposed for III-F/IV-B.
\subsection{Exact finite CAR influence from the actual one-particle owner}
\label{p3f:sec:v247-car-influence-determinant}

\begin{theorem}[Exact determinant parent on a supplied quasifree source branch]
\label{p3f:thm:v247-car-influence-determinant}
Let $\mathfrak h\cong\mathbb C^d$ be the actual retained complete one-particle
band and let the complete source Hamiltonian on $\bigwedge\mathfrak h$ be
\[
 H_q(t)=d\Gamma(h_q(t))+c_q(t)I,\qquad h_q(t)=h_q(t)^\dagger.
\]
Retain the source clock, frame, central contact and fixed preparation.
Suppose the supplied initial density is gauge-invariant quasifree with
one-particle covariance $0\le n\le I$, defined by
$\langle a_j^\dagger a_i\rangle=n_{ij}$. This is a hypothesis on the actual
preparation, not a replacement for one. Let $u_\pm$ solve the actual
one-particle evolution $i\hbar\dot u_\pm=h_{q_\pm}u_\pm$,
$u_\pm(t_i)=I$, and set
\[
 W=u_-^\dagger(t_f)u_+(t_f),\qquad
 \Delta c=\int_{t_i}^{t_f}(c_{q_+}-c_{q_-})dt,\qquad
 K=I-n+nW.
\]
Then the complete finite-source influence is exactly
\begin{equation}
 \boxed{Z[q_+,q_-]=e^{-i\Delta c/\hbar}\det K.}
 \label{p3f:eq:v247-car-influence-determinant}
\end{equation}
On a sufficiently small history neighborhood of the diagonal,
$\|K-I\|<1$, and its normalized parent is
\[
 \boxed{\Gamma_s=-\Delta c-i\hbar\operatorname{tr}\operatorname{Log} K.}
\]
Thus no $2^d$-dimensional Fock propagation or trace is needed on this
branch. The source operator, central coefficient, covariance and history
remain the inherited inputs; none is fitted by the determinant.
\end{theorem}
\begin{proof}
The exterior lift is multiplicative and satisfies
$U_q=e^{-i\int c_q/\hbar}\bigwedge u_q$. Consequently cyclicity gives
$Z=e^{-i\Delta c/\hbar}\langle\bigwedge W\rangle$.
Diagonalize $n$ unitarily. In that mode basis the supplied quasifree
preparation is the product occupation density with probabilities $n_j$.
The diagonal matrix element of $\bigwedge W$ in the occupied subset $S$
is $\det W_{S,S}$. Therefore
\[
 \langle\bigwedge W\rangle
 =\sum_{S\subseteq\{1,\ldots,d\}}
 \left(\prod_{j\in S}n_j\prod_{j\notin S}(1-n_j)\right)\det W_{S,S}
 =\det(I-n+nW),
\]
by the principal-minor expansion of the determinant. This remains valid
when some $n_j$ are $0$ or $1$; unitary covariance returns to the initial
basis. At equal histories $W=I$, $K=I$, and $\Delta c=0$.
Continuity gives the claimed neighborhood and the convergent matrix
logarithm there. Since $\det K=\exp\operatorname{tr}\operatorname{Log} K$, its branch
is precisely the normalized finite influence parent.
\end{proof}

\begin{corollary}[Complete derivative execution by finite matrix inversion]
\label{p3f:cor:v247-car-logdet-jets}
On the stated local branch with the preparation fixed, let $v,w$ be
admitted history variations. Then
\begin{align*}
 D\Gamma_s[v]&=-D\Delta c[v]
          -i\hbar\operatorname{tr}(K^{-1}DK[v]),\\
 D^2\Gamma_s[v,w]&=-D^2\Delta c[v,w]
 -i\hbar\operatorname{tr}\Bigl(
 K^{-1}D^2K[v,w]\\
 &\qquad{}-K^{-1}DK[w]K^{-1}DK[v]\Bigr),\\
 DK&=n\,DW.
\end{align*}
The actual noncommuting one-particle Duhamel variations of $u_\pm$
supply every matrix derivative. Thus these are the full two-branch jets,
including the mixed causal block and the difference-history covariance;
no scalar spectral substitution or separate noise coefficient enters.
\end{corollary}
\begin{proof}
Differentiate $\operatorname{tr}\operatorname{Log} K$ and
$D(K^{-1})[w]=-K^{-1}DK[w]K^{-1}$.
Finite-dimensional cyclicity makes the displayed second derivative
symmetric in $v,w$ without assuming matrix commutation.
\end{proof}

\begin{proposition}[Executed mean, causal response and noise on the same branch]
\label{p3f:prop:v247-car-source-jets}
For a time-local scalar control, put
$b=\partial_qh$, $r=\partial_q^2h$,
$b_H(t)=u_q(t)^\dagger b(t)u_q(t)$, and
$r_H(t)=u_q(t)^\dagger r(t)u_q(t)$.
The actual source mean, direct contact, connected ordered return and
noise are
\begin{align*}
 F(t)&=\operatorname{tr}(n b_H(t))+\partial_qc(t),\\
 Q_{\rm mean}(t)&=\operatorname{tr}(n r_H(t))+\partial_q^2c(t),\\
 C(t,s)&=\operatorname{tr}\bigl(n b_H(t)(I-n)b_H(s)\bigr),\\
 N(t,s)&=\operatorname{Re}C(t,s),\\
 K^{\rm force}(t,s)&=Q_{\rm mean}(t)\delta(t-s)
 -\frac{i}{\hbar}\Theta(t-s)
        \operatorname{tr}\bigl(n[b_H(t),b_H(s)]\bigr).
\end{align*}
For every real test $v$, with $A_v=\int v(t)b_H(t)dt$,
\[
 \boxed{\iint v(t)N(t,s)v(s)dtds
 =\operatorname{tr}(n A_v(I-n)A_v)
 =\|(I-n)^{1/2}A_v n^{1/2}\|_{\rm HS}^2\ge0.}
\]
These evaluate the already proved causal/noise jets of
Theorem~\ref{p4:thm:v2203-causal-mean-response-noise} from the actual
one-particle propagation and covariance. In particular a central contact
enters $F$ and $Q_{\rm mean}$ and cancels from $C,N$ and the commutator.
\end{proposition}
\begin{proof}
The first two mean formulas follow from
$\langle d\Gamma(A)\rangle=\operatorname{tr}(nA)$.
The finite CAR and the supplied quasifree Wick rule give
\[
 \langle d\Gamma(A)d\Gamma(B)\rangle
 -\langle d\Gamma(A)\rangle\langle d\Gamma(B)\rangle
 =\operatorname{tr}(nA(I-n)B).
\]
Their commutator is exactly $d\Gamma([A,B])$, giving the causal term
with its inherited sign. Smearing the covariance and cycling the finite
trace gives the Hilbert--Schmidt square. Central multiples of $I$ have
zero connected covariance and commute with every insertion.
\end{proof}

The quasifree assumption is needed for the determinant and the displayed
one-particle noise formula. The one-particle mean and retarded commutator
formulas remain valid for every initial finite state with the same $n$,
by the CAR Lie homomorphism. A non-quasifree or paired preparation retains
its higher occupied-sector correlations in the full trace; it is not
replaced by the gauge-invariant quasifree state having the same $n$.
Source-dependent frames, bands and preparations retain their actual
additional variations. This finite evaluation supplies no continuum
measure or new gravity source-to-metric incidence.

\begin{proposition}[Complete finite frame transport preserves the influence parent]
\label{p3f:prop:v247-car-influence-frame-transport}
Let $V_\pm(t)$ be differentiable unitary changes of the complete
one-particle source frame. Suppose their initial and final frames agree
between branches, $V_\pm(t_i)=V_i$ and $V_\pm(t_f)=V_f$, with the
endpoint frames fixed under the admitted interior source variation.
Transport the full generators and initial covariance by
\[
 h'_\pm=V_\pm h_\pm V_\pm^\dagger
                  +i\hbar\dot V_\pm V_\pm^\dagger,
 \qquad n'=V_i nV_i^\dagger,
 \qquad c'_\pm=c_\pm.
\]
Then $Z'=Z$ and the normalized local $\Gamma'_s=\Gamma_s$, including
all admitted source jets. In particular a time- or source-dependent
frame cannot be varied by conjugating only $h$ and omitting its
connection and endpoint data.
\end{proposition}
\begin{proof}
The transformed exact propagator is
$u'_\pm(t)=V_\pm(t)u_\pm(t)V_i^\dagger$, by direct differentiation.
Thus $W'=V_iWV_i^\dagger$ and
$K'=V_iKV_i^\dagger$. The actual central phase is unchanged and the
determinants agree. Normalization at equal histories fixes the same
local logarithm. Exact equality of these functionals preserves every
source derivative on the admitted fixed-endpoint domain. This is the
finite number-conserving exterior-lift transport, not a continuum
measure or anomaly assertion. If branch endpoint frames differ, their
actual endpoint intertwiners must first be retained in the trace.
\end{proof}


% v2.247 proposed source evaluation; uses existing general finite quasifree determinant theorem.
\begin{theorem}[Native complete-band source factorization and a global zero-free response chart]
\label{p4:thm:v247-native-source-factorization}
Use $\hbar=1$ on the actual fixed complete conjugate-sign free source
of Theorem~\ref{p3f:thm:v2216-source-response}, with no zero modes
of $D_\gamma$ on that band, as for its native round-sphere band.
Let $J$ be a connected admitted control interval containing $q_0$
on which the fixed band, frame, clock, constant mass and central
allocation are defined. Suppose its admitted one-particle operator
is invertible at the reference $q_0$ and the
supplied initial state is its filled-negative stationary Slater ground.
Retain the inherited central allocation $c(q)I$. For every positive
Dirac value $\ell$, let $g_\ell$ be its actual positive-sign
multiplicity, counted once per conjugate pair. Then the complete
source influence factorizes as
\[
 Z[q_+,q_-]=
 e^{-i\int(c(q_+)-c(q_-))\,d\eta}
 \prod_{\ell>0}
 \langle g_{\ell,0},u_{\ell,-}^\dagger u_{\ell,+}g_{\ell,0}\rangle^{g_\ell},
\]
where $g_{\ell,0}$ is the normalized lower eigenvector at $q_0$ and
\[
 i\partial_\eta u_\ell
  =(\ell\sigma_3+x(\eta)\sigma_1)u_\ell,
 \qquad x=am.
\]
For constant histories $q_-=q_0$, $q_+=q$ of duration $T$, put
\[
 x_0=a_0m,\quad x=am,\quad E_0=\sqrt{\ell^2+x_0^2},
 \quad E=\sqrt{\ell^2+x^2},\quad
 d_\ell=\frac{\ell^2+x_0x}{E_0E}.
\]
The exact mode factor is
\[
 \boxed{z_\ell=e^{-iE_0T}
 [\cos(ET)+i d_\ell\sin(ET)],\qquad
 |z_\ell|^2=1-
 \frac{\ell^2(x-x_0)^2}{E_0^2E^2}\sin^2(ET).}
\]
For $m>0$ and $a,a_0>0$, $|z_\ell|^2\ge d_\ell^2>0$.
Consequently the full finite-band $Z$ is nonzero on the entire
stationary-reference, constant-comparison chart. More explicitly, put
\[
 t_\ell=E_\ell T,\qquad
 r_\ell=\frac{1+d_\ell}{2}
       +\frac{1-d_\ell}{2}e^{-2it_\ell}.
\]
Then $\operatorname{Re}r_\ell\ge d_\ell>0$, and the right-half-plane
logarithm gives the global anchored formula
\[
 \boxed{\Gamma(q,q_0;T)
 =-[c(q)-c(q_0)]T
 +\sum_\ell g_\ell
 [(E_\ell-E_{\ell,0})T-i\operatorname{Log}r_\ell].}
\]
It is continuous on the full admitted comparison chart
$(q,T)\in J\times[0,\infty)$,
vanishes when $q=q_0$ or $T=0$, and is the normalized continuation
of $-i\log Z$ there. This is an all-order finite-source identity,
without an expansion in $q-q_0$ or $T$.
\end{theorem}
\begin{proof}
On the invariant complete band, $\beta_D$ maps each positive
$\ell$ eigenspace of $D_\gamma$ unitarily to the negative one.
For each positive-sign basis vector $v$, the ordered basis
$(v,\beta_Dv)$ gives $D_\gamma=\ell\sigma_3$ and
$\beta_D=\sigma_1$. These two-dimensional spaces reduce every
$h(q)$ and hence every history propagator. The ground fills precisely
one lower state in each such pair: it is a gauge-invariant Slater state
with covariance $\mathbf1_{(-\infty,0)}(h(q_0))$.
The finite CAR determinant therefore becomes the displayed product
of rank-one occupied-space overlaps; the central phase is retained.
For the constant histories,
\[
 u_q=\cos(ET)I-i\sin(ET)
        (\ell\sigma_3+x\sigma_1)/E.
\]
The reference lower state satisfies
$u_{q_0}g_{\ell,0}=e^{iE_0T}g_{\ell,0}$ and its Bloch vector is
$-(x_0,0,\ell)/E_0$. Taking the matrix element gives $z_\ell$.
The identity
\[
 1-d_\ell^2=
 \frac{\ell^2(x-x_0)^2}{E_0^2E^2}
\]
gives its modulus and strict lower bound. A finite product of these
nonzero factors and a unit central phase is nonzero. Also
$z_\ell=e^{i(E_\ell-E_{\ell,0})T}r_\ell$ and
$\operatorname{Re}r_\ell\ge d_\ell>0$, so the half-plane logarithm
of $r_\ell$ is single valued and continuous. This yields the printed
global logarithm, with the actual central phase, anchored at both
$T=0$ and equal histories.
\end{proof}
If a different band admits zero-Dirac blocks, the displayed product
must additionally include
$Z_0=\exp(i r_0\int(x_+-x_-)d\eta)$, where $r_0$ is the
negative-$\beta_D$ multiplicity on the zero block for $x>0$.
Its constant-comparison contribution to $\Gamma$ is
$r_0(x-x_0)T$. The main theorem uses the stated no-zero-Dirac band. The multiplicity $g_\ell$ includes every
already retained internal copy and introduces no extra factor of two.
The stationary Dirac ground in this theorem is a stated specialization;
QROOT does not prepare every auxiliary matter band, by
Remark~\ref{p4:rem:v2216-qroot-dynamic-scope}. Arbitrary interacting states and two unrestricted histories retain their
complete state and propagator rather than this zero-free specialization.
The constant comparison holds the initial density at $q_0$ fixed even
when $q_+(\eta_i)=q\ne q_0$. It is an admitted independent
source-control comparison, not a curvature perturbation with fixed
initial metric scale and velocity. On a domain requiring those
initial metric conditions, only the arbitrary-history block
factorization applies; the zero-free guarantee for this constant
comparison is not asserted for a ramped metric history.

\begin{corollary}[Exact native work noise and higher cumulants]
\label{p4:cor:v247-native-source-work-cumulants}
Freeze the source and ground of the theorem at one $q$. For the actual
complete insertion $B=\partial_qH_\eta$, define
$\Delta B=B-\langle B\rangle$ and
$A_\ell=x^2\ell^2/E_\ell^2$. Its ordered connected correlator,
symmetrized noise and retarded force memory are
\[
 \boxed{C(t,s)=\sum_\ell g_\ell A_\ell
                e^{-2iE_\ell(t-s)},\qquad
 N(t,s)=\sum_\ell g_\ell A_\ell
                \cos[2E_\ell(t-s)],}
\]
\[
 K^{\rm force}_{\rm mem}(t,s)
 =-2\Theta(t-s)\sum_\ell g_\ell A_\ell
                  \sin[2E_\ell(t-s)].
\]
Here $K^{\rm force}$ is the response of $+\langle B\rangle$;
the source-action mixed block is $-K^{\rm force}$.
The full equal-time characteristic function is
\[
 \boxed{\langle e^{ik\Delta B}\rangle
 =\prod_\ell\left\{
 e^{ikx^2/E_\ell}
 [\cos(kx)-i(x/E_\ell)\sin(kx)]\right\}^{g_\ell}.}
\]
In particular the second through fourth cumulants are sums, with the
actual $g_\ell$, of
\[
 \boxed{\kappa_{2,\ell}=\frac{x^2\ell^2}{E_\ell^2},\qquad
 \kappa_{3,\ell}=\frac{2x^4\ell^2}{E_\ell^3},\qquad
 \kappa_{4,\ell}=-\frac{2x^4\ell^2(\ell^2-2x^2)}{E_\ell^4}.}
\]
For $x>0$, the total third cumulant is strictly positive whenever
at least one admitted $\ell>0$ block occurs. Thus the quasifree
fermionic ground fixes a nonGaussian bilinear work source on that
domain, including every higher work cumulant, without an independently
assigned noise amplitude.
\end{corollary}
\begin{proof}
On each occupied pair the noncentral insertion is $b=x\sigma_1$.
Its ground mean is $-x^2/E$ and its squared transition matrix element
is $x^2\ell^2/E^2$. Ground spectral resolution gives its ordered
correlator; taking the real part and twice the retarded imaginary part
gives the stated kernels. Independent Slater pairs have vanishing
connected cross terms. Since $b^2=x^2I$,
$e^{ikb}=\cos(kx)I+i\sin(kx)\sigma_1$.
Its ground expectation and the centering phase give the characteristic
function. Differentiating its logarithm through order four gives the
printed cumulants; independence makes them additive. Central
$c'(q)I$ cancels from every centered cumulant.
\end{proof}
The unshifted source retains the instantaneous force contact
$c''(q)-\sum_\ell g_\ell x^2/E_\ell$.
The mathematical excitation shift in
Theorem~\ref{p4:thm:v2216-zero-excitation-dynamic-response} changes
that contact and yields its already established subgap coefficients;
it is not an additional native vacuum subtraction.







% BEGIN INTEGRATED V2.251 FORWARD PACKAGE
\subsection{Forward charged/current integration and a finite symmetry sieve}
\label{p4:v251:sec:packet}

The actual local finite constructor of Paper~III-F,
Theorem~\ref{p3f:thm:v1148-finite-metric-constructor}, consumes the field
embedding $J$, parent kinetic form $G_p$, skew connection $B_p$,
energy coefficient $K_p$, every inherited interaction, the actual
occupation/constraint projectors and the central allocation. It gives
\[
 G=J^\dagger G_pJ,\quad S=G^{1/2},\quad W=S^{-1},\quad Z=JW.
\]
\[
 \widehat H=d\Gamma(WKW)-i\hbar d\Gamma(\gamma)
           +\widehat V_{\rm inherited,can}+\varepsilon_0I.
\]
This section applies that forward constructor before comparing its output
with terminal mass targets. It does not supply an unprinted physical field
embedding or assume that the primitive member $Y_*=(I,I,I,I)$ makes the
complete rooted action family blind. The same source owner also retains
the actual rooted transport cell $(I,C_1,C_2)$.

\begin{proposition}[Symmetry survives the complete finite constructor]
\label{p4:v251:prop:symmetry-transport}
On a fixed-rank positive-kinetic chart, let $U$ and $u$ be time-independent
parent and coefficient unitaries. Suppose the actual embedding obeys
$Ju=UJ$, the complete parent $G_p,K_p,B_p$ commute with $U$, all inherited
vertices are invariant under the corresponding slot action, and all actual
constraint and occupation projectors commute with the lifted symmetry.
Then the canonical field embedding $Z$ intertwines the same symmetry,
and the complete finite Hamiltonian commutes with its lift.
\end{proposition}
\begin{proof}
$u^\dagger Gu=J^\dagger U^\dagger G_pUJ=G$, and similarly for $K$.
The embedding-derivative terms in the pulled-back $B$ also commute with
$u$, since differentiating $Ju=UJ$ is legitimate for the fixed unitaries.
Positive functional calculus makes $S,W$ commute with $u$, and hence
$Zu=UZ$. The complete normalized connection
$\gamma=WBW+(W\dot S-\dot SW)/2$ commutes with $u$.
Second quantization therefore preserves symmetry for both one-body terms.
The slot rule preserves every invariant inherited vertex, and an invariant
actual occupation/constraint compression preserves this commutation.
The central term commutes as well. No averaging of the Hamiltonian or
deletion of an inherited interaction was performed.
\end{proof}

For the native symmetric Gibbs/full-ground-support branches, every actual
Hamiltonian symmetry also preserves the state. Thus the state invariance
needed below follows there without an additional vacuum-selection axiom.
For another admitted state, its actual symmetry/support test remains needed.

\begin{proposition}[Finite cyclic response certificate]
\label{p4:v251:prop:finite-response-symmetry}
Let $T=T^\dagger$ act on an actual finite positive-pairing response carrier
$\mathcal V$, and let $J:G\to\mathcal V$ be the actual normalized tag map,
$J^\dagger J=I_3$. Put $d=\dim\mathcal V$ and
$\mathcal G(z)=J^\dagger(zI-T)^{-1}J$. For the rooted cyclic unitary
$X$ on $G$, the following are equivalent:
\[
 \begin{aligned}
 &[\mathcal G(z),X]=0\quad\hbox{for all }z\notin\operatorname{Spec}T,\\
 &\hspace{18mm}\Longleftrightarrow
 [J^\dagger T^nJ,X]=0\quad(0\le n<d).
 \end{aligned}
\]
A sufficient upstream certificate is a unitary $\mathcal U$ on
$\mathcal V$ with $[T,\mathcal U]=0$ and $\mathcal UJ=JX$.
\end{proposition}
\begin{proof}
Commuting a Laurent expansion at infinity proves necessity for every moment.
Conversely Cayley--Hamilton reduces every $T^n$ to the first $d$ powers.
All moments then commute with $X$, so the convergent Laurent series of
$\mathcal G$ does, and rational continuation proves the identity on its
whole resolvent set. The sufficient certificate follows by intertwining
the resolvent. This proves response symmetry without requiring that every
dark complementary degree of freedom share the symmetry.
\end{proof}

The actual normalized positive response may include addition and removal
channels together. Its supplied state/support and pairing are retained.
An arbitrary filled Slater/free boson state is not substituted for the
stationary support: IV-B already evaluates its positive weak leakage on
that particular comparison branch.

Write
\[
 X^3=I,\quad C_1=X+X^\dagger=3P_0-I,\quad C_2=-i(X-X^\dagger).
\]
\[
 \Pi_0=P_0,\quad \Pi_\pm=\tfrac12(Q_\perp\pm C_2/\sqrt3).
\]
On a response satisfying the finite certificate,
\[
 \boxed{\mathcal G(z)=g_0(z)\Pi_0+g_+(z)\Pi_++g_-(z)\Pi_-.}
\]
Each $g_k$ is the resolvent of a positive finite spectral measure of total
mass one, obtained from the corresponding actual normalized source tag.
No character coefficients are fitted or selected.

\begin{theorem}[Exact joint charged-lepton response on the cyclic branch]
\label{p4:v251:thm:lepton-cyclic-response}
Use the actual negative-Hopf left-current rays of Paper~V-A,
Theorem~\ref{p5:thm:v3108-lepton-hopf-incidence-parent}, and Paper~V-B,
Corollary~\ref{p6:cor:v597-zeta-e-closure}. With the representatives
\[
 g_e=\tfrac12(-1-i,i,1)^T,\quad c=\mathbf1/\sqrt3,
\]
\[
 g_\mu=(2,-1+3i,-1-3i)^T/(2\sqrt6),
\]
put $a=(2+\sqrt3)/4$, $b=(2-\sqrt3)/4$ and
$\chi=(1-i)/(4\sqrt2)$. If the actual complete response satisfies the
finite cyclic certificate, its matrix in the current-role frame is exactly
\[
 \boxed{
 \begin{pmatrix}
 ag_++bg_-&\chi(g_+-g_-)&0\\
 \bar\chi(g_+-g_-)&bg_++ag_-&0\\
 0&0&g_0
 \end{pmatrix}.}
\]
Independent role-line rephasings transform the mixed entries covariantly;
the displayed $\chi$ uses exactly the stated representatives. The spectral
measure weights and support conclusions below are phase independent.
The positive electron and muon spectral measures satisfy
\[
 \mu_e=a\mu_++b\mu_-,\quad
 \mu_\mu=b\mu_++a\mu_-,\qquad
 \boxed{r\mu_\mu\le\mu_e\le r^{-1}\mu_\mu,\quad r=7-4\sqrt3>0.}
\]
In particular they have the same visible pole support. A pair of distinct
exact single-pole charged-role claims cannot hold on this branch.
\end{theorem}
\begin{proof}
The native current orientation gives
$\operatorname{Tr}(P_eC_2)=3/2$ and
$\operatorname{Tr}(P_\mu C_2)=-3/2$.
Both rays lie in $Q_\perp G$, so their overlaps with $\Pi_\pm$ are
$(a,b)$ and $(b,a)$ respectively. Direct evaluation gives
$g_e^\dagger\Pi_+g_\mu=\chi$ and the opposite sign for $\Pi_-$;
the singlet cross terms vanish. This proves the complete matrix, including
its mixed role response. Spectral decomposition proves the measure formulas.
Since $a>b>0$ and $r=b/a$, both coefficients of
$\mu_e-r\mu_\mu$ and $r^{-1}\mu_\mu-\mu_e$ are nonnegative.
The same inequalities hold for every Borel subset and hence for each finite
pole residue. They imply equal supports. Distinct single-pole supports
therefore contradict the actual cyclic certificate and tag hypotheses.
\end{proof}

This is a conditional statement about a complete interacting response,
not a no-go for the framework, a requirement of noncentral tree Yukawa
coefficients, or a rejection of a weaker resonance/readout claim. It narrows
the forward alternatives: an actual noncyclic consumed return, different
source/support incidence, or an explicitly different observable must be
exhibited if that strong pair of pole claims is retained.

The common response also obeys the exact phase-sensitive relation
\[
 \sqrt6\,\mathcal G_{e\mu}
   =\frac{1-i}{2}(\mathcal G_{ee}-\mathcal G_{\mu\mu}).
\]
Thus diagonal role responses and their mixed current-sensitive response
are not independent targets on this branch.

\begin{corollary}[Quantified noncyclic content of the desired free return]
\label{p4:v251:cor:required-noncyclic-content}
For the desired unequal free left return
$B=e^2P_e+m^2P_\mu+t^2P_0$,
\[
 \boxed{\|[B,C_2]\|_{\rm HS}^2=\tfrac32(m^2-e^2)^2,\qquad
        \|[B,X]\|_{\rm HS}^2=\tfrac38(m^2-e^2)^2.}
\]
Consequently a cyclic-only action return cannot equal $B$ when $e\ne m$.
\end{corollary}
\begin{proof}
$P_0$ commutes with $C_2$ and $P_e+P_\mu=Q_\perp$.
Thus $[B,C_2]=(e^2-m^2)[P_e,C_2]$. In the two current rays,
$|g_e^\dagger C_2g_\mu|^2=3/4$, giving the first norm after summing
the two off-diagonal entries. Since
$X=P_0-Q_\perp/2+iC_2/2$, the second norm is one quarter of the first.
\end{proof}

\subsubsection{Complete character blocks and the same weak vertex}
On an actual tensor block $\mathcal K\otimes G$, cyclic invariance has the
fully computable equivalent form
\[
 T=T_0\otimes I+T_1\otimes C_1+T_2\otimes C_2,
\]
\[
 T_0=\tfrac13\operatorname{Tr}_G T,\quad
 T_1=\tfrac16\operatorname{Tr}_G[(I\otimes C_1)T],\quad
 T_2=\tfrac16\operatorname{Tr}_G[(I\otimes C_2)T].
\]
The complete environment Hamiltonians are
\[
 T^{(0)}=T_0+2T_1,\qquad
 T^{(\pm)}=T_0-T_1\pm\sqrt3T_2.
\]
The outer matrices need not commute. All bosonic and constraint returns
retained in this tensor block remain in them.

For two such actual charged sectors, suppose their actual canonical weak
insertion satisfies the corresponding cyclic intertwining condition.
Then $V_{ud}=\sum_k B_{ud}^{(k)}\otimes\Pi_k$, with each coefficient
extracted from that same pulled-back action. The complete current-resolved
response between actual factorized tags is
\[
 \begin{aligned}
 &(j_u\otimes g)^\dagger R_u(z)V_{ud}R_d(w)(j_d\otimes h)\\
 &\qquad=\sum_{k=0,\pm}(g^\dagger\Pi_kh)
 j_u^\dagger(z-T_u^{(k)})^{-1}B_{ud}^{(k)}
                  (w-T_d^{(k)})^{-1}j_d.
 \end{aligned}
\]
This follows by inserting the three reducing character projectors. It
is an exact common-action vertex formula, rather than a numerical CKM
assignment or a new mass-eigenframe ontology. Moving sources, pairings,
embeddings and stationary supports retain their corresponding derivative
contacts; the displayed fixed-chart insertion is not their complete
replacement.

The parent generation-blind gauge coefficient is $I_G$. If the actual
finite embeddings intertwine the cyclic symmetry, the slot rule proves
the displayed canonical cyclic condition. It does not automatically make
all $B_{ud}^{(k)}$ equal: different actual kinetic embeddings can give
different character pullbacks. When the actual canonical insertion is
independently proved to be $B_{ud}\otimes I_G$, the three coefficients
are the same $B_{ud}$. In either case they are dependent action outputs,
not three additional interaction parameters.

\subsubsection{Finite forward tests and remaining scope}
The actual primitive $Y_*=I$ alone does not settle any of these symmetry
tests, because the full rooted source retains the noncentral transport
ports. The forward calculation is to extract the actual finite coefficient
word, apply the symmetry-transport certificate or its weaker finite response
moment test, and evaluate the three complete character responses if it
passes. If it fails, the actual noncommuting source coefficient and its
joint source/current/metric jets are retained rather than projected away.

The family symmetry hierarchy is explicit: identity-only coefficients
have commutant $U(3)$; $C_1$ gives $U(1)\times U(2)$; $(C_1,C_2)$ gives
$U(1)^3$; adjoining the actual rooted $R_-$ to $C_1$ gives the full
$M_3(\mathbb C)$ algebra and scalar commutant. The last assertion follows
from $P_0=(I+C_1)/3$ and the rooted matrix units generated by $(R_-,P_0)$.
Availability of that last algebra does not by itself assert which of its
coefficients the action consumes.

Central absolute energy remains in the parent action throughout. Its
cancellation from a fixed gap response does not justify deleting its metric
or source derivatives. No global, continuum, or actual physical hierarchy
closure is inferred from this finite calculation alone.

% END INTEGRATED V2.251 FORWARD PACKAGE


% BEGIN V2.254 PASTERSKI-OBSERVABLE-TRANSFER (no holographic postulate)
\section{Action-generated detector/Ward covariance and finite zero-mode retention}
\label{p4b:sec:v254-detector-ward}
The detector-charge and memory comparisons of
\cite{OertelMoultPasterski2026Detectors,MoultNarayananPasterski2026MemoryCorrelators,PasterskiStromingerZhiboedov2016Memory}
are used here as \emph{observable targets}, not as a holographic dictionary or
as an assumption of asymptotic flatness. This section draws a finite algebraic
consequence of the existing complete-action Ward identity,
Theorem~\ref{p4:thm:v4106-canonical-stress-ward}. Its full proof and the
spectral/angular calculations are in Paper~IV-C,
Section~\ref{p4c:sec:v254-detector-technical}.

\begin{theorem}[Finite in-in charge/detector covariance with explicit Ward residual]
\label{p4b:thm:v254-in-in-charge-detector-covariance}
On an admitted finite quantum/state chart let $\rho\ge0$ with
$\operatorname{Tr}\rho=1$. Suppose that the \emph{actual action} has
independently supplied incoming, outgoing hard-detector and outgoing
soft-radiative self-adjoint charge insertions $Q^-_f,Q^{H,+}_f,Q^{S,+}_f$
for each of two chosen real angular tests $f,g$. All operators are pulled
back to the same finite incoming-picture Hilbert space; no radiation
operator is defined by conservation. Put
\[
R_f:=Q^{H,+}_f+Q^{S,+}_f-Q^-_f,\quad
\Delta A:=A-\operatorname{Tr}(\rho A)I,\quad
\operatorname{Cov}_\rho(A,B):=\tfrac12\operatorname{Tr}\rho\{\Delta A,\Delta B\}.
\]
If both incoming charges are sharp in $\rho$, i.e.
$\operatorname{Cov}_\rho(Q^-_h,Q^-_h)=0$ for $h=f,g$, then
\begin{align}
 &\operatorname{Cov}_\rho(Q^{S,+}_f,Q^{S,+}_g)
   -\operatorname{Cov}_\rho(Q^{H,+}_f,Q^{H,+}_g) \nonumber\\
 &\quad=\operatorname{Cov}_\rho(R_f,R_g)
       -\operatorname{Cov}_\rho(Q^{H,+}_f,R_g)
       -\operatorname{Cov}_\rho(R_f,Q^{H,+}_g).
\label{p4b:eq:v254-charge-covariance}
\end{align}
Writing $\epsilon_h=\sqrt{\operatorname{Cov}_\rho(R_h,R_h)}$
and $\sigma_h=\sqrt{\operatorname{Cov}_\rho(Q^{H,+}_h,Q^{H,+}_h)}$,
\[
\left|\operatorname{Cov}_\rho(Q^{S,+}_f,Q^{S,+}_g)
-\operatorname{Cov}_\rho(Q^{H,+}_f,Q^{H,+}_g)\right|
\le\epsilon_f\epsilon_g+\sigma_f\epsilon_g+\epsilon_f\sigma_g.
\]
Thus exact action-generated charge transport ($R_f=R_g=0$) implies
exact equality of these \emph{connected} covariance kernels. If the
incoming charges are not sharp, the corresponding incoming covariances
and cross terms must be retained; the equality is not automatic.
\end{theorem}
\begin{proof}
Rewrite $Q^{S,+}_h=-Q^{H,+}_h+Q^-_h+R_h$ and expand the centered,
symmetrized covariance. An operator with zero $\rho$-variance has
zero $\rho$-covariance with every other Hermitian operator by
Cauchy--Schwarz for the positive $\rho$-weighted inner product.
The same inequality applied to the three residual terms proves the
bound. No use of a celestial or BMS representation occurs.
\end{proof}

\begin{corollary}[Ward defect controls a generated endpoint-charge residual]
\label{p4b:cor:v254-ward-operator-residual}
Suppose in addition that the independently constituted total
outgoing and incoming charges are exactly the endpoint representatives
of the \emph{same} canonical Ward generator $G_f$:
$Q^{H,+}_f+Q^{S,+}_f=U_f^\dagger G_f(t_f)U_f$ and
$Q^-_f=G_f(t_i)$, with $U_f=U(t_f,t_i)$ the actual evolution.
On an admitted bounded-source chart, the IV-B Ward defect
$\mathfrak D_f=\delta H[R_f^{\rm sym}]-\mathfrak d_tG_f$
then gives
\[
\boxed{\ \epsilon_f\le\|R_f\|
\le \int_{t_i}^{t_f}
\bigl(\|\delta H[R_f^{\rm sym}](t)\|+\|\mathfrak D_f(t)\|\bigr)\,dt.\ }
\]
Here $R_f^{\rm sym}$ is the \emph{symmetry tangent}, distinct from the
charge residual $R_f$ above. All moving-frame, constraint, endpoint,
source-history and physical-mark terms of the actual canonical transfer
remain included. This bound is conditional on the displayed
same-charge identification; it does not generate it.
\end{corollary}
\begin{proof}
Differentiate $U(t,t_i)^\dagger G_f(t)U(t,t_i)$ and use
$\mathfrak d_tG_f=\delta H[R_f^{\rm sym}]-\mathfrak D_f$;
then integrate, take the operator norm, and use
$\epsilon_f\le\|R_f\|$. The defect does not vanish solely because
$G_f$ is a mathematically available generator.
\end{proof}

\begin{proposition}[Same-action collar-smeared matter/gauge flux insertion]
\label{p4b:prop:v254-same-action-angular-flux}
On an admitted smooth \emph{finite four-volume collar} $\mathcal W$
with physical observer $t^\mu$, outward spatial unit direction
$s^\mu$, with $t\cdot s=0$, and smooth compactly supported real test function $w$, suppose the retained
matter/gauge action has the Hilbert insertion
$\delta S_s[v]=\tfrac12\int_{\mathcal W}T^{\mu\nu}v_{\mu\nu}\,d\mu_{\rm act}$
and admits $v_{\mu\nu}=-w(t_\mu s_\nu+s_\mu t_\nu)$.
Then
\[
\boxed{\ \mathcal F_s[w]:=-\int_{\mathcal W}
 w T_{\mu\nu}t^\mu s^\nu\,d\mu_{\rm act}
 =\delta S_s[v].\ }
\]
Its corresponding Hermitian quantum insertion, \emph{if licensed by the
same functional source family}, inherits the complete finite in-in
mean and connected covariance, including all action contacts and
moving-frame derivatives. This is a collar \emph{flux-density test},
not automatically a surface-energy operator, gravitational-wave flux,
asymptotic supertranslation charge or a detector POVM.
\end{proposition}
\begin{proof}
Insert the indicated symmetric metric tangent into the actual
Hilbert first variation. The two symmetric contractions are equal.
\end{proof}

\begin{remark}[No soft-sector or memory shortcut]
The finite stationary zero-frequency noise theorem and the round
angular spin-two inverse are proved in
Paper~IV-C, Theorem~\ref{p4c:thm:v254-zero-mode-no-omission}
and Proposition~\ref{p4c:prop:v254-angular-spin-two-inverse}.
The soft-omission example of \cite{JorstadPasterski2026SoftOmissions}
is a warning not to replace the full two-branch quantum covariance by
retarded response alone. The asymptotic Noether construction of
\cite{Pasterski2023NoetherCurrents} is a useful comparison only
\emph{after} action-generated physical charges, their boundary terms and
the requisite radiative domain are established. None of these cited
papers proves the missing Born--Hopf parent-action coefficients,
physical BMS charges, $f(R)$ scalaron mixing or infrared limits.
\end{remark}
% END V2.254 PASTERSKI-OBSERVABLE-TRANSFER

% BEGIN FRAMEWORK v2.266: CANONICAL FINITE INSTRUMENT EXISTENCE
\subsection{Abstract repeated localization instruments and action firewall}
\label{p4b:sec:v266-instrument-existence}
\begin{theorem}[Sector-preserving finite instrument and dilation exist]
\label{p4b:thm:v266-native-instrument}
For each of the native exact-position five-, six-, and seven-effect
POVMs of Paper~II-C, the operators $K_i=E_i^{1/2}$ give a
sector-preserving finite completely positive instrument
$\mathcal I_i(\rho)=K_i\rho K_i$, with nonselective trace-preserving
channel $\Phi=\sum_i\mathcal I_i$. Its geometric and protected
populations are conserved for every input state. It admits a finite
unitary pointer/ancilla realization preserving the two sectors.
Neither this existence theorem nor the specific L\"uders choice
proves that the common physical action consumes or performs the
instrument.
\end{theorem}
\begin{proof}
Each effect lies in the native algebra
$M_{11}(\mathbb C)\oplus M_4(\mathbb C)$, so functional calculus gives
$K_i\in M_{11}\oplus M_4$ and $[K_i,P_S]=0$.
Since $\sum_iK_i^\dagger K_i=\sum_iE_i=I$, the maps form a CP
instrument and $\Phi$ is trace preserving; commutation with $P_S$
gives $\operatorname{Tr}(P_S\Phi(\rho))=\operatorname{Tr}(P_S\rho)$.
The isometry
$V\psi=\sum_iK_i\psi\otimes|i\rangle$ from $\mathcal H$
into $\mathcal H\otimes\mathbb C^N$ obeys $V^\dagger V=I$.
In finite dimension it extends to a unitary with a prepared ancilla;
the extension can be performed independently on
$\mathcal H_G\otimes\mathbb C^N$ and
$\mathcal H_S\otimes\mathbb C^N$.
This constructs one possible instrument, not the unique instrument
associated with the effects or an action-incidence proof.
\end{proof}

\begin{corollary}[Quadratic fingerprint as a conditional repeated event]
\label{p4b:cor:v266-quadratic-repeat}
For the five-outcome regular-simplex POVM at radius $r\ge2\Rmom$,
sequential repetition of the \emph{same canonical L\"uders instrument}
has same-label probability
\[
\boxed{p_{\rm same}(\rho)=\frac15+
\frac{4}{5r^2}\operatorname{Tr}\left(\rho\sum_a(X_a^D)^2\right).}
\]
At $r=2\Rmom$ this gives the protected benchmark probabilities
$7/20$ (intrinsic) and $2/5$ (cone), for the corresponding specified
protected maximally mixed preparations. First-outcome probabilities
are uniform in both benchmarks, but two-event correlations differ.
\end{corollary}
\begin{proof}
The effect identity and its state/instrument qualifications are
Proposition~\ref{p2c:prop:v266-repeat-fingerprint}.
\end{proof}
The original physical response/detector obligations, timing, pointer
preparation, sector source state and same-action action-consumer
incidences remain independent: an abstract Stinespring extension is
not evidence that any such selected process occurs.
% END FRAMEWORK v2.266: CANONICAL FINITE INSTRUMENT EXISTENCE


% BEGIN FRAMEWORK v2.320: PAPER IV-B PETERSEN-HOPF QUANTUM RESPONSE
\section{Petersen mirror transfer, finite pole visibility and loop interference}
\label{p4b:sec:v320-petersen-response}
\noindent\textit{Scope.} These are exact conditional quantum comparisons on
the marked Paper~I-A Petersen modules. They produce neither a new term
in the complete generated action nor an actual new one-particle mass,
source pole, time variable, or matter multiplet. All physical responses
remain controlled by the present Paper~IV-B same-owner source and
complete finite-action requirements.

\begin{theorem}[Conditional two-core chiral block and its source-visible poles]
\label{p4b:thm:v320-chiral-auxiliary-poles}
Fix a pair of cores $i,j$, with $T=T_{ij}:\mathbb C^3\to\mathbb C^{3'}$
and positive pair metric $T^*T=4P_a+\frac14P_t$ from
Paper~I-A. On the \emph{auxiliary} direct sum
$\mathcal H_{\rm cmp}=\mathbb C^3\oplus\mathbb C^{3'}$, define
\[
 H_T=\begin{pmatrix}0&T^*\\T&0\end{pmatrix},
 \qquad Z=\begin{pmatrix}I_3&0\\0&-I_3\end{pmatrix}.
\]
Then
\[
 \det(zI_6-H_T)=(z^2-4)(z^2-\tfrac14)^2,
 \quad \operatorname{Spec}(H_T)
 =\{+2,-2,+\tfrac12^{\,(2)},-\tfrac12^{\,(2)}\}.
\]
For an initial unit vector $\psi\in\mathbb C^3$, set
$a=\|P_a\psi\|^2\in[0,1]$. The diagonal source resolvent is
\[
 \boxed{\langle(\psi,0),(z-H_T)^{-1}(\psi,0)\rangle
 =z\left(\frac{a}{z^2-4}
 +\frac{1-a}{z^2-\tfrac14}\right).}
\]
The unitary comparison evolution $e^{-itH_T}$ transfers probability
to the conjugate summand according to
\[
 \boxed{P_{\mathbf3\to\mathbf3'}(t;\psi)
 =a\sin^2(2t)+(1-a)\sin^2(t/2).}
\]
Thus an axial source ($a=1$) sees only the pair of poles $\pm2$,
whereas a transverse source ($a=0$) sees only $\pm1/2$. These
numbers are dimensionless geometric comparison eigenvalues, not masses,
frequencies with physical units, or action-produced resonances.
\end{theorem}
\begin{proof}
The spectrum of a self-adjoint off-diagonal block consists of the
signed nonzero singular values of $T$. The Schur complement gives the
upper-left resolvent block
\[
 [(z-H_T)^{-1}]_{11}
 =z(z^2-T^*T)^{-1}
 =\frac{z}{z^2-4}P_a+\frac{z}{z^2-\tfrac14}P_t.
\]
Taking its quadratic form in $\psi$ yields the source-resolvent
formula. Functional calculus of the two singular values gives the
transfer sine-squares. The coefficient $a$ is a source-preparation
mark; its value and the existence of $H_T$ in the physical action
are not selected by the geometric construction.
\end{proof}

\begin{corollary}[Complementary sign is a relative phase, not a separate pole]
\label{p4b:cor:v320-complement-phase-gauge}
For the complementary three-core selection, $T_{S^c}=-T$ and
\[
 H_{-T}=ZH_TZ,
 \quad (z-H_{-T})^{-1}=Z(z-H_T)^{-1}Z.
\]
Every initial state and registered effect commuting with $Z$ has
identical evolution and finite resolvent probabilities for $T$ and
$-T$. In particular the two constructions have identical characteristic
polynomials and identical diagonal source-visible poles. A
complement-odd distinction requires an independently referenced
inter-sector coherence or an additional coupling not transformed
away by $Z$; the determinant sign of the bare real map is not itself
a new quantum mass or species label.
\end{corollary}
\begin{proof}
Conjugation by $Z$ changes the sign of both off-diagonal blocks.
Conjugation covariance holds for every spectral function and
resolvent. For $[Z,\rho]=[Z,E]=0$, cyclicity of trace gives
$\operatorname{Tr}(E e^{-itH_T}\rho e^{itH_T})
=\operatorname{Tr}(E e^{-itH_{-T}}\rho e^{itH_{-T}})$.
\end{proof}

\begin{theorem}[Finite signed Petersen loops as a conditional phase
comparison]
\label{p4b:thm:v320-z2-loop-interference}
Choose unoriented Hopf-axis representatives and give a Petersen
edge $v\sim w$ the real line transport
$s_{vw}=\operatorname{sgn}(u_v\cdot u_w)\in\{\pm1\}$.
For each closed edge loop $\gamma$ define the gauge-invariant
Wilson sign
\[
 W(\gamma)=\prod_{(v,w)\in\gamma}s_{vw}.
\]
The twelve simple Petersen pentagons divide exactly into six with
$W=+1$ and six with $W=-1$ (Paper~I-A). A mathematical two-path
interferometer coherently comparing two based loops with calibrated
relative phase $\delta$ has one normalized port probability
\[
 \boxed{p_\delta
 =\frac14\left|W(\gamma_1)+e^{i\delta}W(\gamma_2)\right|^2
 =\frac12\left(1+W(\gamma_1)W(\gamma_2)\cos\delta\right).}
\]
Loops of opposite Wilson sign reverse the fringes. The ratio
$W(\gamma_1)W(\gamma_2)$, rather than either arbitrary edge sign, is
the relevant relative phase. Realizing the two paths as controlled
quantum processes and calibrating their amplitudes requires
additional action, preparation and readout incidence.
\end{theorem}
\begin{proof}
A gauge change $u_v\mapsto\epsilon_vu_v$ changes
$s_{vw}\mapsto\epsilon_v\epsilon_ws_{vw}$, whose factors cancel
around any closed loop. The $6+6$ classification is the exact finite
cycle certificate in Paper~I-A. Interference of two equal-normalized
complex amplitudes $W(\gamma_1)/\sqrt2$ and
$e^{i\delta}W(\gamma_2)/\sqrt2$ at a balanced recombiner gives
the displayed port probability. No physical identification of this
real sign with a measured Berry phase is used.
\end{proof}

\begin{proposition}[Source-action visibility and the obstruction to a
Petersen mass claim]
\label{p4b:prop:v320-source-incidence-boundary}
The finite operators $S$, $T_{ij}$, $H_T$ and the Wilson signs are
not automatically constituents of the actual Paper~IV-B clock
Hamiltonian or its source-cyclic quotient. A proposed physical
Petersen pole requires (i) an occurrence-typed map from the actual
one-particle carrier to the marked module, (ii) the complete
same-action first and second response jets and their source/receiver
ports, and (iii) a nonzero residue in the resulting physical
source-cyclic resolvent. The bare polynomial
$(z^2-4)(z^2-\frac14)^2$ satisfies none of these incidence tests by
itself. Its mere existence does not modify $E_{32}$, $E_{36}$,
$E_{40}$, particle masses, the continuum limit or the accepted
finite action.
\end{proposition}
\begin{proof}
The comparison operator was defined on an independently marked
six-dimensional auxiliary module and lacks the action-consumed
one-particle and source ownership maps. The exact finite-source
visibility definitions and pole/receiver theorems of this paper
require those maps before an eigenvalue can be promoted to a
physical pole. The conditional resolvent in
Theorem~\ref{p4b:thm:v320-chiral-auxiliary-poles} explicitly
illustrates that even a supplied finite operator need not expose
all of its eigenvalues through a given source.
\end{proof}
% END FRAMEWORK v2.320: PAPER IV-B PETERSEN-HOPF QUANTUM RESPONSE



% BEGIN FRAMEWORK v2.347: EXACT INTERACTING RESPONSE AND EVEN-READOUT WITNESSES
\section{Exact limits of one-body response and finite interaction hierarchies}
\label{p4b:sec:v347-response-closure}
The unrestricted many-body response and source-jet contracts already
established above remain the primary statements.  The new criterion below
specifies exactly when an observable-specific reduction to one-body
densities is possible, even when the Hamiltonian has interactions.

\begin{theorem}[Necessary and sufficient one-body information criterion]
\label{p4b:thm:v347-onebody-criterion}
Let \(K\) be a complex \(d\)-mode fermionic one-particle space and
\(\mathscr H=\mathcal F_-(K)\).  In the full finite CAR algebra define the
unital, star-closed linear space
\[
 \mathscr W_1=\operatorname{span}_{\mathbb C}
 \{I,c_i^\dagger c_j:1\le i,j\le d\}
 =\mathbb CI+d\Gamma(\operatorname{End}(K)).
\]
For an arbitrary operator \(O\in\operatorname{End}(\mathscr H)\), its
expectation \(\operatorname{Tr}(\rho O)\) is determined by the
one-particle density matrix \(n(\rho)\) for \emph{every} density matrix
\(\rho\) if and only if \(O\in\mathscr W_1\).  Equivalently,
\begin{equation}\label{p4b:eq:v347-onebody-defect}
 \boxed{\eta(O):=\|O-\mathsf P_1O\|_{\rm HS}=0,}
\end{equation}
where \(\mathsf P_1\) is the Hilbert--Schmidt orthogonal projection onto
\(\mathscr W_1\).  This is a criterion for one observable, not an
assumption that a full interacting Heisenberg algebra is quadratic.
\end{theorem}
\begin{proof}
Every element of \(\mathscr W_1\) has expectation equal to a
constant plus a linear combination of entries of \(n(\rho)\), giving
sufficiency.  Conversely suppose first that \(O=O^\dagger\), and write
\(Y=O-\mathsf P_1O\).  Because \(\mathscr W_1\) is star closed,
\(Y=Y^\dagger\), \(\operatorname{Tr}Y=0\), and
\(\operatorname{Tr}(Yc_i^\dagger c_j)=0\).  If \(Y\ne0\), put
\(D=2^d\) and choose
\(0<\epsilon<(D\|Y\|_{\rm op})^{-1}\).  The states
\(\rho_\pm=D^{-1}I\pm\epsilon Y\) are positive, normalized and have
identical one-body matrices, but
\[
 \operatorname{Tr}((\rho_+-\rho_-)O)
 =2\epsilon\operatorname{Tr}(Y^2)>0.
\]
This contradiction proves necessity.  For non-Hermitian \(O\), apply
the same argument to its Hermitian real and imaginary parts; the
star-closedness of \(\mathscr W_1\) gives the result.
\end{proof}

\begin{corollary}[Sharp criterion for an arbitrary-state retarded CAR source]
\label{p4b:cor:v347-retarded-criterion}
Let \(U(t,s)\) be the \emph{complete} finite many-body propagator and
let \(A=d\Gamma(A_1)\), \(B=d\Gamma(B_1)\) be actual bilinear source
operators.  Define
\[
 K_{AB}(t,s)=[U(t,s)^\dagger A U(t,s),B],\qquad
 \eta_{AB}(t,s)=\eta(K_{AB}(t,s)).
\]
The retarded commutator at fixed \(t,s\) is determined by the
one-particle density matrix at \(s\), for every preparation, precisely
when \(\eta_{AB}(t,s)=0\).  If this condition fails, two states with the
same one-body matrix have different responses.  For any two such states,
\[
 \big|\operatorname{Tr}((\rho-\sigma)K_{AB})\big|
 \le\|\rho-\sigma\|_{\rm HS}\,\eta_{AB}(t,s).
\]
A bilinear-invariant Heisenberg dynamics is sufficient but not
necessary for this \emph{source-pair-specific} property.
\end{corollary}
\begin{proof}
The theorem applies to \(K_{AB}(t,s)\).  For the bound, the state
difference is orthogonal to \(\mathscr W_1\) by hypothesis; apply the
Hilbert--Schmidt Cauchy--Schwarz inequality to its pairing with
\((I-\mathsf P_1)K_{AB}\).
\end{proof}

\begin{proposition}[Explicit three-mode interacting obstruction]
\label{p4b:prop:v347-three-mode-response-obstruction}
On three CAR modes take \(H=Un_1n_2\),
\(A=c_1^\dagger c_3\), \(B=A^\dagger\) and
\(\theta=U(t-s)/\hbar\).  Then
\[
 A(t,s)=e^{i\theta n_2}A,\qquad
 [A(t,s),B]=e^{i\theta n_2}(n_1-n_3).
\]
The two normalized states
\[
 \rho_1=\tfrac12(|110\rangle\langle110|+|001\rangle\langle001|),
 \quad
 \rho_2=\tfrac12(|100\rangle\langle100|+|011\rangle\langle011|)
\]
have the identical one-body matrix \(n=I_3/2\), while
\[
 \langle[A(t,s),B]\rangle_{\rho_1}
 =\tfrac12(e^{i\theta}-1),\qquad
 \langle[A(t,s),B]\rangle_{\rho_2}
 =\tfrac12(1-e^{i\theta}).
\]
Thus the universal arbitrary-state one-body response reduction is
false under generic interacting evolution.
\end{proposition}
\begin{proof}
The number operator \(n_2\) commutes with \(A,B\), and
\([n_1n_2,A]=n_2A\), giving the first identity.  The second follows
from \([c_1^\dagger c_3,c_3^\dagger c_1]=n_1-n_3\).  Direct diagonal
evaluation on the four displayed Fock basis vectors gives the two
expectations and identical one-particle marginals.
\end{proof}

\begin{theorem}[Finite terminating interacting Majorana moment hierarchy]
\label{p4b:thm:v347-majorana-hierarchy}
Let \(\gamma_1,\ldots,\gamma_{2d}\) be Majorana generators with
\(\{\gamma_a,\gamma_b\}=2\delta_{ab}I\).  Define
\(\mathscr C_{\le r}^{\rm even}\) as the span of even ordered
Majorana monomials of degree at most \(r\), including the identity.
For a time-dependent even Hamiltonian
\(H(t)=H_0(t)+H_2(t)+H_4(t)\), whose nonconstant terms are sums of
quadratic and quartic monomials,
\[
 [H_2(t),\mathscr C_{\le2k}^{\rm even}]
 \subseteq\mathscr C_{\le2k}^{\rm even},\qquad
 [H_4(t),\mathscr C_{\le2k}^{\rm even}]
 \subseteq\mathscr C_{\le2k+2}^{\rm even}.
\]
Thus the exact differential equations for even moments through degree
\(2k\) involve at most moments of degree \(2k+2\).  The sequence
terminates at degree \(2d\), giving a finite exact interacting closure;
it need not yield a small or experimentally licensed dynamical system.
\end{theorem}
\begin{proof}
Every even Majorana monomial squares to a scalar sign up to its phase.
For a degree-\(2r\) and a degree-\(2s\) ordered monomial, the
commutator vanishes when the number of common generators is even; when
it is odd, the resulting monomial has degree \(2r+2s-2j\), where
\(j\ge1\) is that odd number.  A quadratic factor therefore retains
or lowers the filtered degree, while a quartic factor raises it by at
most two.  Linear extension gives the displayed inclusions.  The
complete Clifford matrix algebra has only \(2d\) generators, so its
even word filtration terminates.
\end{proof}

\begin{proposition}[Short-time odd dynamics witnessed by even projectors]
\label{p4b:prop:v347-odd-transfer-witness}
On a physically admitted self-adjoint finite generator, let the
\(11\oplus4\) central decomposition have block form
\[
 H=\begin{pmatrix}H_+&T\\T^\dagger&H_-\end{pmatrix},
 \qquad T:\mathcal H_-\longrightarrow\mathcal H_+.
\]
For every initial state \(\rho_+=P_+\rho_+P_+\), the \emph{even}
projector \(P_-\) has short-time transition probability
\begin{equation}\label{p4b:eq:v347-odd-transfer}
 \boxed{\operatorname{Tr}(P_-e^{-itH/\hbar}\rho_+e^{itH/\hbar})
 =\frac{t^2}{\hbar^2}\operatorname{Tr}(\rho_+TT^\dagger)
 +O(t^3).}
\end{equation}
Taking the supremum over all \(\rho_+\), the \(t^2\) coefficient times
\(\hbar^2\) is \(\|T\|_{\rm op}^2\).  The witness only detects an
\emph{actually occurring} odd dynamical bridge: an auxiliary odd matrix
or field-typed \(\mathsf{FENV}\) algebra alone does not certify this
Hamiltonian, the preparation or the measured even transition.
\end{proposition}
\begin{proof}
The off-diagonal propagator block obeys
\(P_-e^{-itH/\hbar}P_+=-itT^\dagger/\hbar+O(t^2)\).
Taking its adjoint product gives the probability and coefficient.
Maximizing a positive quadratic form over density matrices selects the
largest eigenvalue of \(TT^\dagger\), equal to \(\|T\|_{\rm op}^2\).
\end{proof}
% END FRAMEWORK v2.347



% BEGIN SOURCE-INTEGRATED CROSS-PAPER MATHEMATICS v3_017, 15 November 2026.
%  Derived geometric comparison modules; physical source occurrence not promoted.
%  E3/E4 independent audit of these NEW results remains pending; inherited source statements retained.
\section{Additional exact mathematical consequences from the verified source geometry}

% -------- Research module Paper4B-DirectedNonetLudersChannel-v1_1 (latest supplied version) --------
% Born--Hopf Geometry CC17 / proposed Paper IV-B auxiliary quantum extension.
% Consumes PROPOSED (not canonical) cc12 source:
%   p1:thm:cc12-directed-coherent-gram,
%   p2c:thm:cc12-twenty-complex-povm,
%   p2c:thm:cc16-directed-tomography-stability.
% This module is conditional on those mathematical constructions and
% does not establish an occurring meson bath, GKLS coefficient, or source.

\subsection{The twenty directed core pairs define an exact anisotropic quantum channel}
\label{p4b:sec:cc17-nonet-luders-semigroup}
\begin{theorem}[A covariant entanglement-breaking qutrit channel with exact 3/5 rates]
\label{p4b:thm:cc17-directed-nonet-luders}
On the proposed marked twenty-effect auxiliary coherent nonet
$H_{ij}=H_{ij}^2=H_{ij}^*$ with
$\sum_{i\ne j}H_{ij}=\frac{20}{3}I_3$, define
\[
\boxed{\Phi(X)=\sum_{i\ne j}\frac3{20}H_{ij}XH_{ij}
 =\frac3{20}\sum_{i\ne j}\operatorname{Tr}(XH_{ij})H_{ij}.}
\]
Then $\Phi$ is completely positive, trace preserving, unital,
self-adjoint for the Hilbert--Schmidt metric and
entanglement breaking. It is equivariant under the same marked
$A_5$ representation as the core-pair projectors.
In the source's orthogonal decomposition
\[\operatorname{Herm}_3(\mathbb C)
 =\mathbb RI\oplus i\mathfrak{so}(3)\oplus\operatorname{Sym}^2_0(\mathbb R^3)
 \cong\mathbf1\oplus\mathbf3\oplus\mathbf5,\]
one has the \emph{exact} operator spectrum
\[\boxed{\operatorname{spec}(\Phi)=
\{1^{(1)},\ (1/2)^{(3)},\ (1/10)^{(5)}\}.}\]
For a trace-one $\rho=I/3+Q+iA(v)$,
\[\boxed{\Phi^k(\rho)=I/3+10^{-k}Q+2^{-k}iA(v),\qquad k\ge0.}\]
Its only fixed density matrix is $I_3/3$; the sharp
Hilbert--Schmidt contraction factor on traceless matrices is $1/2$.
\end{theorem}
\begin{proof}
The Kraus operators $K_{ij}=\sqrt{3/20}\,H_{ij}$ obey
$\sum K_{ij}^*K_{ij}=I_3$ and likewise
$\sum K_{ij}K_{ij}^*=I_3$. Since each $H_{ij}$ is rank one,
$H_{ij}XH_{ij}=\operatorname{Tr}(XH_{ij})H_{ij}$, making
$\Phi$ a rank-one measure-and-prepare channel; such a channel
is entanglement breaking. Covariance follows by permuting the
projectors. As a real Hilbert--Schmidt superoperator,
$\Phi=(3/20)SS^*$ for the previous twenty-projector synthesis
$Se_{ij}=H_{ij}$. Multiply the source's nonzero synthesis Gram
eigenvalues $(20/3,10/3,2/3)$ by $3/20$ to obtain
$1,1/2,1/10$ with multiplicities $1,3,5$.
Strict positivity of the last two rates gives the fixed-point
and iteration statements.
\end{proof}

\begin{corollary}[An exact auxiliary GKLS semigroup and purity law]
\label{p4b:cor:cc17-nonet-gkls-purity}
For any \emph{chosen comparison} rate $\gamma>0$, the generator
\[\boxed{\mathcal L_\gamma=\gamma(\Phi-I)}\]
is in Lindblad form with jump operators
$L_{ij}=\sqrt{3\gamma/20}\,H_{ij}$. Its unital CPTP semigroup
acts on the triplet and quintet with distinct rates
$\gamma/2$ and $9\gamma/10$:
\[
 e^{t\mathcal L_\gamma}\rho
 =I/3+e^{-9\gamma t/10}Q+e^{-\gamma t/2}iA(v).
\]
For integer channel iterations the purity is exactly
\[\boxed{\operatorname{Tr}(\Phi^k\rho)^2
=\frac13+100^{-k}\operatorname{Tr}Q^2
       +4^{-k}\operatorname{Tr}(iA(v))^2.}\]
The semigroup purity obeys the analogous formula with exponential
factors $e^{-9\gamma t/5}$ and $e^{-\gamma t}$.
The von Neumann entropy cannot decrease under either the unital
channel or its semigroup; no physical decoherence rate is fixed.
\end{corollary}
\begin{proof}
Using $\sum L_{ij}^*L_{ij}=\gamma I$, the Lindblad expression
$\sum L_{ij}X L_{ij}^*-\frac12\{\gamma I,X\}$ equals
$\gamma(\Phi(X)-X)$. The eigenspace formula gives both semigroup
and discrete powers. Orthogonality of the three Hilbert--Schmidt
summands yields the purity equations. Quantum relative-entropy
data processing for a unital channel with fixed state $I/3$
proves entropy monotonicity.
\end{proof}

\begin{remark}[Physical source firewall]
This is a mathematically valid quantum channel on the auxiliary
$\mathbb C^3$ comparison. Its existence does not identify a physical
meson nonet measurement, a hadronic dissipative source, a decay pole,
QCD branching fractions or the rate $\gamma$. In particular
no known hadronic lifetime is predicted by the eigenvalues $1/2,1/10$.
\end{remark}
% END SOURCE-INTEGRATED CROSS-PAPER MATHEMATICS v3_017.

\printbibliography

% BEGIN FRAMEWORK v2.356: exact current-visible cyclic spectral solver
\section{Exact current-visible cyclic response reduction}
\label{p4b:sec:v356-current-krylov}

\begin{theorem}[Cyclic Krylov evaluation of the inverse-cubic response]
\label{p4b:thm:v356-current-visible-krylov}
Let $H$ be finite-dimensional and Hermitian with ground energy $E_0$
and normalized ground vector $\lvert 0\rangle$.  For a Hermitian
source/current $J$, put $v=(J-\langle J\rangle_0)\lvert0\rangle$.
Assume $v$ is orthogonal to the \emph{entire} ground eigenspace.
On the current-visible cyclic subspace
\[
 \mathcal K_v=\operatorname{span}\{v,Av,A^2v,\ldots\},
 \qquad A=(H-E_0)|_{\mathcal K_v}>0,
\]
let $Q:\mathbb C^r\to\mathcal K_v$ be an isometry from a
\emph{completed} Lanczos/Krylov basis with
$Qe_1=v/\|v\|$ and $T_r=Q^*AQ$.  For $v\ne0$,
\[
 \boxed{Z_{\rm spec}(0)=2\langle v,A^{-3}v\rangle
     =2\|v\|^2e_1^*T_r^{-3}e_1
     =\int_0^\infty t^2\langle v,e^{-tA}v\rangle\,dt.}
\]
The value is unchanged by a common unitary transport of $H,J$ and
the vacuum.  An incomplete Krylov basis gives a projection or
approximation, not this exact identity.
\end{theorem}
\begin{proof}
By construction $\mathcal K_v$ is $A$-invariant; finite-dimensional
Hermiticity makes it reducing.  Completion of the Krylov process
therefore gives $AQ=QT_r$, hence
$Q^*A^{-3}Q=T_r^{-3}$ and the first equality follows.
For every positive eigenvalue $\Delta$,
$\int_0^\infty t^2 e^{-t\Delta}\,dt=2\Delta^{-3}$;
finite spectral resolution proves the last equality.
Simultaneous unitary conjugation transports the spectral measure
without changing its weights.  A truncated Krylov span is not
necessarily invariant and cannot justify replacing $A^{-3}$ by a
smaller inverse cube without an additional error certificate.
\end{proof}

\begin{remark}[Contact term and physical occurrence]
\label{p4b:rem:v356-current-krylov-scope}
The theorem evaluates only the positive-frequency spectral part
of an already owned current-response packet.  It neither determines
the same-action local temporal contact $Z_{\rm loc}(0)$ nor
constructs a physical QCD Hamiltonian, current, or occurring vacuum.
The complete electromagnetic coefficient is
$Z_J^{\rm tot}(0)=Z_{\rm loc}(0)+Z_{\rm spec}(0)$ on the licensed
temporal source chart.
\end{remark}
% END FRAMEWORK v2.356


% BEGIN SOURCE-INTEGRATED DAG CROSS-CONNECTIONS CC18 (16 November 2026)
\section{Symmetry ownership and Haar twirl of auxiliary rank-one channels}
\label{p4b:sec:cc18-su3-nonet-twirl}
\begin{theorem}[Sharp full-$SU(3)$ covariance obstruction of the marked nonet channel]
\label{p4b:thm:cc18-nonet-su3-obstruction}
The $A_5$-covariant auxiliary twenty-effect qutrit channel
$\Phi$ of Theorem~\ref{p4b:thm:cc17-directed-nonet-luders}
acts on $\operatorname{Herm}_3(\mathbb C)=\mathbf1\oplus\mathbf3
\oplus\mathbf5$ by $1,1/2,1/10$ respectively.  It is \emph{not}
covariant under the full $SU(3)$ conjugation action.  In fact for
$X=i(E_{12}-E_{21})$ and
$U=e^{-i\pi/6}\operatorname{diag}(i,1,1)\in SU(3)$,
\[
\boxed{\Phi(UXU^*)=\tfrac1{10}UXU^*,\qquad
U\Phi(X)U^*=\tfrac12 UXU^*.}
\]
The Hilbert--Schmidt covariance defect normalized by $\|UXU^*\|$
is exactly $2/5$.  This excludes full $SU(3)$ covariance for this
same \emph{auxiliary} channel; it does not identify that channel with
physical color-gauge transport or physical meson evolution.
\end{theorem}
\begin{proof}
The matrix $X$ belongs to the imaginary antisymmetric triplet,
whereas $UXU^*=-(E_{12}+E_{21})$ lies in the real symmetric
traceless quintet.  The exact source eigenvalues give the two
formulas and $|1/2-1/10|=2/5$.  The same conclusion follows from the
irreducibility of the traceless $SU(3)$ conjugation module: a fully
covariant linear channel must act by a single scalar there.
\end{proof}

\begin{theorem}[Universal full-unitary twirl of a rank-one L\"uders channel]
\label{p4b:thm:cc18-rankone-haar-twirl}
Let $d\ge2$ and $P_k$ be rank-one orthogonal projectors in
$\mathbb C^d$ with positive coefficients $w_k$ such that
$\sum_kw_kP_k=I_d$.  Define the normalized L\"uders channel
$\Phi(X)=\sum_kw_kP_kXP_k$.  Its full $SU(d)$ conjugation twirl is
\[
\boxed{\overline\Phi(X)=\frac{X+\operatorname{Tr}(X)I_d}{d+1},
\qquad p_{\mathrm{traceless}}=\frac1{d+1}.}
\]
For $d=3$ this maps the preceding nonet channel to
$(X+\operatorname{Tr}(X)I_3)/4$, losing its distinct triplet/quintet
rates.  In the positive-depolarizing parameter family, $1/(d+1)$
is exactly the entanglement-breaking boundary.
\end{theorem}
\begin{proof}
The superoperator Hilbert--Schmidt trace of
$X\mapsto P_kXP_k$ is $(\operatorname{Tr}P_k)^2=1$.
Thus $\operatorname{Tr}_{\mathrm{super}}\Phi=\sum_kw_k=d$
by tracing the POVM normalization.  The twirled channel fixes
$I_d$ and, by irreducibility of the complex traceless adjoint
module, is scalar $p$ on its $d^2-1$ dimensions.
Therefore $1+(d^2-1)p=d$, yielding $p=1/(d+1)$ and the boxed
formula.  Each rank-one L\"uders summand is measure-and-prepare,
so its Haar twirl is entanglement breaking.  Conversely, for a
depolarizing channel with positive parameter $p$, the normalized
Choi partial transpose has antisymmetric eigenvalue
$[1-(d+1)p]/d^2$, negative for $p>1/(d+1)$; hence the boundary
is sharp.
\end{proof}
\begin{remark}[Auxiliary algebra versus physical color]
The nonet Jordan/Lie comparison algebra and a physical
$SU(3)_{\mathrm c}$ gauge representation are not automatically
identified.  This theorem is a covariance diagnostic \emph{if}
that same finite auxiliary channel is ever proposed as fully
$SU(3)$-covariant; it imposes no additional demand on its source's
actual $A_5$ covariance.
\end{remark}
% END SOURCE-INTEGRATED DAG CROSS-CONNECTIONS CC18



% BEGIN SOURCE-INTEGRATED CROSS-DAG MATHEMATICS CC20 (2026-11-16)
% BEGIN MODULE: Paper4B-DirectedPairComparisonMarkov-v1_1.tex
% Mathematical comparison channel, not a physical QCD decay model.
\section{The directed-core coherent Gram as a finite reversible Markov operator}
\label{p4b:sec:cc20-core-pair-markov}
\begin{proposition}[Exact directed-pair stochastic kernel]
\label{p4b:prop:cc20-directed-pair-markov}
For the existing twenty auxiliary rank-one core-pair projectors
$H_{ij}$ with oriented axes $n_{ij}$, define
\[
 K_{ij,kl}=\frac3{20}\operatorname{Tr}(H_{ij}H_{kl})
 =\frac3{80}(1+n_{ij}\cdot n_{kl})^2.
\]
Then $K$ is symmetric, doubly stochastic, positive semidefinite, and
\[
 \boxed{\operatorname{spec}K
 =\{1^{(1)},(1/2)^{(3)},(1/10)^{(5)},0^{(11)}\}.}
\]
For \emph{every} probability vector $p$ on the twenty directed pairs,
including vectors outside the quantum-state image,
\[
 \sum_{j\ne i}(Kp)_{ij}=\sum_{j\ne i}(Kp)_{ji}=\frac15,
\qquad i=1,\ldots,5.
\]
Furthermore, for each directed-pair label $a$ and $k\ge1$,
\[
 (K^k)_{aa}=\frac{1+3\,2^{-k}+5\,10^{-k}}{20}.
\]
The spectrum is a mathematical discrete-step comparison, not a
physical meson decay rate or an action-generated interaction.
\end{proposition}
\begin{proof}
The twenty-projector Gram has spectrum
$\{(20/3)^{(1)},(10/3)^{(3)},(2/3)^{(5)},0^{(11)}\}$.
Rescaling by $3/20$ yields the displayed stochastic spectrum.
Every four-projector outgoing or incoming core sum equals
$(4/3)I_3$; hence the corresponding $K$-row sum equals $1/5$
independently of the incoming probability law.  Transitivity on the
twenty labels makes every diagonal entry of each spectral projector
its rank divided by twenty, proving the return formula.  The
positive Gram form and row sums establish all remaining properties.
\end{proof}

% END MODULE: Paper4B-DirectedPairComparisonMarkov-v1_1.tex
% END SOURCE-INTEGRATED CROSS-DAG MATHEMATICS CC20


% BEGIN FRAMEWORK v3.027: EXACT HIGHER-ORDER FINITE-SECTOR TRANSITIONS
\section{Higher-order finite-sector transition order and all-time exclusions}
\label{p4b:sec:v3027-high-order-sector-transitions}
\begin{theorem}[First nonzero Hamiltonian moment determines short-time order]
\label{p4b:thm:v3027-high-order-sector-transition}
Let $H=H^*$ on a finite Hilbert space of dimension $D$, and let
$P_\alpha,P_\beta$ be mutually orthogonal nonzero sector projectors,
$d_\alpha=\operatorname{rank}P_\alpha$.
For the uniform preparation $\rho_\alpha=P_\alpha/d_\alpha$, set
\[
 p_{\alpha\to\beta}(t)=\operatorname{tr}
 \left[P_\beta e^{-itH/\hbar}\rho_\alpha
  e^{itH/\hbar}P_\beta\right].
\]
If $n_*\ge1$ is the smallest integer with
$P_\beta H^{n_*}P_\alpha\neq0$, then
\[
\boxed{p_{\alpha\to\beta}(t)
=\frac{\|P_\beta H^{n_*}P_\alpha\|_{\rm HS}^2}
{d_\alpha(n_*!)^2\hbar^{2n_*}}t^{2n_*}
+O(t^{2n_*+1}).}
\]
The transition is absent at \emph{every} time if and only if
$P_\beta H^nP_\alpha=0$ for $n=0,\ldots,D-1$, equivalently
$P_\beta\Pi_E P_\alpha=0$ for each distinct spectral projector
$\Pi_E$ of $H$.
\end{theorem}
\begin{proof}
Write the finite exponential as a convergent power series. For
orthogonal sectors its zeroth off-block vanishes, and all lower
terms through $n_*-1$ vanish by definition. The first off-block
amplitude is
$(-it/\hbar)^{n_*}P_\beta H^{n_*}P_\alpha/n_*!$.
The probability is its squared Hilbert--Schmidt norm divided by
$d_\alpha$, giving the formula and remainder. If all moments vanish,
the exponential block vanishes. Conversely, if the exponential block
vanishes for all times, all its derivatives at zero vanish.
Cayley--Hamilton reduces the moment test to $0,\ldots,D-1$;
the spectral decomposition gives the equivalent projective test.
A zero \emph{direct} block alone does not forbid higher-order
indirect paths, and finite unitary recurrences are not irreversible
physical widths.
\end{proof}
\begin{remark}[Geometric path multiplicity is conditional]
\label{p4b:rem:v3027-hurwitz-path-comparison}
If $H=gA_{600}$ is deliberately chosen as a finite scalar-graph
\emph{comparison}, first nonzero powers enumerate weighted paths.
For vertex projectors, the graph distance gives a lower bound on
$n_*$, and for equal positive weights the shortest-path coefficient
is the squared number of shortest walks. No equality between this
comparison and a proper-action QCD/Fock decay Hamiltonian is asserted.
\end{remark}
% END FRAMEWORK v3.027 IV-B MODULE


% BEGIN PHYSICAL OCCUPATION TO SHEAR INTERFACE v1.17
\section{State actualization and metric-source response}
\label{bh:berger:v117:p4b:state-shear}
\begin{definition}[Same-state shape-source packet]
\label{bh:berger:v117:p4b:packet}
On a \emph{physically licensed} finite state $\rho_L$ and a metric
variation $H_L(B)$ derived from its own complete quantum action,
with $B=\sum_{a=1}^5b_aE_a$ on
$\operatorname{Sym}^2_0(\mathbb R^3)$, define
\[
 \mathcal F_a(\rho_L)=
 \left.\operatorname{Tr}(\rho_L\,\partial_{b_a}H_L)\right|_0,
 \qquad
 \mathcal K_{ab}(\rho_L)=\text{complete second source response}.
\]
The second response includes the direct second derivative of the
Hamiltonian, retarded/thermal correlations, finite CAR/Fock contacts,
metric-dependent physical projection/embedding and existing
clock, scalar and boundary terms. These data are not replaced by a
freely chosen angular comparison density.
\end{definition}
\begin{remark}[Distinguish four quantum-state types]
\label{bh:berger:v117:p4b:occupation-gates}
The already proved neutral, unprepared Fock origin carries no
independently occupied higher angular band. An admitted filled-negative
Dirac reference ground may have nonzero tensor \emph{susceptibility}
while its mean shear vanishes; its occurrence on a given physical
band is a separate condition. A fully rotationally invariant Gibbs
state has zero rank-two mean source. Only a sourced, prepared,
measured or dynamically selected anisotropic state may produce a
nonzero first-order homogeneous Berger force.
Spin-three/spin-five geometrical comparisons and their icosahedral
couplings do not establish their physical occupation, transition
widths, or particle pole identifications.
\end{remark}
% END PHYSICAL OCCUPATION TO SHEAR INTERFACE v1.17


\end{document}
